% =====================================================================
%  Advanced Structural Theory -- Lecture Note (LaTeX edition)
%  Compile:  tectonic main.tex
% =====================================================================
\documentclass[a4paper,12pt,oneside]{book}

% ---------- encoding & fonts ----------
\usepackage[T1]{fontenc}
\usepackage[utf8]{inputenc}
\usepackage{lmodern}
\usepackage{microtype}
\usepackage[english]{babel}

% ---------- layout ----------
\usepackage[a4paper,scale=0.8]{geometry}
\usepackage{fancyhdr}
\setlength{\headheight}{15pt}
\usepackage{indentfirst}
\usepackage{setspace}
\setlength{\parindent}{2em}
\setlength{\parskip}{0.35em}
\setstretch{1.08}

% ---------- math ----------
\usepackage{amsmath,amssymb,amsthm,mathtools}
\usepackage{bm}
\usepackage{siunitx}
\usepackage{cancel}
\allowdisplaybreaks
\setcounter{MaxMatrixCols}{20}

% ---------- tables / lists ----------
\usepackage{booktabs}
\usepackage{array}
\usepackage{multirow}
\usepackage{makecell}
\usepackage{longtable}
\usepackage{diagbox}
\usepackage{enumitem}
\setlist{itemsep=0.15em,topsep=0.3em}
\setlist[itemize]{label=\textbullet}

% ---------- figures ----------
\usepackage{graphicx}
\usepackage{float}
\usepackage{caption}
\usepackage{subcaption}
\usepackage{tikz}
\usepackage{pgfplots}
\pgfplotsset{compat=1.18}
\usetikzlibrary{arrows.meta,calc,patterns,patterns.meta,positioning,
  decorations.pathmorphing,decorations.markings,decorations.pathreplacing,
  shapes.misc,shapes.geometric,angles,quotes,matrix,fit,backgrounds,intersections,
  bending,shadings}
\captionsetup{font=small,labelfont=bf}

% ---------- common TikZ styles & pics for structural sketches ----------
\tikzset{
  >=Latex,
  member/.style={line width=1.6pt,black!75},
  thinmember/.style={line width=1.0pt,black!75},
  dim/.style={<->,thin},
  load/.style={->,thick},
  hatch/.style={pattern={Lines[angle=45,distance=3pt]},pattern color=black!70},
  ground/.style={pattern={Lines[angle=45,distance=3pt]},pattern color=black!70},
  node dot/.style={circle,fill=white,draw=black,inner sep=0pt,minimum size=4pt},
  % pin support at the pic origin (member attaches at origin)
  pics/pin/.style={code={
    \draw[fill=black!60] (0,0) -- (-0.22,-0.35) -- (0.22,-0.35) -- cycle;
    \fill[black!25] (-0.35,-0.35) rectangle (0.35,-0.47);
    \fill[white,draw=black] (0,0) circle (0.05);}},
  % roller support
  pics/roller/.style={code={
    \draw[fill=black!60] (0,0) -- (-0.22,-0.35) -- (0.22,-0.35) -- cycle;
    \draw[fill=white] (-0.13,-0.43) circle (0.07); \draw[fill=white] (0.13,-0.43) circle (0.07);
    \fill[black!25] (-0.35,-0.5) rectangle (0.35,-0.62);
    \fill[white,draw=black] (0,0) circle (0.05);}},
  % fixed support: a small hatched block under the origin
  pics/fixed/.style={code={
    \fill[hatch] (-0.35,-0.25) -- (0.35,-0.25) -- (0.35,0) -- (-0.35,0) -- cycle; \draw[thick] (-0.35,0) -- (0.35,0);}},
  % fixed wall support on the left (member goes to the right); polygons so that they rotate properly
  pics/wallL/.style={code={
    \fill[hatch] (-0.25,-0.4) -- (0,-0.4) -- (0,0.4) -- (-0.25,0.4) -- cycle; \draw[thick] (0,-0.4) -- (0,0.4);}},
  pics/wallR/.style={code={
    \fill[hatch] (0,-0.4) -- (0.25,-0.4) -- (0.25,0.4) -- (0,0.4) -- cycle; \draw[thick] (0,-0.4) -- (0,0.4);}},
  % counter-clockwise moment arrow of given radius
  pics/ccw/.style={code={\draw[->,thick] (-30:#1) arc (-30:210:#1);}},
  pics/cw/.style={code={\draw[->,thick] (210:#1) arc (210:-30:#1);}},
}

% ---------- boxes ----------
\usepackage[most]{tcolorbox}
\tcbset{
  notebox/.style={colback=black!3,colframe=black!60,boxrule=0.5pt,arc=1.5pt,
    left=6pt,right=6pt,top=4pt,bottom=4pt,fonttitle=\bfseries,breakable},
}
\newtcolorbox{remark}[1][]{notebox,title={Remark},#1}
\newtcolorbox{keypoint}[1][]{notebox,colframe=blue!50!black,colback=blue!3,title={Key point},#1}

% ---------- code ----------
\usepackage{listings}
\usepackage{xcolor}
\definecolor{dkgreen}{rgb}{0,0.5,0}
\definecolor{mgray}{rgb}{0.5,0.5,0.5}
\definecolor{mauve}{rgb}{0.58,0,0.82}
\lstset{frame=single,language=Fortran,aboveskip=3mm,belowskip=3mm,
  showstringspaces=false,columns=flexible,basicstyle={\small\ttfamily},
  numbers=left,numberstyle=\tiny\color{mgray},keywordstyle=\color{blue},
  commentstyle=\color{dkgreen},stringstyle=\color{mauve},breaklines=true,
  breakatwhitespace=true,tabsize=4}

% ---------- table of contents: wider number boxes ----------
\makeatletter
\renewcommand*\l@section{\@dottedtocline{1}{1.5em}{3.0em}}
\renewcommand*\l@subsection{\@dottedtocline{2}{4.5em}{3.6em}}
\makeatother

% ---------- hyperref (last) ----------
\usepackage[hidelinks,bookmarksnumbered]{hyperref}

% ---------- headers ----------
\pagestyle{fancy}
\fancyhf{}
\fancyhead[L]{\nouppercase{\leftmark}}
\fancyhead[R]{\textsc{Advanced Structural Theory}}
\fancyfoot[C]{\thepage}
\renewcommand{\headrulewidth}{0.8pt}
\renewcommand{\footrulewidth}{0pt}
\renewcommand{\chaptermark}[1]{\markboth{\thechapter.\ #1}{}}
\fancypagestyle{plain}{\fancyhf{}\fancyfoot[C]{\thepage}\renewcommand{\headrulewidth}{0pt}}

% ---------- theorem-like ----------
\theoremstyle{definition}
\newtheorem{example}{Example}[chapter]
\newtheorem{definition}{Definition}[chapter]
\theoremstyle{remark}
\newtheorem*{note}{Note}
\newtheorem*{solution}{Solution}

% =====================================================================
%  Macros
% =====================================================================
\newcommand{\HRule}{\rule{\linewidth}{2pt}}
\newcommand{\mat}[1]{\mathbf{#1}}                 % matrix
\newcommand{\vect}[1]{\mathbf{#1}}                % vector
\newcommand{\bs}[1]{\bm{#1}}                      % bold symbol (greek)
\newcommand{\tr}{\mathsf{T}}                      % transpose
\newcommand{\dd}{\mathrm{d}}
\newcommand{\pd}[2]{\frac{\partial #1}{\partial #2}}
\newcommand{\od}[2]{\frac{\dd #1}{\dd #2}}
\newcommand{\bmat}[1]{\begin{bmatrix}#1\end{bmatrix}}
\newcommand{\pmat}[1]{\begin{pmatrix}#1\end{pmatrix}}
\newcommand{\Bmat}[1]{\begin{Bmatrix}#1\end{Bmatrix}}   % column vector with braces
\newcommand{\vmat}[1]{\begin{vmatrix}#1\end{vmatrix}}
\newcommand{\diag}{\operatorname{diag}}
\newcommand{\sym}{\text{sym.}}
\newcommand{\F}{\mat{F}}
\newcommand{\Kmat}{\mat{K}}
\newcommand{\kmat}{\mat{k}}
\newcommand{\Tmat}{\mat{T}}
\newcommand{\dvec}{\mat{d}}
\newcommand{\Dvec}{\mat{D}}
\newcommand{\Pvec}{\mat{P}}
\newcommand{\pvec}{\mat{p}}
\newcommand{\evec}{\mat{e}}
\newcommand{\dW}{\delta W}
\newcommand{\dU}{\delta U}
\newcommand{\Dl}{\Delta}
\newcommand{\lb}{\left}
\newcommand{\rb}{\right}
\newcommand{\eps}{\varepsilon}
\newcommand{\Var}{\text{Var}}
\newcommand{\NN}{\mathbf{N}}
\newcommand{\Bm}{\mathbf{B}}

% run-in style headings (safe before figures)
\newcommand{\topic}[1]{\par\medskip\noindent\textbf{#1}\par\nopagebreak\smallskip\noindent\ignorespaces}
\newcommand{\topicnote}[2]{\par\medskip\noindent\textbf{#1}\ #2\par\nopagebreak\smallskip\noindent\ignorespaces}
% non-floating centered figure
\newenvironment{fig}{\par\medskip\begin{center}}{\end{center}\par\medskip}

\graphicspath{{./}{./figures/}}

% =====================================================================
\begin{document}

% ---------------------------- Title page ----------------------------
\begin{titlepage}
\begin{center}
\includegraphics[width=0.25\textwidth]{logo.pdf}\\[1cm]
\textsc{\LARGE National Taiwan University}\\[1.5cm]
\textsc{\Large Advanced Structural Theory}\\[0.5cm]
\HRule\\[0.5cm]
{\huge\bfseries Lecture Note}\\[0.5cm]
\HRule\\[2.5cm]
\begin{minipage}{0.45\textwidth}
\begin{flushleft}\large
\emph{Author:}\\
Chun-Yu Ke
\end{flushleft}
\end{minipage}
\begin{minipage}{0.45\textwidth}
\begin{flushright}\large
\emph{Typeset in \LaTeX{} by:}\\
Gauss Chang
\end{flushright}
\end{minipage}
\vfill
{\large\today}
\end{center}
\end{titlepage}

\frontmatter
\tableofcontents

\mainmatter
\chapter{Introduction}

\section{Purpose of this course}
To calculate the responses, including internal forces (axial, shear, moment) distribution and displacements (deflections), of framed structures.

\section{Five elements of structural analysis}
\begin{enumerate}
  \item \textbf{Basic mechanics:}
    \begin{enumerate}[label=(\alph*)]
      \item stress--strain relationships
      \item equilibrium equations
      \item compatibility equations
    \end{enumerate}
  \item \textbf{Finite element mechanics}, or member force--deformation relationships
    \begin{fig}
      % Simply supported beam with end moments M1(alpha1), M2(alpha2)
\begin{tikzpicture}[scale=1.1,>=Latex]
  \def\L{6}
  % beam
  \draw[line width=2.2pt,black!70] (0,0) -- (\L,0);
  \node[above] at (\L/2,0.05) {$EI$};
  % pin support (left)
  \draw[fill=black!60] (0,0) -- (-0.25,-0.4) -- (0.25,-0.4) -- cycle;
  \fill[white,draw=black] (0,0) circle (0.06);
  \fill[black!25] (-0.4,-0.4) rectangle (0.4,-0.55);
  % roller support (right)
  \draw[fill=black!60] (\L,0) -- (\L-0.25,-0.4) -- (\L+0.25,-0.4) -- cycle;
  \fill[white,draw=black] (\L,0) circle (0.06);
  \draw[fill=white] (\L-0.15,-0.5) circle (0.09);
  \draw[fill=white] (\L+0.15,-0.5) circle (0.09);
  \fill[black!25] (\L-0.4,-0.6) rectangle (\L+0.4,-0.75);
  % end moments
  \draw[->,thick] (-0.05,0.55) arc (150:-30:0.35);
  \node[above] at (0.2,0.75) {$M_1(\alpha_1)$};
  \draw[->,thick] (\L+0.05,0.55) arc (30:210:0.35);
  \node[above] at (\L-0.2,0.75) {$M_2(\alpha_2)$};
  % length dimension
  \draw[<->] (0,-1.1) -- (\L,-1.1) node[midway,fill=white,inner sep=1pt] {$L$};
  \draw (0,-0.85) -- (0,-1.3);
  \draw (\L,-0.85) -- (\L,-1.3);
\end{tikzpicture}

    \end{fig}
    \begin{itemize}
      \item \emph{Stiffness methods}: forces in terms of displacements through stiffness coefficients
        \begin{equation}
          \Bmat{M_1\\ M_2} = \frac{2EI}{L}\bmat{2 & 1\\ 1 & 2}\Bmat{\alpha_1\\ \alpha_2}
        \end{equation}
      \item \emph{Flexibility methods}: deformations in terms of forces through flexibility coefficients
        \begin{equation}
          \Bmat{\alpha_1\\ \alpha_2} = \frac{L}{6EI}\bmat{2 & -1\\ -1 & 2}\Bmat{M_1\\ M_2}
        \end{equation}
    \end{itemize}
  \item \textbf{Formulation of structural (global) equations}\\
    Meaning of structural equations:
    \begin{itemize}
      \item Stiffness methods: equilibrium equations
      \item Flexibility methods: compatibility equations (consistent deformation)
    \end{itemize}
  \item \textbf{Equation solving}\\
    Gauss elimination, decomposition methods, etc.
  \item \textbf{Solution interpretation}
    \begin{itemize}
      \item Member force diagrams such as axial, shear and moment diagrams
      \item Deformed shapes
    \end{itemize}
\end{enumerate}

\section{Types of analysis}
(See also Sec.~8.1 of the textbook.)

\subsection*{Linear analysis}

\subsection*{Nonlinear analysis}

\topicnote{Geometrically nonlinear analysis}{(second-order analysis, $P$-$\Delta$ effects)}
\begin{fig}
  % P-Delta effect on a cantilever column: three stages
\begin{tikzpicture}[scale=1.0,>=Latex,font=\small]
  \def\h{3.2}
  % ---------- (a) undeformed ----------
  \begin{scope}
    \fill[black!45] (-0.55,-0.35) rectangle (0.55,0);
    \draw[fill=white] (-0.09,0) rectangle (0.09,\h);
    \draw[->,thick] (0,\h+0.75) -- (0,\h+0.08) node[pos=0,above] {$P$};
    \draw[->,thick] (-1.0,\h) -- (-0.15,\h) node[pos=0,left] {$V$};
    \draw[<->] (-1.3,0) -- (-1.3,\h) node[midway,left] {$h$};
    % reactions
    \draw[->,thick] (1.1,-0.18) -- (0.6,-0.18) node[pos=0,right] {$V$};
    \draw[->,thick] (0,-0.8) -- (0,-0.4) node[pos=0,below] {$P$};
    \draw[->,thick] (-0.75,-0.75) arc (-120:20:0.3);
    \node[left] at (-0.8,-0.6) {$Vh$};
  \end{scope}
  \draw[-{Triangle[width=6pt,length=6pt]},line width=1.6pt] (1.6,\h/2) -- (2.3,\h/2);
  % ---------- (b) first-order deformed ----------
  \begin{scope}[xshift=4.2cm]
    \fill[black!45] (-0.55,-0.35) rectangle (0.55,0);
    \draw[thick] (0,0) .. controls (0,\h*0.55) and (0.15,\h*0.85) .. (0.55,\h);
    \draw[->,thick] (-0.9,\h) -- (0.4,\h) node[pos=0,left] {$V$};
    \draw[->,thick] (1.1,-0.18) -- (0.6,-0.18) node[pos=0,right] {$V$};
    \draw[->,thick] (-0.75,-0.75) arc (-120:20:0.3);
    \node[left] at (-0.8,-0.6) {$Vh$};
  \end{scope}
  % ---------- (c) second-order deformed ----------
  \begin{scope}[xshift=8.8cm]
    \fill[black!45] (-0.55,-0.35) rectangle (0.55,0);
    \draw[densely dashed] (0,0) .. controls (0,\h*0.55) and (0.15,\h*0.85) .. (0.55,\h);
    \draw[thick] (0,0) .. controls (0,\h*0.5) and (0.2,\h*0.85) .. (0.85,\h);
    \draw[->,thick] (0.85,\h+0.75) -- (0.85,\h+0.08) node[pos=0,above] {$P$};
    \draw[->,thick] (-0.9,\h) -- (0.65,\h) node[pos=0,left] {$V$};
    \draw[->,thick] (1.1,-0.18) -- (0.6,-0.18) node[pos=0,right] {$V$};
    \draw[->,thick] (0,-0.8) -- (0,-0.4) node[pos=0,below] {$P$};
    \draw[->,thick] (-0.75,-0.75) arc (-120:20:0.3);
    \node[left] at (-0.8,-0.6) {$Vh+P\Delta$};
    \draw[<->] (0,\h+0.35) -- (0.85,\h+0.35) node[midway,above] {$\Delta$};
  \end{scope}
  % ---------- formulas ----------
  \node[anchor=west,align=left] at (11.0,\h*0.55) {%
    $\Delta \cong \Delta_L \times \dfrac{1}{1-\dfrac{P}{P_E}}$\\[10pt]
    $\Delta_L=\dfrac{Vh^3}{3EI}$\\[8pt]
    $P_E=\dfrac{\pi^2EI}{4h^2}$};
\end{tikzpicture}

\end{fig}

\topic{Materially nonlinear analysis}
\begin{fig}
  % Materially nonlinear analysis: nonlinear elastic vs. plastic
\begin{tikzpicture}[scale=1.0,>=Latex,font=\small]
  % ---------- nonlinear elastic ----------
  \begin{scope}
    \draw[->] (0,0) -- (3.4,0) node[right] {$\varepsilon$};
    \draw[->] (0,0) -- (0,2.6) node[above] {$\sigma$};
    \draw[thick] (0,0) .. controls (0.6,1.7) and (1.6,2.2) .. (3.0,2.3);
    % loading / unloading arrows along the same curve
    \draw[->,thin] (0.35,0.85) .. controls (0.5,1.2) .. (0.85,1.55);
    \draw[<-,thin] (1.3,1.83) .. controls (1.5,1.95) .. (1.9,2.05);
    \node[left] at (-0.2,1.3) {\itshape elastic};
  \end{scope}
  % ---------- plastic ----------
  \begin{scope}[xshift=6.8cm]
    \draw[->] (0,0) -- (3.4,0) node[right] {$\varepsilon$};
    \draw[->] (0,0) -- (0,2.6) node[above] {$\sigma$};
    % hysteresis parallelogram, hatched
    \fill[pattern={Lines[angle=90,distance=3pt]},pattern color=black!70]
      (0,0) -- (0.8,1.9) -- (2.4,1.9) -- (1.6,0) -- cycle;
    \draw[thick] (0,0) -- (0.8,1.9) -- (2.4,1.9) -- (1.6,0) -- cycle;
    \draw[->,thin] (0.35,0.83) -- (0.5,1.19);
    \draw[->,thin] (1.4,1.9) -- (1.7,1.9);
    \draw[->,thin] (2.05,1.07) -- (1.9,0.71);
    \node[left] at (-0.2,1.1) {\itshape plastic};
  \end{scope}
\end{tikzpicture}

\end{fig}

\topicnote{Nonlinear analysis due to change in boundary conditions}{(such as contact problems)}
\begin{fig}
  % Contact problem: a ball pressed onto a rigid surface
\begin{tikzpicture}[scale=1.0,>=Latex]
  % ground
  \fill[pattern={Lines[angle=45,distance=4pt]},pattern color=black!70] (-1.6,-0.35) rectangle (1.6,0);
  \draw[thick] (-1.6,0) -- (1.6,0);
  % ball (shaded)
  \shade[ball color=black!45] (0,1.0) circle (1.0);
  % load
  \draw[->,very thick] (0,3.0) -- (0,2.1);
  \node[right] at (0.1,2.75) {$P(\Delta)$};
\end{tikzpicture}

\end{fig}
Considering a ball subjected to a load $P$, the contact area will increase as the load increases.

\subsection*{Example of geometrically nonlinear analysis}
A simple example, the \emph{two-member truss}, illustrating the geometrically nonlinear analysis is as follows.
\begin{fig}
  % Two-member (von Mises) truss: geometrically nonlinear example
\begin{tikzpicture}[scale=1.0,>=Latex]
  \def\Lx{5.0} \def\h{2.0} \def\dl{0.7}
  \coordinate (A) at (0,0);
  \coordinate (B) at (2*\Lx,0);
  \coordinate (C0) at (\Lx,\h);
  \coordinate (C1) at (\Lx,\h-\dl);
  % undeformed (dotted)
  \draw[densely dotted] (A) -- (C0) -- (B);
  % baseline (dashed)
  \draw[dashed] (A) -- (B);
  % deformed (thick)
  \draw[line width=2.5pt,black!70] (A) -- (C1) -- (B);
  % supports
  \foreach \p in {A,B} {
    \draw[fill=black!60] (\p) -- ++(-0.25,-0.4) -- ++(0.5,0) -- cycle;
    \fill[black!25] ($(\p)+(-0.4,-0.4)$) rectangle ($(\p)+(0.4,-0.55)$);
    \fill[white,draw=black] (\p) circle (0.06);
  }
  \fill[white,draw=black] (C1) circle (0.07);
  % load
  \draw[->,very thick] ($(C0)+(0,1.0)$) -- ($(C1)+(0,0.1)$);
  \node[above] at ($(C0)+(0,1.0)$) {$P$};
  % labels
  \node[above left,inner sep=1pt] at ($(A)!0.55!(C0)$) {$L$};
  \node[above,inner sep=1pt] at ($(A)!0.62!(C1)$) {$L'$};
  % angles
  \draw[->] ($(A)+(1.4,0)$) arc (0:{atan(\h/\Lx)}:1.4);
  \node at ($(A)+(1.65,0.22)$) {$\theta$};
  \draw[->] ($(A)+(2.3,0)$) arc (0:{atan((\h-\dl)/\Lx)}:2.3);
  \node at ($(A)+(2.55,0.22)$) {$\theta'$};
  % dimension lines
  \draw[thin] (C0) -- ($(B)+(1.0,\h)$);
  \draw[thin] (C1) -- ($(B)+(1.0,\h-\dl)$);
  \draw[thin] (B) -- ($(B)+(1.0,0)$);
  \draw[<->] ($(B)+(0.6,\h)$) -- ($(B)+(0.6,\h-\dl)$) node[midway,right] {$\Delta$};
  \draw[<->] ($(B)+(0.9,\h-\dl)$) -- ($(B)+(0.9,0)$) node[midway,right] {$h-\Delta$};
\end{tikzpicture}

\end{fig}

\begin{solution}
For a displacement $\Delta$, we have
\begin{equation}
  L' = \sqrt{(h-\Delta)^2 + (L\cos\theta)^2}.
\end{equation}
Knowing $\eps = \dfrac{L-L'}{L}$ and the stress--strain relationship $\sigma = f(\eps)$, we can first obtain the stress corresponding to this displacement and then the axial force $F = \sigma A$. Finally, the equilibrium equation states that
\begin{equation}
  P = 2F\sin\theta', \qquad \text{where } \sin\theta' = (h-\Delta)/L'.
\end{equation}
It shall be noted that the equilibrium equation is written for the \emph{deformed configuration}, which is the main feature of geometric nonlinear analysis. With different stress--strain relationships, the $P$--$\Delta$ curves are of different shapes, but all the curves pass through $(0,0)$, $(h,0)$ and $(2h,0)$.
\end{solution}
\begin{fig}
  % Load-displacement curve of the two-member truss (snap-through)
\begin{tikzpicture}[>=Latex]
  \begin{axis}[
    width=12cm,height=6cm,
    axis lines=middle,
    xlabel={$\Delta$},ylabel={$P$},
    xlabel style={at={(ticklabel* cs:1)},anchor=west},
    ylabel style={at={(ticklabel* cs:1)},anchor=south},
    xmin=0,xmax=2.6,ymin=-1.4,ymax=1.9,
    xtick={0.1,1,2},xticklabels={$\Delta_i$,$h$,$2h$},
    ytick=\empty,
    clip=false,
    samples=120,domain=0:2.35,
  ]
    % P proportional to (h-D)(D)(2h-D) (shape of the exact curve)
    \addplot[thick,black] {2.6*(1-x)*x*(2-x)};
    % P_lim and P_i
    \pgfmathsetmacro{\plim}{2.6*(1-0.4226)*0.4226*(2-0.4226)}
    \pgfmathsetmacro{\pI}{2.6*(1-0.1)*0.1*(2-0.1)}
    \draw[densely dotted] (axis cs:0,\plim) -- (axis cs:2.2,\plim);
    \draw[densely dotted] (axis cs:0,\pI) -- (axis cs:0.1,\pI) -- (axis cs:0.1,0);
    \node[left] at (axis cs:0,\plim) {$P_{\lim}$};
    \node[left] at (axis cs:0,\pI) {$P_i$};
  \end{axis}
\end{tikzpicture}

\end{fig}

Another famous example related to geometric nonlinear analysis is a horizontal cable fixed at both ends and with a pretension $T$. Assume the length of the cable is $2L$ and the cable is subjected to a vertical load $P$ applied at the midpoint. Assuming $T$ is much larger than $P$, calculate the vertical deflection at the loaded point.
\begin{fig}
  % Pretensioned cable with a concentrated load at midspan
\begin{tikzpicture}[scale=1.0,>=Latex]
  \def\L{4.0} \def\v{0.8}
  \coordinate (A) at (0,0); \coordinate (B) at (2*\L,0); \coordinate (C) at (\L,-\v);
  % undeformed cable (dotted)
  \draw[densely dotted] (A) -- (B);
  % deformed cable (double line)
  \draw[double,double distance=1.5pt] (A) -- (C) -- (B);
  % supports
  \foreach \p in {A,B} {
    \draw[fill=black!60] (\p) -- ++(-0.25,-0.4) -- ++(0.5,0) -- cycle;
    \fill[black!25] ($(\p)+(-0.4,-0.4)$) rectangle ($(\p)+(0.4,-0.55)$);
    \fill[white,draw=black] (\p) circle (0.06);
  }
  \draw[fill=white] (C) circle (0.12); \fill (C) circle (0.05);
  % tension arrows (cut)
  \draw[->,very thick] ($(A)+(-0.3,0)$) -- ++(-1.1,0) node[left] {$T$};
  \draw[->,very thick] ($(B)+(0.3,0)$) -- ++(1.1,0) node[right] {$T$};
  % load
  \draw[->,very thick] (C) -- ++(0,-1.0) node[right] {$P$};
  % labels
  \node[below left] at ($(A)!0.5!(C)$) {$L'$};
  \draw[->] ($(A)+(2.2,0)$) arc (0:{-atan(\v/\L)}:2.2);
  \node[right] at ($(A)+(2.25,-0.2)$) {$\theta$};
  % dimensions
  \draw[<->] (-0.7,0) -- (-0.7,-\v) node[midway,left] {$v$};
  \draw[thin] (-0.45,0) -- (-0.9,0); \draw[thin] (C) -- (-0.9,-\v);
  \draw[<->] (0,-1.4) -- (\L,-1.4) node[midway,fill=white,inner sep=1pt] {$L$};
  \draw[thin] (0,-1.2) -- (0,-1.6); \draw[thin] (\L,-1.2) -- (\L,-1.6);
\end{tikzpicture}

\end{fig}

\begin{solution}
The equilibrium equation for the deformed configuration is
\begin{equation}
  2(T+\Delta F)\sin\theta = P,
\end{equation}
where $\Delta F$ is the increment of the tension. Since $T$ is much larger than $P$, $\Delta F$ must be very small compared with $T$. Also, $\sin\theta \approx v/L$; hence,
\begin{equation}
  2T\frac{v}{L} = P \quad\Longrightarrow\quad v = \frac{PL}{2T}.
\end{equation}
The relation $v \propto \dfrac{P}{T}$ is what we would expect.
\end{solution}

\section{Different levels of analysis}
\begin{figure}[H]
  \centering
  % Levels of analysis (after McGuire, Gallagher & Ziemian, Fig. 8.1)
\begin{tikzpicture}[>=Latex,font=\footnotesize]
  \begin{axis}[
    width=13.5cm,height=9.5cm,
    axis lines=left,
    xlabel={Lateral deflection, $\Delta$},
    ylabel={Lateral load, $H$},
    xmin=0,xmax=10,ymin=0,ymax=10,
    xtick=\empty,
    ytick={8.6,7.6,5.2,3.9,2.8},
    yticklabels={$H_{ec}$,$H_{es}$,$H_{ic}$,$H_p$,$H_{is}$},
    clip=false,
  ]
    % horizontal limit lines
    \draw[dashed] (axis cs:0,8.6) -- (axis cs:9.6,8.6);
    \node[above right,inner sep=1pt] at (axis cs:6.2,8.6) {Elastic critical load};
    \draw[dashed] (axis cs:0,7.6) -- (axis cs:9.6,7.6);
    \node[above right,inner sep=1pt] at (axis cs:6.2,7.6) {Elastic stability limit};
    \draw[dashed] (axis cs:0,5.2) -- (axis cs:9.6,5.2);
    \node[above right,inner sep=1pt] at (axis cs:6.2,5.2) {Inelastic critical load};
    \draw[dashed] (axis cs:0,3.9) -- (axis cs:9.6,3.9);
    \node[above right,inner sep=1pt] at (axis cs:6.2,3.9) {Plastic limit load};
    \draw[dashed] (axis cs:0,2.8) -- (axis cs:9.6,2.8);
    \node[above right,inner sep=1pt] at (axis cs:6.2,2.8) {Inelastic stability limit};
    % first-order elastic: straight line
    \draw[thick] (axis cs:0,0) -- (axis cs:3.4,10);
    \node[right,align=left,fill=white,inner sep=1pt] at (axis cs:3.5,9.4) {First-order\\elastic analysis};
    % second-order elastic: softening, asymptote to elastic stability limit
    \draw[thick] (axis cs:0,0) .. controls (axis cs:1.4,4.2) and (axis cs:3.0,7.4) .. (axis cs:6.2,7.6)
                 .. controls (axis cs:7.0,7.6) .. (axis cs:7.6,7.3);
    % bifurcation branch at H_ec (elastic)
    \draw[thick] (axis cs:0.9,8.6) -- (axis cs:3.7,8.6);
    \node[left] at (axis cs:2.6,8.95) {Bifurcation};
    % bifurcation on 2nd-order elastic curve
    \draw[thick] (axis cs:3.9,6.6) .. controls (axis cs:4.6,6.4) .. (axis cs:5.4,6.1);
    \node[right] at (axis cs:4.0,6.95) {Bifurcation};
    \node[right,align=left,fill=white,inner sep=1pt] at (axis cs:4.3,5.4) {Second-order\\elastic analysis};
    % first-order inelastic: follows first-order elastic then flattens to H_p
    \draw[thick] (axis cs:0,0) -- (axis cs:1.05,3.1) .. controls (axis cs:1.6,3.8) .. (axis cs:2.6,3.9)
                 -- (axis cs:5.8,3.9);
    \node[right,align=left,fill=white,inner sep=1pt] at (axis cs:3.2,4.5) {First-order\\inelastic analysis};
    % bifurcation at H_ic (inelastic)
    \draw[thick] (axis cs:0.6,5.2) -- (axis cs:2.6,5.2);
    \node[left] at (axis cs:1.6,5.5) {Bifurcation};
    % second-order inelastic: peaks at H_is then descends
    \draw[thick] (axis cs:0,0) -- (axis cs:0.85,2.5) .. controls (axis cs:1.2,2.8) and (axis cs:1.8,2.9) .. (axis cs:2.6,2.8)
                 .. controls (axis cs:3.4,2.6) .. (axis cs:4.4,1.7);
    \node[right,align=left] at (axis cs:1.7,1.5) {Second-order\\inelastic analysis};
    % small inset frame
    \begin{scope}[shift={(axis cs:7.4,0.4)},x=0.22cm,y=0.22cm]
      \draw[thick] (0,0) -- (0,7) -- (5,7) -- (5,0);
      \draw[thick] (0,3.5) -- (5,3.5);
      \fill[pattern={Lines[angle=45,distance=2pt]}] (-0.6,-0.5) rectangle (0.6,0);
      \fill[pattern={Lines[angle=45,distance=2pt]}] (4.4,-0.5) rectangle (5.6,0);
      \draw[->] (-1.8,7) -- (-0.1,7) node[pos=0,left] {\tiny $H=\alpha P$};
      \draw[->] (-1.8,3.5) -- (-0.1,3.5) node[pos=0,left] {\tiny $H=\alpha P$};
      \draw[->] (0.6,8.6) -- (0.6,7.2) node[pos=0,above] {\tiny $P$};
      \draw[->] (4.4,8.6) -- (4.4,7.2);
      \draw[->] (0.6,5.1) -- (0.6,3.7) node[pos=0,above] {\tiny $P$};
      \draw[->] (4.4,5.1) -- (4.4,3.7);
      \draw[<->] (5.3,0.2) -- (5.3,7) ; \node[right] at (5.3,3.5) {\tiny $\Delta$};
    \end{scope}
  \end{axis}
\end{tikzpicture}

  \caption{Levels of analysis (after Fig.~8.1 of the textbook).}
\end{figure}

\chapter{Definitions and Concepts}

\section{Mathematical model}
A framed structure is composed of \emph{members} and \emph{joints}. Joints can be classified as rigid, pin (simple), and semi-rigid. Joints are commonly referred to as connections; they are the structural components that tie the members together. For an analysis mathematical model, we like to say that the model is composed of \emph{elements} and \emph{nodes} instead of members and joints.

\section{Coordinate systems}
\begin{itemize}
  \item Member (local) coordinate system $(xyz)$
  \item Structural (global) coordinate system $(XYZ)$
\end{itemize}
\begin{fig}
  % Gable frame: local (xyz) vs. global (XYZ) coordinate systems
\begin{tikzpicture}[scale=1.0]
  \coordinate (A) at (0,0); \coordinate (B) at (0,2.2);
  \coordinate (C) at (1.6,3.3); \coordinate (D) at (3.2,2.2); \coordinate (E) at (3.2,0);
  \draw[member] (A) -- (B) -- (C) -- (D) -- (E);
  \pic at (A) {fixed}; \pic at (E) {fixed};
  % global axes
  \draw[->] (A) -- ++(1.1,0) node[right] {$X$};
  \draw[->] (A) -- ++(0,1.1) node[above] {$Y$};
  % local axes on member BC
  \draw[->] (B) -- ($(B)!0.55!(C)$) node[above left] {$x$};
  \draw[->] (B) -- ($(B)+(-0.6,0.87)$) node[above] {$y$};
\end{tikzpicture}

\end{fig}

\section{Degrees of freedom (DOFs)}
DOFs are kinematic parameters such as displacements and rotations used to describe the configuration (deformed shape) of the structure. For an analysis model, each node will be assigned appropriate DOFs according to its restraints.

\begin{remark}
\begin{enumerate}
  \item The number of DOFs is infinity for structures theoretically.
  \item Different types of analysis such as static, dynamic, and buckling may need different numbers of DOFs to achieve the same level of accuracy. For example, for buckling and dynamic analyses, a member usually needs to be divided into four elements. However, for linear static analysis, no discretization of a member is needed in general.
\end{enumerate}
\end{remark}

\section{External loads, nodal displacements, and member displacements (forces)}
\topic{External (applied) loads}
\begin{itemize}
  \item Nodal loads: concentrated; referred to the $XYZ$ coordinate system
  \item Member loads: distributed; referred to the $xyz$ coordinate system
\end{itemize}

\topic{Nodal displacements}
Referred to the $XYZ$ coordinate system.

\topic{Member displacements (forces)}
Referred to the $xyz$ coordinate system.

\section{Static and kinematic indeterminacy}
Please refer to Sec.~3.4 and Fig.~3.7 (shown below) of the textbook. Note that:
\begin{itemize}
  \item Degree of \emph{static} indeterminacy $=$ the number of unknowns in force (flexibility) methods.
  \item Degree of \emph{kinematic} indeterminacy $=$ the number of unknowns in displacement (stiffness) methods.
\end{itemize}
\begin{figure}[H]
  \centering
  % Static and kinematic indeterminacy (after Fig. 3.7 of the textbook)
\begin{tikzpicture}[scale=0.75,font=\footnotesize,
  redund/.style={->,thick},
  dofm/.style={->,thick,black!60}]
  % column headers
  \node at (1.6,1.1) {Structure};
  \node at (6.6,1.1) {Forces};
  \node[align=center] at (11.6,1.3) {Degrees of\\Freedom};
  \node[align=center] at (15.0,1.3) {Indeterminacy\\[-1pt]Static \quad Kinematic};
  % helper: a load arrow pointing down-right onto a point
  \newcommand{\ld}[1]{\draw[load] ($#1+(-0.45,0.6)$) -- ($#1+(-0.05,0.07)$);}
  % helper: a small rotational dof symbol
  \newcommand{\rot}[1]{\draw[dofm] ($#1+(0.3,0.05)$) arc (0:180:0.3);}
  \newcommand{\transv}[1]{\draw[dofm] ($#1+(0,-0.05)$) -- ($#1+(0,-0.6)$);}
  \newcommand{\transh}[1]{\draw[dofm] ($#1+(0.05,0)$) -- ($#1+(0.6,0)$);}
  %
  % ---------------- (a) simply supported beam ----------------
  \begin{scope}[yshift=0cm]
    \node at (-0.6,0) {$(a)$};
    \draw[member] (0,0) -- (3.2,0); \pic at (0,0) {pin}; \pic at (3.2,0) {roller};
    \ld{(1.4,0)}
    % forces (free body)
    \begin{scope}[xshift=5cm]
      \draw[member] (0,0) -- (3.2,0); \ld{(1.4,0)}
      \draw[redund] (-0.7,0) -- (-0.1,0); \draw[redund] (0,-0.7) -- (0,-0.1); \draw[redund] (3.2,-0.7) -- (3.2,-0.1);
    \end{scope}
    % dofs
    \begin{scope}[xshift=10cm]
      \draw[member] (0,0) -- (3.2,0); \pic at (0,0) {pin}; \pic at (3.2,0) {roller}; \ld{(1.4,0)}
      \rot{(0,0)} \rot{(3.2,0)}
    \end{scope}
    \node at (14.3,0) {0}; \node at (15.8,0) {2};
  \end{scope}
  % ---------------- (b) two-span continuous beam ----------------
  \begin{scope}[yshift=-2.2cm]
    \node at (-0.6,0) {$(b)$};
    \draw[member] (0,0) -- (3.2,0); \pic at (0,0) {pin}; \pic at (1.6,0) {roller}; \pic at (3.2,0) {roller};
    \ld{(0.8,0)} \ld{(2.4,0)}
    \begin{scope}[xshift=5cm]
      \draw[member] (0,0) -- (3.2,0); \ld{(0.8,0)} \ld{(2.4,0)}
      \draw[redund] (-0.7,0) -- (-0.1,0); \draw[redund] (0,-0.7) -- (0,-0.1);
      \draw[redund] (1.6,-0.7) -- (1.6,-0.1); \draw[redund] (3.2,-0.7) -- (3.2,-0.1);
    \end{scope}
    \begin{scope}[xshift=10cm]
      \draw[member] (0,0) -- (3.2,0); \pic at (0,0) {pin}; \pic at (1.6,0) {roller}; \pic at (3.2,0) {roller};
      \ld{(0.8,0)} \ld{(2.4,0)}
      \rot{(0,0)} \rot{(1.6,0)} \rot{(3.2,0)}
    \end{scope}
    \node at (14.3,0) {1}; \node at (15.8,0) {3};
  \end{scope}
  % ---------------- (c) fixed-fixed beam ----------------
  \begin{scope}[yshift=-4.4cm]
    \node at (-0.6,0) {$(c)$};
    \draw[member] (0,0) -- (3.2,0); \pic at (0,0) {wallL}; \pic at (3.2,0) {wallR};
    \ld{(1.4,0)}
    \begin{scope}[xshift=5cm]
      \draw[member] (0,0) -- (3.2,0); \ld{(1.4,0)}
      \draw[redund] (-0.7,0) -- (-0.1,0); \draw[redund] (0,-0.7) -- (0,-0.1); \draw[redund] (0.05,-0.35) arc (-90:150:0.35);
      \draw[redund] (3.2,-0.7) -- (3.2,-0.1); \draw[redund] (3.9,0) -- (3.3,0); \draw[redund] (3.15,-0.35) arc (270:30:0.35);
    \end{scope}
    \begin{scope}[xshift=10cm]
      \draw[member] (0,0) -- (3.2,0); \pic at (0,0) {wallL}; \pic at (3.2,0) {wallR}; \ld{(1.4,0)}
    \end{scope}
    \node at (14.3,0) {3}; \node at (15.8,0) {0};
  \end{scope}
  % ---------------- (d) portal frame ----------------
  \begin{scope}[yshift=-7.4cm]
    \node at (-0.6,0.6) {$(d)$};
    \draw[member] (0,0) -- (0,2) -- (3,2) -- (3,0); \pic at (0,0) {fixed}; \pic at (3,0) {fixed};
    \ld{(1.3,2)}
    \begin{scope}[xshift=5cm]
      \draw[member] (0,0) -- (0,2) -- (3,2) -- (3,0); \ld{(1.3,2)}
      \foreach \x in {0,3}{
        \draw[redund] (\x-0.6,0) -- (\x-0.1,0); \draw[redund] (\x,-0.7) -- (\x,-0.1);
        \draw[redund] (\x+0.35,-0.3) arc (-40:200:0.35);}
    \end{scope}
    \begin{scope}[xshift=10cm]
      \draw[member] (0,0) -- (0,2) -- (3,2) -- (3,0); \pic at (0,0) {fixed}; \pic at (3,0) {fixed}; \ld{(1.3,2)}
      \transh{(3,2)} \draw[dofm] (0,2) -- (0.6,2);
      \draw[dofm] (0.4,2.3) arc (0:200:0.35); \draw[dofm] (3.4,2.3) arc (0:200:0.35);
      \draw[dofm] (0,1.95) -- (0,1.4);
    \end{scope}
    \node at (14.3,0.6) {3}; \node at (15.8,0.6) {3};
  \end{scope}
  % ---------------- (e) multi-storey frame ----------------
  \begin{scope}[yshift=-15.4cm]
    \node at (-0.6,2.5) {$(e)$};
    % 3 bays x 7 storeys
    \newcommand{\frameE}{%
      \foreach \x in {0,0.8,...,2.4}{\draw[thinmember] (\x,0) -- (\x,4.9);}
      \foreach \y in {0.7,1.4,...,4.9}{\draw[thinmember] (0,\y) -- (2.4,\y);}
      \fill[hatch] (-0.3,-0.25) rectangle (2.7,0); \draw (-0.3,0) -- (2.7,0);}
    \frameE \ld{(0.6,4.9)} \ld{(2.0,4.9)}
    \begin{scope}[xshift=5cm]
      \frameE
      % cut members: show internal forces at midpoint of every member
      \foreach \x in {0,0.8,...,2.4}{\foreach \y in {0.35,1.05,...,4.55}{
        \draw[redund,thin] (\x-0.15,\y) -- (\x+0.15,\y); \draw[thin] (\x,\y-0.1) -- (\x,\y+0.1);}}
      \foreach \x in {0.4,1.2,2.0}{\foreach \y in {0.7,1.4,...,4.9}{
        \draw[thin] (\x-0.1,\y) -- (\x+0.1,\y); \draw[redund,thin] (\x,\y-0.15) -- (\x,\y+0.15);}}
      % detail bubble
      \draw[rounded corners=8pt] (-0.3,5.5) rectangle (2.7,6.4);
      \draw[thin] (0.7,5.5) -- (0.6,4.95);
      \draw[redund] (0.3,5.95) -- (0.9,5.95); \draw[redund] (1.2,5.6) -- (1.2,6.3);
      \draw[redund] (1.95,5.65) arc (-60:240:0.32);
      \node at (1.6,5.95) {$\odot$};
    \end{scope}
    \begin{scope}[xshift=10cm]
      \frameE
      % dofs at the top two storeys (3 per node): shown as arrows + rotation
      \foreach \x in {0,0.8,...,2.4}{
        \draw[dofm] (\x+0.3,4.95) arc (0:180:0.3); \draw[dofm] (\x,4.9) -- (\x,4.5);
        \draw[dofm] (\x,4.2) -- (\x+0.4,4.2); \draw[dofm] (\x+0.3,4.25) arc (0:180:0.3);}
      \draw[dofm] (2.4,4.9) -- (2.85,4.9);
    \end{scope}
    \node at (14.3,2.5) {63}; \node at (15.8,2.5) {35};
  \end{scope}
\end{tikzpicture}

  \caption{Static and kinematic indeterminacy (after Fig.~3.7 of the textbook).}
\end{figure}

\chapter[Basic Equations]{Basic Equations: Equilibrium, Compatibility, and Constitutive Equations}

\section{Equilibrium equations}
We shall distinguish the following three levels of equilibrium equations.

\topicnote{Material level:}{e.g.\ $\sigma_{ij,j}+f_i=0$}
\begin{fig}
  % 2D differential element with stresses and body forces
\begin{tikzpicture}[scale=1.0,font=\small]
  \def\a{3.0}
  \draw[thick] (0,0) rectangle (\a,\a);
  \node[anchor=north west,inner sep=2pt] at (0,\a*0.55) {$dy$};
  \node[anchor=south west,inner sep=2pt] at (\a*0.55,0) {$dx$};
  % body forces
  \draw[->] (\a/2,\a/2) -- ++(0.8,0) node[below] {$f_x$};
  \draw[->] (\a/2,\a/2) -- ++(0,0.8) node[right] {$f_y$};
  % left face
  \draw[->] (0,\a/2) -- ++(-1.1,0) node[left] {$\sigma_{xx}$};
  \draw[->] (-0.45,\a-0.2) -- (-0.45,0.2);
  \node[left] at (-0.5,0.55) {$\sigma_{xy}$};
  % right face
  \draw[->] (\a,\a/2) -- ++(1.1,0) node[right] {$\sigma_{xx}+\dfrac{\partial\sigma_{xx}}{\partial x}dx$};
  \draw[->] (\a+0.45,0.2) -- (\a+0.45,\a-0.2);
  \node[anchor=north west] at (\a+0.3,-0.1) {$\sigma_{xy}+\dfrac{\partial\sigma_{xy}}{\partial x}dx$};
  % bottom face
  \draw[->] (\a/2,0) -- ++(0,-1.2) node[right] {$\sigma_{yy}$};
  \draw[->] (\a-0.2,-0.45) -- (0.2,-0.45);
  \node[below] at (0.7,-0.5) {$\sigma_{xy}$};
  % top face
  \draw[->] (\a/2,\a) -- ++(0,1.3) node[above] {$\sigma_{yy}+\dfrac{\partial\sigma_{yy}}{\partial y}dy$};
  \draw[->] (0.2,\a+0.45) -- (\a-0.2,\a+0.45);
  \node[anchor=south east] at (\a*0.3,\a+0.5) {$\sigma_{xy}+\dfrac{\partial\sigma_{xy}}{\partial y}dy$};
\end{tikzpicture}

\end{fig}
\begin{itemize}
  \item $\sum F_x = 0$ (assuming thickness $=1$):
  \begin{multline*}
    -\sigma_{xx}\,(dy\cdot 1) + \Bigl(\sigma_{xx}+\pd{\sigma_{xx}}{x}dx\Bigr)(dy\cdot 1)
    - \sigma_{xy}\,(dx\cdot 1) \\
    + \Bigl(\sigma_{xy}+\pd{\sigma_{xy}}{y}dy\Bigr)(dx\cdot 1) + f_x\,(dx\cdot dy\cdot 1) = 0
  \end{multline*}
  \begin{equation}
    \Longrightarrow\quad \pd{\sigma_{xx}}{x}+\pd{\sigma_{xy}}{y}+f_x = 0
  \end{equation}
  \item $\sum F_y = 0$:
  \begin{multline*}
    -\sigma_{yy}\,(dx\cdot 1) + \Bigl(\sigma_{yy}+\pd{\sigma_{yy}}{y}dy\Bigr)(dx\cdot 1)
    - \sigma_{xy}\,(dy\cdot 1) \\
    + \Bigl(\sigma_{xy}+\pd{\sigma_{xy}}{x}dx\Bigr)(dy\cdot 1) + f_y\,(dx\cdot dy\cdot 1) = 0
  \end{multline*}
  \begin{equation}
    \Longrightarrow\quad \pd{\sigma_{yy}}{y}+\pd{\sigma_{xy}}{x}+f_y = 0
    \quad\xRightarrow{\;\sigma_{xy}=\sigma_{yx}\;}\quad
    \pd{\sigma_{yx}}{x}+\pd{\sigma_{yy}}{y}+f_y = 0
  \end{equation}
\end{itemize}

\topicnote{Nodal level:}{$\sum \mat{f}_i = \Pvec_i$}
This means: summation of internal forces $\mat{f}_i$ (including moments) $=$ externally applied loads at node $i$, $\Pvec_i$.

\noindent$*$ There are six nodes in this cantilever but only one joint (node 1).
\begin{fig}
  % A cantilever modelled with six nodes but only one joint (node 1)
\begin{tikzpicture}[scale=1.0,font=\small]
  \fill[left color=black!60,right color=black!15] (-0.35,-0.45) rectangle (0,0.45);
  \draw[thick] (0,-0.2) rectangle (6,0.2);
  \foreach \i in {1,...,6}{
    \fill ({(\i-1)*1.2},0) circle (0.07);
    \node[above] at ({(\i-1)*1.2},0.3) {\i};}
\end{tikzpicture}

\end{fig}

\topic{Structural level:}
\begin{equation*}
  \Kmat\mat{U} = \Pvec \qquad \text{(obtained by assembling the nodal equilibrium equations)}
\end{equation*}

\section{Compatibility equations}
\topic{Material level}
\begin{enumerate}[label=e.g.\,\arabic*.,leftmargin=*]
  \item $\mat{u}(\mat{x}') = \mat{u}(\mat{x})$ as $\mat{x}'\to\mat{x}$
  \item 2D elasticity problems:
    \begin{equation}
      2\frac{\partial^2\eps_{xy}}{\partial x\,\partial y} = \frac{\partial^2\eps_x}{\partial y^2}+\frac{\partial^2\eps_y}{\partial x^2},
      \qquad\text{where }
      \eps_x = \pd{u}{x},\quad \eps_y=\pd{v}{y},\quad \eps_{xy} = \frac12\Bigl(\pd{u}{y}+\pd{v}{x}\Bigr).
    \end{equation}
    Mathematically, this is called the \emph{integrable condition}. Once this equation is satisfied, we can determine the displacement field from a given strain field.
\end{enumerate}

We all know the integrable condition associated with a conservative force as shown below.
\begin{fig}
  % Work done by a force along a path from point 1 to point 2
\begin{tikzpicture}[scale=1.0,font=\small]
  \draw[->] (0,0) -- (5.0,0) node[right] {$x$};
  \draw[->] (0,0) -- (0,3.4) node[above] {$y$};
  \draw[thick] (0.3,0.8) .. controls (0.5,2.3) and (1.2,2.4) .. (2.0,1.7)
               .. controls (2.6,1.2) and (3.2,1.2) .. (3.9,2.2);
  \fill (0.3,0.8) circle (0.05) node[below right] {1};
  \fill (3.9,2.2) circle (0.05) node[above right] {2};
  \coordinate (Q) at (1.0,2.0);
  \fill (Q) circle (0.04);
  \draw[->] (Q) -- ++(1.0,0) node[right] {$P(x,y)$};
  \draw[->] (Q) -- ++(0,1.0) node[above] {$Q(x,y)$};
\end{tikzpicture}

\end{fig}
\begin{gather}
  \mat{F} = P(x,y)\,\mat{i} + Q(x,y)\,\mat{j}, \qquad d\mat{r} = dx\,\mat{i} + dy\,\mat{j} \\
  \int_1^2 \mat{F}\cdot d\mat{r} = \int_1^2 P(x,y)\,dx + Q(x,y)\,dy = \int_1^2 dW = -\int_1^2 dV = V_1 - V_2
\end{gather}
Note that the potential energy $V$ exists when $\mat{F}$ is conservative. The work done by a force needs to be evaluated by integration along the path of the particle. However, for a conservative force, no integration is needed. The work done can be evaluated using the change in the potential energy.
\begin{gather}
  \Longrightarrow\quad -dV = -\pd{V}{x}dx - \pd{V}{y}dy = P(x,y)\,dx + Q(x,y)\,dy \\
  \Longrightarrow\quad P(x,y) = -\pd{V}{x},\quad Q(x,y) = -\pd{V}{y}
  \qquad\Longrightarrow\quad \pd{P}{y} = \pd{Q}{x}
\end{gather}

Consider two well-known conservative forces and their potentials.
\begin{enumerate}[label=(\alph*)]
  \item \textbf{Gravity force}
    \begin{fig}
      % Gravity force acting on a particle
\begin{tikzpicture}[scale=1.0,font=\small]
  \draw[->] (0,0) -- (4.5,0) node[right] {$x$};
  \draw[->] (0,0) -- (0,3.2) node[above] {$y$};
  \shade[ball color=black!40] (1.2,1.8) circle (0.22);
  \draw[->] (1.42,1.8) -- ++(0.9,0) node[right] {$P(x,y)$};
  \draw[->] (1.2,2.02) -- ++(0,0.9) node[above] {$Q(x,y)$};
\end{tikzpicture}

    \end{fig}
    \begin{gather*}
      P(x,y) = 0, \qquad Q(x,y) = -mg \\
      \text{Check: } \pd{P}{y} = \pd{Q}{x} \;\Rightarrow\; 0 = 0 \\
      dW = (-mg)\,dy \;\Rightarrow\; dV = -dW = mg\,(dy) \;\Rightarrow\; V = mgy + c
    \end{gather*}
  \item \textbf{Spring (along the $x$-direction) force}
    \begin{gather*}
      P(x,y) = -kx, \qquad Q(x,y) = 0 \\
      \text{Check: } \pd{P}{y} = \pd{Q}{x} \;\Rightarrow\; 0 = 0 \\
      dW = (-kx)\,dx \;\Rightarrow\; dV = -dW = kx\,(dx) \;\Rightarrow\; V = \tfrac12 kx^2 + c
    \end{gather*}
\end{enumerate}

\topic{Nodal level}
E.g.\ for a plane truss, $(u_1,u_2)_1 = (u_1,u_2)_2 = (u_1,u_2)_3$, where a subscript represents the member number.
\begin{fig}
  % Compatibility at a truss node: three members share the same nodal displacements
\begin{tikzpicture}[scale=1.0,font=\small]
  \coordinate (N) at (0,0);
  \draw[member] (N) -- (-1.3,-2.0); \draw[member] (N) -- (0.4,-2.0); \draw[member] (N) -- (1.9,-1.7);
  \fill (N) circle (0.06);
  \draw[->,thick] (N) -- (1.4,0) node[right] {$u_1$};
  \draw[->,thick] (N) -- (0,1.2) node[right] {$u_2$};
  \node[circle,draw,inner sep=1pt] at (-1.1,-0.85) {\scriptsize 1};
  \node[circle,draw,inner sep=1pt] at (-0.2,-1.3) {\scriptsize 2};
  \node[circle,draw,inner sep=1pt] at (1.3,-0.8) {\scriptsize 3};
\end{tikzpicture}

\end{fig}

Skew support: no penetration, i.e.\ the displacement perpendicular to the skew support equals zero.
\begin{fig}
  % Skew (inclined) roller support: no penetration, v = 0
\begin{tikzpicture}[scale=1.0,font=\small,rotate=35]
  % ground surface
  \draw[thick] (-1.6,0) -- (1.8,0);
  \fill[hatch] (-1.6,-0.3) rectangle (1.8,0);
  % roller
  \draw[fill=white] (0,0.3) circle (0.3);
  % attached wedge (member end)
  \draw[thick] (-0.9,0.75) -- (0.6,0.6) -- (-0.5,1.5) -- cycle;
  % tangential displacement u
  \draw[->,thick] (0.4,0.85) -- (1.6,0.85) node[above] {$u$};
  % normal displacement v = 0
  \draw[->,thick] (0.6,-1.3) -- (0.6,-0.35) node[pos=0,below] {$v=0$};
\end{tikzpicture}

\end{fig}

\section{Constitutive equations}
\topic{Material level}
Stress--strain relationships; for linearly elastic material $\bs{\sigma} = \mat{C}\bs{\eps}$ or $\bs{\eps} = \mat{S}\bs{\sigma}$, where $\mat{C} =$ elasticity matrix and $\mat{S} =$ compliance matrix.

\topic{3D problems:}
\begin{equation}
  \mat{C} = \frac{E}{1+\nu}
  \begin{bmatrix}
    \dfrac{1-\nu}{1-2\nu} & \dfrac{\nu}{1-2\nu} & \dfrac{\nu}{1-2\nu} & 0 & 0 & 0\\[8pt]
    \dfrac{\nu}{1-2\nu} & \dfrac{1-\nu}{1-2\nu} & \dfrac{\nu}{1-2\nu} & 0 & 0 & 0\\[8pt]
    \dfrac{\nu}{1-2\nu} & \dfrac{\nu}{1-2\nu} & \dfrac{1-\nu}{1-2\nu} & 0 & 0 & 0\\[8pt]
    0 & 0 & 0 & 1 & 0 & 0\\
    0 & 0 & 0 & 0 & 1 & 0\\
    0 & 0 & 0 & 0 & 0 & 1
  \end{bmatrix}
\end{equation}
\begin{equation}
  \mat{S} = \frac{1}{E}
  \begin{bmatrix}
    1 & -\nu & -\nu & 0 & 0 & 0\\
    -\nu & 1 & -\nu & 0 & 0 & 0\\
    -\nu & -\nu & 1 & 0 & 0 & 0\\
    0 & 0 & 0 & 1+\nu & 0 & 0\\
    0 & 0 & 0 & 0 & 1+\nu & 0\\
    0 & 0 & 0 & 0 & 0 & 1+\nu
  \end{bmatrix}
\end{equation}

\topic{2D problems:}
\begin{description}[style=nextline,leftmargin=1.5em]
  \item[\underline{Case 1: plane stress}]
    Assumption: $\sigma_{zz} = \sigma_{zx} = \sigma_{zy} = 0$ and in-plane ($xy$-plane) loading.
    \begin{fig}
      % Plane stress: thin plate loaded in its own plane
\begin{tikzpicture}[scale=1.0,font=\small]
  % thin disc
  \fill[black!8] (0,0) ellipse (1.6 and 0.45);
  \draw (0,0) ellipse (1.6 and 0.45);
  \draw (-1.6,0) -- (-1.6,-0.22) arc (180:360:1.6 and 0.45) -- (1.6,0);
  \fill[black!15] (-1.6,-0.22) arc (180:360:1.6 and 0.45) -- (1.6,0) arc (360:180:1.6 and 0.45) -- cycle;
  % axes
  \draw[->,thick] (0,0) -- (0,2.4) node[above] {$z$};
  \draw[->,thick] (0,0) -- (-2.3,-0.9) node[left] {$x$};
  \draw[->,thick] (0,0) -- (2.3,-0.9) node[right] {$y$};
\end{tikzpicture}

    \end{fig}
    \begin{equation}
      \mat{C} = \frac{E}{1-\nu^2}\bmat{1 & \nu & 0\\ \nu & 1 & 0\\ 0 & 0 & 1-\nu}, \qquad
      \mat{S} = \frac{1}{E}\bmat{1 & -\nu & 0\\ -\nu & 1 & 0\\ 0 & 0 & 1+\nu}
    \end{equation}
  \item[\underline{Case 2: plane strain}]
    Assumption: $\eps_{zz} = \eps_{zx} = \eps_{zy} = 0$ and in-plane ($xy$-plane) loading.
    \begin{fig}
      % Plane strain: long prismatic body loaded uniformly along its length
\begin{tikzpicture}[scale=1.0,font=\small]
  % cylinder along a diagonal direction
  \begin{scope}[rotate=25]
    \fill[black!6] (0,-0.55) -- (3.8,-0.55) arc (270:450:0.22 and 0.55) -- (0,0.55) arc (90:270:0.22 and 0.55) -- cycle;
    \draw (0,-0.55) -- (3.8,-0.55); \draw (0,0.55) -- (3.8,0.55);
    \draw (3.8,-0.55) arc (270:450:0.22 and 0.55);
    \draw[densely dashed] (3.8,0.55) arc (90:270:0.22 and 0.55);
    \shade[left color=black!55,right color=black!15] (0,0) ellipse (0.22 and 0.55);
    \draw (0,0) ellipse (0.22 and 0.55);
  \end{scope}
  % axes
  \draw[->,thick] (0,0) -- (0,2.2) node[above] {$y$};
  \draw[->,thick] (0,0) -- (1.9,-0.9) node[right] {$x$};
  \draw[->,thick] (0,0) -- (-1.9,-0.9) node[left] {$z$};
\end{tikzpicture}

    \end{fig}
    \begin{equation}
      \mat{C} = \frac{E}{1+\nu}
      \begin{bmatrix}
        \dfrac{1-\nu}{1-2\nu} & \dfrac{\nu}{1-2\nu} & 0\\[8pt]
        \dfrac{\nu}{1-2\nu} & \dfrac{1-\nu}{1-2\nu} & 0\\[8pt]
        0 & 0 & 1
      \end{bmatrix}, \qquad
      \mat{S} = \frac{1}{E}
      \bmat{1-\nu^2 & -\nu(1+\nu) & 0\\ -\nu(1+\nu) & 1-\nu^2 & 0\\ 0 & 0 & 1+\nu}
    \end{equation}
\end{description}
\noindent$*$ Note that the $\mat{C}$ and $\mat{S}$ matrices for 2D problems can be derived from those for 3D problems using Gauss elimination; we shall do it in class.

\topic{Cross-section level}
Usually called \emph{generalized stress--strain relationships}, e.g.\ axial force--axial strain relationship ($F = EA\eps$), moment--curvature relationship ($M = EI\phi$), and torsion--rate of twist relationship ($T = GJ\theta'$).

\chapter{Axial Element}

\section{Stiffness matrix for an axial (bar, or truss) element}
\begin{fig}
  % Axial (bar/truss) element with end displacements and forces
\begin{tikzpicture}[scale=1.0,font=\small]
  \def\L{5}
  \draw[thick] (0,-0.15) rectangle (\L,0.15);
  \node[above] at (\L/2,0.15) {$EA$};
  \draw[->,thick] (-1.5,0) -- (-0.1,0);
  \node[above] at (-0.8,0.1) {$\delta_1\,(f_1)$};
  \draw[->,thick] (\L+0.1,0) -- (\L+1.5,0);
  \node[above] at (\L+0.8,0.1) {$\delta_2\,(f_2)$};
  \draw[dim] (0,-0.55) -- (\L,-0.55) node[midway,fill=white,inner sep=1pt] {$L$};
\end{tikzpicture}

\end{fig}
\begin{equation}
  \mat{f} = \kmat\bs{\delta} \quad\text{where}\quad
  \mat{f} = \Bmat{f_1\\ f_2},\quad \bs{\delta} = \Bmat{\delta_1\\ \delta_2}
  \quad\text{and}\quad \kmat = \frac{EA}{L}\bmat{1 & -1\\ -1 & 1}.
\end{equation}
\begin{equation*}
  \delta = \delta_2-\delta_1,\qquad f_2 = k\delta = \frac{EA}{L}(\delta_2-\delta_1)
  \quad\text{and}\quad f_1 = -f_2 = \frac{EA}{L}(\delta_1-\delta_2).
\end{equation*}

\section{Methods for deriving the element stiffness matrix \texorpdfstring{$\kmat$}{k}}
\begin{enumerate}[leftmargin=*]
  \item \textbf{Direct method}\\
    By definition, the $i$-th column of $\kmat$ is the $\mat{f}$ vector corresponding to $\delta_i = 1$, and other $\delta_j = 0$ ($j\neq i$).

    For example, take $\delta_1 = 1$, $\delta_2 = 0$; then the first column of the stiffness matrix can be obtained as shown.
    \begin{fig}
      % Direct method: impose delta_1 = 1, delta_2 = 0
\begin{tikzpicture}[scale=1.0,font=\small]
  \def\L{5}
  % undeformed bar
  \draw[thick] (0,0.35) rectangle (\L,0.65);
  % wall at the right end
  \fill[left color=black!15,right color=black!60] (\L,-0.5) rectangle (\L+0.3,1.0);
  % deformed bar (dotted), shifted 1 to the right at the left end
  \draw[densely dotted,thick] (0.6,-0.15) rectangle (\L,0.15);
  \draw[dim] (0,-0.15) -- (0.6,-0.15); \node[below] at (0.3,-0.2) {1};
  \draw (0,-0.3) -- (0,0.3);
  % end forces
  \draw[->,thick] (-2.0,0) -- (-0.1,0);
  \node[left] at (-2.1,0) {$k_{11}=\dfrac{EA}{L}$};
  \draw[->,thick] (\L+0.4,0) -- (\L+1.9,0);
  \node[right] at (\L+2.0,0) {$k_{21}=-\dfrac{EA}{L}$};
\end{tikzpicture}

    \end{fig}
  \item \textbf{Energy method}\\
    The strain energy $U$ in terms of $\bs{\delta}$ for a linearly elastic member is a quadratic form. Hence
    \begin{equation}
      U = \tfrac12\,\bs{\delta}^\tr\kmat\bs{\delta}.
    \end{equation}
    This is equivalent to
    \begin{equation}
      k_{ij} = \frac{\partial^2 U}{\partial\delta_i\,\partial\delta_j}.
    \end{equation}
    For a truss element, $U$ can be easily shown to be
    \begin{equation}
      U = \frac12\frac{EA}{L}(\Delta\delta)^2 = \frac12\frac{EA}{L}(\delta_2-\delta_1)^2.
    \end{equation}
  \item \textbf{Principle of virtual displacement}, or finite element approach employing shape functions (we shall discuss this method in Ch.~\ref{ch:pvd2}).
\end{enumerate}

\section{Global stiffness matrix for an axial element}
\emph{Stiffness matrix ready for assembling.}
\begin{equation}
  \Bmat{F_{x1}\\ F_{y1}\\ F_{x2}\\ F_{y2}} = \frac{EA}{L}
  \bmat{c^2 & sc & -c^2 & -sc\\ sc & s^2 & -sc & -s^2\\ -c^2 & -sc & c^2 & sc\\ -sc & -s^2 & sc & s^2}
  \Bmat{u_1\\ v_1\\ u_2\\ v_2}
  \label{eq:truss-global-k}
\end{equation}
\begin{minipage}{0.68\textwidth}
where $c = \cos\phi$, $s = \sin\phi$, and $\phi$ is the angle between the local ($x$) and the global ($X$) axes, which is positive if the $x$ axis is rotated counterclockwise with respect to the $X$ axis. Methods for deriving the local stiffness matrix can also be used to obtain the above global stiffness matrix.
\end{minipage}\hfill
\begin{minipage}{0.28\textwidth}
  \centering
  % Local (x,y) axes rotated by phi from global (X,Y) axes
\begin{tikzpicture}[scale=1.0,font=\small]
  \draw[->,thick] (0,0) -- (3.2,0) node[right] {$X$};
  \draw[->,thick] (0,0) -- (0,3.2) node[above] {$Y$};
  \coordinate (O) at (1.2,0.7);
  \draw[dashed] (O) -- ++(1.6,0);
  \draw[->,thick] (O) -- ++(35:2.0) node[above right] {$x$};
  \draw[->,thick] (O) -- ++(125:1.4) node[above left] {$y$};
  \draw[->] ($(O)+(1.1,0)$) arc (0:35:1.1);
  \node at ($(O)+(1.35,0.35)$) {$\phi$};
\end{tikzpicture}

\end{minipage}

Note that the strain energy is
\begin{equation}
  U = \frac12\frac{EA}{L}(\delta_2-\delta_1)^2
    = \frac12\frac{EA}{L}\langle\delta_1\;\;\delta_2\rangle\bmat{1 & -1\\ -1 & 1}\Bmat{\delta_1\\ \delta_2}.
\end{equation}
The member displacements in the local axes, $\delta_1$ and $\delta_2$, and the member displacements in the global axes, $u_1$, $v_1$, $u_2$ and $v_2$, are shown in the figure below.
\begin{fig}
  % Member displacements in local (delta) and global (u,v) axes
\begin{tikzpicture}[scale=1.0,font=\small]
  \coordinate (A) at (0,0); \coordinate (B) at (35:4.2);
  \draw[member] (A) -- (B);
  \draw[dashed] (A) -- ++(215:0.9); \draw[dashed] (B) -- ++(35:0.9);
  % node 1
  \draw[->,thick] (A) -- ++(1.4,0) node[below right] {$u_1$};
  \draw[->,thick] (A) -- ++(0,1.4) node[above] {$v_1$};
  \draw[->,thick] ($(A)+(215:0.9)$) -- ($(A)+(215:0.05)$);
  \node[left] at ($(A)+(215:0.7)$) {$\delta_1$};
  \draw[->] ($(A)+(0.9,0)$) arc (0:35:0.9); \node at ($(A)+(1.15,0.32)$) {$\phi$};
  % node 2
  \draw[->,thick] (B) -- ++(1.4,0) node[right] {$u_2$};
  \draw[->,thick] (B) -- ++(0,1.4) node[above] {$v_2$};
  \draw[->,thick] (B) -- ++(35:1.0) node[above right] {$\delta_2$};
\end{tikzpicture}

\end{fig}

Relationship between the displacements in the local axes and the displacements in the global axes:
\begin{equation}
  \bs{\delta} = \Bmat{\delta_1\\ \delta_2} = \bmat{c & s & 0 & 0\\ 0 & 0 & c & s}\Bmat{u_1\\ v_1\\ u_2\\ v_2} = \Tmat\bs{\Delta}.
\end{equation}
Writing the strain energy $U$ in terms of $\bs{\delta}$ and $\bs{\Delta}$:
\begin{equation}
  U = \frac12\bs{\delta}^\tr\kmat_l\bs{\delta} = \frac12\bs{\Delta}^\tr\Tmat^\tr\kmat_l\Tmat\bs{\Delta}
    = \frac12\bs{\Delta}^\tr\kmat_g\bs{\Delta},
  \qquad\text{where } \kmat_g = \Tmat^\tr\kmat_l\Tmat.
\end{equation}
Therefore,
\begin{equation}
  U = \frac12\Bmat{u_1\\ v_1\\ u_2\\ v_2}^{\!\tr}
  \bmat{c & 0\\ s & 0\\ 0 & c\\ 0 & s}\frac{EA}{L}\bmat{1 & -1\\ -1 & 1}\bmat{c & s & 0 & 0\\ 0 & 0 & c & s}
  \Bmat{u_1\\ v_1\\ u_2\\ v_2}
  = \frac12\Bmat{u_1\\ v_1\\ u_2\\ v_2}^{\!\tr}[k_g]\Bmat{u_1\\ v_1\\ u_2\\ v_2},
\end{equation}
where $[k_g]$ is the global stiffness matrix given above.

\section{General features of an element stiffness matrix}
\begin{enumerate}[leftmargin=*]
  \item \textbf{Symmetric}
  \item \textbf{Singular}\\
    Number of zero eigenvalues $=$ number of rigid-body modes.\\
    E.g.\ for the local stiffness matrix of an axial element, the number is 1 (one translation), and for the global one, the number is 3 (two translations plus one rotation). For example, let's take a look at the rigid-body rotation (w.r.t.\ node 1) mode, $\bs{\Delta}_r$, where
    \begin{fig}
      \begin{minipage}{0.35\textwidth}
        \begin{equation*}
          \bs{\Delta}_r = \Bmat{0\\ 0\\ -L\theta_r\sin\phi\\ L\theta_r\cos\phi}
        \end{equation*}
      \end{minipage}
      \begin{minipage}{0.55\textwidth}
        \centering% Rigid-body rotation mode about node 1
\begin{tikzpicture}[scale=1.0,font=\small]
  \coordinate (A) at (0,0); \coordinate (B) at (30:4.2); \coordinate (Br) at (55:4.2);
  \draw[member] (A) -- (B); \draw[member] (A) -- (Br);
  \draw[dashed] (B) -- (Br);
  \node[left] at ($(B)!0.5!(Br)+(-0.05,0)$) {$L\theta_r$};
  \draw[->,thick] (A) -- ++(1.4,0) node[below right] {$u_1$};
  \draw[->,thick] (A) -- ++(0,1.4) node[left] {$v_1$};
  \draw[->,thick] (B) -- ++(1.2,0) node[right] {$u_2$};
  \draw[->,thick] (B) -- ++(0,1.2) node[above] {$v_2$};
  \draw[->] ($(A)+(1.0,0)$) arc (0:30:1.0); \node at ($(A)+(1.25,0.28)$) {$\phi$};
  \draw[->] ($(A)+(30:1.7)$) arc (30:55:1.7); \node at ($(A)+(42:2.0)$) {$\theta_r$};
\end{tikzpicture}

      \end{minipage}
    \end{fig}
  \item Each column of the element stiffness matrix constitutes a set of member end forces that are in \textbf{equilibrium}. Therefore, $\F = \kmat\bs{\Delta}$ is also a set of forces in equilibrium.

    For example, $v_2 = 1$ and others equal 0. The axial force is equal to $(EA/L)\sin\phi$. With the geometric relationship, we get
    \begin{equation}
      \Bmat{F_{x1}\\ F_{y1}\\ F_{x2}\\ F_{y2}} = \frac{EA}{L}\Bmat{-sc\\ -s^2\\ sc\\ s^2}.
    \end{equation}
    This vector is the 4th column of the element stiffness matrix.
    \begin{fig}
      % Fourth column of the global stiffness matrix: v2 = 1, others 0
\begin{tikzpicture}[scale=1.0,font=\small]
  % --- left: geometry ---
  \begin{scope}
    \coordinate (A) at (0,0); \coordinate (B) at (25:4.0);
    \draw[member] (A) -- (B);
    \draw[dashed,thick] (A) -- ($(B)+(0,1.0)$);
    \draw[->,thick] (B) -- ($(B)+(0,1.0)$);
    \node[left] at ($(B)+(-0.05,0.55)$) {1};
    % component of the unit displacement along the member
    \draw[densely dotted] ($(B)+(0,1.0)$) -- ($(B)+(25:{sin(25)})$);
    \draw[dim] ($(B)+(-0.1,0)$) -- ($(B)+(25:{sin(25)})+(-0.1,0)$);
    \node[below right,inner sep=1pt] at ($(B)+(0.05,-0.05)$) {\scriptsize $1\cdot\sin\phi$};
    \draw[->] ($(A)+(1.4,0)$) arc (0:25:1.4); \node at ($(A)+(1.7,0.3)$) {$\phi$};
    \draw (A) -- ++(1.8,0);
  \end{scope}
  % --- right: end forces ---
  \begin{scope}[xshift=7.5cm]
    \coordinate (A) at (0,0); \coordinate (B) at (25:4.0);
    \draw[member] (A) -- (B);
    \draw[->,thick] (A) -- ++(0,1.3) node[right] {$-\dfrac{EA}{L}\cdot s^2$};
    \draw[->,thick] (A) -- ++(1.2,0) node[below] {$-\dfrac{EA}{L}\cdot sc$};
    \draw[->,thick] (B) -- ++(0,1.3) node[right] {$\dfrac{EA}{L}\cdot s^2$};
    \draw[->,thick] (B) -- ++(1.3,0) node[below right] {$\dfrac{EA}{L}\cdot sc$};
    \draw[->] ($(A)+(1.0,0)$) arc (0:25:1.0); \node at ($(A)+(1.3,0.35)$) {$\phi$};
  \end{scope}
\end{tikzpicture}

    \end{fig}
\end{enumerate}

\chapter{Direct Stiffness Method}
\label{ch:dsm}

\section{Direct stiffness method}
The equilibrium equations for a structural system can be written as
\begin{equation}
  \Kmat\bs{\Delta} = \Pvec
\end{equation}
where $\Kmat =$ structural stiffness matrix, $\bs{\Delta} =$ displacement vector, and $\Pvec =$ load vector. The stiffness matrix $\Kmat$ can be obtained by assembling the element stiffness matrices using the member degree-of-freedom table (the \texttt{LM} matrix in our program).

In carrying out the assembling procedure, compatibility conditions of nodes need to be taken into account. For example, for a rigid joint, the members sharing this joint all have the same degrees of freedom; however, for an internal hinge, the members sharing this joint have the same translational degrees of freedom but have different rotational degrees of freedom.

The meaning of $\Kmat\bs{\Delta} = \Pvec$ is nothing but \emph{summation of internal forces $=$ external forces at nodes}, because
\begin{equation*}
  \Kmat\bs{\Delta} = \sum_{i=1}^{n}\kmat_i\bs{\Delta}_i = \sum_{i=1}^{n}\F_i = \text{summation of internal forces},
\end{equation*}
where $\kmat_i =$ element stiffness matrix, and $\bs{\Delta}_i =$ element displacement vector of the $i$-th element, both being referred to the global axes.
\begin{fig}
  % Nodal equilibrium: sum of member end forces = applied nodal loads
\begin{tikzpicture}[scale=0.9,font=\small]
  \coordinate (N) at (0,0);
  % node with member stubs and applied loads
  \draw[member] (N) -- (-0.9,-1.0); \draw[member] (N) -- (0.45,-1.1);
  \fill (N) circle (0.06);
  \draw[load] (N) -- ++(1.3,0) node[right] {$P_x$};
  \draw[load] (N) -- ++(0,1.3) node[above] {$P_y$};
  % forces exerted by the members on the node
  \draw[load] (-2.2,-0.75) -- (-1.1,-0.75) node[pos=0,left] {$F^1_{x2}$};
  \draw[load] (-0.75,-2.1) -- (-0.75,-1.1) node[pos=0,below] {$F^1_{y2}$};
  \draw[load] (2.0,-0.75) -- (0.9,-0.75) node[pos=0,right] {$F^2_{x2}$};
  \draw[load] (0.75,-2.1) -- (0.75,-1.1) node[pos=0,below] {$F^2_{y2}$};
  % ---- member 1 free body ----
  \coordinate (T1) at (-2.3,-3.1); \coordinate (A) at (-5.0,-6.1);
  \draw[member] (T1) -- (A);
  \draw[load] (T1) -- ++(-1.1,0) node[left] {$F^1_{x2}$};
  \draw[load] (T1) -- ++(0,-1.0) node[below] {$F^1_{y2}$};
  \draw[load] (A) -- ++(1.1,0) node[right] {$F^1_{x1}$};
  \draw[load] (A) -- ++(0,1.0) node[left] {$F^1_{y1}$};
  % ---- member 2 free body ----
  \coordinate (T2) at (1.8,-3.3); \coordinate (B) at (2.9,-6.6);
  \draw[member] (T2) -- (B);
  \draw[load] (T2) -- ++(1.1,0) node[right] {$F^2_{x2}$};
  \draw[load] (T2) -- ++(0,-1.0) node[below] {$F^2_{y2}$};
  \draw[load] (B) -- ++(1.1,0) node[right] {$F^2_{x1}$};
  \draw[load] (B) -- ++(0,1.0) node[midway,right] {$F^2_{y1}$};
  % ---- equations ----
  \node[anchor=west,align=left] at (4.8,-2.8) {$F^1_{x2}+F^2_{x2}=P_x$\\[6pt]$F^1_{y2}+F^2_{y2}=P_y$};
\end{tikzpicture}

\end{fig}

\begin{figure}[tb]
  \centering
  \begin{subfigure}[b]{0.52\textwidth}\centering% Representative interior joint of a plane truss (after McGuire Fig. 3.1)
\begin{tikzpicture}[scale=0.72,font=\small]
  \coordinate (q) at (0,0);
  \coordinate (n1) at (-2.4,-2.9);  % lower-left node (D)
  \coordinate (n2) at (-2.0,0.6);   % upper-left node (C,G,F)
  \coordinate (n3) at (2.2,2.0);    % upper-right node (B,F,E)
  \coordinate (n4) at (2.2,-2.0);   % lower-right node (A,E)
  \draw[thinmember] (q) -- (n1) node[pos=0.45,above left] {(D)};
  \draw[thinmember] (q) -- (n2) node[pos=0.72,above] {(C)};
  \draw[thinmember] (q) -- (n3) node[pos=0.6,below right] {(B)};
  \draw[thinmember] (q) -- (n4) node[midway,below left] {(A)};
  \draw[thinmember] (n1) -- (n2) node[midway,left] {(G)};
  \draw[thinmember] (n2) -- (n3) node[midway,above] {(F)};
  \draw[thinmember] (n3) -- (n4) node[midway,right] {(E)};
  \fill (q) circle (0.05) node[below right] {$q$};
  % loads / dofs at q
  \draw[load] (-1.3,0) -- (-0.1,0) node[pos=0,above] {$P_i$};
  \draw[load] (q) -- (1.2,0) node[right] {$\Delta_i$};
  \draw[load] (q) -- (0,1.2) node[right] {$\Delta_2$};
  % dofs at the other nodes
  \draw[load] ($(n1)+(-1.2,0)$) -- (n1) node[pos=0,above] {$\Delta_1$};
  \draw[load] ($(n1)+(0,-1.0)$) -- (n1) node[pos=0,right] {$\Delta_9$};
  \draw[load] ($(n2)+(-1.2,0)$) -- (n2) node[pos=0,above] {$\Delta_8$};
  \draw[load] (n2) -- ++(0,1.0) node[right] {$\Delta_7$};
  \draw[load] (n3) -- ++(1.1,0) node[above] {$\Delta_5$};
  \draw[load] (n3) -- ++(0,1.0) node[right] {$\Delta_6$};
  \draw[load] (n4) -- ++(1.1,0) node[above] {$\Delta_4$};
  \draw[load] ($(n4)+(0,-1.0)$) -- (n4) node[pos=0,right] {$\Delta_3$};
  % axes
  \draw[->] (-3.4,-3.4) -- (-2.6,-3.4) node[below] {$x$};
  \draw[->] (-3.4,-3.4) -- (-3.4,-2.6) node[left] {$y$};
\end{tikzpicture}

    \caption{Representative interior joint of a plane truss.}\end{subfigure}\hfill
  \begin{subfigure}[b]{0.44\textwidth}\centering% Study of equilibrium in the P_i - Delta_i direction (after McGuire Fig. 3.2)
\begin{tikzpicture}[scale=0.72,font=\small]
  % member ends with their x-force components
  \draw[thinmember] (-2.6,1.8) -- (-1.4,1.1) node[pos=0,above left] {(C)};
  \fill (-1.4,1.1) circle (0.04); \draw[load] (-1.4,1.1) -- ++(1.0,0) node[above] {$F_i^C$};
  \draw[thinmember] (2.6,1.9) -- (1.6,1.1) node[pos=0,above right] {(B)};
  \fill (1.6,1.1) circle (0.04); \draw[load] (1.6,1.1) -- ++(1.0,0) node[above] {$F_i^B$};
  \draw[thinmember] (-2.6,-1.9) -- (-1.4,-1.1) node[pos=0,below left] {(D)};
  \fill (-1.4,-1.1) circle (0.04); \draw[load] (-1.4,-1.1) -- ++(1.0,0) node[below] {$F_i^D$};
  \draw[thinmember] (2.6,-1.9) -- (1.6,-1.1) node[pos=0,below right] {(A)};
  \fill (1.6,-1.1) circle (0.04); \draw[load] (1.6,-1.1) -- ++(1.0,0) node[above] {$F_i^A$};
  % the joint
  \fill (0,0) circle (0.06) node[below] {$q$};
  \draw[load] (-1.8,0) -- (-0.15,0) node[pos=0,above] {$P_i$};
  \foreach \y in {0.3,0.1,-0.1,-0.3}{\draw[load] (1.2,\y) -- (0.2,\y);}
  \draw[decorate,decoration={brace,amplitude=4pt}] (1.35,0.45) -- (1.35,-0.45);
  \node[right] at (1.5,0) {$F_i^A,\,F_i^B,\,F_i^C,\,F_i^D$};
  \node[align=center,font=\scriptsize] at (0,-2.5) {Note: $y$ components of bar forces\\are omitted for clarity.};
\end{tikzpicture}

    \caption{Study of equilibrium in the $P_i$--$\Delta_i$ directions.}\end{subfigure}\\[1em]
  \begin{subfigure}[b]{0.52\textwidth}\centering% Effect of adding a member (H) to a joint (after McGuire Fig. 3.3)
\begin{tikzpicture}[scale=0.72,font=\small]
  \coordinate (q) at (0,0);
  \coordinate (n1) at (-2.4,-2.9);
  \coordinate (n2) at (-2.0,0.6);
  \coordinate (n3) at (2.2,2.0);
  \coordinate (n4) at (2.2,-2.0);
  \coordinate (n5) at (-0.6,-2.6);   % new node (H,L,M)
  \draw[thinmember] (q) -- (n1) node[pos=0.45,above left] {(D)};
  \draw[thinmember] (q) -- (n2) node[pos=0.72,above] {(C)};
  \draw[thinmember] (q) -- (n3) node[pos=0.6,below right] {(B)};
  \draw[thinmember] (q) -- (n4) node[midway,above right] {(A)};
  \draw[thinmember] (n1) -- (n2) node[midway,left] {(G)};
  \draw[thinmember] (n2) -- (n3) node[midway,above] {(F)};
  \draw[thinmember] (n3) -- (n4) node[midway,right] {(E)};
  \draw[thinmember] (q) -- (n5) node[midway,right] {(H)};
  \draw[thinmember] (n5) -- (n4) node[midway,above] {(L)};
  \draw[thinmember] (n1) -- (n5) node[midway,above] {(M)};
  \fill (q) circle (0.05) node[below right] {$q$};
  \draw[load] (-1.3,0) -- (-0.1,0) node[pos=0,above] {$P_i$};
  \draw[load] (q) -- (1.2,0) node[right] {$\Delta_i$};
  \draw[load] (q) -- (0,1.2) node[right] {$\Delta_2$};
  \draw[load] ($(n1)+(-1.2,0)$) -- (n1) node[pos=0,above] {$\Delta_1$};
  \draw[load] ($(n1)+(0,-1.0)$) -- (n1) node[pos=0,right] {$\Delta_9$};
  \draw[load] (n4) -- ++(1.1,0) node[above] {$\Delta_4$};
  \draw[load] ($(n4)+(0,-1.0)$) -- (n4) node[pos=0,right] {$\Delta_3$};
  \draw[load] (n5) -- ++(1.1,0) node[below] {$\Delta_{10}$};
  \draw[load] ($(n5)+(0,-1.0)$) -- (n5) node[pos=0,right] {$\Delta_{11}$};
\end{tikzpicture}

    \caption{Effect of adding a member to a joint.}\end{subfigure}\hfill
  \begin{subfigure}[b]{0.44\textwidth}\centering% Equilibrium in the P_i - Delta_i direction, member added (after McGuire Fig. 3.4)
\begin{tikzpicture}[scale=0.72,font=\small]
  \draw[thinmember] (-2.6,1.8) -- (-1.4,1.1) node[pos=0,above left] {(C)};
  \fill (-1.4,1.1) circle (0.04); \draw[load] (-1.4,1.1) -- ++(1.0,0) node[above] {$F_i^C$};
  \draw[thinmember] (2.6,1.9) -- (1.6,1.1) node[pos=0,above right] {(B)};
  \fill (1.6,1.1) circle (0.04); \draw[load] (1.6,1.1) -- ++(1.0,0) node[above] {$F_i^B$};
  \draw[thinmember] (-2.6,-1.9) -- (-1.4,-1.1) node[pos=0,below left] {(D)};
  \fill (-1.4,-1.1) circle (0.04); \draw[load] (-1.4,-1.1) -- ++(1.0,0) node[below] {$F_i^D$};
  \draw[thinmember] (2.6,-1.9) -- (1.6,-1.1) node[pos=0,below right] {(A)};
  \fill (1.6,-1.1) circle (0.04); \draw[load] (1.6,-1.1) -- ++(1.0,0) node[above] {$F_i^A$};
  \draw[thinmember] (-0.2,-2.2) -- (0,-1.4) node[pos=0,below left] {(H)};
  \fill (0,-1.4) circle (0.04); \draw[load] (0,-1.4) -- ++(0.9,0) node[above] {$F_i^H$};
  \fill (0,0) circle (0.06) node[below] {$q$};
  \draw[load] (-1.8,0) -- (-0.15,0) node[pos=0,above] {$P_i$};
  \foreach \y in {0.4,0.2,0,-0.2,-0.4}{\draw[load] (1.2,\y) -- (0.2,\y);}
  \draw[decorate,decoration={brace,amplitude=4pt}] (1.35,0.55) -- (1.35,-0.55);
  \node[right] at (1.5,0) {$F_i^A,\,F_i^B,\,F_i^C,\,F_i^D,\,F_i^H$};
\end{tikzpicture}

    \caption{Equilibrium in the $P_i$--$\Delta_i$ direction, member added.}\end{subfigure}
  \caption{Assembly of the global stiffness equation at a joint (after Figs.~3.1--3.4 of the textbook).}
  \label{fig:mcguire-joint}
\end{figure}

With reference to Fig.~\ref{fig:mcguire-joint}, equilibrium at joint $q$ in the $\Delta_i$ direction requires
\begin{equation}
  P_i = F_i^A + F_i^B + F_i^C + F_i^D. \tag{3.2}
\end{equation}
Writing each bar force in terms of the bar's own stiffness coefficients and displacements,
\begin{equation}
  \begin{aligned}
    P_i ={}& (k_{ii}^A\Delta_i^A + k_{i2}^A\Delta_2^A + k_{i3}^A\Delta_3^A + k_{i4}^A\Delta_4^A)\\
          &+ (k_{ii}^B\Delta_i^B + k_{i2}^B\Delta_2^B + k_{i5}^B\Delta_5^B + k_{i6}^B\Delta_6^B)\\
          &+ (k_{ii}^C\Delta_i^C + k_{i2}^C\Delta_2^C + k_{i7}^C\Delta_7^C + k_{i8}^C\Delta_8^C)\\
          &+ (k_{ii}^D\Delta_i^D + k_{i1}^D\Delta_1^D + k_{i2}^D\Delta_2^D + k_{i9}^D\Delta_9^D).
  \end{aligned} \tag{3.3}
\end{equation}
Since compatibility requires $\Delta_i^A = \Delta_i^B = \Delta_i^C = \Delta_i^D = \Delta_i$, it follows that
\begin{equation}
  \begin{aligned}
    P_i ={}& (k_{ii}^A + k_{ii}^B + k_{ii}^C + k_{ii}^D)\Delta_i + k_{i1}^D\Delta_1\\
          &+ (k_{i2}^A + k_{i2}^B + k_{i2}^C + k_{i2}^D)\Delta_2 + k_{i3}^A\Delta_3 + \cdots + k_{i9}^D\Delta_9
  \end{aligned} \tag{3.4}
\end{equation}
or
\begin{equation}
  P_i = K_{ii}\Delta_i + K_{i1}\Delta_1 + K_{i2}\Delta_2 + K_{i3}\Delta_3 + \cdots + K_{i9}\Delta_9. \tag{3.5}
\end{equation}
This is the final form of the desired equations. The capitalized terms $K_{ii}, K_{i1}, K_{i2},\dots,K_{i9}$ are \emph{global stiffness coefficients} and Equation~(3.5) is a \emph{global stiffness equation}.

\section{System equations for supported structures}
The system equations of $\Kmat\bs{\Delta} = \Pvec$ can be rewritten as
\begin{equation}
  \bmat{\Kmat_{ff} & \Kmat_{fs}\\ \Kmat_{sf} & \Kmat_{ss}}\Bmat{\bs{\Delta}_f\\ \bs{\Delta}_s} = \Bmat{\Pvec_f\\ \Pvec_s}
\end{equation}
where all quantities pertinent to the supports are assigned the subscript $s$ and those relating to the remaining degrees of freedom have the subscript $f$. Note that $\Kmat_{fs} = \Kmat_{sf}^\tr$.

For the case where $\bs{\Delta}_s = \mat{0}$, we have
\begin{align}
  \Kmat_{ff}\bs{\Delta}_f &= \Pvec_f\\
  \Kmat_{sf}\bs{\Delta}_f &= \Pvec_s
\end{align}
We can solve the first equation to obtain $\bs{\Delta}_f$, and substitute $\bs{\Delta}_f$ into the second equation to obtain the reaction vector $\Pvec_s$.

For the case where $\bs{\Delta}_s = \overline{\bs{\Delta}}_s$, a given vector, we have
\begin{align}
  \Kmat_{ff}\bs{\Delta}_f &= \Pvec_f - \Kmat_{fs}\overline{\bs{\Delta}}_s\\
  \Kmat_{sf}\bs{\Delta}_f + \Kmat_{ss}\overline{\bs{\Delta}}_s &= \Pvec_s
\end{align}
Similarly, we can solve the first equation to obtain $\bs{\Delta}_f$, and substitute $\bs{\Delta}_f$ and $\overline{\bs{\Delta}}_s$ into the second equation to obtain the reaction vector $\Pvec_s$.

In our program, we only formulate $\Kmat_{ff}$ and we use the fixed-end forces to calculate the equivalent nodal forces, which are exactly the same as $(-\Kmat_{fs}\overline{\bs{\Delta}}_s)$ of the first equation. By doing so, we can still obtain the same $\bs{\Delta}_f$.

After obtaining $\bs{\Delta}_f$, the member (element) forces can be calculated by $\kmat_i\bs{\Delta}_i = \F_i$ for the $i$-th element, where note that $\F_i$ is referred to the global (local) axes if $\kmat_i$ and $\bs{\Delta}_i$ are referred to the global (local) axes. Of course, when talking about internal forces such as axial and shear forces, and bending moment, $\F_i$ should always be referred to the local (member) axes.

\section{Example: assembling of stiffness matrix}
\begin{fig}
  % Example: 3-4-5 triangular truss
\begin{tikzpicture}[scale=1.0,font=\small]
  \coordinate (n1) at (0,0); \coordinate (n2) at (3,0); \coordinate (n3) at (0,4);
  \draw[member] (n1) -- (n2) -- (n3) -- cycle;
  \foreach \p in {n1,n2,n3}{\node[node dot] at (\p) {};}
  \node[below left] at (n1) {1}; \node[right] at (n2) {2}; \node[above] at (n3) {3};
  \node[circle,draw,inner sep=1.5pt] at (1.5,0.45) {\scriptsize 1};
  \node[circle,draw,inner sep=1.5pt] at (1.95,2.35) {\scriptsize 2};
  \node[circle,draw,inner sep=1.5pt] at (0.45,1.6) {\scriptsize 3};
  % dimensions
  \draw[dim] (-0.8,0) -- (-0.8,4) node[midway,fill=white,inner sep=1pt] {4\,m};
  \draw (-1.0,0) -- (-0.6,0); \draw (-1.0,4) -- (-0.6,4);
  \draw[dim] (0,-0.6) -- (3,-0.6) node[midway,fill=white,inner sep=1pt] {3\,m};
  \draw (0,-0.8) -- (0,-0.4); \draw (3,-0.8) -- (3,-0.4);
  % axes
  \draw[->] (-0.9,-1.1) -- (-0.3,-1.1) node[right] {$x$};
  \draw[->] (-0.9,-1.1) -- (-0.9,-0.5) node[left] {$y$};
  % data
  \node[anchor=west,align=left] at (4.6,2.6) {Areas\quad $\begin{aligned}A_1 &= 600~\mathrm{mm^2}\\ A_2 &= 1000~\mathrm{mm^2}\\ A_3 &= 800~\mathrm{mm^2}\end{aligned}$};
  \node[anchor=west] at (4.6,1.0) {Young's modulus $E = 200$~GPa};
\end{tikzpicture}

\end{fig}
For all elements, $\dfrac{EA}{L} = 40{,}000$~kN/m.
\begin{fig}
  % Global (u,v,F) and local (u',v',F') member displacements and forces
\begin{tikzpicture}[scale=0.9,font=\small]
  \def\ang{40}
  % ---- global displacements ----
  \begin{scope}
    \coordinate (A) at (0,0); \coordinate (B) at (\ang:3);
    \draw[member] (A) -- (B);
    \draw[load] (A) -- ++(1.2,0) node[right] {$u_1$}; \draw[load] (A) -- ++(0,1.2) node[left] {$v_1$};
    \draw[load] (B) -- ++(1.2,0) node[right] {$u_2$}; \draw[load] (B) -- ++(0,1.2) node[left] {$v_2$};
  \end{scope}
  % ---- global forces ----
  \begin{scope}[xshift=5.5cm]
    \coordinate (A) at (0,0); \coordinate (B) at (\ang:3);
    \draw[member] (A) -- (B);
    \draw[load] (A) -- ++(1.2,0) node[right] {$F_{x1}$}; \draw[load] (A) -- ++(0,1.2) node[left] {$F_{y1}$};
    \draw[load] (B) -- ++(1.2,0) node[right] {$F_{x2}$}; \draw[load] (B) -- ++(0,1.2) node[left] {$F_{y2}$};
  \end{scope}
  \node[anchor=west] at (-5.4,1.2) {Global\quad $\kmat\dvec = \mat{f}$};
  % ---- local displacements ----
  \begin{scope}[yshift=-4.2cm]
    \coordinate (A) at (0,0); \coordinate (B) at (\ang:3);
    \draw[member] (A) -- (B);
    \draw[load] ($(A)+(\ang+180:1.2)$) -- ($(A)+(\ang+180:0.05)$) node[pos=0,below] {$u_1'$};
    \draw[load] (A) -- ++(\ang+90:1.1) node[above left] {$v_1'$};
    \draw[load] (B) -- ++(\ang:1.2) node[right] {$u_2'$};
    \draw[load] (B) -- ++(\ang+90:1.1) node[above] {$v_2'$};
  \end{scope}
  % ---- local forces ----
  \begin{scope}[xshift=5.5cm,yshift=-4.2cm]
    \coordinate (A) at (0,0); \coordinate (B) at (\ang:3);
    \draw[member] (A) -- (B);
    \draw[load] ($(A)+(\ang+180:1.2)$) -- ($(A)+(\ang+180:0.05)$) node[pos=0,below] {$F_{x1}'$};
    \draw[load] (A) -- ++(\ang+90:1.1) node[above left] {$F_{y1}'$};
    \draw[load] (B) -- ++(\ang:1.2) node[right] {$F_{x2}'$};
    \draw[load] (B) -- ++(\ang+90:1.1) node[above] {$F_{y2}'$};
  \end{scope}
  \node[anchor=west] at (-5.4,-3.0) {Local\quad $\kmat'\dvec' = \mat{f}'$};
\end{tikzpicture}

\end{fig}
\begin{equation}
  \boxed{\kmat = \Tmat^\tr\kmat'\Tmat}
\end{equation}
where
\begin{equation}
  \Tmat = \bmat{c & s & & \\ -s & c & & \\ & & c & s\\ & & -s & c}, \qquad c = \cos\phi,\quad s = \sin\phi
\end{equation}
and
\begin{equation}
  \kmat' = \frac{EA}{L}\;
  \begin{array}{c}
    \begin{array}{cccc} u_1' & v_1' & u_2' & v_2' \end{array}\\
    \left[\begin{array}{@{\;}r@{\quad}r@{\quad}r@{\quad}r@{\;}}
      1 & 0 & -1 & 0\\ 0 & 0 & 0 & 0\\ -1 & 0 & 1 & 0\\ 0 & 0 & 0 & 0
    \end{array}\right]
  \end{array}
\end{equation}
Therefore, the element stiffness matrices in the global coordinate system can be obtained as:

\noindent\textbf{Mem \textcircled{\scriptsize 1}:} $\phi = 0^\circ$, $c = 1.0$, $s = 0.0$
\begin{equation}
  \kmat_1 = \kmat' = 40{,}000\;
  \begin{array}{c}
    \begin{array}{cccc} \;u_1\; & \;v_1\; & \;u_2\; & \;v_2\; \end{array}\\
    \left[\begin{array}{@{\;}r@{\quad}r@{\quad}r@{\quad}r@{\;}}
      1 & 0 & -1 & 0\\ & 0 & 0 & 0\\ & & 1 & 0\\ \multicolumn{2}{l}{\textit{Sym.}} & & 0
    \end{array}\right]
  \end{array}\;\text{kN/m}
\end{equation}

\noindent\textbf{Mem \textcircled{\scriptsize 2}:} $\phi = 126.9^\circ$, $c = -0.6$, $s = 0.8$
\begin{fig}
  % Member 2 of the example truss: phi = 126.9 degrees
\begin{tikzpicture}[scale=1.0,font=\small]
  \coordinate (A) at (0,0); \coordinate (B) at (126.87:2.5);
  \draw[member] (A) -- (B);
  \draw[dashed] (A) -- (2.0,0);
  \draw[->] (1.0,0) arc (0:126.87:1.0);
  \node at (0.55,1.05) {$\phi$};
  \node[below] at (A) {1}; \node[above left] at (B) {2};
\end{tikzpicture}

\end{fig}
\begin{equation}
  \kmat_2 = 40{,}000\;
  \begin{array}{c}
    \begin{array}{cccc} \;\;u_2\;\; & \;\;v_2\;\; & \;\;u_3\;\; & \;\;v_3\;\; \end{array}\\
    \left[\begin{array}{@{\;}r@{\quad}r@{\quad}r@{\quad}r@{\;}}
      0.36 & -0.48 & -0.36 & 0.48\\ & 0.64 & 0.48 & -0.64\\ & & 0.36 & -0.48\\ \multicolumn{2}{l}{\textit{Sym.}} & & 0.64
    \end{array}\right]
  \end{array}\;\text{kN/m}
\end{equation}

\noindent\textbf{Mem \textcircled{\scriptsize 3}:} $\phi = 90^\circ$, $c = 0.0$, $s = 1.0$
\begin{equation}
  \kmat_3 = 40{,}000\;
  \begin{array}{c}
    \begin{array}{cccc} \;u_1\; & \;v_1\; & \;u_3\; & \;v_3\; \end{array}\\
    \left[\begin{array}{@{\;}r@{\quad}r@{\quad}r@{\quad}r@{\;}}
      0 & 0 & 0 & 0\\ & 1 & 0 & -1\\ & & 0 & 0\\ \multicolumn{2}{l}{\textit{Sym.}} & & 1
    \end{array}\right]
  \end{array}\;\text{kN/m}
\end{equation}

After assembling,
\begin{equation}
  \Kmat = 40{,}000\;
  \begin{array}{c}
    \begin{array}{cccccc} \;u_1\; & \;\;v_1\;\; & \;\;u_2\;\; & \;\;v_2\;\; & \;\;u_3\;\; & \;\;v_3\;\; \end{array}\\
    \left[\begin{array}{@{\;}r@{\quad}r@{\quad}r@{\quad}r@{\quad}r@{\quad}r@{\;}}
      1 & 0 & -1.00 & 0.00 & 0.00 & 0.00\\
        & 1 &  0.00 & 0.00 & 0.00 & -1.00\\
        &   &  1.36 & -0.48 & -0.36 & 0.48\\
        &   &       & 0.64 & 0.48 & -0.64\\
        &   &       &      & 0.36 & -0.48\\
      \multicolumn{3}{l}{\textit{Sym.}} & & & 1.64
    \end{array}\right]
  \end{array}\;\text{kN/m}
\end{equation}
\begin{fig}
  % Assembly pattern of the structural stiffness matrix (which element contributes where)
\begin{tikzpicture}[scale=0.9,font=\small,
  e1/.style={pattern={Lines[angle=45,distance=3pt]},pattern color=black!70},
  e2/.style={pattern={Lines[angle=90,distance=3pt]},pattern color=black!70},
  e3/.style={pattern={Lines[angle=0,distance=3pt]},pattern color=black!70}]
  % rows/cols: 1=u1 2=v1 3=u2 4=v2 5=u3 6=v3 ; cell (i,j) at x=j-1, y=6-i
  % element 1: dofs 1-4 ; element 2: dofs 3-6 ; element 3: dofs 1,2,5,6
  \fill[e1] (0,2) rectangle (4,6);
  \fill[e2] (2,0) rectangle (6,4);
  \fill[e3] (0,4) rectangle (2,6); \fill[e3] (4,0) rectangle (6,2);
  \fill[e3] (0,0) rectangle (2,2); \fill[e3] (4,4) rectangle (6,6);
  % grid lines separating nodes
  \draw[black!50] (0,0) rectangle (6,6);
  \draw[black!50] (2,0) -- (2,6); \draw[black!50] (4,0) -- (4,6);
  \draw[black!50] (0,2) -- (6,2); \draw[black!50] (0,4) -- (6,4);
  % brackets
  \draw[thick] (-0.15,-0.1) -- (-0.3,-0.1) -- (-0.3,6.1) -- (-0.15,6.1);
  \draw[thick] (6.15,-0.1) -- (6.3,-0.1) -- (6.3,6.1) -- (6.15,6.1);
  \node[left] at (-0.6,3) {$\Kmat=$};
  % column labels
  \foreach \x/\t in {0.5/u_1,1.5/v_1,2.5/u_2,3.5/v_2,4.5/u_3,5.5/v_3}{\node[above] at (\x,6.2) {$\t$};}
  % legend
  \begin{scope}[xshift=7.6cm]
    \fill[e1] (0,4.6) rectangle (1,5.6); \draw (0,4.6) rectangle (1,5.6);
    \node[right] at (1.2,5.1) {Element \textcircled{\scriptsize 1}};
    \fill[e2] (0,2.9) rectangle (1,3.9); \draw (0,2.9) rectangle (1,3.9);
    \node[right] at (1.2,3.4) {Element \textcircled{\scriptsize 2}};
    \fill[e3] (0,1.2) rectangle (1,2.2); \draw (0,1.2) rectangle (1,2.2);
    \node[right] at (1.2,1.7) {Element \textcircled{\scriptsize 3}};
  \end{scope}
\end{tikzpicture}

\end{fig}
It can be found that the eigenvalues of $\Kmat$ are: $0, 0, 0, 40000, 80000$, and $120000$. The three zero eigenvalues correspond to the rigid body modes in the plane (two translations plus one rotation).

\noindent$\Kmat\bs{\Delta}_r = \mat{0}$: find the rigid body modes
\begin{equation}
  \bs{\Delta}_{r1} = \Bmat{1\\0\\1\\0\\1\\0},\qquad
  \bs{\Delta}_{r2} = \Bmat{0\\1\\0\\1\\0\\1},\qquad
  \bs{\Delta}_{r3} = \Bmat{0\\0\\0\\3\\-4\\0}
\end{equation}
\begin{fig}
  % Three rigid-body modes of the unsupported example truss
\begin{tikzpicture}[scale=0.7,font=\small]
  \def\tri{(0,0) -- (3,0) -- (0,4) -- cycle}
  % ---- mode 1: translation in x ----
  \begin{scope}
    \draw[thick] \tri;
    \draw[dashed] (1,0) -- (4,0) -- (1,4) -- cycle;
    \foreach \p in {(0,0),(3,0),(0,4),(1,0),(4,0),(1,4)}{\node[node dot] at \p {};}
    \draw[dim] (0,4.7) -- (1,4.7) node[midway,above] {1};
    \draw[dim] (0,-0.7) -- (1,-0.7) node[midway,below] {1};
    \draw[dim] (3,-0.7) -- (4,-0.7) node[midway,below] {1};
  \end{scope}
  % ---- mode 2: translation in y ----
  \begin{scope}[xshift=8cm]
    \draw[thick] \tri;
    \draw[dashed] (0,1) -- (3,1) -- (0,5) -- cycle;
    \foreach \p in {(0,0),(3,0),(0,4),(0,1),(3,1),(0,5)}{\node[node dot] at \p {};}
    \draw[dim] (-0.7,0) -- (-0.7,1) node[midway,left] {1};
    \draw[dim] (3.7,0) -- (3.7,1) node[midway,right] {1};
    \draw[dim] (0.7,4) -- (0.7,5) node[midway,right] {1};
  \end{scope}
  % ---- mode 3: rotation about node 1 (theta = 1) ----
  \begin{scope}[xshift=16cm]
    \draw[thick] \tri;
    \draw[dashed] (0,0) -- (3,3) -- (-4,4) -- cycle;
    \foreach \p in {(0,0),(3,0),(0,4),(3,3),(-4,4)}{\node[node dot] at \p {};}
    \draw[dim] (-4,4.7) -- (0,4.7) node[midway,above] {4};
    \draw[dim] (3.7,0) -- (3.7,3) node[midway,right] {3};
    \draw[->] (1.6,0) arc (0:45:1.6); \node[right] at (1.7,0.55) {$\theta=1$};
  \end{scope}
\end{tikzpicture}

\end{fig}

Let's support the truss and apply two loads as shown below. Now we would like to solve for the displacements.
\begin{fig}
  % Supported example truss with two loads at node 2
\begin{tikzpicture}[scale=1.0,font=\small]
  \coordinate (n1) at (0,0); \coordinate (n2) at (3,0); \coordinate (n3) at (0,4);
  \draw[member] (n1) -- (n2) -- (n3) -- cycle;
  \fill[hatch] (-0.85,3.5) rectangle (-0.65,4.5); \draw (-0.65,3.5) -- (-0.65,4.5);
  \foreach \p in {n1,n2,n3}{\node[node dot] at (\p) {};}
  \node[below right] at ($(n1)+(0.05,-0.1)$) {1}; \node[below] at ($(n2)+(0,-0.05)$) {2}; \node at (-0.3,3.55) {3};
  \node[circle,draw,inner sep=1.5pt] at (1.5,0.45) {\scriptsize 1};
  \node[circle,draw,inner sep=1.5pt] at (1.95,2.35) {\scriptsize 2};
  \node[circle,draw,inner sep=1.5pt] at (-0.45,1.6) {\scriptsize 3};
  % supports: pin at 1, horizontal roller at 3
  \pic at (n1) {pin};
  \pic[rotate=-90] at (n3) {roller};
  % dof arrows
  \foreach \p/\i in {n1/1,n2/2,n3/3}{
    \draw[->,thin] ($(\p)+(0.15,0.15)$) -- ++(0.55,0) node[right,inner sep=1pt] {$u_\i$};
    \draw[->,thin] ($(\p)+(0.15,0.15)$) -- ++(0,0.55) node[above,inner sep=1pt] {$v_\i$};}
  % loads
  \draw[load] (n2) -- ++(1.4,0) node[right] {72\,kN};
  \draw[load] (n2) -- ++(0,1.4) node[above] {90\,kN};
\end{tikzpicture}

\end{fig}
Because $u_1 = v_1 = u_3 = 0$ $\Rightarrow$ delete the columns and rows corresponding to $u_1$, $v_1$ and $u_3$. Knowing that $P_{x2} = 72$~kN and $P_{y2} = 90$~kN, therefore
\begin{equation}
  40{,}000\;
  \begin{array}{c}
    \begin{array}{ccc} \;\;u_2\;\; & \;\;v_2\;\; & \;\;v_3\;\; \end{array}\\
    \left[\begin{array}{@{\;}r@{\quad}r@{\quad}r@{\;}}
      1.36 & -0.48 & 0.48\\ & 0.64 & -0.64\\ \textit{Sym.} & & 1.64
    \end{array}\right]
  \end{array}
  \Bmat{u_2\\ v_2\\ v_3} = \Bmat{72\\ 90\\ 0}
  \qquad (\Kmat\bs{\Delta} = \Pvec)
\end{equation}
\begin{equation}
  \Longrightarrow\quad \bs{\Delta} = \Bmat{u_2\\ v_2\\ v_3} = 10^{-3}\Bmat{3.49\\ 8.38\\ 2.25}\text{ m} = \Bmat{3.49\\ 8.38\\ 2.25}\text{ mm}
\end{equation}

\subsection*{Calculation of member forces (force recovery)}
Member \textcircled{\scriptsize 2}, for example.

\topic{Member (local) displacements}
\begin{fig}
  % Member 2: global nodal displacements -> local displacements
\begin{tikzpicture}[scale=0.85,font=\small]
  \def\ang{126.87}
  \begin{scope}
    \coordinate (A) at (0,0); \coordinate (B) at (\ang:3);
    \draw[member] (A) -- (B);
    \node[below] at (A) {1}; \node[left] at (B) {2};
    \node[circle,draw,inner sep=1pt] at ($(A)!0.5!(B)+(-0.45,0)$) {\scriptsize 2};
    \draw[load] (A) -- ++(1.0,0) node[right] {$u_2$}; \draw[load] (A) -- ++(0,1.0) node[right] {$v_2$};
    \draw[load] (B) -- ++(1.0,0) node[right] {$u_3=0$}; \draw[load] (B) -- ++(0,1.0) node[left] {$v_3$};
  \end{scope}
  \node at (4.2,1.2) {$\Longrightarrow$};
  \begin{scope}[xshift=8.5cm]
    \coordinate (A) at (0,0); \coordinate (B) at (\ang:3);
    \draw[member] (A) -- (B);
    \draw[load] ($(A)+(\ang+180:1.0)$) -- ($(A)+(\ang+180:0.05)$) node[pos=0,right] {$u_1'$};
    \draw[load] (A) -- ++(\ang+90:1.0) node[left] {$v_1'$};
    \draw[load] (B) -- ++(\ang:1.0) node[above] {$u_2'$};
    \draw[load] (B) -- ++(\ang+90:1.0) node[below left] {$v_2'$};
  \end{scope}
\end{tikzpicture}

\end{fig}
\begin{equation}
  \dvec' = \Tmat\dvec \quad\Longrightarrow\quad
  \Bmat{u_1'\\ v_1'\\ u_2'\\ v_2'} =
  \bmat{c & s & & \\ -s & c & & \\ & & c & s\\ & & -s & c}
  \Bmat{u_2\\ v_2\\ u_3\\ v_3} = \Bmat{4.61\\ -7.82\\ 1.80\\ -1.35}\text{ mm}
\end{equation}
($c = -0.6$, $s = 0.8$ $\because$ $\phi = 126.9^\circ$)

\topic{Member forces}
\begin{equation}
  \mat{f}' = \kmat'\dvec' = \frac{EA}{L}
  \begin{array}{c}
    \begin{array}{cccc} u_1' & v_1' & u_2' & v_2' \end{array}\\
    \left[\begin{array}{@{\;}r@{\quad}r@{\quad}r@{\quad}r@{\;}}
      1 & 0 & -1 & 0\\ 0 & 0 & 0 & 0\\ -1 & 0 & 1 & 0\\ 0 & 0 & 0 & 0
    \end{array}\right]
  \end{array}
  \Bmat{4.61\\ -7.82\\ 1.80\\ -1.35}\times 10^{-3}\text{ m}
  = \Bmat{112.4\\ 0.\\ -112.4\\ 0.}\text{ kN}
\end{equation}
$\bigl(\frac{EA}{L} = 40{,}000\text{ kN/m}\bigr)$
\begin{fig}
  % Recovered member forces of the example truss
\begin{tikzpicture}[scale=1.0,font=\small]
  \def\ang{126.87}
  % member 2
  \begin{scope}
    \coordinate (A) at (0,0); \coordinate (B) at (\ang:3);
    \draw[member] (A) -- (B);
    \node[circle,draw,inner sep=1pt] at ($(A)!0.5!(B)+(-0.45,0)$) {\scriptsize 2};
    \draw[load] ($(A)+(\ang+180:1.1)$) -- ($(A)+(\ang+180:0.1)$) node[pos=0,right] {112.4\,kN};
    \draw[load] ($(B)+(\ang:1.1)$) -- ($(B)+(\ang:0.1)$) node[pos=0,right] {112.4\,kN};
  \end{scope}
  \node at (3.6,1.6) {Similarly,};
  % member 3
  \begin{scope}[xshift=6.2cm,yshift=-0.6cm]
    \draw[member] (0,0) -- (0,3.2);
    \node[circle,draw,inner sep=1pt] at (-0.45,1.6) {\scriptsize 3};
    \draw[load] (0,3.3) -- (0,4.3) node[right] {90\,kN};
    \draw[load] (0,-0.1) -- (0,-1.1) node[right] {90\,kN};
  \end{scope}
  % member 1
  \begin{scope}[xshift=8.6cm,yshift=1.5cm]
    \draw[member] (0,0) -- (2.6,0);
    \node[circle,draw,inner sep=1pt] at (1.3,0.4) {\scriptsize 1};
    \draw[load] (-0.1,0) -- (-1.0,0) node[below,pos=0.5] {139.5\,kN};
    \draw[load] (2.7,0) -- (3.6,0) node[below,pos=0.5] {139.5\,kN};
  \end{scope}
\end{tikzpicture}

\end{fig}
\noindent$***$ You can verify that the nodal equilibrium equations are satisfied.

\section{Example: adding a member to a structure}
\begin{minipage}[t]{0.5\textwidth}
The truss shown is subjected to a force at node $d$; the cross-sectional areas of members are as follows:
\begin{align*}
  A_{ab} = A_{cd} &= 10\times10^3~\mathrm{mm^2}\\
  A_{bc} = A_{bd} = A_{ac} &= 15\times10^3~\mathrm{mm^2}
\end{align*}
By the direct stiffness method described above, we can calculate the nodal displacements:
\begin{equation*}
  \{\Delta\} = \Bmat{u_b\\ u_c\\ v_c\\ u_d\\ v_d} = \Bmat{0\\ 2.2\\ -0.5\\ 2.45\\ -1.41}\text{ mm}
\end{equation*}
\end{minipage}\hfill
\begin{minipage}[t]{0.45\textwidth}
\vspace{0pt}\centering
% Truss subjected to a load at node d
\begin{tikzpicture}[scale=0.8,font=\small]
  \coordinate (a) at (0,0); \coordinate (b) at (2,0); \coordinate (c) at (2,4); \coordinate (d) at (4,4);
  \draw[thinmember] (a) -- (b) -- (c) -- (d); \draw[thinmember] (a) -- (c); \draw[thinmember] (b) -- (d);
  \node[above left] at (a) {$a$}; \node[right] at ($(b)+(0.1,0.2)$) {$b$}; \node[above left] at (c) {$c$}; \node[above] at (d) {$d$};
  \pic at (a) {pin}; \pic at (b) {roller};
  \draw[load] (d) -- ++(1.3,0) node[right] {100\,N};
\end{tikzpicture}

\end{minipage}

\begin{minipage}[t]{0.5\textwidth}
Now, add one member to the structure at node $d$; the same load of 100~N is applied at node $d$. If we want to determine the displacement at node $d$, there is a simple way to solve it.
\begin{equation*}
  A_{de} = 5\times10^3~\mathrm{mm^2}\ \text{(given)}
\end{equation*}
\end{minipage}\hfill
\begin{minipage}[t]{0.45\textwidth}
\vspace{0pt}\centering
% The same truss with an additional member de
\begin{tikzpicture}[scale=0.8,font=\small]
  \coordinate (a) at (0,0); \coordinate (b) at (2,0); \coordinate (c) at (2,4); \coordinate (d) at (4,4); \coordinate (e) at (6,4);
  \draw[thinmember] (a) -- (b) -- (c) -- (d); \draw[thinmember] (a) -- (c); \draw[thinmember] (b) -- (d);
  \draw[thinmember] (d) -- (e);
  \node[above left] at (a) {$a$}; \node[right] at ($(b)+(0.1,0.2)$) {$b$}; \node[above left] at (c) {$c$}; \node[above] at (d) {$d$}; \node[above] at (e) {$e$};
  \pic at (a) {pin}; \pic at (b) {roller};
  \pic at (e) {wallR};
\end{tikzpicture}

\end{minipage}

This new structure is like two springs in parallel. The stiffness of the original structure is $k_1 = \dfrac{P_0}{u_d} = \dfrac{100}{2.45}$, and that of member $de$ is $k_2 = \dfrac{EA_{de}}{L_{de}}$. Then we have
\begin{equation}
  u_d = \frac{100}{k_1 + k_2}.
\end{equation}

\section{General features of the structural stiffness matrix}
\begin{enumerate}[leftmargin=*]
  \item The matrix is \textbf{singular if not supported}; the number of zero eigenvalues $=$ the number of rigid body modes plus the number of mechanisms. For example, the structural stiffness matrix of Example~3.1 of the textbook has three eigenvalues that are zero. They correspond to the translation along the $x$- and $y$-axes and the rotation about the $z$-axis.

    \emph{Note: meaning of eigenvalues.}
    \begin{equation}
      \Kmat\bs{\varphi} = \lambda\bs{\varphi} \;\Rightarrow\; \bs{\varphi}^\tr\Kmat\bs{\varphi} = \lambda\bs{\varphi}^\tr\bs{\varphi}
      \;\Rightarrow\; \bs{\varphi}^\tr\Kmat\bs{\varphi} = \lambda\bs{\varphi}^\tr\bs{\varphi}
      \;\Rightarrow\; \bs{\varphi}^\tr\Kmat\bs{\varphi} = \lambda \quad\text{if } \bs{\varphi}^\tr\bs{\varphi} = 1
    \end{equation}
    This means that $\lambda_i =$ twice the strain energy associated with $\bs{\varphi}_i$ if $\bs{\varphi}_i$ is normalized. Therefore, $\lambda_i = 0$ corresponds to rigid body modes or mechanisms.

    These eigenvalues also indicate whether the solution for $\Kmat\bs{\Delta} = \Pvec$ exists or not. If there is no eigenvalue that equals 0, $\Kmat$ is nonsingular. Then the solution exists and is unique undoubtedly. However, if there is any eigenvalue that equals 0, it does not mean that there is no solution. If $\bs{\varphi}$ is the eigenvector corresponding to $\lambda = 0$ and if $\bs{\varphi}\cdot\Pvec = 0$, then we will get an infinite number of solutions. Why? If $\bs{\varphi}\cdot\Pvec\neq 0$, there will be no solution at all. Problem~3.12 of McGuire et al.\ (2000) is a good example for the above discussion. Also, please refer to the paper reproduced at the end of this chapter (Sec.~\ref{sec:leu-paper}).

  \item The $k$-th column of the structural stiffness matrix of an unsupported structure represents a set of forces in \textbf{equilibrium} due to $\Delta_k = 1$, and other $\Delta$'s $= 0$. For example, we can verify that each column of $\Kmat$ of Example~3.1 of the textbook satisfies the following three equilibrium equations:
    \begin{equation*}
      \sum F_x = 0,\qquad \sum F_y = 0,\qquad\text{and}\qquad \sum M_z = 0.
    \end{equation*}

  \item $\Kmat$ is \textbf{banded} if the nodes are properly numbered; usually we number the nodes along the side whose number of nodes is less than that of the other side(s). For the structure shown below (left), the stiffness matrix is full since nodes are not properly numbered.
    \begin{fig}
      % Node numbering and bandwidth: improper (left) vs. proper (right)
\begin{tikzpicture}[scale=0.8,font=\small]
  % ---- tall frame, numbered along the long side: full matrix ----
  \begin{scope}
    \foreach \x in {0,2,4}{\draw[member] (\x,0) -- (\x,5);}
    \foreach \y in {1,2,3,4,5}{\draw[member] (0,\y) -- (4,\y);}
    \foreach \x in {0,2,4}{\pic at (\x,0) {fixed};}
    \foreach \y/\r in {0/0,1/1,2/2,3/3,4/4,5/5}{
      \pgfmathtruncatemacro{\a}{3*\r+1}\pgfmathtruncatemacro{\b}{3*\r+2}\pgfmathtruncatemacro{\c}{3*\r+3}
      \node[above left,inner sep=1pt,fill=white] at (0,\y) {\a};
      \node[above left,inner sep=1pt,fill=white] at (2,\y) {\b};
      \node[above left,inner sep=1pt,fill=white] at (4,\y) {\c};}
  \end{scope}
  % ---- wide frame, numbered along the short side: banded ----
  \begin{scope}[xshift=7cm]
    \foreach \x in {0,2,4,6,8}{\draw[member] (\x,0) -- (\x,2);}
    \foreach \y in {1,2}{\draw[member] (0,\y) -- (8,\y);}
    \foreach \x in {0,2,4,6,8}{\pic at (\x,0) {fixed};}
    \foreach \x/\c in {0/0,2/1,4/2,6/3,8/4}{
      \pgfmathtruncatemacro{\a}{3*\c+1}\pgfmathtruncatemacro{\b}{3*\c+2}\pgfmathtruncatemacro{\d}{3*\c+3}
      \node[above left,inner sep=1pt,fill=white] at (\x,0) {\a};
      \node[above left,inner sep=1pt,fill=white] at (\x,1) {\b};
      \node[above left,inner sep=1pt,fill=white] at (\x,2) {\d};}
  \end{scope}
\end{tikzpicture}

    \end{fig}
    \begin{fig}
      % Continuous beam with an improper DOF numbering (K not banded)
\begin{tikzpicture}[scale=1.0,font=\small]
  \draw[member] (0,0) -- (9,0);
  \pic at (0,0) {pin}; \pic at (9,0) {pin};
  \foreach \x in {1.5,3,4.5,6,7.5}{\pic at (\x,0) {roller};}
  \foreach \x/\n in {0/1,1.5/7,3/2,4.5/6,6/3,7.5/5,9/4}{
    \draw[->,thick] ($(\x,0)+(0.35,0.2)$) arc (20:160:0.35);
    \node[above] at (\x,0.55) {$D\n$};}
\end{tikzpicture}

    \end{fig}
    \begin{equation*}
      [\Kmat] =
      \begin{array}{c}
        \begin{array}{ccccccc} 1 & 2 & 3 & 4 & 5 & 6 & 7 \end{array}\\
        \left[\begin{array}{ccccccc}
          X & O & O & O & O & O & X\\
          O & X & O & O & O & X & X\\
          O & O & X & O & X & X & O\\
          O & O & O & X & X & O & O\\
          O & O & X & X & X & O & O\\
          O & X & X & O & O & X & O\\
          X & X & O & O & O & O & X
        \end{array}\right]
      \end{array}
    \end{equation*}
    However, if we number the nodes as shown below, $\Kmat$ is banded with the semi-band width $=2$.
    \begin{fig}
      % Continuous beam with a proper DOF numbering (semi-bandwidth = 2)
\begin{tikzpicture}[scale=1.0,font=\small]
  \draw[member] (0,0) -- (9,0);
  \pic at (0,0) {pin}; \pic at (9,0) {pin};
  \foreach \x in {1.5,3,4.5,6,7.5}{\pic at (\x,0) {roller};}
  \foreach \x/\n in {0/1,1.5/2,3/3,4.5/4,6/5,7.5/6,9/7}{
    \draw[->,thick] ($(\x,0)+(0.35,0.2)$) arc (20:160:0.35);
    \node[above] at (\x,0.55) {$D\n$};}
\end{tikzpicture}

    \end{fig}
    \begin{equation*}
      [\Kmat] =
      \begin{array}{c}
        \begin{array}{ccccccc} 1 & 2 & 3 & 4 & 5 & 6 & 7 \end{array}\\
        \left[\begin{array}{ccccccc}
          X & X & O & O & O & O & O\\
          X & X & X & O & O & O & O\\
          O & X & X & X & O & O & O\\
          O & O & X & X & X & O & O\\
          O & O & O & X & X & X & O\\
          O & O & O & O & X & X & X\\
          O & O & O & O & O & X & X
        \end{array}\right]
      \end{array}
    \end{equation*}
    \begin{figure}[H]
      \centering
      \begin{subfigure}[t]{0.48\textwidth}\centering% Joint stiffness matrix for a continuous beam (after Weaver & Gere, Fig. 4-9)
\begin{tikzpicture}[scale=0.38,font=\scriptsize]
  % ---- (a) beam with DOF numbering: translation odd, rotation even ----
  \begin{scope}[yshift=2cm,xscale=1.0]
    \draw[member] (0,0) -- (18,0);
    \foreach \i in {0,...,6}{
      \pgfmathtruncatemacro{\t}{2*\i+1}\pgfmathtruncatemacro{\r}{2*\i+2}
      \draw[fill=black!50] (3*\i,0) -- ++(-0.3,-0.55) -- ++(0.6,0) -- cycle;
      \draw[->,thick] (3*\i,0.1) -- (3*\i,1.6) node[above] {\t};
      \draw[->,thick] (3*\i-1.3,-1.1) -- (3*\i-0.1,-0.1) node[pos=0,below left] {\r};}
    \foreach \i in {1,...,6}{\node[above] at (3*\i-1.5,0.05) {\i};}
    \node at (9,-2.6) {(a)};
  \end{scope}
  % ---- (b) matrix pattern ----
  \begin{scope}[yshift=-19cm]
    \foreach \m in {1,...,6}{
      \pgfmathsetmacro{\x}{2*(\m-1)} \pgfmathsetmacro{\y}{14-2*(\m-1)}
      \fill[pattern={Lines[angle=45,distance=2.5pt]},pattern color=black!70] (\x,\y-4) rectangle (\x+4,\y);
      \node[fill=white,inner sep=0.5pt] at (\x+3.5,\y-0.5) {\m};}
    \draw[step=1,black!60,thin] (0,0) grid (14,14);
    \draw[thick] (0,0) rectangle (14,14);
    \foreach \i in {1,...,14}{
      \node[above] at (\i-0.5,14) {\i}; \node[left] at (0,14.5-\i) {\i};}
    % alternate numbering shown on the right and bottom
    \foreach \i/\n in {1/8,2/1,3/9,4/2,5/10,6/3,7/11,8/4,9/12,10/5,11/13,12/6,13/14,14/7}{
      \node[right] at (14,14.5-\i) {\n}; \node[below] at (\i-0.5,0) {\n};}
    \node at (7,-1.6) {(b)};
  \end{scope}
\end{tikzpicture}

        \caption{Joint stiffness matrix for a continuous beam (Fig.~4-9).}\end{subfigure}\hfill
      \begin{subfigure}[t]{0.48\textwidth}\centering% Rearranged joint stiffness matrix for a continuous beam (after Weaver & Gere, Fig. 4-10)
\begin{tikzpicture}[scale=0.38,font=\scriptsize]
  % ---- (a) beam: rotations 1..7, translations 8..14 ----
  \begin{scope}[yshift=2cm]
    \draw[member] (0,0) -- (18,0);
    \foreach \i in {0,...,6}{
      \pgfmathtruncatemacro{\t}{\i+8}\pgfmathtruncatemacro{\r}{\i+1}
      \draw[fill=black!50] (3*\i,0) -- ++(-0.3,-0.55) -- ++(0.6,0) -- cycle;
      \draw[->,thick] (3*\i,0.1) -- (3*\i,1.6) node[above] {\t};
      \draw[->,thick] (3*\i-1.3,-1.1) -- (3*\i-0.1,-0.1) node[pos=0,below left] {\r};}
    \foreach \i in {1,...,6}{\node[above] at (3*\i-1.5,0.05) {\i};}
    \node at (9,-2.6) {(a)};
  \end{scope}
  % ---- (b) matrix pattern: four tridiagonal 7x7 blocks ----
  \begin{scope}[yshift=-19cm]
    \foreach \bx in {0,7}{\foreach \by in {0,7}{
      \foreach \i in {1,...,7}{
        \pgfmathsetmacro{\x}{\bx+\i-1} \pgfmathsetmacro{\y}{\by+7-\i}
        \fill[pattern={Lines[angle=45,distance=2.5pt]},pattern color=black!70] (\x,\y) rectangle (\x+1,\y+1);
        \ifnum\i<7
          \fill[pattern={Lines[angle=45,distance=2.5pt]},pattern color=black!70] (\x+1,\y) rectangle (\x+2,\y+1);
          \fill[pattern={Lines[angle=45,distance=2.5pt]},pattern color=black!70] (\x,\y-1) rectangle (\x+1,\y);
        \fi}}}
    \draw[step=1,black!60,thin] (0,0) grid (14,14);
    \draw[thick] (0,0) rectangle (14,14);
    \draw[dashed,thick] (7,0) -- (7,14); \draw[dashed,thick] (0,7) -- (14,7);
    \foreach \i in {1,...,14}{\node[above] at (\i-0.5,14) {\i}; \node[left] at (0,14.5-\i) {\i};}
    \node at (7,-1.6) {(b)};
  \end{scope}
\end{tikzpicture}

        \caption{Rearranged joint stiffness matrix for a continuous beam (Fig.~4-10).}\end{subfigure}
      \caption{After Weaver, Jr.\ and Gere (1990), \emph{Matrix Analysis of Framed Structures}.}
    \end{figure}
\end{enumerate}

\section{Generalized stiffness and loads}
\topic{Original system:}
\begin{equation}
  \Kmat_{n\times n}\bs{\Delta}_{n\times 1} = \Pvec_{n\times 1} \qquad\text{(assume properly supported; w/o rigid body modes)}
\end{equation}
\topic{Eigenvalues and eigenvectors:} $(\lambda_i,\bs{\varphi}_i)$, $i = 1\sim n$
\begin{equation}
  \Kmat\bs{\varphi}_i = \lambda_i\bs{\varphi}_i
\end{equation}
\topic{Property of $\bs{\varphi}_i$ (orthogonality):}
\begin{equation}
  \begin{cases}
    i\neq j:\; \bs{\varphi}_j^\tr\Kmat\bs{\varphi}_i = \lambda_i\bs{\varphi}_j^\tr\bs{\varphi}_i = 0\\[4pt]
    i = j:\; \bs{\varphi}_i^\tr\Kmat\bs{\varphi}_i = \lambda_i\bs{\varphi}_i^\tr\bs{\varphi}_i = \lambda_i\|\bs{\varphi}_i\|^2 = \lambda_i \quad\text{if } \|\bs{\varphi}_i\| = 1
  \end{cases}
\end{equation}
Without loss of generality, we let $\|\bs{\varphi}_i\| = 1$, $i = 1\sim n$.

The above orthogonal condition can be written in the following matrix form
\begin{equation}
  \bs{\Phi}^\tr\Kmat\bs{\Phi} = \bs{\Lambda}
\end{equation}
where
\begin{equation}
  \bs{\Phi}_{n\times n} = [\bs{\varphi}_1,\bs{\varphi}_2,\bs{\varphi}_3,\dots,\bs{\varphi}_n],\qquad
  \bs{\Lambda} = \bmat{\lambda_1 & & & \\ & \lambda_2 & & \\ & & \ddots & \\ & & & \lambda_n},\ \text{a diagonal matrix}.
\end{equation}
Note: $\bs{\Phi}^\tr\bs{\Phi} = \mat{I}$ (i.e.\ $\bs{\varphi}_i^\tr\bs{\varphi}_i = 1$, $i = 1\sim n$; $\bs{\varphi}_i^\tr\bs{\varphi}_j = 0$, $i\neq j$), where $\mat{I}$ is an identity matrix.

\topic{Stiffness equations in the eigenspace:}
Write the displacement vector in terms of the eigenvectors:
\begin{equation}
  \bs{\Delta} = \sum_{i=1}^{n}c_i\bs{\varphi}_i = \bs{\Phi}\mat{c};\qquad\text{we would like to determine }\mat{c}.
\end{equation}
Express the load vector $\Pvec$ in the eigenspace:
\begin{equation}
  \Pvec = \sum_{i=1}^{n}d_i\bs{\varphi}_i = \bs{\Phi}\mat{d}
\end{equation}
where the generalized load coefficients $d_i$ can be obtained by premultiplying the above equation by $\bs{\Phi}^\tr$. Hence,
\begin{gather}
  \bs{\Phi}^\tr\Pvec = \bs{\Phi}^\tr\bs{\Phi}\mat{d} = \mat{d}\quad(\because\ \bs{\Phi}^\tr\bs{\Phi} = \mat{I})\\
  \text{i.e.}\quad \bs{\varphi}_i^\tr\Pvec = d_i,\ i = 1\sim n\qquad(\bs{\varphi}_i^\tr\Pvec = \bs{\varphi}_i\bullet\Pvec)
\end{gather}
This means $d_i$ is the projection of $\Pvec$ on $\bs{\varphi}_i$; $d_i$ can be considered a generalized load.

Substituting $\bs{\Delta} = \bs{\Phi}\mat{c}$ and $\Pvec = \bs{\Phi}\mat{d}$ into the equilibrium equations $\Kmat\bs{\Delta} = \Pvec$ yields
\begin{equation}
  \Kmat(\bs{\Phi}\mat{c}) = \bs{\Phi}\mat{d}
\end{equation}
Premultiply it by $\bs{\Phi}^\tr$: $\bs{\Phi}^\tr\Kmat\bs{\Phi}\mat{c} = \bs{\Phi}^\tr\bs{\Phi}\mat{d}$; $\because\ \bs{\Phi}^\tr\Kmat\bs{\Phi} = \bs{\Lambda}$ and $\bs{\Phi}^\tr\bs{\Phi} = \mat{I}$, we have
\begin{equation}
  \bs{\Lambda}\mat{c} = \mat{d},\quad\text{i.e.}\quad
  \bmat{\lambda_1 & & & \\ & \lambda_2 & & \\ & & \ddots & \\ & & & \lambda_n}\Bmat{c_1\\ c_2\\ \vdots\\ c_n} = \Bmat{d_1\\ d_2\\ \vdots\\ d_n}
\end{equation}
$\Longrightarrow$ Uncoupled system; can be called the ``Generalized System''.
\begin{itemize}
  \item $\bs{\Lambda} =$ generalized stiffness matrix
  \item $\mat{d} =$ generalized load vector
\end{itemize}

\topic{Existence of solutions}
\begin{enumerate}[leftmargin=*]
  \item All $\lambda_i\neq 0$ $\Leftrightarrow$ $\Kmat$ is nonsingular $\Leftrightarrow$ solution is unique:
    \begin{equation*}
      c_i = \frac{d_i}{\lambda_i},\ i = 1\sim n. \qquad\Bigl(\text{all }\lambda_i > 0,\text{ then }\Kmat = \text{positive definite}\Bigr)
    \end{equation*}
  \item Some $\lambda_i = 0$; say $\lambda_1 = 0$, $\lambda_2 = 0$, and other $\lambda_i\neq 0$
    \begin{equation*}
      \Longrightarrow\quad
      \begin{cases}
        (0)c_1 = d_1\\ (0)c_2 = d_2\\ \quad\vdots\\ \lambda_n c_n = d_n
      \end{cases}
    \end{equation*}
    \begin{description}[leftmargin=3.5em,style=nextline]
      \item[\fbox{2(a)}] If $d_1 = d_2 = 0$ (i.e.\ $\bs{\varphi}_1\bullet\Pvec = 0$, $\bs{\varphi}_2\bullet\Pvec = 0$)
        $\ \Updownarrow\ $ infinite number of solutions; $c_1, c_2 =$ arbitrary, and other $c_i = d_i/\lambda_i$.
      \item[\fbox{2(b)}] 1.~$d_1\neq 0$;\quad 2.~$d_2\neq 0$;\quad 3.~$d_1\neq 0$ and $d_2\neq 0$
        $\quad\Leftrightarrow\quad$ no solution.
    \end{description}
\end{enumerate}

\clearpage
% =====================================================================
\section[Appendix: Existence of a solution for a singular system (Leu, 2004)]{Appendix: Condition for existence of a solution and its physical interpretation for a singular structural system with applications to a bar chain}
\label{sec:leu-paper}
\begin{center}
  \emph{Liang-Jenq Leu}\\
  Department of Civil Engineering, National Taiwan University, Taipei 10617, Taiwan\\[2pt]
  {\small The Seventeenth KKCNN Symposium on Civil Engineering, December 13--15, 2004, Thailand}
\end{center}

\subsection*{Abstract}
This paper is concerned with the condition for the existence of a solution to a singular structural system and the corresponding physical meanings. A three-bar chain is discussed first. It is shown that this three-bar chain has one mechanism and the condition for the existence of a solution to this structure is such that the external load vector is orthogonal to the mechanism. Then, two methods are proposed to find the mechanism. A physical interpretation of the condition for the existence of a solution for the three-bar chain is also given. On the basis of the physical interpretation, a third method is proposed to obtain the mechanism. The advantage of the three proposed methods is that there is no need to form the structural stiffness matrix. Finally, it is shown that the proposed methods for finding a mechanism and the physical interpretation of the condition for the existence of a solution can be readily extended to the general case of an $N$-bar chain. To this end, a five-bar chain with three mechanisms is illustrated.

\subsection*{Introduction}
Structural systems taught in courses like matrix structural analysis usually have nonsingular stiffness matrices. However, singular stiffness matrices exist for some real problems. For example, the stiffness matrix for the three-bar chain shown in Fig.~\ref{fig:leu1} (McGuire et al.\ 2000) is singular. ``Can a useful solution be obtained when $P_b = P_c$?'' and ``what happens when $P_b\neq P_c$?'' were asked in McGuire et al.\ 2000. There are several ways to answer the above questions as discussed below. For the convenience of discussion, a more general loading case is considered as depicted in Fig.~\ref{fig:leu2}. Note that the effect of geometric nonlinearity is not considered in this study.

\begin{figure}[H]
  \centering
  \begin{subfigure}[b]{0.48\textwidth}\centering% Three-bar chain subject to two loads (paper Fig. 1)
\begin{tikzpicture}[scale=0.45,font=\small]
  \coordinate (a) at (0,3); \coordinate (b) at (4,0); \coordinate (c) at (8,0); \coordinate (d) at (12,3);
  \draw[member] (a) -- (b) -- (c) -- (d);
  \pic[rotate=-36.87] at (a) {wallL};
  \pic[rotate=36.87] at (d) {wallR};
  \node[above right] at (a) {$a$}; \node[above] at (b) {$b$}; \node[above] at (c) {$c$}; \node[above left] at (d) {$d$};
  \draw[load] (b) -- ++(0,-1.6) node[below] {$P_b$};
  \draw[load] (c) -- ++(0,-1.6) node[below] {$P_c$};
  \draw[dashed] (a) -- (0,0) -- (b); \draw[dashed] (d) -- (12,0) -- (c);
  \draw[dim] (-1.2,0) -- (-1.2,3) node[midway,left] {$3m$};
  \draw[dim] (12.9,0) -- (12.9,3) node[midway,right] {$3m$};
  \draw[dim] (0,-0.8) -- (4,-0.8) node[midway,fill=white,inner sep=1pt] {$4m$};
  \draw[dim] (4,-0.8) -- (8,-0.8) node[midway,fill=white,inner sep=1pt] {$4m$};
  \draw[dim] (8,-0.8) -- (12,-0.8) node[midway,fill=white,inner sep=1pt] {$4m$};
\end{tikzpicture}

    \caption{Three-bar chain subject to two loads.}\label{fig:leu1}\end{subfigure}\hfill
  \begin{subfigure}[b]{0.48\textwidth}\centering% Three-bar chain subject to four loads (paper Fig. 2)
\begin{tikzpicture}[scale=0.55,font=\small]
  \coordinate (n0) at (0,3); \coordinate (n1) at (4,0); \coordinate (n2) at (8,0); \coordinate (n3) at (12,3);
  \draw[member] (n0) -- (n1) -- (n2) -- (n3);
  \pic[rotate=-36.87] at (n0) {wallL};
  \pic[rotate=36.87] at (n3) {wallR};
  \foreach \p in {n1,n2}{\node[node dot] at (\p) {};}
  \node[below right] at (n0) {0}; \node[below] at (n1) {1}; \node[below] at (n2) {2}; \node[below left] at (n3) {3};
  \node[circle,draw,inner sep=1pt] at (1.5,2.0) {\scriptsize 1};
  \node[circle,draw,inner sep=1pt] at (6,-0.9) {\scriptsize 2};
  \node[circle,draw,inner sep=1pt] at (10.5,2.0) {\scriptsize 3};
  \draw[load] (n1) -- ++(0,1.8) node[above] {$v_1\,(P_2)$};
  \draw[load] (n2) -- ++(0,1.8) node[above] {$v_2\,(P_4)$};
  \draw[load] (n1) -- ++(-2.2,0) node[above] {$u_1\,(P_1)$};
  \draw[load] (n2) -- ++(2.2,0) node[above] {$u_2\,(P_3)$};
\end{tikzpicture}

    \caption{Three-bar chain subject to four loads.}\label{fig:leu2}\end{subfigure}
\end{figure}

\subsection*{Mathematical background}
First let's write the stiffness equations of a typical structural system in the usual way:
\begin{equation}
  \Kmat\mat{U} = \Pvec \tag{1}\label{leu:1}
\end{equation}
where $\Kmat_{n\times n}$ is the elastic stiffness matrix that is symmetric, $\mat{U}_{n\times 1}$ is the displacement vector, $\Pvec_{n\times 1}$ is the load vector, and $n$ is the number of degrees of freedom. It is well known that the solution of Eq.~\eqref{leu:1} exists and is unique when $\Kmat$ is nonsingular.

This paper mainly investigates the existence of a solution for the case where $\Kmat$ is singular and more importantly the corresponding physical meanings. Mathematically, a solution to Eq.~\eqref{leu:1} exists if and only if the rank of $\Kmat$ is the same as that of the augmented matrix, which is obtained by adjoining the vector $\Pvec$ to the matrix $\Kmat$ as the $(n+1)$th column. If this rank condition is satisfied, the number of solutions is unique when $\Kmat$ is nonsingular and infinite when singular. It is also known that the rank of $\Kmat$ is $n$ when $\Kmat$ is nonsingular and less than $n$ when singular.

Although the concept of rank is simple, it may not be convenient to use. The conditions for the existence of an infinite number of solutions can be stated in another way. According to Greenberg (1978), if the homogeneous form of Eq.~\eqref{leu:1} has a nontrivial solution, denoted as $\bs{\varphi}$, and in addition the solution is orthogonal to the load vector $\Pvec$, then there are an infinite number of solutions.
\begin{equation}
  \Kmat\bs{\varphi} = \mat{0}\ \text{ and }\ \bs{\varphi}^\tr\Pvec = 0 \quad\Longleftrightarrow\quad \text{infinitely many solutions} \tag{2}\label{leu:2}
\end{equation}

\subsection*{Generalized system}
To appreciate the meaning of Eq.~\eqref{leu:2}, the concept of generalized systems is employed. The generalized system is defined here as a system written in the space spanned by the eigenvectors of $\Kmat$. The same idea is frequently used in structural dynamic analysis when carrying out a modal superposition to obtain the dynamic response. However, this idea is seldom used in static analysis.

Let the eigenvalues and eigenvectors of $\Kmat$ be written as: $(\lambda_i,\bs{\varphi}_i)$, $i = 1,2,\dots,n$. Each eigenpair satisfies the equation:
\begin{equation}
  \Kmat\bs{\varphi}_i = \lambda_i\bs{\varphi}_i \tag{3}
\end{equation}
According to the orthogonality of eigenvectors, one has
\begin{equation}
  \begin{cases}
    i\neq j:\; \bs{\varphi}_j^\tr\Kmat\bs{\varphi}_i = \lambda_i\bs{\varphi}_j^\tr\bs{\varphi}_i = 0\\[4pt]
    i = j:\; \bs{\varphi}_i^\tr\Kmat\bs{\varphi}_i = \lambda_i\bs{\varphi}_i^\tr\bs{\varphi}_i = \lambda_i\|\bs{\varphi}_i\|^2 = \lambda_i\quad\text{if }\|\bs{\varphi}_i\|\equiv\sqrt{\bs{\varphi}_i^\tr\bs{\varphi}_i} = 1
  \end{cases} \tag{4a,b}
\end{equation}
Note that without loss of generality, the eigenvectors in Eq.~(4b) have been normalized.

Equations (4a,b) can be written in matrix form as
\begin{equation}
  \bs{\Phi}^\tr\Kmat\bs{\Phi} = \bs{\Lambda} \tag{5}\label{leu:5}
\end{equation}
where
\begin{equation}
  \bs{\Phi}_{n\times n} = [\bs{\varphi}_1,\bs{\varphi}_2,\dots,\bs{\varphi}_n]\quad\text{and}\quad
  \bs{\Lambda} = \bmat{\lambda_1 & & & \\ & \lambda_2 & & \\ & & \ddots & \\ & & & \lambda_n}. \tag{6a,b}
\end{equation}
From Eqs.~(4a,b), one also obtains
\begin{equation}
  \bs{\Phi}^\tr\bs{\Phi} = \mat{I} \tag{7}
\end{equation}
Let the displacement vector $\mat{U}$ be expressed as
\begin{equation}
  \mat{U} = \sum_{i=1}^{n}c_i\bs{\varphi}_i = \bs{\Phi}\mat{c} \tag{8}\label{leu:8}
\end{equation}
Substituting Eq.~\eqref{leu:8} into \eqref{leu:1} results in
\begin{equation}
  \Kmat\bs{\Phi}\mat{c} = \Pvec \tag{9}\label{leu:9}
\end{equation}
Premultiplying Eq.~\eqref{leu:9} by $\bs{\Phi}^\tr$ and using Eq.~\eqref{leu:5}, one obtains
\begin{equation}
  \bs{\Lambda}\mat{c} = \bs{\Phi}^\tr\Pvec \tag{10}\label{leu:10}
\end{equation}
Note that the generalized system of Eq.~\eqref{leu:10} is uncoupled as expected. When written in component form, Eq.~\eqref{leu:10} becomes
\begin{equation}
  \lambda_i c_i = \overline{P}_i,\quad i = 1,2,\dots,n \tag{11}\label{leu:11}
\end{equation}
where the generalized loads are defined as
\begin{equation}
  \overline{P}_i = \bs{\varphi}_i^\tr\Pvec,\quad i = 1,2,\dots,n \tag{12}
\end{equation}
From Eq.~\eqref{leu:10} or \eqref{leu:11}, it can be seen that the eigenvalues $\lambda_i$'s can be interpreted as the generalized stiffnesses and the coefficients $c_i$'s generalized displacements.

\subsection*{Existence of a solution}
The generalized system of Eq.~\eqref{leu:11} can be used conveniently to discuss the existence of a solution to the original system of Eq.~\eqref{leu:1}. Obviously, if all $\lambda_i$'s are nonzero, the solution is unique since in this case $c_i = \overline{P}_i/\lambda_i$ and the solution $\mat{U}$ can be obtained using Eq.~\eqref{leu:8}.

If some $\lambda_i$'s are zero, then there are an infinite number of solutions if the corresponding generalized loads $\overline{P}_i$ are also zero; this is exactly what Eq.~\eqref{leu:2} says. To prove this, without loss of generality, assume the first $r$ eigenvalues are zero. Then the generalized displacements $c_i$, $i = 1,2,\dots,r$ are arbitrary if $\overline{P}_i = 0$, $i = 1,2,\dots,r$ according to Eq.~\eqref{leu:11}. Also, because of $\Kmat\bs{\varphi}_i = \mat{0}$, $i = 1,2,\dots,r$ the linear combination of these eigenvectors satisfies the equation
\begin{equation}
  \Kmat\mat{U}_c = \mat{0} \tag{13}
\end{equation}
where
\begin{equation}
  \mat{U}_c = c_1\bs{\varphi}_1 + c_2\bs{\varphi}_2 + \dots + c_r\bs{\varphi}_r \tag{14}\label{leu:14}
\end{equation}
The complete solution to Eq.~\eqref{leu:1} is
\begin{equation}
  \mat{U} = \mat{U}_p + \mat{U}_c \tag{15}\label{leu:15}
\end{equation}
where $\mat{U}_p$ is a particular solution. It can be obtained from Eqs.~\eqref{leu:11} and \eqref{leu:8} as
\begin{equation}
  \mat{U}_p = \frac{\overline{P}_{r+1}}{\lambda_{r+1}}\bs{\varphi}_{r+1} + \frac{\overline{P}_{r+2}}{\lambda_{r+2}}\bs{\varphi}_{r+2} + \dots + \frac{\overline{P}_n}{\lambda_n}\bs{\varphi}_n \tag{16}
\end{equation}
The number of solutions given in Eq.~\eqref{leu:15} is infinite since the coefficients $c_i$ ($i = 1,2,\dots,r$) in Eq.~\eqref{leu:14} are arbitrary.

What happens when an eigenvalue $\lambda_i = 0$ and its corresponding generalized load $\overline{P}_i$ is nonzero? From Eq.~\eqref{leu:11}, the generalized displacement $c_i$ will be infinity since its stiffness is zero and therefore the solution won't be meaningful. Mathematically speaking, the resulting equation is inconsistent since $0 = \overline{P}_i$ and hence there is no solution.

\subsection*{Mechanism of three-bar chain}
Given a properly supported structure (without rigid body modes) such as the one shown in Fig.~\ref{fig:leu2}, how does one find the eigenvector corresponding to a zero eigenvalue if its $\Kmat$ is singular? Of course, this can be done easily by using an eigen-solver if $\Kmat$ has been formed. However, the associated physical meaning may not be appreciated if done in this manner. A displacement vector that satisfies $\Kmat\bs{\varphi} = \mat{0}$ is actually a mechanism and therefore the eigenvector corresponding to a zero eigenvalue is a mechanism.

Two methods are proposed in this section to find a mechanism for the three-bar chain shown in Fig.~\ref{fig:leu2}. The first method uses the condition that the axial deformation of each bar is zero since its axial force is zero. With reference to Fig.~\ref{fig:leu3}, the axial deformation of a typical bar can be calculated using
\begin{equation}
  \delta_i = (u_i\cos\theta_i + v_i\sin\theta_i) - (u_{i-1}\cos\theta_i + v_{i-1}\sin\theta_i) \tag{17}\label{leu:17}
\end{equation}
By making use of Eq.~\eqref{leu:17}, the condition of three zero axial deformations for the three-bar chain of Fig.~\ref{fig:leu2} can be expressed as
\begin{equation}
  \begin{aligned}
    \delta_1 &= \tfrac45 u_1 - \tfrac35 v_1 = 0\\
    \delta_2 &= u_2 - u_1 = 0\\
    \delta_3 &= -\tfrac45 u_2 - \tfrac35 v_2 = 0
  \end{aligned} \tag{18a,b,c}\label{leu:18}
\end{equation}
There are four unknowns and three equations in Eqs.~\eqref{leu:18}; the solution is
\begin{equation}
  \bs{\varphi} = \frac{1}{\sqrt{50}}\langle 3, 4, 3, -4\rangle^\tr \tag{19}\label{leu:19}
\end{equation}
where $\bs{\varphi}$ has been normalized and its displacement components are listed in the order of $u_1$, $v_1$, $u_2$, and $v_2$. It can be verified that the displacement vector of Eq.~\eqref{leu:19} satisfies $\Kmat\bs{\varphi} = \mat{0}$ and therefore it corresponds to a mechanism.

\begin{figure}[H]
  \centering
  \begin{subfigure}[b]{0.4\textwidth}\centering% A typical bar (paper Fig. 3)
\begin{tikzpicture}[scale=1.0,font=\small]
  \coordinate (A) at (0,0); \coordinate (B) at (40:2.6);
  \draw[member] (A) -- (B);
  \draw[load] (A) -- ++(1.3,0) node[below] {$u_{i-1}$};
  \draw[load] (A) -- ++(0,1.3) node[left] {$v_{i-1}$};
  \draw[load] (B) -- ++(1.3,0) node[right] {$u_i$};
  \draw[load] (B) -- ++(0,1.3) node[left] {$v_i$};
  \draw[->] ($(A)+(0.8,0)$) arc (0:40:0.8); \node at ($(A)+(1.05,0.35)$) {$\theta_i$};
\end{tikzpicture}

    \caption{A typical bar.}\label{fig:leu3}\end{subfigure}\hfill
  \begin{subfigure}[b]{0.55\textwidth}\centering% Compatible rigid-body mode of the three-bar chain (paper Fig. 4)
\begin{tikzpicture}[scale=0.48,font=\small]
  \coordinate (n0) at (0,3); \coordinate (n1) at (4,0); \coordinate (n2) at (8,0); \coordinate (n3) at (12,3);
  \coordinate (a) at (6,-1.5);
  \draw[member] (n0) -- (n1) -- (n2) -- (n3);
  \pic[rotate=-36.87] at (n0) {wallL};
  \pic[rotate=36.87] at (n3) {wallR};
  \foreach \p in {n1,n2}{\node[node dot] at (\p) {};}
  \node[below right] at (n0) {0}; \node[above right] at (n1) {1}; \node[above left] at (n2) {2}; \node[below left] at (n3) {3};
  % axes of bars 1 and 3 meet at a
  \draw[dashed] (n1) -- (a) -- (n2);
  \fill (a) circle (0.08); \node[below] at (a) {$a$};
  % normal displacements
  \draw[load] (n1) -- ++(53.13:1.8) node[above] {$\delta_{n1}\mat{n}_1$};
  \draw[load] (n2) -- ++(-53.13:1.8) node[right] {$\delta_{n2}\mat{n}_2$};
  % moment arms
  \draw[dim] ($(n1)+(-0.9,-0.9)$) -- ($(a)+(-0.9,-0.9)$) node[midway,below left] {$d_1$};
  \draw[dim] ($(a)+(0.9,-0.9)$) -- ($(n2)+(0.9,-0.9)$) node[midway,below right] {$d_2$};
\end{tikzpicture}

    \caption{Compatible rigid body mode.}\label{fig:leu4}\end{subfigure}
\end{figure}

The second way to obtain the mechanism is to draw a compatible rigid body mode for each bar. Fig.~\ref{fig:leu4} shows a compatible rigid body mode for the three-bar chain of Fig.~\ref{fig:leu2}. The rigid body modes for bars 1 and 3 result from a rotation with respect to points 0 and 3, respectively. As shown in Fig.~\ref{fig:leu4}, let the resulting displacement vector of points 1 and 2 be denoted as $\delta_{n1}\mat{n}_1$ and $\delta_{n2}\mat{n}_2$, where $\mat{n}_1$ and $\mat{n}_2$ are the unit vectors normal to the axes of bars 1 and 3, respectively. Since $\mat{n}_1$ and $\mat{n}_2$ are known, the remaining problem is to determine the ratio of the normal displacements $\delta_{n1}$ to $\delta_{n2}$ so that bar 2 also moves as a rigid body.

By employing the concept of ``instantaneous center of rotation'' in conventional rigid body dynamics, one can find the center of rotation of bar 2, point $a$, is actually the intersection of axes of bars 1 and 3 as depicted in Fig.~\ref{fig:leu4}. Let the distances between the center of rotation and points 1 and 2 be represented as $d_1$ and $d_2$, respectively. Then, the displacements of points 1 and 2 of bar 2 due to a rotation with respect to the center of rotation are proportional to their distances to it, i.e.
\begin{equation}
  \frac{\delta_{n1}}{\delta_{n2}} = \frac{d_1}{d_2} \tag{20}\label{leu:20}
\end{equation}
For the particular case of Fig.~\ref{fig:leu2}, $\delta_{n1} = \delta_{n2}$ since $d_1 = d_2$. Therefore, one can obtain the same mechanism as that expressed in Eq.~\eqref{leu:19}.

Having obtained the mechanism of Eq.~\eqref{leu:19}, the condition of Eq.~\eqref{leu:2} for infinitely many solutions can be written explicitly as
\begin{equation}
  \bs{\varphi}^\tr\Pvec = \frac{1}{\sqrt{50}}(3P_1 + 4P_2 + 3P_3 - 4P_4) = 0 \tag{21}\label{leu:21}
\end{equation}
According to Eq.~\eqref{leu:21}, it is easy to answer the questions concerned with the three-bar chain of Fig.~\ref{fig:leu1} given in McGuire et al.\ (2000). Obviously, there will be a solution and actually infinitely many if $P_b = P_c$ and there is no solution when $P_b\neq P_c$. However, Eq.~\eqref{leu:21} is very general and can be used to answer whether a solution exists for the general loading case of Fig.~\ref{fig:leu2}.

\subsection*{Physical interpretation of zero generalized load associated with a mechanism}
That the generalized load associated with a mechanism is zero as stated in Eq.~\eqref{leu:21} is the condition for a solution to exist. This was derived from Eq.~\eqref{leu:11} and was discussed before. In addition to this derivation, an interesting and useful physical interpretation of Eq.~\eqref{leu:21} is given below and on the basis of interpretation, a third method for finding a mechanism is proposed.

The free-body diagram of the three-bar chain is shown in Fig.~\ref{fig:leu5}, where $f_1$ and $f_3$ are the axial forces of bars 1 and 3, respectively. The existence of a solution to the three-bar chain is equivalent to saying that the equilibrium condition is satisfied for the free-body diagram of Fig.~\ref{fig:leu5}. Hence, the moment about point $a$, the center of rotation of bar 2, due to $f_1$, $f_3$, and $P_i$, $i = 1\sim 4$ must be zero. The contribution of $f_1$ and $f_3$ is zero since these two axial forces pass through point $a$. The contribution of the applied loads at points 1 and 2 can be calculated by multiplying their normal components by the moment arms. Hence,
\begin{equation}
  \textstyle\sum M_a = (\mat{n}_1^\tr\Pvec_1)d_1 + (\mat{n}_2^\tr\Pvec_2)d_2 = 0 \tag{22}\label{leu:22}
\end{equation}
where $\Pvec_1 = \langle P_1, P_2\rangle^\tr$ and $\Pvec_2 = \langle P_3, P_4\rangle^\tr$ are the applied loads at points 1 and 2, respectively. The normal vectors $\mat{n}_1$ and $\mat{n}_2$ and the moment arms $d_1$ and $d_2$ were defined in Fig.~\ref{fig:leu4} before. According to Eq.~\eqref{leu:20}, Eq.~\eqref{leu:22} is equivalent to
\begin{equation}
  (\mat{n}_1^\tr\Pvec_1)\delta_{n1} + (\mat{n}_2^\tr\Pvec_2)\delta_{n2} = 0 \tag{23}\label{leu:23}
\end{equation}
From Fig.~\ref{fig:leu4}, the displacements of points 1 and 2 can be written, respectively, as $\bs{\Delta}_1 = \delta_{n1}\mat{n}_1$ and $\bs{\Delta}_2 = \delta_{n2}\mat{n}_2$; they are actually the first two and the last two components of the mechanism $\bs{\varphi}$. Hence Eq.~\eqref{leu:23} can be rewritten as
\begin{equation}
  \bs{\Delta}_1^\tr\Pvec_1 + \bs{\Delta}_2^\tr\Pvec_2 = \bs{\varphi}^\tr\Pvec = 0 \tag{24}
\end{equation}
since $\Pvec_1$ and $\Pvec_2$ are also the first two and the last two components of $\Pvec$. Equation~(24) shows that Eq.~\eqref{leu:22} is equivalent to Eq.~\eqref{leu:21}.

The main finding here is that Eq.~\eqref{leu:21} is equivalent to stating that the resulting moment of the applied loads at points 1 and 2 about the intersection of the axes of bars 1 and 3 is zero. By using this finding, one can also find the mechanism for the three-bar chain of Fig.~\ref{fig:leu2}. According to Fig.~\ref{fig:leu5}, this zero-moment condition yields
\begin{equation}
  P_1\Bigl(\frac32\Bigr) + P_2(2) + P_3\Bigl(\frac32\Bigr) - P_4(2) = 0 \tag{25}\label{leu:25}
\end{equation}
Clearly, this is equivalent to Eq.~\eqref{leu:21} and therefore the same mechanism can be found. This is called the third proposed method for finding a mechanism. It is emphasized that there is no need to form $\Kmat$ when using the three proposed methods to find a mechanism.

\begin{figure}[H]
  \centering
  \begin{subfigure}[b]{0.42\textwidth}\centering% Free-body diagram of the three-bar chain (paper Fig. 5)
\begin{tikzpicture}[scale=0.48,font=\small]
  \coordinate (n0) at (0,3); \coordinate (n1) at (4,0); \coordinate (n2) at (8,0); \coordinate (n3) at (12,3);
  \coordinate (a) at (6,-1.5);
  \draw[member] (n0) -- (n1) -- (n2) -- (n3);
  \foreach \p in {n1,n2}{\node[node dot] at (\p) {};}
  \node[below right] at (n0) {0}; \node[above right] at (n1) {1}; \node[above left] at (n2) {2}; \node[below left] at (n3) {3};
  \draw[dashed] (n1) -- (a) -- (n2); \fill (a) circle (0.08); \node[below] at (a) {$a$};
  % axial forces at the supports
  \draw[load] (n0) -- ++(143.13:1.3) node[above left,inner sep=1pt] {$f_1$};
  \draw[load] (n3) -- ++(36.87:1.3) node[above right,inner sep=1pt] {$f_3$};
  % applied loads
  \draw[load] ($(n1)+(-2.2,0)$) -- (n1) node[pos=0,left] {$P_1$};
  \draw[load] (n1) -- ++(0,1.8) node[above] {$P_2$};
  \draw[load] (n2) -- ++(2.2,0) node[right] {$P_3$};
  \draw[load] (n2) -- ++(0,1.8) node[above] {$P_4$};
  \draw[dim] ($(n1)+(-0.9,-0.9)$) -- ($(a)+(-0.9,-0.9)$) node[midway,below left] {$d_1$};
  \draw[dim] ($(a)+(0.9,-0.9)$) -- ($(n2)+(0.9,-0.9)$) node[midway,below right] {$d_2$};
\end{tikzpicture}

    \caption{Free-body diagram of the three-bar chain.}\label{fig:leu5}\end{subfigure}\hfill
  \begin{subfigure}[b]{0.55\textwidth}\centering% Five-bar chain (paper Fig. 6)
\begin{tikzpicture}[scale=0.62,font=\scriptsize]
  \coordinate (n0) at (0,5);
  \coordinate (n1) at ($(n0)+(-60:3.4)$);
  \coordinate (n2) at ($(n1)+(-30:3.4)$);
  \coordinate (n3) at ($(n2)+(0:3.4)$);
  \coordinate (n4) at ($(n3)+(30:3.4)$);
  \coordinate (n5) at ($(n4)+(60:3.4)$);
  \draw[member] (n0) -- (n1) -- (n2) -- (n3) -- (n4) -- (n5);
  \pic[rotate=-60] at (n0) {wallL};
  \pic[rotate=60] at (n5) {wallR};
  \foreach \p in {n1,n2,n3,n4}{\node[node dot] at (\p) {};}
  \node[right] at ($(n0)+(0.1,0)$) {0}; \node[left] at (n1) {1}; \node[below left] at (n2) {2}; \node[below right] at (n3) {3}; \node[above left] at (n4) {4}; \node[left] at ($(n5)+(-0.1,0)$) {5};
  % rotation centres
  \coordinate (a) at (intersection of n0--n1 and n2--n3);
  \coordinate (b) at (intersection of n1--n2 and n3--n4);
  \coordinate (c) at (intersection of n2--n3 and n4--n5);
  \draw[dashed] (n1) -- (a) -- (n2); \draw[dashed] (n2) -- (b) -- (n3); \draw[dashed] (n3) -- (c) -- (n4);
  \fill (a) circle (0.07) node[left] {a}; \fill (b) circle (0.07) node[below] {b}; \fill (c) circle (0.07) node[right] {c};
  % dofs / loads
  \draw[load] (n1) -- ++(1.1,0) node[below right,inner sep=0pt] {$u_1(P_1)$}; \draw[load] (n1) -- ++(0,1.1) node[above,inner sep=1pt] {$v_1(P_2)$};
  \draw[load] (n2) -- ++(1.1,0) node[above right,inner sep=0pt] {$u_2(P_3)$}; \draw[load] (n2) -- ++(0,1.1) node[above left,inner sep=0pt] {$v_2(P_4)$};
  \draw[load] (n3) -- ++(1.1,0) node[below right,inner sep=1pt] {$u_3(P_5)$}; \draw[load] (n3) -- ++(0,1.1) node[above left,inner sep=0pt] {$v_3(P_6)$};
  \draw[load] (n4) -- ++(1.1,0) node[right,inner sep=1pt] {$u_4(P_7)$}; \draw[load] (n4) -- ++(0,1.1) node[above,inner sep=1pt] {$v_4(P_8)$};
  % angles
  \draw[dashed] (n0) -- ++(0,-1.4); \draw[->] ($(n0)+(0,-1.1)$) arc (-90:-60:1.1); \node[below,inner sep=1pt] at ($(n0)+(0.35,-1.15)$) {$30^\circ$};
  \draw[dashed] (n2) -- ++(-1.3,0); \draw[->] ($(n2)+(-1.1,0)$) arc (180:150:1.1); \node[left,inner sep=1pt] at ($(n2)+(-1.0,0.3)$) {$60^\circ$};
  \draw[dashed] (n4) -- ++(1.3,0); \draw[->] ($(n4)+(0.8,0)$) arc (0:60:0.8); \node[right,inner sep=1pt] at ($(n4)+(0.75,0.45)$) {$60^\circ$};
  \draw[dashed] (n3) -- ++(1.6,0); \draw[->] ($(n3)+(1.35,0)$) arc (0:30:1.35); \node[right,inner sep=1pt] at ($(n3)+(1.35,0.28)$) {$30^\circ$};
\end{tikzpicture}

    \caption{Five-bar chain.}\label{fig:leu6}\end{subfigure}
\end{figure}

\subsection*{\texorpdfstring{$N$}{N}-bar chain and an example of five-bar chain}
The above three proposed methods for finding the mechanism for the three-bar chain can be extended to an $N$-bar chain easily, where $N\geq 3$. The number of displacement degrees of freedom for the $N$-bar chain is $(2N-2)$. There will be $N$ equations if the zero axial deformation condition such as Eqs.~\eqref{leu:18} is used. Therefore, there are $(N-2)$ mechanisms since the difference between $(2N-2)$ and $N$ is $(N-2)$. By solving these $N$ equations of zero axial deformation, one can obtain the $(N-2)$ mechanisms.

The graphic (compatible rigid body mode) method can also be used to find the $(N-2)$ mechanisms. The idea is to take a three-bar free body from the $N$-bar chain. Hence, totally there will be $(N-2)$ three-bar free bodies; the first one consists of bars 1, 2, and 3 and the last one consists of bars $N-2$, $N-1$, and $N$. Each three-bar free body can be analyzed using the same method used for the three-bar chain of Fig.~\ref{fig:leu2} to find a mechanism. Note that each of the obtained mechanisms is local since only four displacement components associated with the middle bar are involved.

The local mechanisms can also be obtained conveniently by applying the moment equilibrium condition to each three-bar free body similarly to that was done in Eq.~\eqref{leu:22} or \eqref{leu:25}. Clearly, there are $(N-2)$ centers of rotation because there are $(N-2)$ three-bar free bodies and each has a center of rotation. By applying the condition of moment equilibrium about each center of rotation, one can obtain the corresponding mechanism.

Finally, consider the five-bar chain of Fig.~\ref{fig:leu6} for illustration. According to the above discussion, there are three three-bar free bodies and therefore three mechanisms can be found. By applying any of the three proposed methods, one can obtain the following three normalized mechanisms:
\begin{align}
  \bs{\varphi}_1 &= \frac{1}{\sqrt 8}\langle\sqrt3, 1, 0, -2, 0, 0, 0, 0\rangle^\tr \tag{26a}\\
  \bs{\varphi}_2 &= \frac{1}{\sqrt 8}\langle 0, 0, 1, \sqrt3, 1, -\sqrt3, 0, 0\rangle^\tr \tag{26b}\\
  \bs{\varphi}_3 &= \frac{1}{\sqrt 8}\langle 0, 0, 0, 0, 0, -2, -\sqrt3, 1\rangle^\tr \tag{26c}
\end{align}
where the displacement components in $\bs{\varphi}_1$, $\bs{\varphi}_2$ and $\bs{\varphi}_3$ are listed in the order of $u_1$, $v_1$, $u_2$, $v_2$, $u_3$, $v_3$, $u_4$, and $v_4$. Among the three methods, perhaps the third one is most convenient to use although the three methods obtain the same mechanisms. Points $a$, $b$, and $c$ in Fig.~\ref{fig:leu6} are the rotation centers when the third method is applied. It is also interesting to know that the obtained mechanisms are not affected by the lengths of the bars but are affected by the inclined angles of the bars. This is why the bar lengths are not given in Fig.~\ref{fig:leu6}. This remark is true for an $N$-bar chain as well.

\subsection*{Conclusion}
Three methods for finding the mechanism for a three-bar chain were first presented. The mechanism is useful for determining whether a solution exists for this kind of singular structure. If the external load vector is orthogonal to the mechanism, there will be an infinite number of solutions. Otherwise, there is no solution. This orthogonal condition was shown to be equivalent to a moment equilibrium condition about the center of rotation of the middle bar of the three-bar chain. The advantage of the three proposed methods is that there is no need to form the structural stiffness matrix $\Kmat$. Discussions are then given on how to extend the three methods to find the $(N-2)$ mechanisms of an $N$-bar chain, where $N\geq 3$. Finally, a five-bar chain having three mechanisms was used to demonstrate the extension.

\subsection*{References}
\begin{itemize}[leftmargin=*]
  \item Greenberg, M.D., \emph{Foundations of Applied Mathematics}, Prentice-Hall, Englewood Cliffs, N.J., USA, 1978, pp.~342--348.
  \item McGuire, W., Gallagher, R.~H., and Ziemian, R.D., \emph{Matrix Structural Analysis}, John Wiley and Sons, New York, USA, 2000, p.~53.
\end{itemize}

\chapter{Programming for FRAME21}
\label{ch:frame21}

\section{Introduction}
\topic{Purpose}
Develop \textsc{FRAME21} --- a linear analysis program for framed structures.

\topic{Procedure}
\begin{enumerate}
  \item Input data.
  \item Set up global degrees of freedom (equation numbering).
  \item Assemble external loads into the load vector.
  \item Form the global stiffness matrix and add the equivalent forces (arising from the fixed-end forces) to the load vector.
  \item Solve for the nodal unknowns.
  \item Calculate member displacements and member end forces (force recovery).
\end{enumerate}

\section{Input data}
Functions: \texttt{INPUT} and \texttt{OUTPUT1}; read in and write out the following data.

\begin{enumerate}[leftmargin=*]
  \item \textbf{Structural data}
    \begin{description}[font=\ttfamily,leftmargin=6em,style=sameline,itemsep=0pt]
      \item[TITLE] title of the problem analyzed; e.g.\ Example 5.3 of McGuire et al.\ (2000).
      \item[FUNIT] unit of force such as kN, kips.
      \item[LUNIT] unit of length such as mm, in.
      \item[NNOD] number of nodes
      \item[NBC] number of elements (number of beam columns)
      \item[NMAT] number of sets of material properties
      \item[NSEC] number of sets of sectional properties
      \item[ITP] frame type; 1: beam, 2: 2D truss, 3: 2D frame, 4: grid, 5: 3D truss, 6: 3D frame.
      \item[NCO] number of coordinates per node
      \item[NDN] number of d.o.f.\ per node
      \item[NNE] number of nodes per element
      \item[NDE] number of d.o.f.\ per element
      \item[IFORCE] indicator for inputting the fixed-end forces (FEF). When \texttt{IFORCE}=1, no FEF is needed; when \texttt{IFORCE}=2, FEF needs to be input.
    \end{description}

  \item \textbf{Nodal data}
    \begin{enumerate}[label=(\alph*)]
      \item Coordinates: \texttt{COOR(NCO,NNOD)} --- $x$, $y$, and $z$.
      \item Boundary conditions: \texttt{NFIX(NDN,NNOD)} --- $u$, $v$, $w$, $\theta_x$, $\theta_y$, and $\theta_z$ (3D frame).

        Specify ``$-1$'' if the d.o.f.\ is constrained (or fixed); specify ``0'' if free. Specify the smaller node number for the d.o.f.\ which is the same as the d.o.f.\ of the larger node number if the \emph{double node technique} is used. For example, if nodes 3 and 4 are double nodes (they have the same coordinates) used to model an internal pin for a beam, then the \texttt{NFIX} values are 0, 0 for node 3, and 3, 0 for node 4.
        \begin{fig}
          % Beam with an internal pin modelled by double nodes 3 and 4
\begin{tikzpicture}[scale=1.0,font=\small]
  \draw[member] (0,0) -- (7,0);
  \pic at (0,0) {wallL}; \pic at (7,0) {wallR};
  \pic at (2.3,0) {roller};
  \node[node dot] at (4.6,0) {};
  \node[above right,inner sep=1pt] at (0,0) {1};
  \node[above right,inner sep=1pt] at (2.3,0.02) {2};
  \node[above left,inner sep=1pt] at (4.5,0.05) {3};
  \node[above right,inner sep=1pt] at (4.7,0.05) {4};
  \node[above left,inner sep=1pt] at (7,0) {5};
  % rotations
  \foreach \x/\l in {2.3/2,4.35/3,4.95/4}{
    \draw[->,thick] ($(\x,0)+(0.3,0.35)$) arc (20:160:0.3);
    \node[above] at (\x,0.6) {$\theta_\l$};}
  \draw[->,thick] (4.6,-0.9) -- (4.6,-0.15) node[pos=0,right] {$v$};
\end{tikzpicture}

        \end{fig}
        \begin{center}
          \begin{minipage}{0.47\linewidth}\centering
            \texttt{NFIX(NDN,NNOD)}\\[4pt]
            \begin{tabular}{|c|c|c|c|c|c|}\hline
                       & 1    & 2    & 3 & 4 & 5   \\\hline
              $v$      & $-1$ & $-1$ & 0 & 3 & $-1$\\\hline
              $\theta$ & $-1$ & 0    & 0 & 0 & $-1$\\\hline
            \end{tabular}
          \end{minipage}
          \begin{minipage}{0.5\linewidth}\centering
            \texttt{IDND(NDN,NNOD)} (discussed later)\\[4pt]
            \begin{tabular}{|c|c|c|c|c|c|}\hline
                       & 1    & 2    & 3 & 4 & 5   \\\hline
              $v$      & $-1$ & $-3$ & 2 & 2 & $-4$\\\hline
              $\theta$ & $-2$ & 1    & 3 & 4 & $-5$\\\hline
            \end{tabular}
          \end{minipage}
        \end{center}
      \item External ``concentrated'' nodal loads: \texttt{EXLD(NDN,NNOD)}
        \begin{fig}
          % Portal frame with concentrated nodal loads
\begin{tikzpicture}[scale=1.0,font=\small]
  \draw[member] (0,0) -- (0,3) -- (3.2,3) -- (3.2,0);
  \pic at (0,0) {fixed}; \pic at (3.2,0) {pin};
  \node[right,inner sep=2pt] at (0,0.3) {1}; \node[below right,inner sep=2pt] at (0,3) {2};
  \node[below left,inner sep=2pt] at (3.2,3) {3}; \node[above left,inner sep=2pt] at (3.2,0.05) {4};
  \draw[load] (-1.0,3) -- (-0.1,3) node[pos=0,left] {2};
  \draw[load] (0,4.0) -- (0,3.1) node[pos=0,above] {10};
  \draw[load] (3.2,3.1) -- (3.2,4.0) node[above] {5};
  \draw[->,thick] ($(3.2,3)+(0.45,0.2)$) arc (30:-170:0.45); \node[right] at (3.75,2.95) {3};
  \draw[load] (4.2,0) -- (3.4,0) node[pos=0,right] {7};
\end{tikzpicture}

        \end{fig}
        \begin{center}
          \begin{minipage}{0.47\linewidth}\centering
            \texttt{NFIX(NDN,NNOD)}\\[4pt]
            \begin{tabular}{|c|c|c|c|c|}\hline
                       & 1    & 2 & 3 & 4   \\\hline
              $u$      & $-1$ & 0 & 0 & 0   \\\hline
              $v$      & $-1$ & 0 & 0 & $-1$\\\hline
              $\theta$ & $-1$ & 0 & 0 & 0   \\\hline
            \end{tabular}
          \end{minipage}
          \begin{minipage}{0.5\linewidth}\centering
            \texttt{EXLD(NDN,NNOD)}\\[4pt]
            \begin{tabular}{|c|c|c|c|c|}\hline
                    & 1 & 2     & 3    & 4   \\\hline
              $F_x$ & 0 & 2     & 0    & $-7$\\\hline
              $F_y$ & 0 & $-10$ & 5    & 0   \\\hline
              $M$   & 0 & 0     & $-3$ & 0   \\\hline
            \end{tabular}
          \end{minipage}
        \end{center}
    \end{enumerate}

  \item \textbf{Elemental data}
    \begin{enumerate}[label=(\alph*)]
      \item Identification matrices (connectivities):\\
        \texttt{IDBC(5,NBC)}: node 1, node 2, material type, sectional type, blank.
      \item Direction cosines of the local $y$ axis: \texttt{VECTY(3,NBC)}. (\texttt{ITP}=6 only)
        \begin{fig}
          % Global (x,y,z) and local (x',y',z') axes of a 3D frame member
\begin{tikzpicture}[scale=1.0,font=\small]
  \begin{scope}
    \draw[->,thick] (0,0) -- (0,2.0) node[above] {$z$};
    \draw[->,thick] (0,0) -- (-1.5,-0.9) node[left] {$x$};
    \draw[->,thick] (0,0) -- (1.5,-0.9) node[right] {$y$};
  \end{scope}
  \begin{scope}[xshift=6cm]
    % a small box (member cross-section sketch)
    \draw (0,0) -- (0.5,0) -- (0.5,1.3) -- (0,1.3) -- cycle;
    \draw (0,1.3) -- (-0.4,1.55) -- (0.1,1.55) -- (0.5,1.3);
    \draw (0,0) -- (-0.4,0.25) -- (-0.4,1.55);
    \draw (-0.4,1.55) -- (-1.4,2.1); \draw (0.1,1.55) -- (-0.9,2.1);
    \draw[->,thick] (0,0) -- (0,2.4) node[above] {$y'$};
    \draw[->,thick] (0,0) -- (-1.6,-0.7) node[left] {$z'$};
    \draw[->,thick] (0,0) -- (1.6,-0.7) node[right] {$x'$};
    \node[below,font=\footnotesize] at (0,-1.1) {$0,0,1$ for this example};
  \end{scope}
\end{tikzpicture}

        \end{fig}
      \item Fixed-end forces (in local axis) \texttt{FEF(NDE,NBC)} $= \F_l$. (\texttt{IFORCE}=2 only)
    \end{enumerate}

  \item \textbf{Material and sectional properties}
    \begin{enumerate}[label=(\alph*)]
      \item \texttt{PROP(5,NMAT)}: Young's modulus $E$, Poisson's ratio $\nu$, blank, blank, blank.\\
        For dynamic problems, $\rho$ is required; for plastic analysis problems, $\sigma_y$ is needed.
      \item \texttt{SECT(5,NSEC)}: cross-sectional area $A$, moments of inertia $I_z$ and $I_y$, and torsional constant $J$, blank.
    \end{enumerate}
\end{enumerate}

\subsection*{Example of input files}
\begin{fig}
  % 2-D truss example for the input file (8 nodes, 14 members)
\begin{tikzpicture}[scale=0.55,font=\small]
  \coordinate (n1) at (0,0); \coordinate (n2) at (6,0); \coordinate (n3) at (12,0); \coordinate (n4) at (18,0); \coordinate (n5) at (24,0);
  \coordinate (n6) at (6,8); \coordinate (n7) at (12,8); \coordinate (n8) at (18,8);
  \draw[thinmember] (n1) -- (n5); \draw[thinmember] (n6) -- (n8);
  \draw[thinmember] (n1) -- (n6) -- (n2); \draw[thinmember] (n6) -- (n3) -- (n7); \draw[thinmember] (n3) -- (n8) -- (n4);
  \draw[thinmember] (n4) -- (n7); \draw[thinmember] (n5) -- (n8);
  \foreach \p in {n1,...,n8}{\node[node dot] at (\p) {};}
  \pic at (n1) {pin}; \pic at (n5) {roller};
  \node[above right] at (n1) {1}; \node[above right] at (n2) {2}; \node[below] at (n3) {3}; \node[below] at (n4) {4}; \node[above left] at (n5) {5};
  \node[above] at (n6) {6}; \node[above] at (n7) {7}; \node[above] at (n8) {8};
  \newcommand{\mem}[2]{\node[circle,draw,fill=white,inner sep=1pt,font=\scriptsize] at #1 {#2};}
  \mem{(3,0)}{1} \mem{(9,0)}{2} \mem{(15,0)}{3} \mem{(21,0)}{4}
  \mem{(3,4)}{5} \mem{(6,4)}{6} \mem{(9,4)}{7} \mem{(12,4)}{8}
  \mem{(13.7,5.7)}{9} \mem{(16.3,5.7)}{10} \mem{(18,4)}{11} \mem{(21,4)}{12}
  \mem{(9,8)}{13} \mem{(15,8)}{14}
  \draw[load,very thick] (n2) -- ++(0,-1.6) node[right] {64\,kN};
  \draw[dashed] (n6) -- (-1.0,8);
  \draw[dim] (-1.0,0) -- (-1.0,8) node[midway,fill=white,inner sep=1pt,rotate=90] {32\,m};
  \draw[dim] (0,-2.3) -- (24,-2.3) node[midway,fill=white,inner sep=1pt] {4@24\,m $=$ 96\,m};
  \draw (0,-2.6) -- (0,-2.0); \draw (24,-2.6) -- (24,-2.0);
\end{tikzpicture}

\end{fig}
\begin{itemize}
  \item $E = 300~\mathrm{kN/cm^2} = 3{,}000{,}000~\mathrm{kN/m^2}$ for members 1, 2, 3, 4, 13, and 14
  \item $E = 200~\mathrm{kN/cm^2} = 2{,}000{,}000~\mathrm{kN/m^2}$ for members 5, 6, 7, 8, 9, 10, 11, 12
  \item Cross section area $= 0.10~\mathrm{m^2}$ for members 1, 2, 3, 4, 13, and 14
  \item Cross section area $= 0.20~\mathrm{m^2}$ for members 6, 8, and 11
  \item Cross section area $= 0.30~\mathrm{m^2}$ for members 5, 7, 9, 10 and 12
\end{itemize}
The contents of \texttt{2Dtruss.ipt}:
\begin{lstlisting}[language={},numbers=none,basicstyle=\small\ttfamily]
2-D truss example
kN
m

* NNOD  NBC  NMAT  NSEC  ITP  NNE  IFORCE
   8    14    2     3     2    2     1

* COOR (m)
1    0.   0.
2   24.   0.
3   48.   0.
4   72.   0.
5   96.   0.
6   24.  32.
7   48.  32.
8   72.  32.

* NFIX (u, v)
1   -1  -1
2    0   0
3    0   0
4    0   0
5   -1  -1
6    0   0
7    0   0
8    0   0

* External Load (kN): Px, Py
1   0.    0.
2   0.  -64.
3   0.    0.
4   0.    0.
5   0.    0.
6   0.    0.
7   0.    0.
8   0.    0.

* IDBC
1   1 2 1 1 0
2   2 3 1 1 0
3   3 4 1 1 0
4   4 5 1 1 0
5   1 6 2 3 0
6   2 6 2 2 0
7   3 6 2 3 0
8   3 7 2 2 0
9   4 7 2 3 0
10  3 8 2 3 0
11  4 8 2 2 0
12  5 8 2 3 0
13  6 7 1 1 0
14  7 8 1 1 0

* PROP
1   3000000. 0. 0. 0. 0.
2   2000000. 0. 0. 0. 0.

* SECT
1   0.10 0. 0. 0. 0.
2   0.20 0. 0. 0. 0.
3   0.30 0. 0. 0. 0.
\end{lstlisting}

\section{Set up global degrees of freedom (equation numbering)}
Functions: \texttt{IDMAT}, \texttt{MEMDOF} and \texttt{SEMIBAND}.
\begin{enumerate}[leftmargin=*]
  \item \texttt{IDMAT}: Equation numbering --- transform \texttt{NFIX(NDN,NNOD)} into \texttt{IDND(NDN,NNOD)} and obtain \texttt{NEQ}.

    Example: 2-D truss (\texttt{NDN}=2)
    \begin{fig}
      % 2-D truss example for IDMAT (6 nodes, 9 members)
\begin{tikzpicture}[scale=0.8,font=\small]
  \coordinate (n1) at (0,0); \coordinate (n2) at (0,3); \coordinate (n3) at (3.5,0); \coordinate (n4) at (3.5,3);
  \coordinate (n5) at (7,0); \coordinate (n6) at (7,3);
  \draw[member] (n1) -- (n2) -- (n4) -- (n3) -- (n1); \draw[member] (n2) -- (n3); \draw[member] (n3) -- (n5) -- (n6) -- (n4); \draw[member] (n3) -- (n6);
  \foreach \p in {n1,...,n6}{\node[node dot] at (\p) {};}
  \pic[rotate=90] at (n2) {pin}; \pic[rotate=90] at (n1) {roller};
  \node[below right] at (n1) {1}; \node[above right] at (n2) {2}; \node[below] at (n3) {3}; \node[above] at (n4) {4}; \node[below] at (n5) {5}; \node[above] at (n6) {6};
  \newcommand{\mem}[2]{\node[circle,draw,fill=white,inner sep=1.5pt,font=\scriptsize] at #1 {#2};}
  \mem{(-0.5,1.5)}{1} \mem{(1.75,-0.5)}{2} \mem{(2.3,1.8)}{3} \mem{(1.75,3.5)}{4}
  \mem{(4.0,1.5)}{5} \mem{(5.25,-0.5)}{6} \mem{(5.8,1.2)}{7} \mem{(5.25,3.5)}{8} \mem{(7.5,1.5)}{9}
\end{tikzpicture}

    \end{fig}
    \begin{center}
      \begin{minipage}{0.47\linewidth}\centering
        \texttt{NFIX(NDN,NNOD)}\\[4pt]
        \begin{tabular}{|c|c|c|c|c|c|c|}\hline
              & 1    & 2    & 3 & 4 & 5 & 6\\\hline
          $u$ & $-1$ & $-1$ & 0 & 0 & 0 & 0\\\hline
          $v$ & 0    & $-1$ & 0 & 0 & 0 & 0\\\hline
        \end{tabular}
      \end{minipage}
      \begin{minipage}{0.47\linewidth}\centering
        \texttt{IDND(NDN,NNOD)}\\[4pt]
        \begin{tabular}{|c|c|c|c|c|c|c|}\hline
              & 1    & 2    & 3 & 4 & 5 & 6\\\hline
          $u$ & $-1$ & $-2$ & 2 & 4 & 6 & 8\\\hline
          $v$ & 1    & $-3$ & 3 & 5 & 7 & 9\\\hline
        \end{tabular}
      \end{minipage}
    \end{center}

  \item \texttt{MEMDOF}: create member d.o.f.\ table, \texttt{LM(NDE,NBC)}.
    \begin{itemize}
      \item Use connectivities \texttt{IDBC(1,IB)} and \texttt{IDBC(2,IB)} (the \texttt{IB}th element) and the nodal d.o.f.\ table, \texttt{IDND(NDN,NNOD)}, to create the member d.o.f.\ table, \texttt{LM(NDE,NBC)}.
    \end{itemize}
    \begin{center}
      \texttt{LM(NDE,NBC)}\\[4pt]
      \begin{tabular}{|c|c|c|c|c|c|c|c|c|c|}\hline
                 & 1    & 2    & 3 & 4 & 5 & 6 & 7 & 8 & 9 (\texttt{NBC})\\\hline
        1        & $-1$ & $-1$ &   &   &   &   & 2 &   & \\\hline
        2        & 1    & 1    &   &   &   &   & 3 &   & \\\hline
        3        & $-2$ & 2    &   &   &   &   & 8 &   & \\\hline
        4 (\texttt{NDE}) & $-3$ & 3 &   &   &   &   & 9 &   & \\\hline
      \end{tabular}
    \end{center}

  \item \texttt{SEMIBAND}: calculate the semi-band width (including diagonal) \texttt{NSBAND}. $[\Kmat]$ can be stored in banded form.
    \begin{center}
      \begin{minipage}{0.5\linewidth}\centering
        $\Kmat$ (\texttt{NEQ},\texttt{NEQ})\\[4pt]
        \renewcommand{\arraystretch}{1.3}
        \begin{tabular}{|c|c|c|c|c|c|}\hline
          $K_{11}$ & $K_{12}$ & $K_{13}$ &          &          &         \\\hline
          $K_{12}$ & $K_{22}$ & $K_{23}$ &          &          &         \\\hline
          $K_{13}$ & $K_{23}$ & $K_{33}$ & $K_{34}$ &          &         \\\hline
                   &          & $K_{34}$ & $K_{44}$ & $K_{45}$ & $K_{46}$\\\hline
                   &          &          & $K_{45}$ & $K_{55}$ & $K_{56}$\\\hline
                   &          &          & $K_{46}$ & $K_{56}$ & $K_{66}$\\\hline
        \end{tabular}
      \end{minipage}
      \begin{minipage}{0.42\linewidth}\centering
        $\widetilde{\Kmat}$ (\texttt{NEQ},\texttt{NSBAND})\\[4pt]
        \renewcommand{\arraystretch}{1.3}
        \begin{tabular}{|c|c|c|}\hline
          $K_{11}$ & $K_{12}$ & $K_{13}$\\\hline
          $K_{22}$ & $K_{23}$ &         \\\hline
          $K_{33}$ & $K_{34}$ &         \\\hline
          $K_{44}$ & $K_{45}$ & $K_{46}$\\\hline
          $K_{55}$ & $K_{56}$ &         \\\hline
          $K_{66}$ &          &         \\\hline
        \end{tabular}
      \end{minipage}
    \end{center}
    $K(i,j)$, $j\geq i$ is stored in $\widetilde{K}(i,\,j-i+1)$; for example,
    \begin{align*}
      K(2,3) &\;\rightarrow\; \widetilde{K}(2,2)\\
      K(4,6) &\;\rightarrow\; \widetilde{K}(4,3)
    \end{align*}
    Calculate \texttt{NSBAND} from the \texttt{LM(NDE,NBC)} matrix columnwise using the algorithm below:
    \begin{equation}
      \texttt{NSBAND} = \max_{i=1,\texttt{NBC}}\Bigl(\max_{j=1,\texttt{NDE}}LM(j,i) - \min_{j=1,\texttt{NDE}}LM(j,i) + 1\Bigr)\quad\text{for } LM(j,i) > 0
    \end{equation}
\end{enumerate}

\section{Assemble external loads into the load vector}
Function \texttt{LOAD}: assemble \texttt{EXLD(NDN,NNOD)} into \texttt{GLOAD(NEQ)}.

\topic{Algorithm:}
\begin{lstlisting}[language=Matlab,numbers=none]
GLOAD = zeros(NEQ,1);
for j = 1:NNOD
    for i = 1:NDN
        if IDND(i,j) > 0
            GLOAD(IDND(i,j)) = GLOAD(IDND(i,j))+EXLD(i,j);
        end
    end
end
\end{lstlisting}

\section{Form the global stiffness matrix and add the equivalent member forces to the load vector}
Functions: \texttt{FORMKP} (\texttt{ROTATION}, \texttt{ELKE}).

For each element, perform the following steps to add its stiffness matrix to the banded global stiffness matrix.
\begin{enumerate}
  \item Calculate element transformation (rotation) matrix $\Tmat_{\texttt{NDE}\times\texttt{NDE}}$. (function \texttt{ROTATION})
  \item Calculate element stiffness matrix referred to the local (elemental) coordinate system, $[k_l]_{\texttt{NDE}\times\texttt{NDE}} = \kmat_l$. (function \texttt{ELKE})
  \item Calculate element stiffness matrix $\kmat$ referred to the global (structural) coordinate system using $\kmat = \Tmat^\tr\kmat_l\Tmat$. ($\kmat_l =$ matrix \texttt{EE}, $\kmat =$ matrix \texttt{ELK})
  \item Assemble $\kmat_{\texttt{NDE}\times\texttt{NDE}}$ into the stiffness matrix $\Kmat_{\texttt{NEQ}\times\texttt{NEQ}}$.
\end{enumerate}
Assembling process: $k_{mn}\;\rightarrow\;K_{IJ}$.

\topic{Algorithm (for the \texttt{IB}th element):}
\begin{lstlisting}[language=Matlab,numbers=none]
% Get element DOFs
LDOF = find(LM(:,IB)>0);
GDOF = LM(LDOF,IB);
% Transform the element stiffness matrix from the local axes
% to global axes : EE --> ELK
ELK = T'*EE*T;
% Assemble the global element stiffness matrix to form the global
% stiffness matrix GLK
GLK(GDOF,GDOF) = GLK(GDOF,GDOF) + ELK(LDOF,LDOF);
\end{lstlisting}

For each element, perform the following three steps to add its equivalent member forces to the load vector.
\begin{enumerate}
  \item Calculate element transformation matrix $\Tmat$. (function \texttt{ROTATION})
  \item Calculate equivalent member forces using $\F_g = -\Tmat^\tr\F_l$, where $\F_l$ (\texttt{FEF(NDE,NBC)}) is the fixed-end forces; $\F_g$ is referred to as equivalent member forces.
    \begin{fig}
      % Fixed-end forces (local) -> equivalent member forces (global)
\begin{tikzpicture}[scale=0.9,font=\footnotesize]
  \def\ang{45}
  % ---- F_l ----
  \begin{scope}
    \coordinate (A) at (0,0); \coordinate (B) at (\ang:3.2);
    \draw[member] (A) -- (B);
    % distributed load
    \foreach \t in {0.1,0.25,...,0.95}{\draw[->] ($(A)!\t!(B)+(\ang+90:0.7)$) -- ($(A)!\t!(B)+(\ang+90:0.1)$);}
    \draw[load] ($(A)+(\ang+90:0.9)$) -- ($(A)+(\ang+90:0.15)$) node[pos=0,left] {10};
    \draw[load] ($(B)+(\ang+90:0.9)$) -- ($(B)+(\ang+90:0.15)$) node[pos=0,right] {$-10$};
    \draw[->,thick] ($(A)+(-0.35,-0.3)$) arc (-120:60:0.35); \node[left] at ($(A)+(-0.3,-0.2)$) {200};
    \draw[->,thick] ($(B)+(0.35,0.3)$) arc (60:240:0.35); \node[right] at ($(B)+(0.25,0.35)$) {$-200$};
    \node at ($(A)+(1.6,-0.7)$) {$\F_l$};
  \end{scope}
  \node at (4.4,2.2) {$\Longrightarrow$};
  % ---- -F_l ----
  \begin{scope}[xshift=5.6cm]
    \coordinate (A) at (0,0); \coordinate (B) at (\ang:3.2);
    \draw[member] (A) -- (B);
    \draw[load] ($(A)+(\ang+90:0.15)$) -- ($(A)+(\ang+90:0.9)$) node[left] {$-10$};
    \draw[load] ($(B)+(\ang+90:0.15)$) -- ($(B)+(\ang+90:0.9)$) node[right] {10};
    \draw[->,thick] ($(A)+(0.35,-0.3)$) arc (-60:-240:0.35); \node[left] at ($(A)+(-0.3,-0.2)$) {$-200$};
    \draw[->,thick] ($(B)+(-0.35,0.3)$) arc (120:-60:0.35); \node[right] at ($(B)+(0.25,0.35)$) {200};
    \node at ($(A)+(1.6,-0.7)$) {$-\F_l$};
  \end{scope}
  \node at (10.2,2.2) {$\Longrightarrow$};
  % ---- -T^T F_l (global components) ----
  \begin{scope}[xshift=11.6cm]
    \coordinate (A) at (0,0); \coordinate (B) at (\ang:3.2);
    \draw[member] (A) -- (B);
    \draw[load] (A) -- ++(0,0.9) node[left] {$-10\frac{1}{\sqrt2}$};
    \draw[load] (A) -- ++(0.9,0) node[right] {$10\frac{1}{\sqrt2}$};
    \draw[load] (B) -- ++(0,0.9) node[left] {$-10\frac{1}{\sqrt2}$};
    \draw[load] (B) -- ++(0.9,0) node[below right] {$10\frac{1}{\sqrt2}$};
    \draw[->,thick] ($(A)+(0.35,-0.3)$) arc (-60:-240:0.35); \node[left] at ($(A)+(-0.35,-0.35)$) {$-200$};
    \draw[->,thick] ($(B)+(-0.35,0.3)$) arc (120:-60:0.35); \node[right] at ($(B)+(0.3,0.45)$) {200};
    \node at ($(A)+(1.6,-1.1)$) {$-\Tmat^\tr\F_l$};
  \end{scope}
\end{tikzpicture}

    \end{fig}
  \item Assemble $\F_g$ into \texttt{GLOAD(NEQ)} using the algorithm below (for the \texttt{IB}th element).
\end{enumerate}
\begin{lstlisting}[language=Matlab,numbers=none]
Steps 2 and 3:
EFEQ = -T'*FEF(:,IB);
GLOAD(GDOF) = GLOAD(GDOF) + EFEQ(LDOF);
\end{lstlisting}

\section{Solve for the nodal unknowns}
(Chapter~\ref{ch:solve}; function \texttt{SOLVE})

\section{Calculate member end forces (force recovery)}
(function \texttt{FORCE})

For each element (member), the member end forces can be calculated as follows:
\begin{enumerate}
  \item Use \texttt{LM(NDE,IB)} (the \texttt{IB}th element) and the global displacement vector \texttt{DISP(NEQ)} to find the element displacement vector in the global coordinate system, $\dvec_g$.
  \item Call function \texttt{ROTATION} to obtain the element transformation matrix $\Tmat$.
  \item Obtain the element displacement vector in the local coordinate system, $\dvec_l$ (dimension: $\texttt{NDE}\times\texttt{NDE}$) using $\dvec_l = \Tmat\dvec_g$.
  \item Call function \texttt{ELKE} to obtain the element stiffness matrix $\kmat_l$.
  \item Calculate the element forces (member end forces) using $\mat{f}_l = \kmat_l\dvec_l + \F_l$.
\end{enumerate}

\chapter[3D Beam-Column Element]{3D Beam-Column Element: Strength of Materials Approach}
\label{ch:beamcolumn}

\section{Stiffness matrix of a 3D beam-column element with bisymmetrical section}
{\footnotesize
\begin{equation}
  \label{eq:k3d}
  \Bmat{F_{x1}\\ F_{y1}\\ F_{z1}\\ M_{x1}\\ M_{y1}\\ M_{z1}\\ F_{x2}\\ F_{y2}\\ F_{z2}\\ M_{x2}\\ M_{y2}\\ M_{z2}}
  = E\,
  \setlength{\arraycolsep}{2.2pt}
  \left[\begin{array}{cccccc|cccccc}
    \frac{A}{L} & 0 & 0 & 0 & 0 & 0 & -\frac{A}{L} & 0 & 0 & 0 & 0 & 0\\[3pt]
    0 & \frac{12I_z}{L^3} & 0 & 0 & 0 & \frac{6I_z}{L^2} & 0 & -\frac{12I_z}{L^3} & 0 & 0 & 0 & \frac{6I_z}{L^2}\\[3pt]
    0 & 0 & \frac{12I_y}{L^3} & 0 & -\frac{6I_y}{L^2} & 0 & 0 & 0 & -\frac{12I_y}{L^3} & 0 & -\frac{6I_y}{L^2} & 0\\[3pt]
    0 & 0 & 0 & \frac{J}{2(1+\nu)L} & 0 & 0 & 0 & 0 & 0 & -\frac{J}{2(1+\nu)L} & 0 & 0\\[3pt]
    0 & 0 & -\frac{6I_y}{L^2} & 0 & \frac{4I_y}{L} & 0 & 0 & 0 & \frac{6I_y}{L^2} & 0 & \frac{2I_y}{L} & 0\\[3pt]
    0 & \frac{6I_z}{L^2} & 0 & 0 & 0 & \frac{4I_z}{L} & 0 & -\frac{6I_z}{L^2} & 0 & 0 & 0 & \frac{2I_z}{L}\\[3pt]\hline
    -\frac{A}{L} & 0 & 0 & 0 & 0 & 0 & \frac{A}{L} & 0 & 0 & 0 & 0 & 0\\[3pt]
    0 & -\frac{12I_z}{L^3} & 0 & 0 & 0 & -\frac{6I_z}{L^2} & 0 & \frac{12I_z}{L^3} & 0 & 0 & 0 & -\frac{6I_z}{L^2}\\[3pt]
    0 & 0 & -\frac{12I_y}{L^3} & 0 & \frac{6I_y}{L^2} & 0 & 0 & 0 & \frac{12I_y}{L^3} & 0 & \frac{6I_y}{L^2} & 0\\[3pt]
    0 & 0 & 0 & -\frac{J}{2(1+\nu)L} & 0 & 0 & 0 & 0 & 0 & \frac{J}{2(1+\nu)L} & 0 & 0\\[3pt]
    0 & 0 & -\frac{6I_y}{L^2} & 0 & \frac{2I_y}{L} & 0 & 0 & 0 & \frac{6I_y}{L^2} & 0 & \frac{4I_y}{L} & 0\\[3pt]
    0 & \frac{6I_z}{L^2} & 0 & 0 & 0 & \frac{2I_z}{L} & 0 & -\frac{6I_z}{L^2} & 0 & 0 & 0 & \frac{4I_z}{L}
  \end{array}\right]
  \Bmat{u_1\\ v_1\\ w_1\\ \theta_{x1}\\ \theta_{y1}\\ \theta_{z1}\\ u_2\\ v_2\\ w_2\\ \theta_{x2}\\ \theta_{y2}\\ \theta_{z2}}
\end{equation}}

The above stiffness matrix contains four parts, one arising from axial deformation, two from bending deformation, and the other from torsional deformation. We shall derive the stiffness matrices for these four parts by solving the corresponding differential equations.
\begin{fig}
  % 3D beam-column element: end forces and displacements (12 DOFs)
\begin{tikzpicture}[scale=1.0,font=\footnotesize,x={(1cm,0cm)},y={(0cm,1cm)},z={(-0.7cm,-0.5cm)}]
  \def\L{7} \def\h{0.45} \def\b{0.9}
  % member as a box
  \coordinate (n1) at (0,0,0); \coordinate (n2) at (\L,0,0);
  \draw (0,\h,\b) -- (\L,\h,\b) -- (\L,\h,-\b) -- (0,\h,-\b) -- cycle;      % top face
  \draw (0,-\h,\b) -- (\L,-\h,\b); \draw (0,-\h,-\b) -- (\L,-\h,-\b);
  \draw (0,\h,\b) -- (0,-\h,\b); \draw (0,\h,-\b) -- (0,-\h,-\b); \draw (0,-\h,\b) -- (0,-\h,-\b);
  \draw (\L,\h,\b) -- (\L,-\h,\b); \draw (\L,\h,-\b) -- (\L,-\h,-\b); \draw (\L,-\h,\b) -- (\L,-\h,-\b);
  \fill (n1) circle (0.06); \fill (n2) circle (0.06);
  \node[below right] at (n1) {1}; \node[above left] at (n2) {2};
  % --- node 1 ---
  \draw[load] (n1) -- ++(0,1.6,0) node[right] {$F_{y1},v_1$};
  \draw[load,double] (0,1.75,0) -- (0,2.9,0) node[right] {$M_{y1},\theta_{y1}$};
  \draw[load] (-1.4,0,0) -- (-0.1,0,0) node[pos=0,above] {$F_{x1},u_1$};
  \draw[load,double] (-3.2,0,0) -- (-1.7,0,0) node[pos=0,above] {$M_{x1},\theta_{x1}$};
  \draw[load] (n1) -- ++(0,0,1.5) node[below left,inner sep=1pt] {$F_{z1},w_1$};
  \draw[load,double] (0,0,1.9) -- (0,0,3.1) node[below left,inner sep=1pt] {$M_{z1},\theta_{z1}$};
  % --- node 2 ---
  \draw[load] (n2) -- ++(0,1.6,0) node[right] {$F_{y2},v_2$};
  \draw[load,double] (\L,1.75,0) -- (\L,2.9,0) node[right] {$M_{y2},\theta_{y2}$};
  \draw[load] (\L+0.1,0,0) -- (\L+1.4,0,0) node[above] {$F_{x2},u_2$};
  \draw[load,double] (\L+1.7,0,0) -- (\L+3.2,0,0) node[above] {$M_{x2},\theta_{x2}$};
  \draw[load] (n2) -- ++(0,0,1.5) node[below right,inner sep=1pt] {$F_{z2},w_2$};
  \draw[load,double] (\L,0,1.9) -- (\L,0,3.1) node[below right,inner sep=1pt] {$M_{z2},\theta_{z2}$};
  % length
  \draw[dim] (0,-2.4,0) -- (\L,-2.4,0) node[midway,fill=white,inner sep=1pt] {$L$};
  \draw (0,-2.2,0) -- (0,-2.6,0); \draw (\L,-2.2,0) -- (\L,-2.6,0);
\end{tikzpicture}

\end{fig}

\section{Axial stiffness matrix}
\begin{fig}
  % Axial bar with distributed load b(x) and a differential element
\begin{tikzpicture}[scale=1.0,font=\small]
  \begin{scope}
    \draw[thick] (0,0) rectangle (5,0.4);
    \foreach \x in {0.2,1.0,...,4.8}{\draw[->] (\x,0.75) -- (\x+0.55,0.75);}
    \node[above] at (2.5,0.85) {$b(x)$};
    \node[align=center,below] at (2.5,-0.1) {Appropriate boundary conditions\\at ends (not shown)};
  \end{scope}
  \begin{scope}[xshift=8cm]
    \draw[thick] (0,0) rectangle (1.0,0.4);
    \draw[->] (0.2,0.75) -- (0.8,0.75); \node[above] at (0.5,0.85) {$b(x)dx$};
    \draw[load] (0,0.2) -- (-0.9,0.2) node[left] {$EAu'$};
    \draw[load] (1.0,0.2) -- (1.9,0.2) node[right] {$EAu'+EAu''dx$};
    \draw[->] (-0.6,-0.4) -- (0,-0.4); \draw[->] (1.6,-0.4) -- (1.0,-0.4);
    \node[below] at (0.5,-0.25) {$dx$};
    \draw (0,-0.6) -- (0,-0.2); \draw (1,-0.6) -- (1,-0.2);
  \end{scope}
\end{tikzpicture}

\end{fig}
\begin{equation*}
  \textstyle\sum F_x = 0:\qquad EAu''\,dx + b(x)\,dx = 0
\end{equation*}
If $b(x) = 0$ then $EAu'' = 0$.
\begin{align*}
  \text{Governing equation (G.E.):}&\quad EAu''(x) = 0\quad(\text{assume } EA = \text{constant}),\ 0 < x < L\\
  \text{Boundary conditions (BCs):}&\quad u(0) = u_1\ \text{ and }\ u(L) = u_2
\end{align*}
Integrating the G.E. twice, we obtain
\begin{equation*}
  u(x) = c_1x + c_2.
\end{equation*}
Using the BCs, we obtain
\begin{equation*}
  c_2 = u_1,\qquad c_1L + c_2 = u_2.
\end{equation*}
Solving the above two equations yields
\begin{equation*}
  c_1 = (u_2-u_1)/L,\qquad c_2 = u_1.
\end{equation*}
Therefore,
\begin{equation}
  u(x) = c_1x + c_2 = u_1\Bigl(1-\frac{x}{L}\Bigr) + u_2\frac{x}{L} = N_1(x)u_1 + N_2(x)u_2
\end{equation}
where
\begin{equation}
  N_1(x) = 1-\frac{x}{L},\qquad N_2(x) = \frac{x}{L}.
\end{equation}
Finally,
\begin{align}
  F_{x1} &= -EAu'(0) = \frac{EA}{L}(u_1-u_2)\\
  F_{x2} &= EAu'(L) = \frac{EA}{L}(-u_1+u_2)
\end{align}
That is,
\begin{equation}
  \Bmat{F_{x1}\\ F_{x2}} = \frac{EA}{L}\bmat{1 & -1\\ -1 & 1}\Bmat{u_1\\ u_2}.
\end{equation}

\topic{Energy method:}
We can also derive the stiffness matrix from the strain energy. Since $u(x) = u_1(1-\frac{x}{L}) + u_2\frac{x}{L} = \NN\bs{\Delta}$, where $\NN(x) = \langle N_1, N_2\rangle$ and $\bs{\Delta} = \langle u_1, u_2\rangle^\tr$, the strain energy can be written as
\begin{equation}
  U = \frac12\int_0^L EA(u')^2\,dx = \frac12 EA\int_0^L (u')^\tr u'\,dx
    = \frac12 EA\int_0^L \bs{\Delta}^\tr\NN'^{\tr}\NN'\bs{\Delta}\,dx = \frac12\bs{\Delta}^\tr\kmat\bs{\Delta}.
\end{equation}
Therefore,
\begin{equation}
  \kmat = EA\int_0^L \NN'^{\tr}\NN'\,dx.
\end{equation}

\section{Flexural stiffness matrix for bending w.r.t.\ the \texorpdfstring{$z$}{z}-axis}
$\theta_z(x) = v'(x)$
\begin{align*}
  \text{Governing equation (G.E.):}&\quad EI_z v^{(iv)}(x) = 0\quad(\text{assume } EI_z = \text{constant}),\ 0 < x < L\\
  \text{Boundary conditions (BCs):}&\quad v(0) = v_1,\ v(L) = v_2,\ v'(0) = \theta_{z1},\ \text{and } v'(L) = \theta_{z2}.
\end{align*}
\begin{fig}
  % Bending about the z-axis: v(x) and v'(x) > 0
\begin{tikzpicture}[scale=1.0,font=\small]
  \draw[->,thick] (0,0) -- (4.2,0) node[right] {$x$};
  \draw[->,thick] (0,0) -- (0,2.8) node[above] {$y$};
  \draw[->,thick] (0,0) -- (-1.6,-1.6) node[below left] {$z$};
  % undeformed and deformed beam
  \draw[fill=black!10] (0,-0.12) rectangle (3.4,0.12);
  \draw[fill=black!10,rotate=12] (0,0.1) rectangle (3.4,0.32);
  \draw[->] (3.2,0.1) -- (3.2,0.66); \node[right] at (3.25,0.4) {$v(x)$};
  \node[right] at (3.3,1.1) {$v'(x)>0$};
\end{tikzpicture}

\end{fig}
Integrating the G.E. four times, we obtain
\begin{equation}
  v(x) = \frac16 c_1x^3 + \frac12 c_2x^2 + c_3x + c_4.
\end{equation}
Using the four boundary conditions, we obtain
\begin{equation}
  \bmat{0 & 0 & 0 & 1\\ L^3/6 & L^2/2 & L & 1\\ 0 & 0 & 1 & 0\\ L^2/2 & L & 1 & 0}\Bmat{c_1\\ c_2\\ c_3\\ c_4} = \Bmat{v_1\\ v_2\\ \theta_{z1}\\ \theta_{z2}}.
\end{equation}
Solving the above system, we obtain
\begin{align*}
  c_1 &= \frac{12}{L^3}(v_1-v_2) + \frac{6}{L^2}(\theta_{z1}+\theta_{z2})\\
  c_2 &= \frac{6}{L^2}(-v_1+v_2) - \frac1L(4\theta_{z1}+2\theta_{z2})\\
  c_3 &= \theta_{z1}\\
  c_4 &= v_1
\end{align*}
Therefore,
\begin{equation}
  \begin{aligned}
    v(x) &= \frac{x^3}{6}\Bigl[\frac{12}{L^3}(v_1-v_2)+\frac{6}{L^2}(\theta_{z1}+\theta_{z2})\Bigr]
          + \frac{x^2}{2}\Bigl[\frac{6}{L^2}(-v_1+v_2)-\frac1L(4\theta_{z1}+2\theta_{z2})\Bigr] + \theta_{z1}x + v_1\\
         &= N_1(x)v_1 + N_2(x)\theta_{z1} + N_3(x)v_2 + N_4(x)\theta_{z2}
  \end{aligned}
\end{equation}
where the shape functions are shown below.
\begin{fig}
  % Hermitian shape functions N1..N4 of a beam element
\begin{tikzpicture}[scale=1.0,font=\small,
  wall/.style={fill=black!50,draw=none},
  declare function={N1(\x)=1-3*\x^2+2*\x^3; N2(\x)=\x*(1-\x)^2; N3(\x)=3*\x^2-2*\x^3; N4(\x)=\x^3-\x^2;}]
  \def\L{5.5}
  % ---- N1 ----
  \begin{scope}
    \node[anchor=west] at (-0.2,1.6) {$N_1(x) = 1-3\bigl(\tfrac{x}{L}\bigr)^2+2\bigl(\tfrac{x}{L}\bigr)^3$};
    \draw[dashed] (0,0) -- (\L,0);
    \fill[wall] (-0.2,-0.25) rectangle (0,0.25); \fill[wall] (-0.2,0.75) rectangle (0,1.25);
    \fill[wall] (\L,-0.25) rectangle (\L+0.2,0.25);
    \draw[member,domain=0:1,samples=60,variable=\t] plot ({\t*\L},{N1(\t)});
    \draw[dim] (-0.5,0) -- (-0.5,1) node[midway,left] {1};
  \end{scope}
  % ---- N2 ----
  \begin{scope}[yshift=-3.2cm]
    \node[anchor=west] at (-0.2,1.6) {$N_2(x) = x\bigl(1-\tfrac{x}{L}\bigr)^2$};
    \draw[dashed] (0,0) -- (\L,0);
    \pic at (0,0) {pin}; \fill[wall] (\L,-0.25) rectangle (\L+0.2,0.25);
    \draw[member,domain=0:1,samples=60,variable=\t] plot ({\t*\L},{N2(\t)*\L});
    \draw[thin] (0,0) -- (1.3,1.3);
    \draw[->] (1.0,0) arc (0:45:1.0); \node at (1.2,0.35) {1};
  \end{scope}
  % ---- N3 ----
  \begin{scope}[yshift=-6.4cm]
    \node[anchor=west] at (-0.2,1.6) {$N_3(x) = 3\bigl(\tfrac{x}{L}\bigr)^2-2\bigl(\tfrac{x}{L}\bigr)^3$};
    \draw[dashed] (0,0) -- (\L,0);
    \fill[wall] (-0.2,-0.25) rectangle (0,0.25);
    \fill[wall] (\L,-0.25) rectangle (\L+0.2,0.25); \fill[wall] (\L,0.75) rectangle (\L+0.2,1.25);
    \draw[member,domain=0:1,samples=60,variable=\t] plot ({\t*\L},{N3(\t)});
    \draw[dim] (\L+0.5,0) -- (\L+0.5,1) node[midway,right] {1};
  \end{scope}
  % ---- N4 ----
  \begin{scope}[yshift=-9.4cm]
    \node[anchor=west] at (-0.2,1.4) {$N_4(x) = \dfrac{x^3}{L^2}-\dfrac{x^2}{L}$};
    \draw[dashed] (0,0) -- (\L,0);
    \fill[wall] (-0.2,-0.25) rectangle (0,0.25); \pic at (\L,0) {pin};
    \draw[member,domain=0:1,samples=60,variable=\t] plot ({\t*\L},{N4(\t)*\L});
    \draw[thin] (\L,0) -- (\L-1.3,-1.3);
    \draw[->] (\L-1.0,0) arc (180:225:1.0); \node at (\L-1.25,-0.3) {1};
  \end{scope}
\end{tikzpicture}

\end{fig}
Finally,
\begin{align}
  F_{y1} &= V_y(0) = EI_z v'''(0) = \frac{EI_z}{L}\Bigl(\frac{12}{L^2}v_1 + \frac6L\theta_{z1} - \frac{12}{L^2}v_2 + \frac6L\theta_{z2}\Bigr)\\
  M_{z1} &= -M_z(0) = -EI_z v''(0) = \frac{EI_z}{L}\Bigl(\frac6L v_1 + 4\theta_{z1} - \frac6L v_2 + 2\theta_{z2}\Bigr)\\
  F_{y2} &= -V_y(L) = -EI_z v'''(L) = \frac{EI_z}{L}\Bigl(-\frac{12}{L^2}v_1 - \frac6L\theta_{z1} + \frac{12}{L^2}v_2 - \frac6L\theta_{z2}\Bigr)\\
  M_{z2} &= M_z(L) = EI_z v''(L) = \frac{EI_z}{L}\Bigl(\frac6L v_1 + 2\theta_{z1} - \frac6L v_2 + 4\theta_{z2}\Bigr)
\end{align}
In matrix form,
\begin{equation}
  \Bmat{F_{y1}\\ M_{z1}\\ F_{y2}\\ M_{z2}} = \frac{EI_z}{L}
  \begin{bmatrix}
    \dfrac{12}{L^2} & \dfrac6L & -\dfrac{12}{L^2} & \dfrac6L\\[8pt]
    \dfrac6L & 4 & -\dfrac6L & 2\\[8pt]
    -\dfrac{12}{L^2} & -\dfrac6L & \dfrac{12}{L^2} & -\dfrac6L\\[8pt]
    \dfrac6L & 2 & -\dfrac6L & 4
  \end{bmatrix}
  \Bmat{v_1\\ \theta_{z1}\\ v_2\\ \theta_{z2}}
\end{equation}

\topic{Energy method:}
Knowing that $v(x) = \NN\bs{\Delta}$, where $\NN(x) = \langle N_1, N_2, N_3, N_4\rangle$ and $\bs{\Delta} = \langle v_1, \theta_{z1}, v_2, \theta_{z2}\rangle^\tr$, the strain energy is
\begin{equation}
  U = \frac12\int_0^L EI_z(v'')^2\,dx = \frac12 EI_z\int_0^L (v'')^\tr v''\,dx
    = \frac12 EI_z\int_0^L \bs{\Delta}^\tr\NN''^{\tr}\NN''\bs{\Delta}\,dx = \frac12\bs{\Delta}^\tr\kmat\bs{\Delta}.
\end{equation}
Therefore,
\begin{equation}
  \kmat = EI_z\int_0^L \NN''^{\tr}\NN''\,dx.
\end{equation}

\section{Flexural stiffness matrix for bending w.r.t.\ the \texorpdfstring{$y$}{y}-axis}
$\theta_y(x) = -w'(x)$
\begin{fig}
  % Bending about the y-axis: w(x) and w'(x) > 0
\begin{tikzpicture}[scale=1.0,font=\small]
  \draw[->,thick] (0,0) -- (4.2,0) node[right] {$x$};
  \draw[->,thick] (0,0) -- (0,2.8) node[above] {$z$};
  \draw[->,thick,dashed] (0,0) -- (1.6,1.6) node[above right] {$y$};
  \draw[fill=black!10] (0,-0.12) rectangle (3.4,0.12);
  \draw[fill=black!10,rotate=12] (0,0.1) rectangle (3.4,0.32);
  \draw[->] (3.2,0.1) -- (3.2,0.66); \node[right] at (3.25,0.4) {$w(x)$};
  \node[right] at (3.3,1.1) {$w'(x)>0$};
\end{tikzpicture}

\end{fig}
The stiffness matrix for bending w.r.t.\ the $y$-axis can be derived in a way similar to that for the $z$-axis bending, but it is noted that $w'(x) = -\theta_y(x)$; the result is
\begin{equation}
  \Bmat{F_{z1}\\ M_{y1}\\ F_{z2}\\ M_{y2}} = \frac{EI_y}{L}
  \begin{bmatrix}
    \dfrac{12}{L^2} & -\dfrac6L & -\dfrac{12}{L^2} & -\dfrac6L\\[8pt]
    -\dfrac6L & 4 & \dfrac6L & 2\\[8pt]
    -\dfrac{12}{L^2} & \dfrac6L & \dfrac{12}{L^2} & \dfrac6L\\[8pt]
    -\dfrac6L & 2 & \dfrac6L & 4
  \end{bmatrix}
  \Bmat{w_1\\ \theta_{y1}\\ w_2\\ \theta_{y2}}
\end{equation}
It shall be noted that the shape functions $N_1(x)$ and $N_3(x)$ for the $y$-axis bending are the same as those for the $z$-axis bending. However, $N_2(x)$ and $N_4(x)$ have a sign difference. Therefore, the corresponding stiffness matrices are also slightly different.

\section{Torsional stiffness matrix}
The torsional stiffness matrix can be derived in a way similar to that used in deriving the axial stiffness matrix; the result is:
\begin{equation}
  \Bmat{M_{x1}\\ M_{x2}} = \frac{GJ}{L}\bmat{1 & -1\\ -1 & 1}\Bmat{\theta_{x1}\\ \theta_{x2}}.
\end{equation}
The table below is adopted from I.~A.~Macleod (1990), \emph{Analytical Modelling of Structural Systems}, Ellis Horwood Limited, page 203. Note that the torsion constant is written as $J$ in this course.

\begin{table}[H]
  \centering\small
  \caption{St Venant torsion constant $C$ (after Macleod, 1990, Table A3).}
  \renewcommand{\arraystretch}{1.2}
  \begin{tabular}{@{}>{\raggedright\arraybackslash}p{0.3\textwidth}>{\raggedright\arraybackslash}p{0.64\textwidth}@{}}
    \toprule
    Shape & $C$\\
    \midrule
    Circular section & $C =$ second polar moment of area $= \pi R^4/2$\\[4pt]
    Closed thin-walled section & $C = \dfrac{4A_0^2}{\displaystyle\oint\frac{ds}{t}}$\rule[-22pt]{0pt}{44pt}\newline
      $A_0$ is the area enclosed by a line through the mid-thickness.\newline
      The integral is carried out along the perimeter.\\[6pt]
    Narrow rectangle, $\dfrac bt > 5$ & \begin{tikzpicture}[scale=1.0,font=\scriptsize,baseline=(current bounding box.center)]
  \draw (0,0) rectangle (2.4,0.35); \node[above] at (1.2,0.35) {$b$}; \node[right] at (2.4,0.17) {$t$};
\end{tikzpicture}
\qquad $C = \dfrac13 bt^3$\rule[-10pt]{0pt}{26pt}\\[6pt]
    Open section composed of narrow rectangles & $C = \sum\dfrac13 bt^3$\rule[-10pt]{0pt}{24pt}\newline
      $b$ is the larger dimension of a rectangle\newline
      $t$ is the width of the rectangle\newline
      The summation is over all rectangles\\[6pt]
    Closed rectangular thick-walled section & See Marshall (1970)\\[6pt]
    Solid rectangle & \begin{tikzpicture}[scale=1.0,font=\scriptsize,baseline=(current bounding box.center)]
  \draw (0,0) rectangle (2.4,0.7); \node[below] at (1.2,0) {$b$}; \node[right] at (2.4,0.35) {$t$};
\end{tikzpicture}
\rule[-14pt]{0pt}{20pt}\newline
      $C = bt^3\Bigl[\dfrac13 - 0.21\dfrac tb\Bigl(1-\dfrac{t^4}{12b^4}\Bigr)\Bigr]$\rule[-12pt]{0pt}{30pt}\\[6pt]
    Solid square & \begin{tikzpicture}[scale=1.0,font=\scriptsize,baseline=(current bounding box.center)]
  \draw (0,0) rectangle (0.9,0.9); \node[above] at (0.45,0.9) {$b$}; \node[left] at (0,0.45) {$b$};
\end{tikzpicture}
\qquad $C = 0.1406\,b^4$\rule[-14pt]{0pt}{30pt}\\
    \bottomrule
  \end{tabular}
\end{table}

\chapter[Coordinate Transformation]{Coordinate Transformation, Contragredient Principle, and Congruent Transformation}
\label{ch:transform}

\section{Coordinate transformation}
\begin{fig}
  % Vector P in two coordinate systems xyz and x'y'z'
\begin{tikzpicture}[scale=1.0,font=\small]
  \coordinate (O) at (0,0);
  \draw[->,thick] (O) -- (3.2,0) node[right] {$x$};
  \draw[->,thick] (O) -- (0,2.6) node[above] {$y$};
  \draw[->,thick] (O) -- (-1.4,-1.4) node[below left] {$z$};
  \draw[->,dashed,thick] (O) -- (2.6,2.0) node[above right] {$x'$};
  \draw[->,dashed,thick] (O) -- (-1.5,2.0) node[above left] {$y'$};
  \draw[->,dashed,thick] (O) -- (-0.2,-1.8) node[below] {$z'$};
  \draw[-{Latex[length=3mm]},line width=1.6pt] (0.6,1.0) -- (1.5,2.1) node[above right] {$\mat{P}$};
\end{tikzpicture}

\end{fig}
The components of a vector in two coordinate systems $xyz$ and $x'y'z'$ have the following relation:
\begin{equation}
  \Bmat{P_{x'}\\ P_{y'}\\ P_{z'}} = \bmat{\lambda_{x'} & \mu_{x'} & \nu_{x'}\\ \lambda_{y'} & \mu_{y'} & \nu_{y'}\\ \lambda_{z'} & \mu_{z'} & \nu_{z'}}\Bmat{P_x\\ P_y\\ P_z}
  \quad\text{or}\quad
  \Bmat{P_x\\ P_y\\ P_z} = \bmat{\lambda_{x'} & \lambda_{y'} & \lambda_{z'}\\ \mu_{x'} & \mu_{y'} & \mu_{z'}\\ \nu_{x'} & \nu_{y'} & \nu_{z'}}\Bmat{P_{x'}\\ P_{y'}\\ P_{z'}}
\end{equation}
i.e.
\begin{equation}
  \Pvec' = \mat{R}\Pvec \quad\text{or}\quad \Pvec = \mat{R}^\tr\Pvec'
\end{equation}
where $\mat{R}$ is a rotation matrix, consisting of the directional cosines of $x'$, $y'$, and $z'$ with respect to the $xyz$ system; namely $\lambda_{x'}^2+\mu_{x'}^2+\nu_{x'}^2 = 1$ and similar expressions for the directional cosines of the $y'$ and $z'$ axes. $\mat{R}$ is orthogonal because $\mat{R}\mat{R}^\tr = \mat{I}$ or equivalently $\mat{R}^{-1} = \mat{R}^\tr$.
\begin{equation}
  \begin{aligned}
    \Pvec &= P_{x'}\mat{i}' + P_{y'}\mat{j}' + P_{z'}\mat{k}'\\
          &= P_{x'}(\lambda_{x'}\mat{i}+\mu_{x'}\mat{j}+\nu_{x'}\mat{k}) + P_{y'}(\lambda_{y'}\mat{i}+\mu_{y'}\mat{j}+\nu_{y'}\mat{k}) + P_{z'}(\lambda_{z'}\mat{i}+\mu_{z'}\mat{j}+\nu_{z'}\mat{k})\\
          &= (P_{x'}\lambda_{x'}+P_{y'}\lambda_{y'}+P_{z'}\lambda_{z'})\mat{i} + (P_{x'}\mu_{x'}+P_{y'}\mu_{y'}+P_{z'}\mu_{z'})\mat{j} + (P_{x'}\nu_{x'}+P_{y'}\nu_{y'}+P_{z'}\nu_{z'})\mat{k}\\
          &= P_x\mat{i} + P_y\mat{j} + P_z\mat{k}
  \end{aligned}
\end{equation}
That is,
\begin{equation}
  \Bmat{P_x\\ P_y\\ P_z} = \bmat{\lambda_{x'} & \lambda_{y'} & \lambda_{z'}\\ \mu_{x'} & \mu_{y'} & \mu_{z'}\\ \nu_{x'} & \nu_{y'} & \nu_{z'}}\Bmat{P_{x'}\\ P_{y'}\\ P_{z'}}
  \qquad\text{or}\qquad \Pvec = \mat{R}^\tr\Pvec'.
\end{equation}

Let the $x'y'z'$ system represent the local (member) coordinate system, and assume that the force--displacement relation for a 3D beam-column element is of the form:
\begin{equation}
  \F' = \Bmat{F_{x'1}\\ F_{y'1}\\ F_{z'1}\\ M_{x'1}\\ M_{y'1}\\ M_{z'1}\\ F_{x'2}\\ F_{y'2}\\ F_{z'2}\\ M_{x'2}\\ M_{y'2}\\ M_{z'2}}
  = \kmat'\Bmat{u_1'\\ v_1'\\ w_1'\\ \theta_{x'1}\\ \theta_{y'1}\\ \theta_{z'1}\\ u_2'\\ v_2'\\ w_2'\\ \theta_{x'2}\\ \theta_{y'2}\\ \theta_{z'2}} = \kmat'\bs{\Delta}'
\end{equation}
where the stiffness matrix in the local axes has been given in Ch.~\ref{ch:beamcolumn} and is not repeated here.

The components of the forces and moments at nodes 1 and 2 in the local coordinate system are related to their components in the global coordinate ($xyz$) system as follows:
\begin{equation}
  \F' = \bmat{\mat{R} & & & \\ & \mat{R} & & \\ & & \mat{R} & \\ & & & \mat{R}}\F = \Tmat\F
\end{equation}
This is because both nodal forces and moments are vectors. Similarly, the nodal degrees of freedom follow the same transformation; that is,
\begin{equation}
  \bs{\Delta}' = \bmat{\mat{R} & & & \\ & \mat{R} & & \\ & & \mat{R} & \\ & & & \mat{R}}\bs{\Delta} = \Tmat\bs{\Delta}
\end{equation}
With the above two transformations, the force--displacement relation becomes:
\begin{equation}
  \F' = \kmat'\bs{\Delta}' \quad\Longrightarrow\quad \Tmat\F = \kmat'\Tmat\bs{\Delta}
\end{equation}
Therefore, we can obtain the force--displacement relation w.r.t.\ the global coordinate system as follows:
\begin{equation}
  \boxed{\F = \Tmat^{-1}\kmat'\Tmat\bs{\Delta} = \Tmat^\tr\kmat'\Tmat\bs{\Delta} = \kmat\bs{\Delta}}
\end{equation}
Note that use has been made of the fact that $\Tmat^{-1} = \Tmat^\tr$ in the derivation of the stiffness transformation; $\kmat = \Tmat^\tr\kmat'\Tmat$. This relation looks very nice and seems familiar to us. Actually, we can see similar transformations in the transformation of stress or strain components, moments of inertia, etc.

\section{Contragredient principle and congruent transformation}
In deriving the stiffness transformation, $\kmat = \Tmat^\tr\kmat'\Tmat$, we have used the fact that $\Tmat^{-1} = \Tmat^\tr$, which holds only when $\Tmat$ is a square matrix. However, the stiffness transformation actually holds even when $\Tmat$ is a rectangular matrix.

Let a structure system (or a member) be described by two pairs of forces and displacements; one consists of $\F$ and $\bs{\Delta}$ and the other $\F'$ and $\bs{\Delta}'$. Usually we say $\F$ and $\bs{\Delta}$ are \emph{work conjugate} to each other and similarly $\F'$ and $\bs{\Delta}'$ are also a work conjugate pair.

The \textbf{contragredient principle} states that
\begin{equation}
  \bs{\Delta}' = \Tmat\bs{\Delta} \quad\Longleftrightarrow\quad \F = \Tmat^\tr\F'
\end{equation}
To prove this principle, we shall know that the strain energy can be written in terms of the two work conjugate pairs as follows.
\begin{equation}
  U = \tfrac12\bs{\Delta}^\tr\F = \tfrac12\bs{\Delta}'^\tr\F'
\end{equation}
Substituting $\bs{\Delta}' = \Tmat\bs{\Delta}$ into the expression for $U$ yields
\begin{equation}
  U = \tfrac12\bs{\Delta}^\tr\F = \tfrac12(\Tmat\bs{\Delta})^\tr\F' = \tfrac12\bs{\Delta}^\tr\Tmat^\tr\F'
\end{equation}
Obviously, $\F = \Tmat^\tr\F'$. Analogously, it is easy to show that if $\F = \Tmat^\tr\F'$, then $\bs{\Delta}' = \Tmat\bs{\Delta}$.

What the contragredient principle says is that once the relation between two force (displacement) vectors is known, then the relation between their conjugate displacement (force) vectors is also sure.

Now assume that the force--displacement relations for the two conjugate pairs are as follows.
\begin{equation}
  \F = \kmat\bs{\Delta},\qquad \F' = \kmat'\bs{\Delta}'
\end{equation}
Then, the strain energy can be written as
\begin{equation}
  U = \tfrac12\bs{\Delta}^\tr\kmat\bs{\Delta} = \tfrac12\bs{\Delta}'^\tr\kmat'\bs{\Delta}'
\end{equation}
Substituting $\bs{\Delta}' = \Tmat\bs{\Delta}$ into the above equation, we can show that
\begin{equation}
  \kmat = \Tmat^\tr\kmat'\Tmat
\end{equation}
This is the so-called \textbf{congruent transformation}.

Note that it is not required that $\Tmat$ be a square and orthogonal matrix when the contragredient principle and the congruent transformation are derived as shown above.

There are various applications of the contragredient principle and congruent transformation. They are not only applicable for a member but also for a structural system. The latter applications include such subjects as static condensation, structures with constraints, and structures with skew supports, which are to be discussed in Ch.~\ref{ch:special}. For the former applications, three examples are provided herein.

\begin{example}
Consider a 2D truss element; its stiffness matrix referred to the member (local) axes is known as:
\begin{equation*}
  \kmat' = \frac{EA}{L}\bmat{1 & -1\\ -1 & 1}
\end{equation*}
Let the local axes $x'$, $y'$ be obtained by rotating counterclockwise the global axes $x$, $y$ an angle $\phi$. Then the displacement vectors in the local and global coordinate systems have the following relation.
\begin{equation*}
  \bs{\Delta}' = \Bmat{u_1'\\ u_2'} = \bmat{\cos\phi & \sin\phi & 0 & 0\\ 0 & 0 & \cos\phi & \sin\phi}\Bmat{u_1\\ v_1\\ u_2\\ v_2} = \Tmat\bs{\Delta}
\end{equation*}
According to the contragredient principle, we obtain
\begin{equation*}
  \F = \Bmat{F_{x1}\\ F_{y1}\\ F_{x2}\\ F_{y2}} = \Tmat^\tr\F' = \Bmat{F_{x'1}\cos\phi\\ F_{x'1}\sin\phi\\ F_{x'2}\cos\phi\\ F_{x'2}\sin\phi}
\end{equation*}
This relation is obviously true.

The stiffness matrix associated with $\bs{\Delta}$ can be obtained through the congruent transformation as follows.
\begin{equation*}
  \kmat = \Tmat^\tr\kmat'\Tmat = \frac{EA}{L}\bmat{c^2 & sc & -c^2 & -sc\\ sc & s^2 & -sc & -s^2\\ -c^2 & -sc & c^2 & sc\\ -sc & -s^2 & sc & s^2}
\end{equation*}
where $c = \cos\phi$ and $s = \sin\phi$.
\end{example}

\begin{example}[Derive {$[k]$} of an unsupported element from {$[k]$} of a supported element]
Considering a simply-supported beam subjected to end moments $M_1$ and $M_2$, these two end moments can be expressed in terms of the corresponding end rotations as:
\begin{fig}
  % Simply supported beam with end moments M1(alpha1), M2(alpha2)
\begin{tikzpicture}[scale=1.1,>=Latex]
  \def\L{6}
  % beam
  \draw[line width=2.2pt,black!70] (0,0) -- (\L,0);
  \node[above] at (\L/2,0.05) {$EI$};
  % pin support (left)
  \draw[fill=black!60] (0,0) -- (-0.25,-0.4) -- (0.25,-0.4) -- cycle;
  \fill[white,draw=black] (0,0) circle (0.06);
  \fill[black!25] (-0.4,-0.4) rectangle (0.4,-0.55);
  % roller support (right)
  \draw[fill=black!60] (\L,0) -- (\L-0.25,-0.4) -- (\L+0.25,-0.4) -- cycle;
  \fill[white,draw=black] (\L,0) circle (0.06);
  \draw[fill=white] (\L-0.15,-0.5) circle (0.09);
  \draw[fill=white] (\L+0.15,-0.5) circle (0.09);
  \fill[black!25] (\L-0.4,-0.6) rectangle (\L+0.4,-0.75);
  % end moments
  \draw[->,thick] (-0.05,0.55) arc (150:-30:0.35);
  \node[above] at (0.2,0.75) {$M_1(\alpha_1)$};
  \draw[->,thick] (\L+0.05,0.55) arc (30:210:0.35);
  \node[above] at (\L-0.2,0.75) {$M_2(\alpha_2)$};
  % length dimension
  \draw[<->] (0,-1.1) -- (\L,-1.1) node[midway,fill=white,inner sep=1pt] {$L$};
  \draw (0,-0.85) -- (0,-1.3);
  \draw (\L,-0.85) -- (\L,-1.3);
\end{tikzpicture}

\end{fig}
\begin{equation*}
  \Bmat{M_{z1}\\ M_{z2}} = \frac{2EI_z}{L}\bmat{2 & 1\\ 1 & 2}\Bmat{\alpha_1\\ \alpha_2}
  \qquad\text{or}\qquad \F' = \kmat'\bs{\Delta}'
\end{equation*}
Assume the corresponding unsupported beam has the following force--displacement relation:
\begin{equation*}
  \F = \Bmat{F_{y1}\\ M_{z1}\\ F_{y2}\\ M_{z2}} = \kmat\bs{\Delta} = \kmat\Bmat{v_1\\ \theta_{z1}\\ v_2\\ \theta_{z2}}
\end{equation*}
The contragredient principle and congruent transformation allow us to investigate the relations between these two conjugate pairs. First, from the equilibrium conditions, we have
\begin{fig}
  % Unsupported beam element: end shears and moments
\begin{tikzpicture}[scale=1.0,font=\small]
  \def\L{6}
  \draw[member] (0,0) -- (\L,0);
  \draw[load] (0,0.1) -- (0,1.3) node[above] {$F_{y1}$};
  \draw[load] (\L,0.1) -- (\L,1.3) node[above] {$F_{y2}$};
  \draw[->,thick] (-0.05,-0.4) arc (-100:100:0.4); \node[left] at (-0.4,0) {$M_{z1}$};
  \draw[->,thick] (\L+0.05,-0.4) arc (280:80:0.4); \node[right] at (\L+0.4,0) {$M_{z2}$};
  \draw[dim] (0,-0.9) -- (\L,-0.9) node[midway,fill=white,inner sep=1pt] {$L$};
  \draw (0,-0.7) -- (0,-1.1); \draw (\L,-0.7) -- (\L,-1.1);
\end{tikzpicture}

\end{fig}
\begin{equation*}
  F_{y1} = \frac{M_{z1}+M_{z2}}{L},\quad\text{and}\quad F_{y2} = -\frac{M_{z1}+M_{z2}}{L}
\end{equation*}
\begin{equation*}
  \Bmat{F_{y1}\\ M_{z1}\\ F_{y2}\\ M_{z2}} =
  \begin{bmatrix} \dfrac1L & \dfrac1L\\[6pt] 1 & 0\\[3pt] -\dfrac1L & -\dfrac1L\\[6pt] 0 & 1\end{bmatrix}
  \Bmat{M_{z1}\\ M_{z2}},\quad\text{or}\quad \F = \Tmat^\tr\F'
\end{equation*}
According to the contragredient principle,
\begin{equation*}
  \bs{\Delta}' = \Tmat\bs{\Delta},\quad\text{or}\quad
  \Bmat{\alpha_1\\ \alpha_2} = \begin{bmatrix}\dfrac1L & 1 & -\dfrac1L & 0\\[6pt] \dfrac1L & 0 & -\dfrac1L & 1\end{bmatrix}\Bmat{v_1\\ \theta_{z1}\\ v_2\\ \theta_{z2}}
\end{equation*}
That is,
\begin{equation*}
  \alpha_1 = \theta_{z1}-\theta_r = \theta_{z1}-\frac{v_2-v_1}{L}\quad\text{and}\quad
  \alpha_2 = \theta_{z2}-\theta_r = \theta_{z2}-\frac{v_2-v_1}{L}
\end{equation*}
\begin{fig}
  % End rotations w.r.t. the chord: alpha_i = theta_zi - theta_r
\begin{tikzpicture}[scale=1.0,font=\small]
  \coordinate (A) at (0,0); \coordinate (B) at (9,3);
  % chord
  \draw[very thick] (A) -- (B);
  % deflected shape (dashed)
  \draw[dashed,thick] (A) .. controls (2.5,2.6) and (6.0,1.2) .. (B);
  % tangents at the ends
  \draw[thin] (A) -- ++(52:1.6); \draw[thin] (B) -- ++(232:1.9);
  % horizontal references
  \draw[thin] (A) -- ++(2.2,0); \draw[thin] (B) -- ++(-2.2,0);
  % angles at end 1
  \draw[->] ($(A)+(1.6,0)$) arc (0:18.4:1.6); \node at (1.95,0.25) {$\theta_r$};
  \draw[->] ($(A)+(1.1,0)$) arc (0:52:1.1); \node at (1.05,0.9) {$\theta_{z1}$};
  \draw[->] ($(A)+(18.4:1.4)$) arc (18.4:52:1.4); \node at (1.1,1.35) {$\alpha_1$};
  % angles at end 2
  \draw[->] ($(B)+(180:1.6)$) arc (180:198.4:1.6); \node at ($(B)+(-1.9,-0.25)$) {$\theta_{z2}$};
  \draw[->] ($(B)+(198.4:1.3)$) arc (198.4:232:1.3); \node at ($(B)+(-1.35,-0.95)$) {$\alpha_2$};
  % vertical dims
  \draw[dim] (-0.3,-1.5) -- (-0.3,0) node[midway,left] {$v_1$};
  \draw[dim] (9.3,-1.5) -- (9.3,3) node[midway,right] {$v_2$};
  \draw[thin] (0,-1.5) -- (-0.5,-1.5); \draw[thin] (9,-1.5) -- (9.5,-1.5);
  \draw[dim] (0,-1.5) -- (9,-1.5) node[midway,fill=white,inner sep=1pt] {$L$};
\end{tikzpicture}

\end{fig}
\begin{equation*}
  \begin{aligned}
  M_{12} &= \frac{2EI}{L}\bigl[2\alpha_1+\alpha_2\bigr] = \frac{2EI}{L}\bigl[2(\theta_{z1}-\theta_r)+(\theta_{z2}-\theta_r)\bigr]\\
         &= \frac{2EI}{L}(2\theta_{z1}+\theta_{z2}-3\theta_r) = \frac{2EI}{L}\Bigl(2\theta_{z1}+\theta_{z2}-3\frac{\Delta}{L}\Bigr)
  \end{aligned}
\end{equation*}
where $\theta_r = (v_2-v_1)/L = \Delta/L$ is the rotation of the chord connecting nodes 1 and 2. Therefore, $\alpha_1$ and $\alpha_2$ are the end rotations w.r.t.\ the chord; they are often called \emph{natural deformations} because rigid body displacements are excluded. This is \emph{the slope deflection equation}.

In the above, we first establish the force relations, $\F = \Tmat^\tr\F'$, and then obtain the displacement relations according to the contragredient principle. Of course, we can also reverse the sequence.

The stiffness matrix for the unsupported beam can be obtained through the use of the congruent transformation; hence,
\begin{equation*}
  \kmat = \Tmat^\tr\kmat'\Tmat = \frac{EI_z}{L}
  \begin{bmatrix}
    \dfrac{12}{L^2} & \dfrac6L & -\dfrac{12}{L^2} & \dfrac6L\\[8pt]
    \dfrac6L & 4 & -\dfrac6L & 2\\[8pt]
    -\dfrac{12}{L^2} & -\dfrac6L & \dfrac{12}{L^2} & -\dfrac6L\\[8pt]
    \dfrac6L & 2 & -\dfrac6L & 4
  \end{bmatrix}
\end{equation*}
The result is the same as that derived in Ch.~\ref{ch:beamcolumn}.
\end{example}

\begin{example}
Consider an elastic beam of length $L$ with rigid zones at ends, whose lengths are $a$ (left) and $b$ (right), respectively. Denote the degrees of freedom of the beam with rigid zones by $\bs{\Delta}$ and those of the beam by $\bs{\Delta}'$. Then, it is easy to show that
\begin{fig}
  % Beam with rigid zones: DOFs Delta (with rigid zones) and Delta' (elastic beam only)
\begin{tikzpicture}[scale=0.68,font=\small]
  % ---- with rigid zones ----
  \begin{scope}
    \fill[hatch] (0,-0.12) rectangle (1.0,0.12); \draw (0,-0.12) rectangle (1.0,0.12);
    \draw[member] (1.0,0) -- (5.5,0);
    \fill[hatch] (5.5,-0.12) rectangle (6.2,0.12); \draw (5.5,-0.12) rectangle (6.2,0.12);
    \draw[load] (-1.2,0) -- (-0.1,0) node[pos=0,left] {$u_1$};
    \draw[load] (0,0.15) -- (0,1.2) node[above] {$v_1$};
    \draw[->,thick] (0.35,-0.3) arc (-40:200:0.35); \node[below left] at (0,-0.35) {$\theta_{z1}$};
    \draw[load] (6.3,0) -- (7.4,0) node[right] {$u_2$};
    \draw[load] (6.2,0.15) -- (6.2,1.2) node[above] {$v_2$};
    \draw[->,thick] (6.55,-0.3) arc (-40:200:0.35); \node[below right] at (6.2,-0.35) {$\theta_{z2}$};
    \node[below] at (0.5,-0.15) {$a$}; \node[below] at (5.85,-0.15) {$b$};
    \draw[dim] (1.0,-0.9) -- (5.5,-0.9) node[midway,fill=white,inner sep=1pt] {$L$};
  \end{scope}
  % ---- elastic beam only ----
  \begin{scope}[xshift=10cm]
    \draw[member] (0,0) -- (4.5,0);
    \draw[load] (-1.2,0) -- (-0.1,0) node[pos=0,left] {$u_1'$};
    \draw[load] (0,0.15) -- (0,1.2) node[above] {$v_1'$};
    \draw[->,thick] (0.35,-0.3) arc (-40:200:0.35); \node[below left] at (0,-0.35) {$\theta_{z1}'$};
    \draw[load] (4.6,0) -- (5.7,0) node[right] {$u_2'$};
    \draw[load] (4.5,0.15) -- (4.5,1.2) node[above] {$v_2'$};
    \draw[->,thick] (4.85,-0.3) arc (-40:200:0.35); \node[below right] at (4.5,-0.35) {$\theta_{z2}'$};
    \draw[dim] (0,-0.9) -- (4.5,-0.9) node[midway,fill=white,inner sep=1pt] {$L$};
  \end{scope}
\end{tikzpicture}

\end{fig}
\begin{equation*}
  \bs{\Delta}' = \Bmat{u_1'\\ v_1'\\ \theta_{z1}'\\ u_2'\\ v_2'\\ \theta_{z2}'} =
  \bmat{1 & 0 & 0 & 0 & 0 & 0\\ 0 & 1 & a & 0 & 0 & 0\\ 0 & 0 & 1 & 0 & 0 & 0\\ 0 & 0 & 0 & 1 & 0 & 0\\ 0 & 0 & 0 & 0 & 1 & -b\\ 0 & 0 & 0 & 0 & 0 & 1}
  \Bmat{u_1\\ v_1\\ \theta_{z1}\\ u_2\\ v_2\\ \theta_{z2}} = \Tmat\bs{\Delta}
\end{equation*}
The same $\Tmat$ can also be derived from the conditions, $\F = \Tmat^\tr\F'$.
\begin{fig}
  % Beam with rigid zones: forces F and F'
\begin{tikzpicture}[scale=0.68,font=\small]
  \begin{scope}
    \fill[hatch] (0,-0.12) rectangle (1.0,0.12); \draw (0,-0.12) rectangle (1.0,0.12);
    \draw[member] (1.0,0) -- (5.5,0);
    \fill[hatch] (5.5,-0.12) rectangle (6.2,0.12); \draw (5.5,-0.12) rectangle (6.2,0.12);
    \draw[load] (-1.2,0) -- (-0.1,0) node[pos=0,left] {$F_{x1}$};
    \draw[load] (0,0.15) -- (0,1.2) node[above] {$F_{y1}$};
    \draw[->,thick] (0.35,-0.3) arc (-40:200:0.35); \node[below left] at (0,-0.35) {$M_{z1}$};
    \draw[load] (6.3,0) -- (7.4,0) node[right] {$F_{x2}$};
    \draw[load] (6.2,0.15) -- (6.2,1.2) node[above] {$F_{y2}$};
    \draw[->,thick] (6.55,-0.3) arc (-40:200:0.35); \node[below right] at (6.2,-0.35) {$M_{z2}$};
    \node[below] at (0.5,-0.15) {$a$}; \node[below] at (5.85,-0.15) {$b$};
    \draw[dim] (1.0,-0.9) -- (5.5,-0.9) node[midway,fill=white,inner sep=1pt] {$L$};
  \end{scope}
  \begin{scope}[xshift=10cm]
    \draw[member] (0,0) -- (4.5,0);
    \draw[load] (-1.2,0) -- (-0.1,0) node[pos=0,left] {$F_{x1}$};
    \draw[load] (0,0.15) -- (0,1.2) node[above] {$F_{y1}$};
    \draw[->,thick] (0.35,-0.3) arc (-40:200:0.35); \node[below left,align=right] at (0.1,-0.35) {$M_{z1}-F_{y1}\cdot a$};
    \draw[load] (4.6,0) -- (5.7,0) node[right] {$F_{x2}$};
    \draw[load] (4.5,0.15) -- (4.5,1.2) node[above] {$F_{y2}$};
    \draw[->,thick] (4.85,-0.3) arc (-40:200:0.35); \node[below right] at (4.4,-0.35) {$M_{z2}+F_{y2}\cdot b$};
    \draw[dim] (0,-0.9) -- (4.5,-0.9) node[midway,fill=white,inner sep=1pt] {$L$};
  \end{scope}
\end{tikzpicture}

\end{fig}
Therefore
\begin{equation*}
  \begin{array}{lcl}
    F_{x1}' = F_{x1} & \rightarrow & F_{x1} = F_{x1}'\\
    F_{y1}' = F_{y1} & \rightarrow & F_{y1} = F_{y1}'\\
    M_{z1}' = M_{z1} - F_{y1}a & \rightarrow & M_{z1} = M_{z1}' + F_{y1}a = M_{z1}' + F_{y1}'a\\
    F_{x2}' = F_{x2} & \rightarrow & F_{x2} = F_{x2}'\\
    F_{y2}' = F_{y2} & \rightarrow & F_{y2} = F_{y2}'\\
    M_{z2}' = M_{z2} + F_{y2}b & \rightarrow & M_{z2} = M_{z2}' - F_{y2}b = M_{z2}' - F_{y2}'b
  \end{array}
\end{equation*}
In matrix form,
\begin{equation*}
  \F = \Bmat{F_{x1}\\ F_{y1}\\ M_{z1}\\ F_{x2}\\ F_{y2}\\ M_{z2}} =
  \bmat{1 & 0 & 0 & 0 & 0 & 0\\ 0 & 1 & 0 & 0 & 0 & 0\\ 0 & a & 1 & 0 & 0 & 0\\ 0 & 0 & 0 & 1 & 0 & 0\\ 0 & 0 & 0 & 0 & 1 & 0\\ 0 & 0 & 0 & 0 & -b & 1}
  \Bmat{F_{x1}'\\ F_{y1}'\\ M_{z1}'\\ F_{x2}'\\ F_{y2}'\\ M_{z2}'} = \Tmat^\tr\F'
\end{equation*}
The stiffness matrix for the beam with rigid zones, $\kmat$, can then be related to the stiffness matrix of the beam without rigid zones, $\kmat'$, through the congruent transformation; that is
\begin{equation*}
  \kmat = \Tmat^\tr\kmat'\Tmat
\end{equation*}
\end{example}

\begin{example}
Consider an elastic beam of length $L$ connecting an axial force member by two rigid links, whose lengths are $e_1$ (left) and $e_2$ (right). The axial force member has the axial rigidity of $EA$ and length of $L_1$, where $L_1 = \sqrt{L^2+(e_2-e_1)^2}$.
\begin{fig}
  % Beam connected to an axial member by two rigid links
\begin{tikzpicture}[scale=0.68,font=\small]
  % ---- (a) DOF Delta ----
  \begin{scope}
    \draw[fill=black!15] (0,0) rectangle (5.5,0.3);
    \node[above] at (2.75,0.3) {$E, I, A, L$};
    \fill[hatch] (0,-0.5) rectangle (0.18,0); \draw (0,-0.5) rectangle (0.18,0);
    \fill[hatch] (5.32,-0.9) rectangle (5.5,0); \draw (5.32,-0.9) rectangle (5.5,0);
    \draw[thinmember] (0.09,-0.5) -- (5.41,-0.9);
    \node[below] at (2.75,-0.75) {$E, A, L_1$};
    \node[right,font=\footnotesize] at (0.2,-0.25) {$e_1$}; \node[left,font=\footnotesize] at (5.3,-0.45) {$e_2$};
    \node[below left] at (0,0) {1}; \node[below right] at (5.5,0) {2};
    \draw[load] (-1.2,0.15) -- (-0.1,0.15) node[pos=0,left] {$u_1$};
    \draw[load] (0,0.4) -- (0,1.4) node[above] {$v_1$};
    \draw[->,thick] (0.3,-0.1) arc (-30:200:0.35); \node[below left] at (-0.1,-0.35) {$\theta_{z1}$};
    \draw[load] (5.6,0.15) -- (6.7,0.15) node[right] {$u_2$};
    \draw[load] (5.5,0.4) -- (5.5,1.4) node[above] {$v_2$};
    \draw[->,thick] (5.8,-0.1) arc (-30:200:0.35); \node[below right] at (5.6,-0.35) {$\theta_{z2}$};
    \node[below] at (2.75,-1.7) {(a) Degrees of freedom $\bs{\Delta}$};
  \end{scope}
  % ---- (b) DOF Delta' ----
  \begin{scope}[xshift=9.5cm]
    \draw[thinmember] (0,-0.5) -- (5.5,-0.9);
    \node[below] at (2.75,-0.85) {$E, A, L_1$};
    \draw[load] (-1.2,-0.5) -- (-0.1,-0.5) node[pos=0,left] {$d_1$};
    \draw[load] (0,-0.4) -- (0,0.8) node[above] {$d_2$};
    \draw[load] (5.6,-0.9) -- (6.7,-0.9) node[right] {$d_3$};
    \draw[load] (5.5,-0.8) -- (5.5,0.4) node[above] {$d_4$};
    \node[below] at (2.75,-1.7) {(b) Degrees of freedom $\bs{\Delta}'$};
  \end{scope}
\end{tikzpicture}

\end{fig}
The strain energy of the system is
\begin{equation*}
  U = U_b + U_{truss} = \tfrac12\bs{\Delta}^\tr\kmat_b\bs{\Delta} + \tfrac12\bs{\Delta}^\tr\kmat_t\bs{\Delta}
  \quad\Longrightarrow\quad \kmat = \kmat_b + \kmat_t
\end{equation*}
The stiffness matrix of the truss $\kmat'$ corresponding to $\bs{\Delta}'$ is known, but what is the stiffness matrix of the truss corresponding to $\bs{\Delta}$? This can be easily done with the aid of the congruent transformation. It is not difficult to show that
\begin{equation*}
  \bs{\Delta}' = \Bmat{d_1\\ d_2\\ d_3\\ d_4} =
  \bmat{1 & 0 & e_1 & 0 & 0 & 0\\ 0 & 1 & 0 & 0 & 0 & 0\\ 0 & 0 & 0 & 1 & 0 & e_2\\ 0 & 0 & 0 & 0 & 1 & 0}
  \Bmat{u_1\\ v_1\\ \theta_1\\ u_2\\ v_2\\ \theta_2},\quad\text{or}\quad \bs{\Delta}' = \Tmat\bs{\Delta}.
\end{equation*}
Then, according to the congruent transformation $\kmat_t = \Tmat^\tr\kmat'\Tmat$, we have
\begin{equation*}
  \kmat_t = \frac{EA}{L_1}
  \begin{bmatrix}
    c^2 & sc & e_1c^2 & -c^2 & -sc & -e_2c^2\\
    sc & s^2 & e_1sc & -sc & -s^2 & -e_2sc\\
    e_1c^2 & e_1sc & e_1^2c^2 & -e_1c^2 & -e_1sc & -e_1e_2c^2\\
    -c^2 & -sc & -e_1c^2 & c^2 & sc & e_2c^2\\
    -sc & -s^2 & -e_1sc & sc & s^2 & e_2sc\\
    -e_2c^2 & -e_2sc & -e_1e_2c^2 & e_2c^2 & e_2sc & e_2^2c^2
  \end{bmatrix}
\end{equation*}
The same result can be obtained if two degrees of freedom along the member axis are used for the axial force member. Try it.
\end{example}

\begin{example}
Use the congruent transformation to derive the stiffness matrix for a beam element bent about the $y$-axis from that bent about the $z$-axis.
\begin{equation*}
  \kmat_y = \frac{EI_y}{L}
  \begin{bmatrix}
    \dfrac{12}{L^2} & -\dfrac6L & -\dfrac{12}{L^2} & -\dfrac6L\\[8pt]
    -\dfrac6L & 4 & \dfrac6L & 2\\[8pt]
    -\dfrac{12}{L^2} & \dfrac6L & \dfrac{12}{L^2} & \dfrac6L\\[8pt]
    -\dfrac6L & 2 & \dfrac6L & 4
  \end{bmatrix}\Bmat{w_1\\ \theta_{y1}\\ w_2\\ \theta_{y2}};\qquad
  \kmat_z = \frac{EI_z}{L}
  \begin{bmatrix}
    \dfrac{12}{L^2} & \dfrac6L & -\dfrac{12}{L^2} & \dfrac6L\\[8pt]
    \dfrac6L & 4 & -\dfrac6L & 2\\[8pt]
    -\dfrac{12}{L^2} & -\dfrac6L & \dfrac{12}{L^2} & -\dfrac6L\\[8pt]
    \dfrac6L & 2 & -\dfrac6L & 4
  \end{bmatrix}
\end{equation*}
\begin{equation*}
  \kmat_y = \Tmat^\tr\kmat_z\Tmat;\qquad \Tmat = ?
\end{equation*}
\end{example}

\section{Transformation matrices for all types of framed elements in FRAME21}
\begin{enumerate}[leftmargin=*]
  \item \underline{Beam element, \texttt{ITP}=1 (lies in the $x$-$y$ plane)}\\
    Degrees of freedom: $\bs{\Delta}^\tr = \langle v_1, \theta_{z1}, v_2, \theta_{z2}\rangle = \langle 2, 6, 8, 12\rangle$
    \begin{equation*}
      \Tmat = \bmat{\mat{R} & \\ & \mat{R}} = \mat{I}_{4\times4},\quad\text{with } \mat{R} = \mat{I}_{2\times2}.
    \end{equation*}

  \item \underline{2D truss element, \texttt{ITP}=2 (lies in the $x$-$y$ plane)}\\
    Degrees of freedom: $\bs{\Delta}^\tr = \langle u_1, v_1, u_2, v_2\rangle = \langle 1, 2, 7, 8\rangle$
    \begin{equation*}
      \Tmat = \bmat{\mat{R} & \\ & \mat{R}},\quad\text{with } \mat{R} = \bmat{\cos\phi & \sin\phi\\ -\sin\phi & \cos\phi}
    \end{equation*}
    where $\cos\phi = \dfrac{x_2-x_1}{L}$, $\sin\phi = \dfrac{y_2-y_1}{L}$. Note that
    $\kmat' = \dfrac{EA}{L}\bmat{1 & 0 & -1 & 0\\ 0 & 0 & 0 & 0\\ -1 & 0 & 1 & 0\\ 0 & 0 & 0 & 0}$.

  \item \underline{2D frame element, \texttt{ITP}=3 (lies in the $x$-$y$ plane)}\\
    Degrees of freedom: $\bs{\Delta}^\tr = \langle u_1, v_1, \theta_{z1}, u_2, v_2, \theta_{z2}\rangle = \langle 1, 2, 6, 7, 8, 12\rangle$
    \begin{equation*}
      \Tmat = \bmat{\mat{R} & \\ & \mat{R}},\quad\text{with } \mat{R} = \bmat{\cos\phi & \sin\phi & 0\\ -\sin\phi & \cos\phi & 0\\ 0 & 0 & 1}
    \end{equation*}

  \item \underline{Grid element, \texttt{ITP}=4 (lies in the $x$-$z$ plane)}\\
    \begin{minipage}[t]{0.6\linewidth}
      Degrees of freedom:
      $\bs{\Delta}^\tr = \langle v_1, \theta_{x1}, \theta_{z1}, v_2, \theta_{x2}, \theta_{z2}\rangle = \langle 2, 4, 6, 8, 10, 12\rangle$
      \begin{enumerate}[label=(\arabic*)]
        \item bend about the strong ($z$-) axis.
        \item can take torsion.
      \end{enumerate}
    \end{minipage}\hfill
    \begin{minipage}[t]{0.36\linewidth}
      \vspace{0pt}\centering% Grid element (ITP=4) lying in the x-z plane
\begin{tikzpicture}[scale=1.0,font=\small]
  \coordinate (O) at (0,0);
  \draw[->,thick] (O) -- (0,3.0) node[above] {$z$};
  \draw[->,thick] (O) -- (4.0,0) node[right] {$x$};
  \draw[->,dashed,thick] (O) -- (2.0,2.2) node[above right] {$y$};
  \draw[->,dashed,thick] (O) -- (1.4,1.6) node[right] {$z'$};
  \draw[->,dashed,thick] (O) -- (2.6,1.4) node[right] {$y'$};
  % member along x'
  \draw[line width=3pt,black!60] (O) -- (-35:2.8);
  \draw[->,thick] (-35:2.8) -- (-35:3.6) node[right] {$x'$};
  \node[below left,font=\footnotesize] at (-35:0.6) {1}; \node[below,font=\footnotesize] at (-35:2.8) {2};
  \draw[->] (1.6,0) arc (0:-35:1.6); \node at (1.9,-0.45) {$\phi$};
\end{tikzpicture}

    \end{minipage}
    \begin{equation*}
      \bs{\Delta}' = \Bmat{v_1'\\ \theta_{x1}'\\ \theta_{z1}'\\ v_2'\\ \theta_{x2}'\\ \theta_{z2}'} = \Tmat\bs{\Delta} = \Tmat\Bmat{v_1\\ \theta_{x1}\\ \theta_{z1}\\ v_2\\ \theta_{x2}\\ \theta_{z2}}
    \end{equation*}
    where $\Tmat = \bmat{\mat{R} & \\ & \mat{R}}$, and $\mat{R} = \bmat{1 & 0 & 0\\ 0 & \cos\phi & -\sin\phi\\ 0 & \sin\phi & \cos\phi}$. The values of $\sin\phi$ and $\cos\phi$ can be evaluated by
    \begin{equation*}
      \cos\phi = \frac{x_2-x_1}{L},\qquad \sin\phi = -\frac{z_2-z_1}{L}\quad\text{(different from 2D truss)}
    \end{equation*}

    \begin{example}[Grid structure; Sack (1989), \emph{Matrix Structural Analysis}, pp.~198--200, Example~5.4]
    Analyze the simple grid shown below. For all elements $I_y = 350$~in$^4$, $I_z = 700$~in$^4$; $J = 800$~in$^4$; $E = 29\times10^3$~kips/in$^2$; and $G = 11\times10^3$~kips/in$^2$.
    \begin{fig}
      % Simple grid (after Sack, 1989, Example 5.4): grid with loading (isometric sketch)
\begin{tikzpicture}[scale=0.9,font=\footnotesize,x={(1cm,0cm)},y={(0cm,1cm)},z={(-0.55cm,-0.42cm)}]
  \coordinate (a) at (0,0,0);
  \coordinate (b) at (4.5,0,0); \coordinate (d) at (-3.6,0,0);
  \coordinate (e) at (0,0,-3.6); \coordinate (c) at (0,0,5.0);
  % members
  \draw[line width=2.4pt,black!65] (d) -- (b);
  \draw[line width=2.4pt,black!65] (e) -- (c);
  % fixed ends
  \foreach \p in {b,d}{\fill[hatch] ($(\p)+(-0.15,-0.25,0)$) -- ($(\p)+(0.15,-0.25,0)$) -- ($(\p)+(0.15,0.55,0)$) -- ($(\p)+(-0.15,0.55,0)$) -- cycle;}
  \foreach \p in {e,c}{\fill[hatch] ($(\p)+(0,-0.25,-0.15)$) -- ($(\p)+(0,-0.25,0.15)$) -- ($(\p)+(0,0.55,0.15)$) -- ($(\p)+(0,0.55,-0.15)$) -- cycle;}
  % distributed load on ad
  \foreach \t in {0.12,0.3,...,0.95}{\draw[->] (-3.6*\t,1.0,0) -- (-3.6*\t,0.1,0);}
  \draw (-3.6*0.12,1.0,0) -- (-3.6*0.95,1.0,0);
  \node[above] at (-1.9,1.0,0) {10 k/ft};
  % distributed load on ac
  \foreach \t in {0.1,0.22,...,0.95}{\draw[->] (0,1.0,5.0*\t) -- (0,0.1,5.0*\t);}
  \draw (0,1.0,5.0*0.1) -- (0,1.0,5.0*0.95);
  \node[left] at (0,1.15,3.0) {10 k/ft};
  % node labels
  \node[above right] at (a) {$a$}; \node[above] at ($(b)+(0,0.6,0)$) {$b$}; \node[right] at ($(c)+(0.2,0,0)$) {$c$};
  \node[above] at ($(d)+(0,0.6,0)$) {$d$}; \node[above] at ($(e)+(0,0.6,0)$) {$e$};
  % axes
  \draw[->] (a) -- (0,2.2,0) node[above] {$y$};
  \draw[->] ($(b)+(0.3,0,0)$) -- (6.0,0,0) node[right] {$x$};
  \draw[->] ($(c)+(0,0,0.3)$) -- (0,0,6.6) node[below left] {$z$};
  % dimensions (below the members)
  \node[below,inner sep=5pt] at (2.25,0,0) {18 ft}; \node[below,inner sep=5pt] at (-1.8,0,0) {12 ft};
  \node[right,inner sep=5pt] at (0,0,2.6) {20 ft}; \node[right,inner sep=5pt] at (0,0,-1.9) {12 ft};
\end{tikzpicture}

    \end{fig}
    \begin{solution}
    The element stiffness matrices are obtained using Eq.~(5.37) of Sack. Assuming axial inextensibility of all elements, we need consider only $V_a$, $\Theta_{xa}$, and $\Theta_{za}$. Note that this simplifies calculation of relevant element stiffness entries (we urge the reader to calculate the appropriate 3 by 3 partition for the four element stiffness matrices). Invoking nodal equilibrium for the structure with support conditions enforced yields (the reader is encouraged to carry out this process)
    \begin{equation*}
      \Bmat{-160\\ 4000\\ 1440} = \bmat{264.35 & 6343.8 & -3263.2\\ 6343.8 & 1116900 & 00000\\ -3263.2 & 00000 & 1049800}\Bmat{V_a\\ \Theta_{xa}\\ \Theta_{za}}
    \end{equation*}
    (with forces in kips and moments in kip$\cdot$in). Hence
    \begin{equation*}
      \Bmat{V_a\\ \Theta_{xa}\\ \Theta_{za}} = \Bmat{-0.8170\text{ in.}\\ 8.222\times10^{-3}\text{ rad}\\ -1.168\times10^{-3}\text{ rad}}
    \end{equation*}
    For elements $ab$ and $ae$ the element forces are obtained using Eq.~(5.42) of Sack; thus
    \begin{equation*}
      \Bmat{\bar p_{\bar ya}\\ \bar m_{\bar xa}\\ \bar m_{\bar za}\\ \bar m_{\bar zb}}^{ab} = \mat{e}\mat{u}^{ab} = \Bmat{-22.8\text{ kips}\\ +335\text{ kip}\cdot\text{in.}\\ -2570\text{ kip}\cdot\text{in.}\\ -2350\text{ kip}\cdot\text{in.}}
      \quad\text{and}\quad
      \Bmat{\bar p_{\bar ya}\\ \bar m_{\bar xa}\\ \bar m_{\bar za}\\ \bar m_{\bar ze}}^{ae} = \mat{e}\mat{u}^{ae} = \Bmat{-45.6\text{ kips}\\ +85.6\text{ kip}\cdot\text{in.}\\ -1350\text{ kip}\cdot\text{in.}\\ -4130\text{ kip}\cdot\text{in.}}
    \end{equation*}
    Elements $ac$ and $ad$ have initial forces applied; therefore using Eq.~(4.65) of Sack we have
    \begin{equation*}
      \Bmat{\bar p_{\bar ya}\\ \bar p_{\bar yc}\\ \bar m_{\bar xa}\\ \bar m_{\bar za}\\ \bar m_{\bar zc}}^{ac} = \mat{e}\mat{u}^{ac}+\bar{\mat{s}}^o
      = \Bmat{-31.8\\ +31.8\\ 42.8\\ -4510\\ -3120} + \Bmat{100\\ 100\\ 0\\ 4000\\ -4000}
      = \Bmat{68.2\text{ kips}\\ 131.8\text{ kips}\\ -42.8\text{ kip}\cdot\text{in.}\\ -510\text{ kip}\cdot\text{in.}\\ -7120\text{ kip}\cdot\text{in.}}
    \end{equation*}
    and
    \begin{equation*}
      \Bmat{\bar p_{\bar ya}\\ \bar p_{\bar yd}\\ \bar m_{\bar xa}\\ \bar m_{\bar za}\\ \bar m_{\bar zd}}^{ad} = \mat{e}\mat{u}^{ad}+\bar{\mat{s}}^o
      = \Bmat{-59.8\\ +59.8\\ -502\\ -4140\\ 4470} + \Bmat{60\\ 60\\ 0\\ 1440\\ -1440}
      = \Bmat{+0.2\text{ kips}\\ 119.8\text{ kips}\\ -502\text{ kip}\cdot\text{in.}\\ -2700\text{ kip}\cdot\text{in.}\\ -5910\text{ kip}\cdot\text{in.}}
    \end{equation*}
    \begin{figure}[H]
      \centering
      % Free-body diagrams of the grid members and node a, drawn in plan view
% (forces in kips, moments in kip-in).  Vertical (y) forces: dot = up (+y), cross = down.
% Double-headed arrows = moment vectors (about x or z).
\begin{tikzpicture}[scale=0.8,font=\scriptsize,
  mom/.style={->,thick,double,double distance=1pt},
  up/.pic={\draw[fill=white] (0,0) circle (0.13); \fill (0,0) circle (0.04);},
  down/.pic={\draw[fill=white] (0,0) circle (0.13); \draw (-0.09,-0.09) -- (0.09,0.09); \draw (-0.09,0.09) -- (0.09,-0.09);}]
  % plan: x to the right, z downwards on the page
  \draw[->] (-7.2,3.6) -- (-6.2,3.6) node[right] {$x$}; \draw[->] (-7.2,3.6) -- (-7.2,2.6) node[below] {$z$};
  % ---------------- node a ----------------
  \draw[line width=1.5pt] (-0.5,0) -- (0.5,0); \draw[line width=1.5pt] (0,-0.5) -- (0,0.5);
  \fill (0,0) circle (0.05); \node[above left] at (0,0) {$a$};
  % from ab (right)
  \pic at (1.4,-0.5) {up}; \node[below] at (1.4,-0.65) {22.8};
  \draw[mom] (0.6,0) -- (2.0,0) node[right] {335};
  \draw[mom] (1.4,0.25) -- (1.4,1.3) node[above] {2570};
  % from ad (left)
  \pic at (-1.4,-0.5) {down}; \node[below] at (-1.4,-0.65) {0.2};
  \draw[mom] (-0.6,0) -- (-2.0,0) node[left] {502};
  \draw[mom] (-1.4,0.25) -- (-1.4,1.3) node[above] {2700};
  % from ae (top, -z)
  \pic at (-0.5,1.9) {up}; \node[left] at (-0.65,1.9) {45.6};
  \draw[mom] (0,0.6) -- (0,2.4) node[above] {85.6};
  \draw[mom] (0.25,2.1) -- (1.4,2.1) node[right] {1350};
  % from ac (bottom, +z)
  \pic at (-0.5,-1.9) {down}; \node[left] at (-0.65,-1.9) {68.2};
  \draw[mom] (0,-0.6) -- (0,-2.4) node[below] {42.8};
  \draw[mom] (0.25,-2.1) -- (1.4,-2.1) node[right] {510};
  % ---------------- member ab ----------------
  \begin{scope}[shift={(4.8,0)}]
    \draw[line width=1.5pt,black!70] (0,0) -- (3.4,0); \node[below] at (1.7,-0.1) {$ab$};
    \pic at (0,0.45) {down}; \node[above] at (0,0.6) {22.8};
    \draw[mom] (-0.15,0) -- (-1.0,0) node[left] {335};
    \draw[mom] (0.3,-0.2) -- (0.3,-1.0) node[below] {2570};
    \pic at (3.4,0.45) {up}; \node[above] at (3.4,0.6) {22.8};
    \draw[mom] (3.55,0) -- (4.4,0) node[right] {335};
    \draw[mom] (3.1,-0.2) -- (3.1,-1.0) node[below] {2350};
  \end{scope}
  % ---------------- member ad ----------------
  \begin{scope}[shift={(-4.8,0)}]
    \draw[line width=1.5pt,black!70] (0,0) -- (-3.4,0); \node[below] at (-1.7,-0.1) {$ad$};
    \pic at (0,0.45) {up}; \node[above] at (0,0.6) {0.2};
    \draw[mom] (0.15,0) -- (1.0,0) node[right] {502};
    \draw[mom] (-0.3,-0.2) -- (-0.3,-1.0) node[below] {2700};
    \pic at (-3.4,0.45) {up}; \node[above] at (-3.4,0.6) {119.8};
    \draw[mom] (-3.55,0) -- (-4.4,0) node[left] {502};
    \draw[mom] (-3.1,-0.2) -- (-3.1,-1.0) node[below] {5910};
  \end{scope}
  % ---------------- member ae (towards -z, i.e. up the page) ----------------
  \begin{scope}[shift={(0,4.6)}]
    \draw[line width=1.5pt,black!70] (0,0) -- (0,2.6); \node[right] at (0.1,1.3) {$ae$};
    \pic at (-0.45,0) {down}; \node[left] at (-0.6,0) {45.6};
    \draw[mom] (0,-0.15) -- (0,-0.9) node[below] {85.6};
    \draw[mom] (0.2,0.3) -- (1.1,0.3) node[right] {1350};
    \pic at (-0.45,2.6) {up}; \node[left] at (-0.6,2.6) {45.6};
    \draw[mom] (0,2.75) -- (0,3.5) node[above] {85.6};
    \draw[mom] (0.2,2.3) -- (1.1,2.3) node[right] {4130};
  \end{scope}
  % ---------------- member ac (towards +z, i.e. down the page) ----------------
  \begin{scope}[shift={(0,-4.8)}]
    \draw[line width=1.5pt,black!70] (0,0) -- (0,-3.2); \node[right] at (0.1,-1.6) {$ac$};
    \pic at (-0.45,0) {up}; \node[left] at (-0.6,0) {68.2};
    \draw[mom] (0,0.15) -- (0,0.9) node[above] {42.8};
    \draw[mom] (0.2,-0.3) -- (1.1,-0.3) node[right] {510};
    \pic at (-0.45,-3.2) {up}; \node[left] at (-0.6,-3.2) {131.8};
    \draw[mom] (0,-3.35) -- (0,-4.1) node[below] {42.8};
    \draw[mom] (0.2,-2.9) -- (1.1,-2.9) node[right] {7120};
  \end{scope}
  % legend
  \node[anchor=west,align=left] at (4.2,-5.4) {\tikz\pic at (0,0) {up}; \ force $+y$ (up)\\[2pt]\tikz\pic at (0,0) {down}; \ force $-y$ (down)\\[2pt]\tikz\draw[mom] (0,0)--(0.5,0); \ moment vector};
\end{tikzpicture}

      \caption{Free-body diagrams with forces in kips and moments in kip-inches (after Fig.~E5.4b of Sack, 1989).}
    \end{figure}
    \emph{Discussion.} Note that only the independent element forces were calculated using Eq.~(5.42). The free-body diagrams for the four elements, plus that for node $a$, are shown above. We urge the reader to check each of these for equilibrium. The stiffness method ensures element and nodal equilibrium, but it is a good idea to perform spot checks to guard against misinterpreting the results.
    \end{solution}
    \end{example}

    \newpage\topic{Same example, with the $y$ and $z$ axes opposite to the direction shown in the textbook.}
    \begin{minipage}{0.5\linewidth}
      \begin{equation*}
        [T] = \bmat{1 & 0 & 0\\ 0 & \cos\phi & -\sin\phi\\ 0 & \sin\phi & \cos\phi}
      \end{equation*}
    \end{minipage}\hfill
    \begin{minipage}{0.4\linewidth}\centering
      % Grid: y and z axes opposite to the textbook (z up, member rotated clockwise by phi)
\begin{tikzpicture}[scale=0.9,font=\small]
  \draw[->,thick] (0,0) -- (0,2.6) node[above] {$z$};
  \draw[->,thick] (0,0) -- (2.6,0) node[right] {$x$};
  \draw[->,dashed,thick] (0,0) -- (1.6,2.1) node[above right] {$z'$};
  \draw[line width=3pt,black!60] (0,0) -- (-37:2.0);
  \draw[->,thick] (-37:2.0) -- (-37:2.6) node[right] {$x'$};
  \draw[->] (1.3,0) arc (0:-37:1.3); \node at (1.55,-0.4) {$\phi$};
\end{tikzpicture}

    \end{minipage}
    \begin{equation*}
      \kmat(E,G,I,J,L,\phi) = [T]^\tr[k_l][T] =
      \bmat{1 & 0 & 0\\ 0 & \cos\phi & \sin\phi\\ 0 & -\sin\phi & \cos\phi}
      \begin{bmatrix}\frac{12EI}{L^3} & 0 & \frac{6EI}{L^2}\\[3pt] 0 & \frac{GJ}{L} & 0\\[3pt] \frac{6EI}{L^2} & 0 & \frac{4EI}{L}\end{bmatrix}
      \bmat{1 & 0 & 0\\ 0 & \cos\phi & -\sin\phi\\ 0 & \sin\phi & \cos\phi}
    \end{equation*}
    \begin{equation*}
      \kmat(E,G,I,J,L,\phi) =
      \begin{bmatrix}
        12E\frac{I}{L^3} & 6sE\frac{I}{L^2} & 6cE\frac{I}{L^2}\\[4pt]
        6sE\frac{I}{L^2} & c^2G\frac JL + 4s^2E\frac IL & -scG\frac JL + 4scE\frac IL\\[4pt]
        6cE\frac{I}{L^2} & -scG\frac JL + 4scE\frac IL & s^2G\frac JL + 4c^2E\frac IL
      \end{bmatrix},\qquad c = \cos\phi,\ s = \sin\phi
    \end{equation*}
    \begin{align*}
      ab:\ \rightarrow x'\quad \kmat(29000, 11000, 700, 800, 216, 0) &= \bmat{24.172 & 0 & 2610.6\\ 0 & 40741. & 0\\ 2610.6 & 0 & 3.7593\times10^5}\\
      ae:\ x'\uparrow\quad \kmat(29000, 11000, 700, 800, 120, -0.5\pi) &= \bmat{140.97 & -8458.3 & 0\\ -8458.3 & 6.7667\times10^5 & 0\\ 0 & 0 & 73333.}\\
      ad:\ x'\leftarrow\quad \kmat(29000, 11000, 700, 800, 144, -\pi) &= \bmat{81.581 & 0 & -5873.8\\ 0 & 61111. & 0\\ -5873.8 & 0 & 5.6389\times10^5}\\
      ac:\ x'\downarrow\quad \kmat(29000, 11000, 700, 800, 240, -1.5\pi) &= \bmat{17.622 & 2114.6 & 0\\ 2114.6 & 3.3833\times10^5 & 0\\ 0 & 0 & 36667.}\\
      \text{sum of }[k] &= \bmat{264.35 & -6343.7 & -3263.2\\ -6343.7 & 1.1169\times10^6 & 0\\ -3263.2 & 0 & 1.0498\times10^6}
    \end{align*}
    Note that the sign for $k_{12}$ is different from that given in the textbook. Why?

    \topic{Same example, with the $y$ and $z$ axes the same as those shown in the textbook.}
    \begin{minipage}{0.5\linewidth}
      \begin{equation*}
        [T] = \bmat{1 & 0 & 0\\ 0 & \cos\phi & -\sin\phi\\ 0 & \sin\phi & \cos\phi}
      \end{equation*}
    \end{minipage}\hfill
    \begin{minipage}{0.4\linewidth}\centering
      % Grid: y and z axes as in the textbook (z down)
\begin{tikzpicture}[scale=0.9,font=\small]
  \draw[->,thick] (0,0) -- (0,-2.6) node[below] {$z$};
  \draw[->,thick] (0,0) -- (2.6,0) node[right] {$x$};
  \draw[->,dashed,thick] (0,0) -- (1.9,-1.6) node[below right] {$z'$};
  \draw[line width=3pt,black!60] (0,0) -- (40:2.0);
  \draw[->,thick] (40:2.0) -- (40:2.6) node[above right] {$x'$};
  \draw[->] (1.3,0) arc (0:40:1.3); \node at (1.55,0.45) {$\phi$};
\end{tikzpicture}

    \end{minipage}
    \begin{equation*}
      \kmat(E,G,I,J,L,\phi) = [T]^\tr[k_l][T] \quad\text{(same expression as above)}
    \end{equation*}
    \begin{align*}
      ab:\ \rightarrow x'\quad \kmat(29000, 11000, 700, 800, 216, 0) &= \bmat{24.172 & 0 & 2610.6\\ 0 & 40741. & 0\\ 2610.6 & 0 & 3.7593\times10^5}\\
      ae:\ x'\uparrow\quad \kmat(29000, 11000, 700, 800, 120, 0.5\pi) &= \bmat{140.97 & 8458.3 & 0\\ 8458.3 & 6.7667\times10^5 & 0\\ 0 & 0 & 73333.}\\
      ad:\ x'\leftarrow\quad \kmat(29000, 11000, 700, 800, 144, \pi) &= \bmat{81.581 & 0 & -5873.8\\ 0 & 61111. & 0\\ -5873.8 & 0 & 5.6389\times10^5}\\
      ac:\ x'\downarrow\quad \kmat(29000, 11000, 700, 800, 240, 1.5\pi) &= \bmat{17.622 & -2114.6 & 0\\ -2114.6 & 3.3833\times10^5 & 0\\ 0 & 0 & 36667.}\\
      \text{sum of }[k] &= \bmat{264.35 & 6343.7 & -3263.2\\ 6343.7 & 1.1169\times10^6 & 0\\ -3263.2 & 0 & 1.0498\times10^6}
    \end{align*}

  \item \underline{3D truss element, \texttt{ITP}=5}\\
    Degrees of freedom: $\bs{\Delta}^\tr = \langle u_1, v_1, w_1, u_2, v_2, w_2\rangle = \langle 1, 2, 3, 7, 8, 9\rangle$
    \begin{equation*}
      \Tmat = \bmat{\mat{R} & \\ & \mat{R}},\quad\text{with } \mat{R} = \Bmat{\mat{x}'\\ \mat{y}'\\ \mat{z}'},
    \end{equation*}
    where $\mat{x}' = \Bigl\langle\dfrac{x_2-x_1}{L}\ \ \dfrac{y_2-y_1}{L}\ \ \dfrac{z_2-z_1}{L}\Bigr\rangle$; $\mat{y}'$ and $\mat{z}'$ can be any vector. Note that
    \begin{equation*}
      \kmat' = \frac{EA}{L}\bmat{1 & 0 & 0 & -1 & 0 & 0\\ 0 & 0 & 0 & 0 & 0 & 0\\ 0 & 0 & 0 & 0 & 0 & 0\\ -1 & 0 & 0 & 1 & 0 & 0\\ 0 & 0 & 0 & 0 & 0 & 0\\ 0 & 0 & 0 & 0 & 0 & 0}.
    \end{equation*}

  \item \underline{3D frame element, \texttt{ITP}=6}\\
    Degrees of freedom: $\bs{\Delta}^\tr = \langle u_1, v_1, w_1, \theta_{x1}, \theta_{y1}, \theta_{z1}, u_2, v_2, w_2, \theta_{x2}, \theta_{y2}, \theta_{z2}\rangle$
    \begin{equation*}
      \Tmat = \bmat{\mat{R} & & & \\ & \mat{R} & & \\ & & \mat{R} & \\ & & & \mat{R}},\quad\text{with } \mat{R} = \Bmat{\mat{x}'\\ \mat{y}'\\ \mat{z}'}
    \end{equation*}
    where $\mat{x}' = \Bigl\langle\dfrac{x_2-x_1}{L}\ \ \dfrac{y_2-y_1}{L}\ \ \dfrac{z_2-z_1}{L}\Bigr\rangle$, $\mat{y}'$ is the unit vector of one of the principal axes $=$ \texttt{VECTY(3,NBC)} in the program, and $\mat{z}' = \mat{x}'\times\mat{y}'$.

    \begin{minipage}[t]{0.55\linewidth}
      \vspace{0pt}
      Consider a space beam element lying along the $Y$ axis as shown.
      \begin{itemize}
        \item node 1: $(0,0,0)$
        \item node 2: $(0,L,0)$
        \item Therefore, $x' = (0,1,0)$.
      \end{itemize}
    \end{minipage}\hfill
    \begin{minipage}[t]{0.42\linewidth}
      \vspace{0pt}\centering% Space beam element lying along the global Y axis, with four choices of y'
\begin{tikzpicture}[scale=0.62,font=\small,x={(-0.75cm,-0.45cm)},y={(1cm,0cm)},z={(0cm,1cm)}]
  % axes
  \draw[->,thick] (0,0,0) -- (2.4,0,0) node[left] {$x$};
  \draw[->,thick] (0,0,0) -- (0,5.6,0) node[right] {$y$};
  \draw[->,thick] (0,0,0) -- (0,0,3.4) node[above] {$z$};
  % member box from y=0.4 to y=5 with section 0.5 (x) by 1.0 (z)
  \def\ya{0.6} \def\yb{4.8} \def\hx{0.35} \def\hz{0.6}
  \draw (\hx,\ya,\hz) -- (\hx,\yb,\hz) -- (\hx,\yb,-\hz) -- (\hx,\ya,-\hz) -- cycle;
  \draw (-\hx,\ya,\hz) -- (-\hx,\yb,\hz); \draw (\hx,\ya,\hz) -- (-\hx,\ya,\hz); \draw (\hx,\yb,\hz) -- (-\hx,\yb,\hz);
  \draw (-\hx,\yb,\hz) -- (-\hx,\yb,-\hz) -- (\hx,\yb,-\hz);
  \draw[dashed] (0,\ya,0) -- (0,\yb,0);
  \node[left] at (0,\ya,\hz+0.2) {1}; \node[below] at (0,\yb,-\hz-0.1) {2};
  \draw[->] (0,\yb+0.3,0) -- (0,\yb+0.3,1.0) node[right,font=\footnotesize] {$I_2 = 2$};
  \draw[->] (0,\yb+0.3,0) -- (-1.2,\yb+0.3,0) node[right,font=\footnotesize] {$I_1 = 10$};
\end{tikzpicture}

    \end{minipage}

    The following four cases can be distinguished with reference to the selection of the $y'$ axis:
    \begin{align*}
      \text{Case 1: } & y' = (-1,0,0);\ z' = (0,0,1) & (I_{y'},I_{z'}) &= (I_1,I_2) = (10,2)\\
      \text{Case 2: } & y' = (1,0,0);\ z' = (0,0,-1) & (I_{y'},I_{z'}) &= (I_1,I_2) = (10,2)\\
      \text{Case 3: } & y' = (0,0,1);\ z' = (1,0,0) & (I_{y'},I_{z'}) &= (I_2,I_1) = (2,10)\\
      \text{Case 4: } & y' = (0,0,-1);\ z' = (-1,0,0) & (I_{y'},I_{z'}) &= (I_2,I_1) = (2,10)
    \end{align*}
    \begin{fig}
      % Four choices of the local y' axis for a beam along the global Y axis
\begin{tikzpicture}[scale=0.75,font=\small]
  \newcommand{\block}{%
    \draw (0,0) rectangle (0.5,1.2); \draw (0,1.2) -- (-0.6,1.6) -- (-0.1,1.6) -- (0.5,1.2); \draw (-0.6,1.6) -- (-0.6,0.4) -- (0,0);}
  \begin{scope}
    \block
    \draw[->,thick] (0.25,0.6) -- (0.25,2.3) node[above] {$z'$};
    \draw[->,thick] (0.25,0.6) -- (1.9,1.4) node[right] {$y'$};
    \draw[->,thick] (0.25,0.6) -- (1.9,0.1) node[right] {$x'$};
    \node at (0.6,-2.4) {Case 1};
  \end{scope}
  \begin{scope}[xshift=5.2cm]
    \block
    \draw[->,thick] (0.25,0.6) -- (0.25,-1.0) node[below] {$z'$};
    \draw[->,thick] (0.25,0.6) -- (-1.4,-0.2) node[left] {$y'$};
    \draw[->,thick] (0.25,0.6) -- (1.9,0.1) node[right] {$x'$};
    \node at (0.6,-2.4) {Case 2};
  \end{scope}
  \begin{scope}[xshift=10.4cm]
    \block
    \draw[->,thick] (0.25,0.6) -- (0.25,2.3) node[above] {$y'$};
    \draw[->,thick] (0.25,0.6) -- (-1.4,-0.2) node[left] {$z'$};
    \draw[->,thick] (0.25,0.6) -- (1.9,0.1) node[right] {$x'$};
    \node at (0.6,-2.4) {Case 3};
  \end{scope}
  \begin{scope}[xshift=15.6cm]
    \block
    \draw[->,thick] (0.25,0.6) -- (0.25,-1.0) node[below] {$y'$};
    \draw[->,thick] (0.25,0.6) -- (1.9,1.4) node[right] {$z'$};
    \draw[->,thick] (0.25,0.6) -- (1.9,0.1) node[right] {$x'$};
    \node at (0.6,-2.4) {Case 4};
  \end{scope}
\end{tikzpicture}

    \end{fig}
    The global stiffness matrix $\kmat = \Tmat^\tr\kmat'\Tmat$ is the same for the above four cases as expected and is given by
    {\scriptsize
    \begin{equation*}
      \kmat(E,\nu,A,I_1,I_2,J,L) = E\,
      \setlength{\arraycolsep}{1.6pt}
      \begin{bmatrix}
        12\frac{I_2}{L^3} & 0 & 0 & 0 & 0 & -6\frac{I_2}{L^2} & -12\frac{I_2}{L^3} & 0 & 0 & 0 & 0 & -6\frac{I_2}{L^2}\\[3pt]
        0 & \frac AL & 0 & 0 & 0 & 0 & 0 & -\frac AL & 0 & 0 & 0 & 0\\[3pt]
        0 & 0 & 12\frac{I_1}{L^3} & 6\frac{I_1}{L^2} & 0 & 0 & 0 & 0 & -12\frac{I_1}{L^3} & 6\frac{I_1}{L^2} & 0 & 0\\[3pt]
        0 & 0 & 6\frac{I_1}{L^2} & 4\frac{I_1}{L} & 0 & 0 & 0 & 0 & -6\frac{I_1}{L^2} & 2\frac{I_1}{L} & 0 & 0\\[3pt]
        0 & 0 & 0 & 0 & \frac12\frac{J}{(1+\nu)L} & 0 & 0 & 0 & 0 & 0 & -\frac12\frac{J}{(1+\nu)L} & 0\\[3pt]
        -6\frac{I_2}{L^2} & 0 & 0 & 0 & 0 & 4\frac{I_2}{L} & 6\frac{I_2}{L^2} & 0 & 0 & 0 & 0 & 2\frac{I_2}{L}\\[3pt]
        -12\frac{I_2}{L^3} & 0 & 0 & 0 & 0 & 6\frac{I_2}{L^2} & 12\frac{I_2}{L^3} & 0 & 0 & 0 & 0 & 6\frac{I_2}{L^2}\\[3pt]
        0 & -\frac AL & 0 & 0 & 0 & 0 & 0 & \frac AL & 0 & 0 & 0 & 0\\[3pt]
        0 & 0 & -12\frac{I_1}{L^3} & -6\frac{I_1}{L^2} & 0 & 0 & 0 & 0 & 12\frac{I_1}{L^3} & -6\frac{I_1}{L^2} & 0 & 0\\[3pt]
        0 & 0 & 6\frac{I_1}{L^2} & 2\frac{I_1}{L} & 0 & 0 & 0 & 0 & -6\frac{I_1}{L^2} & 4\frac{I_1}{L} & 0 & 0\\[3pt]
        0 & 0 & 0 & 0 & -\frac12\frac{J}{(1+\nu)L} & 0 & 0 & 0 & 0 & 0 & \frac12\frac{J}{(1+\nu)L} & 0\\[3pt]
        -6\frac{I_2}{L^2} & 0 & 0 & 0 & 0 & 2\frac{I_2}{L} & 6\frac{I_2}{L^2} & 0 & 0 & 0 & 0 & 4\frac{I_2}{L}
      \end{bmatrix}
    \end{equation*}}
\end{enumerate}

\chapter{Solution of Linear Algebraic Equations}
\label{ch:solve}

\section{Methods for solving linear algebraic equations \texorpdfstring{$\mat{A}\mat{x}=\mat{b}$}{Ax=b}}
\begin{enumerate}
  \item Gauss elimination
  \item Decomposition methods
  \item Iterative methods: such as Gauss--Seidel and Jacobi methods and conjugate gradient method (Sec.~11.3 of textbook).
\end{enumerate}
We shall discuss only the first two methods because they are the most commonly used methods.

\section{Gauss elimination method}
The main idea of Gauss elimination is to reduce the original system $\mat{A}\mat{x} = \mat{b}$ to an upper triangular system through appropriate \emph{row operations}; that is, the desired result is
\begin{equation*}
  \mat{U}\mat{x} = \mat{c}
\end{equation*}
where $\mat{U}$ is an upper triangular matrix. This system can then be solved to obtain $\mat{x}$ using the backward substitution. Before explaining the algorithm, let's consider the following simple example.

\begin{example}
Solving the linear system with Gauss elimination.
\begin{alignat*}{2}
  3x_1 - 6x_2 + 7x_3 &= 3 &\qquad&\text{(a)}\\
  9x_1 + 0 - 5x_3 &= 3 &&\text{(b)}\\
  5x_1 - 8x_2 + 6x_3 &= -4 &&\text{(c)}
\end{alignat*}
First, from (a): $x_1 = (6x_2-7x_3+3)/3 \;\rightarrow\;$ (b), (c). Then, we have
\begin{alignat*}{2}
  18x_2 - 26x_3 &= -6 &\qquad&\text{(b$'$)}\\
  2x_2 - \tfrac{17}{3}x_3 &= -9 &&\text{(c$'$)}
\end{alignat*}
Second, from (b$'$): $x_2 = \frac19(-3+13x_3) \;\rightarrow\;$ (c$'$). Then, we have
\begin{alignat*}{2}
  -\tfrac{25}{9}x_3 &= -\tfrac{25}{3} &\qquad&\text{(c$''$)}
\end{alignat*}
Rewrite these results [(a), (b$'$), and (c$''$)] in matrix form:
\begin{equation*}
  \bmat{3 & -6 & 7\\ 0 & 18 & -26\\ 0 & 0 & -\frac{25}{9}}\Bmat{x_1\\ x_2\\ x_3} = \Bmat{3\\ -6\\ -\frac{25}{3}}
\end{equation*}
The above procedure is called \emph{forward elimination}, which results in an upper triangular system $\mat{U}\mat{x} = \mat{c}$.

The upper triangular system can be solved easily using the \emph{backward substitution} as described below. First, obtain $x_3 = 3$ from (c$''$). Then substituting $x_3 = 3$ into (b$'$) gives $x_2 = 4$. Finally, substitute $x_3 = 3$ and $x_2 = 4$ into (a) to obtain $x_1 = 2$.
\end{example}

\topic{Algorithm for upper triangularization: $\mat{A}\mat{x}=\mat{b}\ \rightarrow\ \mat{U}\mat{x}=\mat{c}$.}
\begin{lstlisting}[language=Fortran,numbers=none,basicstyle=\small\ttfamily\itshape,mathescape=true]
do m=1, n-1
    do i=m+1, n
        r=$a_{im}/a_{mm}$
        $b_i$=$b_i$-r$b_m$
        do j=m, n
            $a_{ij}$=$a_{ij}$ - r $a_{mj}$
        end do
    end do
end do
\end{lstlisting}
It can be represented by the following graph.
\begin{fig}
  % Graphical representation of the upper triangularization step m
\begin{tikzpicture}[scale=0.85,font=\small]
  \def\n{8}
  % matrix frame
  \draw (0,0) rectangle (\n,\n);
  % processed upper-left triangle: X on the diagonal (rows 1..m-1)
  \foreach \i in {1,...,5}{\node at (\i-0.5,\n-\i+0.5) {X};}
  \draw[dashed] (0.6,\n-0.6) -- (2.7,\n-2.7); \node at (3.0,\n-3.0) {$\cdot$}; \node at (3.3,\n-3.3) {$\cdot$};
  % dashed region of zeros below the processed part
  \draw[dashed] (0.3,\n-0.3) -- (0.3,0.3) -- (3.7,0.3) -- (3.7,\n-3.7);
  \node[font=\Large] at (1.6,3.6) {$0$};
  % pivot row m and column m, column j, row i
  \pgfmathsetmacro{\m}{\n-5.5} \pgfmathsetmacro{\ii}{\n-7.0}
  \draw[dotted] (0,\m) -- (\n,\m); \node[left] at (0,\m) {$m$};
  \draw[dotted] (0,\m-0.35) -- (3.7,\m-0.35); \node[left] at (0,\m-0.35) {$m{+}1$};
  \draw[dotted] (0,\ii) -- (\n,\ii); \node[left] at (0,\ii) {$i$};
  \draw[dotted] (3.7,0) -- (3.7,\n); \node[above] at (3.7,\n) {$m$};
  \draw[dotted] (6.3,0) -- (6.3,\n); \node[above] at (6.3,\n) {$j$};
  \node[left] at (0,0.2) {$n$};
  \node[above right,font=\footnotesize] at (3.7,\m) {$a_{mm}$}; \node[above right,font=\footnotesize] at (6.3,\m) {$a_{mj}$};
  \node[above right,font=\footnotesize] at (3.7,\ii) {$a_{im}$}; \node[above right,font=\footnotesize] at (6.3,\ii) {$a_{ij}$};
  \fill (3.7,\m) circle (0.05); \fill (6.3,\m) circle (0.05); \fill (3.7,\ii) circle (0.05); \fill (6.3,\ii) circle (0.05);
  % right-hand side vector
  \draw (\n+0.8,0) rectangle (\n+1.8,\n);
  \node[font=\footnotesize] at (\n+1.3,\m) {$b_m$}; \node[font=\footnotesize] at (\n+1.3,\ii) {$b_i$};
  \draw[->] (\n+1.8,\m) -- (\n+2.6,\m) -- (\n+2.6,\ii) -- (\n+1.85,\ii);
  \node[right,font=\footnotesize] at (\n+2.6,{(\m+\ii)/2}) {$\times\Bigl(-\dfrac{a_{im}}{a_{mm}}\Bigr)$};
\end{tikzpicture}

\end{fig}

\topic{Algorithm for backward substitution: $\mat{U}\mat{x}=\mat{c}\ \rightarrow\ \mat{x}$}
\begin{lstlisting}[language=Fortran,numbers=none,basicstyle=\small\ttfamily\itshape,mathescape=true]
$b_n$= $b_n/a_{nn}$
do i=n-1, 1, -1
    $b_i$=($b_i$- $\sum_{j=i+1}^{n} a_{ij}b_j$ )/$a_{ii}$
end do
\end{lstlisting}

If $\mat{A}$ is symmetric and assume that only its upper triangular coefficients $a_{ij}$ ($j\geq i$) are stored, then the above algorithm for upper triangularization needs to be modified to be of the form:
\begin{lstlisting}[language=Fortran,numbers=none,basicstyle=\small\ttfamily\itshape,mathescape=true]
do m=1, n-1
    do i=m+1, n
        r=$\boldsymbol{a_{mi}}$/$a_{mm}$
        $b_i$= $b_i$-r$b_m$
        do j=$\boldsymbol{i}$, n
            $a_{ij}$=$a_{ij}$ - r $a_{mj}$
        end do
    end do
end do
\end{lstlisting}
Actually, there are only two slight changes; the reasons for these two changes are as follows:
\begin{enumerate}
  \item $a_{im}\ \rightarrow\ a_{mi}$ because $\mat{A}$ is symmetric and its coefficients are available for the upper triangular part.
  \item \texttt{do j=m,n} $\rightarrow$ \texttt{do j=i,n} because $a_{ij} = a_{ji}$, ($j = m, i-1$) when $\mat{A}$ is symmetric even during the upper triangularization process; you may verify this by carrying out an example.
\end{enumerate}
Note that in the process of upper triangularization, if $a_{mm} = 0$, then the above algorithms will fail obviously. In this case, we need to interchange the $m$-th row and a certain row after it, say the $i$-th row, where $a_{im}\neq 0$; there are several ways to do it and this process is called \emph{pivoting}.

\section{Some widely used solvers}
\begin{enumerate}
  \item Banded solver, which is discussed in class and used in our program.
  \item Skyline solver, which will be introduced in the course of Finite Element Methods next semester.
  \item Frontal solver (see Sec.~11.5 of textbook).
\end{enumerate}

\section{Condition number}
$\operatorname{Cond}(\mat{A}) =$ condition number of $\mat{A} = \|\mat{A}^{-1}\|\,\|\mat{A}\|$; when $\mat{A}$ is symmetric such as the structural stiffness matrix $\Kmat$ and $\|\ \| = L_2$ norm, then it can be shown that
\begin{equation}
  \operatorname{cond}(\Kmat) = \frac{\max|\lambda_i|}{\min|\lambda_i|}
\end{equation}
where $\lambda_i$ are eigenvalues of $\Kmat$.

The significance of $\operatorname{cond}(\Kmat)$ lies in the fact that it is an indicator for estimating the relative error of $\bs{\Delta}$ from those of $\Kmat$ and $\Pvec$. Specifically,
\begin{equation}
  \frac{\|\delta\bs{\Delta}\|}{\|\bs{\Delta}\|} \leq
  \frac{\operatorname{cond}(\Kmat)}{1-\operatorname{cond}(\Kmat)\dfrac{\|\delta\Kmat\|}{\|\Kmat\|}}
  \left(\frac{\|\delta\Kmat\|}{\|\Kmat\|} + \frac{\|\delta\Pvec\|}{\|\Pvec\|}\right)
\end{equation}
If $\operatorname{cond}(\Kmat)\|\delta\Kmat\|/\|\Kmat\|$ is much less than 1, the above equation can be simplified to
\begin{equation}
  \frac{\|\delta\bs{\Delta}\|}{\|\bs{\Delta}\|} \leq \operatorname{cond}(\Kmat)\left(\frac{\|\delta\Kmat\|}{\|\Kmat\|} + \frac{\|\delta\Pvec\|}{\|\Pvec\|}\right)
\end{equation}
Notice that $\operatorname{cond}(\Kmat)\geq 1$. When $\operatorname{cond}(\Kmat)$ is very large, say larger than 10000, the system is called \emph{ill-conditioned}. In this case, the accuracy of $\bs{\Delta}$ is usually much less than that of $\Pvec$ and $\Kmat$.

A main cause of ill-conditioning is large differences in the stiffnesses of parts of the structure. For example, a stiff structure on flexible supports can give problems in this direction. However, a flexible structure on stiff supports tends not to suffer from such difficulties. For more examples on ill-conditioned structures, refer to Macleod (1990)\footnote{Macleod, I.A. (1990), \emph{Analytical Modelling of Structural Systems}, Ellis Horwood, London.}, page 117.

\section{Decomposition methods}
If $\mat{A}$ is symmetric (all eigenvalues are real), the two commonly used methods for solving $\mat{A}\mat{x}=\mat{b}$ are the \emph{Cholesky method} (appropriate for only positive definite matrix) and the \emph{triple-factoring method} (called root-free Cholesky method in the textbook). Their algorithms are given in Sec.~\ref{sec:cholesky-triple}.

\begin{example}[Cholesky decomposition]
\begin{equation*}
  \Kmat = \bmat{4 & 6 & 10 & 4\\ 6 & 13 & 13 & 6\\ 10 & 13 & 27 & 2\\ 4 & 6 & 2 & 72}
  = \bmat{L_{11} & 0 & 0 & 0\\ L_{21} & L_{22} & 0 & 0\\ L_{31} & L_{32} & L_{33} & 0\\ L_{41} & L_{42} & L_{43} & L_{44}}
    \bmat{L_{11} & L_{21} & L_{31} & L_{41}\\ 0 & L_{22} & L_{32} & L_{42}\\ 0 & 0 & L_{33} & L_{43}\\ 0 & 0 & 0 & L_{44}} = \mat{L}\mat{L}^\tr
\end{equation*}
\begin{align*}
  4 &= L_{11}^2 \rightarrow L_{11} = 2\\
  6 &= L_{11}L_{21} \rightarrow L_{21} = 3\\
  10 &= L_{11}L_{31} \rightarrow L_{31} = 5\\
  4 &= L_{11}L_{41} \rightarrow L_{41} = 2\\
  13 &= L_{21}^2 + L_{22}^2 \rightarrow L_{22} = 2\\
  13 &= L_{21}L_{31} + L_{22}L_{32} \rightarrow L_{32} = -1\\
  6 &= L_{21}L_{41} + L_{22}L_{42} \rightarrow L_{42} = 0\\
  27 &= L_{31}^2 + L_{32}^2 + L_{33}^2 \rightarrow L_{33} = 1\\
  2 &= L_{31}L_{41} + L_{32}L_{42} + L_{33}L_{43} \rightarrow L_{43} = -8\\
  72 &= L_{41}^2 + L_{42}^2 + L_{43}^2 + L_{44}^2 \rightarrow L_{44} = 2
\end{align*}
Therefore, $\mat{L} = \bmat{2 & 0 & 0 & 0\\ 3 & 2 & 0 & 0\\ 5 & -1 & 1 & 0\\ 2 & 0 & -8 & 2}$.
\end{example}

\begin{example}[Triple-factoring decomposition]
\begin{equation*}
  \Kmat = \bmat{1 & 2 & 3\\ 2 & 2 & 2\\ 3 & 2 & 5} = \mat{L}_1\mat{D}\mat{L}_1^\tr
\end{equation*}
\begin{equation*}
  \mat{D}\mat{L}_1^\tr = \bmat{D_1 & 0 & 0\\ 0 & D_2 & 0\\ 0 & 0 & D_3}\bmat{1 & L_{21} & L_{31}\\ 0 & 1 & L_{32}\\ 0 & 0 & 1}
  = \bmat{D_1 & D_1L_{21} & D_1L_{31}\\ 0 & D_2 & D_2L_{32}\\ 0 & 0 & D_3}
\end{equation*}
\begin{equation*}
  \Kmat = \bmat{1 & 2 & 3\\ 2 & 2 & 2\\ 3 & 2 & 5} = \mat{L}_1(\mat{D}\mat{L}_1^\tr)
  = \bmat{1 & 0 & 0\\ L_{21} & 1 & 0\\ L_{31} & L_{32} & 1}\bmat{D_1 & D_1L_{21} & D_1L_{31}\\ 0 & D_2 & D_2L_{32}\\ 0 & 0 & D_3}
\end{equation*}
\begin{align*}
  1 &= D_1\\
  2 &= D_1L_{21} \rightarrow L_{21} = 2\\
  3 &= D_1L_{31} \rightarrow L_{31} = 3\\
  2 &= D_1L_{21}^2 + D_2 \rightarrow D_2 = -2\\
  2 &= L_{21}D_1L_{31} + D_2L_{32} \rightarrow L_{32} = 2\\
  5 &= D_1L_{31}^2 + D_2L_{32}^2 + D_3 \rightarrow D_3 = 4
\end{align*}
Therefore $\mat{L}_1 = \bmat{1 & 0 & 0\\ 2 & 1 & 0\\ 3 & 2 & 1}$ and $\mat{D} = \bmat{1 & 0 & 0\\ 0 & -2 & 0\\ 0 & 0 & 4}$; $\Kmat$ is indefinite.
\end{example}

\section{Cholesky and triple-factoring methods}
\label{sec:cholesky-triple}
\textbf{Subroutines}: \texttt{DECOM} and \texttt{SOLVER}; solve $\Kmat\mat{U} = \Pvec$ for $\mat{U}$.

\topic{Factorization of $\Kmat$ (\texttt{DECOM}):}
\begin{equation}
  \Kmat = \mat{L}\mat{D}\mat{L}^\tr
\end{equation}
where $\mat{L}$ is a lower triangular matrix, and $\mat{D}$ a diagonal matrix; i.e.,
\begin{equation}
  \mat{L} = \bmat{L_{11} & & & \\ L_{21} & L_{22} & & \\ \vdots & & & \\ L_{n1} & L_{n2} & \cdots & L_{nn}}
  \quad\text{and}\quad
  \mat{D} = \bmat{D_1 & & & \\ & D_2 & & \\ & & \ddots & \\ & & & D_n}
\end{equation}
\begin{equation}
  \Kmat\mat{U} = \mat{L}\underline{\mat{D}\mat{L}^\tr\mat{U}} = \mat{L}\underline{\mat{c}} = \Pvec
\end{equation}

\topic{Two-step solution process (\texttt{SOLVER}):}
\begin{enumerate}
  \item First solve $\mat{L}\mat{c} = \Pvec$ for $\mat{c}$ by forward substitution, i.e.
    \begin{align}
      c_1 &= \frac{1}{L_{11}}P_1\\
      c_i &= \frac{1}{L_{ii}}\Bigl(P_i - \sum_{j=1}^{i-1}L_{ij}c_j\Bigr)\quad(i = 2,\cdots,n)
    \end{align}
  \item Then, solve $\mat{D}\mat{L}^\tr\mat{U} = \mat{c}$ or equivalently $\mat{L}^\tr\mat{U} = \mat{D}^{-1}\mat{c}$ for $\mat{U}$ by backward substitution, i.e.
    \begin{align}
      U_n &= \frac{1}{L_{nn}}\Bigl(\frac{c_n}{D_n}\Bigr)\\
      U_i &= \frac{1}{L_{ii}}\Bigl(\frac{c_i}{D_i} - \sum_{j=i+1}^{n}L_{ji}U_j\Bigr)\quad(i = n-1, n-2,\cdots,1)
    \end{align}
\end{enumerate}

\topic{Methods for obtaining $\mat{L}$ and $\mat{D}$:}
\begin{itemize}
  \item Unknowns: $\frac{n(n+1)}{2}$ entries of $\mat{L}$ and $n$ entries of $\mat{D}$.
  \item Equations: $\frac{n(n+1)}{2}$ equations from $\Kmat = \mat{L}\mat{D}\mat{L}^\tr$ (note that $\Kmat$ is symmetric).
\end{itemize}
We therefore have at least $n$ free choices in assigning the values of $\mat{L}$ and $\mat{D}$. We introduce two commonly used methods herein.

\begin{enumerate}[leftmargin=*]
  \item \textbf{Cholesky method}\\
    Set $D_i = 1$ ($i = 1,2,\cdots,n$); we can then determine $\mat{L}$ column by column using the formula:
    \begin{align}
      L_{jj} &= \Bigl(K_{jj} - \sum_{l=1}^{j-1}L_{jl}L_{jl}\Bigr)^{\frac12}\\
      L_{ij} &= \frac{1}{L_{jj}}\Bigl(K_{ij} - \sum_{l=1}^{j-1}L_{il}L_{jl}\Bigr)\quad(i = j+1, j+2,\cdots,n)
    \end{align}
    Note that for $j = 1$ set the summation terms to zero.

    Drawbacks:
    \begin{enumerate}[label=(\alph*)]
      \item Works efficiently for \textbf{positive definite} matrix only, because otherwise some of the $L_{ij}$ will be imaginary.
      \item Square root requires more computational efforts.
    \end{enumerate}

  \item \textbf{Triple-factoring method}\\
    Set $L_{ii} = 1$ ($i = 1,2,\cdots,n$); we can determine $D_j$ and then $L_{ij}$ columnwise by the formula:
    \begin{align}
      D_j &= K_{jj} - \sum_{l=1}^{j-1}L_{jl}D_lL_{jl}\\
      L_{ij} &= \frac{1}{D_j}\Bigl(K_{ij} - \sum_{l=1}^{j-1}L_{il}D_lL_{jl}\Bigr)\quad(i = j+1, j+2,\cdots,n)
    \end{align}
    Note that:
    \begin{enumerate}[label=(\alph*)]
      \item Works even for $\Kmat$ that is not positive definite.
      \item Determinant of $\Kmat$: $\operatorname{Det}(\Kmat) = \operatorname{Det}(\mat{L}\mat{D}\mat{L}) = D_1D_2\cdots D_n$.
      \item Operation counts including multiplication and division: $O(n^3/6)$ if $\Kmat$ is full, which is the same as that of the Cholesky method. For a banded $\Kmat$ with $\text{MHB} = m$, the operation counts are of $O(m^2n/2)$.
    \end{enumerate}
\end{enumerate}

\chapter[Equivalent Nodal Loads]{Equivalent Nodal Loads; Self-straining Problems}
\label{ch:eqload}

\section{Equivalent nodal loads}
The stiffness equation for a member subjected to intermediate concentrated loads or distributed loads can be written, with respect to the local or member axes, as
\begin{equation}
  \kmat\bs{\Delta}_l + \mat{f}_{fef} = \mat{f} \tag{a}\label{eq10:a}
\end{equation}
where $\mat{f}_{fef}$ is the fixed end force vector. Defining the equivalent force vector as $\mat{f}_{eq} = -\mat{f}_{fef}$, the above equation can be rewritten as
\begin{equation}
  \kmat\bs{\Delta}_l - \mat{f}_{eq} = \mat{f} \tag{b}\label{eq10:b}
\end{equation}
The above member stiffness equation is referred to the member (local) coordinate system. In order to carry out the assembling process, the stiffness equation referred to the structural (global) coordinate system is needed, which is:
\begin{equation}
  \Tmat^\tr\kmat\Tmat\bs{\Delta}_g - \Tmat^\tr\mat{f}_{eq} = \Tmat^\tr\mat{f}
  \quad\text{or}\quad
  \kmat_g\bs{\Delta}_g - \mat{f}_{eqg} = \mat{f}_g \tag{c}\label{eq10:c}
\end{equation}
where $\kmat_g = \Tmat^\tr\kmat\Tmat$, $\mat{f}_{eqg} = \Tmat^\tr\mat{f}_{eq}$, $\mat{f}_g = \Tmat^\tr\mat{f}$, and $\Tmat$ is the transformation matrix (see Ch.~\ref{ch:transform}). Eq.~\eqref{eq10:c} can be derived by using the contragredient principle and the congruent transformation. It can also be derived by premultiplying Eq.~\eqref{eq10:b} by $\Tmat^\tr$ and by using the relation that $\bs{\Delta}_l = \Tmat\bs{\Delta}_g$.

The stiffness equation for a structural system, obtained by assembling Eq.~\eqref{eq10:c}, takes the form
\begin{equation}
  \Kmat\mat{U} - \Pvec_{eq} = \F \tag{d}\label{eq10:d}
\end{equation}
where $\Kmat =$ assemblage of $\kmat_g$, $\mat{U} =$ collection of $\bs{\Delta}_g$, $\Pvec_{eq} =$ assemblage of $\Tmat^\tr\mat{f}_{eq}$, and $\F =$ assemblage of $\mat{f}_g$.

The \emph{nodal equilibrium equation} (considering the free body of a joint) states that
\begin{equation}
  \F = \Pvec_c \tag{e}\label{eq10:e}
\end{equation}
which means the nodal internal forces equal the external forces since $\Pvec_c =$ the applied nodal concentrated load vector. From Eqs.~\eqref{eq10:d} and \eqref{eq10:e}, we have
\begin{equation}
  \Kmat\mat{U} = \Pvec_c + \Pvec_{eq} \tag{f}\label{eq10:f}
\end{equation}
Once $\mat{U}$ is solved from Eq.~\eqref{eq10:f}, the member end forces, referred to the member coordinate system, can then be calculated according to the member stiffness equation, i.e.\ $\kmat\bs{\Delta}_l + \mat{f}_{fef} = \mat{f}$. To this end, we need to find $\bs{\Delta}_g$ first, which is a subset of $\mat{U}$, and then calculate $\bs{\Delta}_l$ by using $\bs{\Delta}_l = \Tmat\bs{\Delta}_g$.

From the above discussions, clearly the determination of $\mat{f}_{fef}$ or $\mat{f}_{eq}$ is the key element when dealing with members with intermediate loads or distributed loads. Table~5.1, on page 111 of the textbook, has listed fixed-end forces for some commonly encountered load types.
\begin{figure}[H]
  \centering
  % Fixed-end forces for commonly encountered load types (after Table 5.1 of the textbook)
\begin{tikzpicture}[scale=0.6,font=\scriptsize]
  \newcommand{\ffbeam}{\draw[member] (0,0) -- (6,0); \pic at (0,0) {wallL}; \pic at (6,0) {wallR};}
  \newcommand{\reactions}[4]{% #1 left moment label, #2 right moment label, #3 left shear, #4 right shear
    \draw[->,thick] (-0.9,0.5) arc (150:-30:0.5); \node[left] at (-1.0,0.55) {#1};
    \draw[->,thick] (6.9,0.5) arc (30:210:0.5); \node[right] at (7.0,0.55) {#2};
    \draw[load] (-0.6,-1.2) -- (-0.6,-0.1) node[pos=0,below] {#3};
    \draw[load] (6.6,-1.2) -- (6.6,-0.1) node[pos=0,below] {#4};}
  % ---------- (1) uniform load ----------
  \begin{scope}
    \ffbeam
    \foreach \x in {0.4,0.9,...,5.7}{\draw[->] (\x,0.9) -- (\x,0.1);} \draw (0.4,0.9) -- (5.7,0.9);
    \node[above] at (3,0.95) {$q$};
    \reactions{$\frac{qL^2}{12}$}{$\frac{qL^2}{12}$}{$\frac{qL}{2}$}{$\frac{qL}{2}$}
    \draw[dim] (0.6,-0.6) -- (5.4,-0.6) node[midway,fill=white,inner sep=1pt] {$L$};
  \end{scope}
  % ---------- (2) central point load ----------
  \begin{scope}[xshift=13cm]
    \ffbeam
    \draw[load] (3,1.3) -- (3,0.1) node[pos=0,right] {$P$};
    \reactions{$\frac{PL}{8}$}{$\frac{PL}{8}$}{$\frac P2$}{$\frac P2$}
    \draw[dim] (0.6,-0.6) -- (3,-0.6) node[midway,fill=white,inner sep=1pt] {$\frac L2$};
    \draw[dim] (3,-0.6) -- (5.4,-0.6) node[midway,fill=white,inner sep=1pt] {$\frac L2$};
    \draw (3,-0.8) -- (3,-0.4);
  \end{scope}
  % ---------- (3) point load at a ----------
  \begin{scope}[yshift=-4.2cm]
    \ffbeam
    \draw[load] (2,1.3) -- (2,0.1) node[pos=0,right] {$P$};
    \reactions{$\frac{Pab^2}{L^2}$}{$\frac{Pa^2b}{L^2}$}{$\frac{Pb^2(3a+b)}{L^3}$}{$\frac{Pa^2(a+3b)}{L^3}$}
    \draw[dim] (0.6,-0.6) -- (2,-0.6) node[midway,fill=white,inner sep=1pt] {$a$};
    \draw[dim] (2,-0.6) -- (5.4,-0.6) node[midway,fill=white,inner sep=1pt] {$b$};
    \draw (2,-0.8) -- (2,-0.4);
  \end{scope}
  % ---------- (4) support settlement ----------
  \begin{scope}[xshift=13cm,yshift=-4.2cm]
    \pic at (0,0) {wallL}; \pic at (6,-0.7) {wallR};
    \draw[member] (0,0) .. controls (2,0) and (4,-0.7) .. (6,-0.7);
    \draw[dim] (6.5,0) -- (6.5,-0.7) node[midway,right] {$\Delta$};
    \draw[->,thick] (-0.9,0.5) arc (150:-30:0.5); \node[left] at (-1.0,0.55) {$\frac{6EI\Delta}{L^2}$};
    \draw[->,thick] (6.9,-0.2) arc (30:-150:0.5); \node[right] at (7.0,0.2) {$\frac{6EI\Delta}{L^2}$};
    \draw[load] (-0.6,-1.2) -- (-0.6,-0.1) node[pos=0,below] {$\frac{12EI\Delta}{L^3}$};
    \draw[load] (6.6,-0.8) -- (6.6,-1.9) node[below] {$\frac{12EI\Delta}{L^3}$};
    \draw[dim] (0.6,-1.4) -- (5.4,-1.4) node[midway,fill=white,inner sep=1pt] {$L$};
  \end{scope}
  % ---------- (5) concentrated moment ----------
  \begin{scope}[yshift=-8.4cm]
    \ffbeam
    \draw[->,thick] (2.4,0.5) arc (100:-100:0.5); \node[right] at (2.6,-0.35) {$M$};
    \reactions{$\frac{Mb(2a-b)}{L^2}$}{$\frac{Ma(2b-a)}{L^2}$}{$\frac{6Mab}{L^3}$}{$\frac{6Mab}{L^3}$}
    \draw[dim] (0.6,-0.6) -- (2.2,-0.6) node[midway,fill=white,inner sep=1pt] {$a$};
    \draw[dim] (2.2,-0.6) -- (5.4,-0.6) node[midway,fill=white,inner sep=1pt] {$b$};
    \draw (2.2,-0.8) -- (2.2,-0.4);
  \end{scope}
  % ---------- (6) triangular load ----------
  \begin{scope}[xshift=13cm,yshift=-8.4cm]
    \ffbeam
    \foreach \x in {0.9,1.4,...,5.7}{\draw[->] (\x,{0.1+0.16*\x}) -- (\x,0.1);} \draw (0.4,0.1) -- (5.7,1.0);
    \node[above] at (4.5,1.0) {$q$};
    \reactions{$\frac{qL^2}{30}$}{$\frac{qL^2}{20}$}{$\frac{3qL}{20}$}{$\frac{7qL}{20}$}
    \draw[dim] (0.6,-0.6) -- (5.4,-0.6) node[midway,fill=white,inner sep=1pt] {$L$};
  \end{scope}
  % ---------- (7) symmetric triangular load ----------
  \begin{scope}[yshift=-12.6cm]
    \ffbeam
    \foreach \x in {0.7,1.2,...,2.9}{\draw[->] (\x,{0.1+0.34*\x}) -- (\x,0.1);}
    \foreach \x in {3.4,3.9,...,5.5}{\draw[->] (\x,{0.1+0.34*(6-\x)}) -- (\x,0.1);}
    \draw (0.3,0.1) -- (3,1.1) -- (5.7,0.1); \node[above] at (3,1.1) {$q$};
    \reactions{$\frac{5qL^2}{96}$}{$\frac{5qL^2}{96}$}{$\frac{qL}{4}$}{$\frac{qL}{4}$}
    \draw[dim] (0.6,-0.6) -- (5.4,-0.6) node[midway,fill=white,inner sep=1pt] {$L$};
  \end{scope}
  % ---------- (8) partial uniform load ----------
  \begin{scope}[xshift=13cm,yshift=-12.6cm]
    \ffbeam
    \foreach \x in {1.8,2.15,...,4.2}{\draw[->] (\x,0.8) -- (\x,0.1);} \draw (1.8,0.8) -- (4.2,0.8);
    \node[above right] at (4.2,0.8) {$q$};
    \draw[dim] (1.8,1.2) -- (3,1.2) node[midway,fill=white,inner sep=1pt] {$\frac d2$};
    \draw[dim] (3,1.2) -- (4.2,1.2) node[midway,fill=white,inner sep=1pt] {$\frac d2$};
    \draw[->,thick] (-0.9,0.5) arc (150:-30:0.5); \node[left,align=right] at (-1.0,0.55) {$\frac{qd}{L^2}\cdot$\\[-2pt]$[ab^2+\frac{(a-2b)d^2}{12}]$};
    \draw[->,thick] (6.9,0.5) arc (30:210:0.5); \node[right,align=left] at (7.0,0.55) {$\frac{qd}{L^2}\cdot$\\[-2pt]$[a^2b+\frac{(b-2a)d^2}{12}]$};
    \draw[load] (-0.6,-1.2) -- (-0.6,-0.1) node[pos=0,below] {$R_1$};
    \draw[load] (6.6,-1.2) -- (6.6,-0.1) node[pos=0,below] {$R_2$};
    \draw[dim] (0.6,-0.6) -- (3,-0.6) node[midway,fill=white,inner sep=1pt] {$a$};
    \draw[dim] (3,-0.6) -- (5.4,-0.6) node[midway,fill=white,inner sep=1pt] {$b$};
    \draw (3,-0.8) -- (3,-0.4);
    \node[anchor=north west,align=left] at (-1.5,-2.0) {$R_1 = \frac{qd}{L^3}\bigl[(2a+L)b^2+\bigl(\frac{a-b}{4}\bigr)d^2\bigr]$\\[4pt]$R_2 = \frac{qd}{L^3}\bigl[(2b+L)a^2-\bigl(\frac{a-b}{4}\bigr)d^2\bigr]$};
  \end{scope}
  \node[anchor=north west,align=left,font=\footnotesize] at (-1.5,-14.6) {Sign convention: all loads and reactive forces\\are positive in the directions shown.};
\end{tikzpicture}

  \caption{Fixed-end forces for common load types (after Table~5.1 of the textbook).}
  \label{fig:fef-table}
\end{figure}

In general, we can determine $\mat{f}_{eq}$ using the principle of virtual displacements. In this way, we let the virtual work done by $\mat{f}_{eq}$ be equal to that done by the intermediate concentrated loads (can be expressed as delta functions) or distributed loads. That is,
\begin{gather}
  \delta\bs{\Delta}_l^\tr\mat{f}_{eq} = \int_0^L \delta v(x)^\tr q(x)\,dx = \int_0^L (\NN\delta\bs{\Delta}_l)^\tr q(x)\,dx
  \;\rightarrow\; \delta\bs{\Delta}_l^\tr\mat{f}_{eq} = \delta\bs{\Delta}_l^\tr\int_0^L \NN^\tr q(x)\,dx\\
  \rightarrow\quad \mat{f}_{eq} = \int_0^L \NN^\tr q(x)\,dx
\end{gather}
where $\NN = \langle N_1, N_2, \dots\rangle$ are shape functions. Written in component form,
\begin{equation}
  (\mat{f}_{eq})_i = \int_0^L N_i\,q(x)\,dx
\end{equation}
Examples illustrating the above equation (will be done in class).

\begin{example}
Consider the beam shown below with the properties: $E = 28$~GPa, $I = 5.8\times10^9$~mm$^4$.
\begin{fig}
  % Example beam: 4 nodes, 3 elements, with loads
\begin{tikzpicture}[scale=0.7,font=\small]
  \draw[member] (0,0) -- (20,0);
  \pic at (0,0) {wallL}; \pic at (10,0) {roller}; \pic at (20,0) {roller};
  \foreach \x/\n in {0/1,6/2,10/3,20/4}{\node[below right,inner sep=2pt] at (\x,0) {\n};}
  \node[node dot] at (6,0) {};
  % triangular load on element 1
  \foreach \x in {0.5,1.2,...,5.6}{\draw[->] (\x,{0.2+0.28*\x}) -- (\x,0.15);} \draw (0,0.2) -- (6,1.9);
  \node[above left] at (1.5,1.6) {30\,kN/m};
  \draw[load] (6,2.8) -- (6,0.2) node[pos=0,above] {200\,kN};
  \draw[->,thick] (11.2,1.3) arc (30:180:0.9); \node[above] at (10.2,1.6) {90\,kNm};
  \draw[load] (15,2.4) -- (15,0.2) node[pos=0,above] {150\,kN};
  \node[below] at (3,-0.4) {$1.5I$}; \node[below] at (15,-0.4) {$I$};
  \draw[dim] (0,-1.8) -- (6,-1.8) node[midway,fill=white,inner sep=1pt] {6\,m};
  \draw[dim] (6,-1.8) -- (10,-1.8) node[midway,fill=white,inner sep=1pt] {4\,m};
  \draw[dim] (10,-1.8) -- (15,-1.8) node[midway,fill=white,inner sep=1pt] {5\,m};
  \draw[dim] (15,-1.8) -- (20,-1.8) node[midway,fill=white,inner sep=1pt] {5\,m};
  \foreach \x in {0,6,10,15,20}{\draw (\x,-1.5) -- (\x,-2.1);}
\end{tikzpicture}

\end{fig}
\begin{fig}
  % Fixed-end forces of elements 1 and 3
\begin{tikzpicture}[scale=0.7,font=\small]
  % element 1
  \begin{scope}
    \draw[member] (0,0) -- (6,0); \pic at (0,0) {wallL}; \pic at (6,0) {wallR};
    \foreach \x in {0.5,1.2,...,5.6}{\draw[->] (\x,{0.2+0.28*\x}) -- (\x,0.15);} \draw (0,0.2) -- (6,1.9);
    \node[above left] at (1.2,1.5) {30\,kN/m};
    \node[below right] at (0,0) {1}; \node[below left] at (6,0) {2};
    \draw[->,thick] (-0.5,-0.6) arc (-100:100:0.6); \node[above left] at (-0.4,0.5) {54};
    \draw[->,thick] (6.5,-0.6) arc (280:80:0.6); \node[above right] at (6.4,0.5) {$-36$};
    \draw[load] (-0.5,-1.4) -- (-0.5,-0.1) node[pos=0,below] {63};
    \draw[load] (6.5,-1.4) -- (6.5,-0.1) node[pos=0,below] {27};
  \end{scope}
  % element 3
  \begin{scope}[xshift=10.5cm]
    \draw[member] (0,0) -- (10,0); \pic at (0,0) {wallL}; \pic at (10,0) {wallR};
    \draw[load] (5,2.2) -- (5,0.15) node[pos=0,above] {150\,kN};
    \node[below right] at (0,0) {3}; \node[below left] at (10,0) {4};
    \draw[->,thick] (-0.5,-0.6) arc (-100:100:0.6); \node[above left] at (-0.4,0.5) {187.5};
    \draw[->,thick] (10.5,-0.6) arc (280:80:0.6); \node[above right] at (10.4,0.5) {$-187.5$};
    \draw[load] (-0.5,-1.4) -- (-0.5,-0.1) node[pos=0,below] {75};
    \draw[load] (10.5,-1.4) -- (10.5,-0.1) node[pos=0,below] {75};
    \node[right] at (11.6,-0.6) {$\mat{f}_{fef}$};
  \end{scope}
\end{tikzpicture}

\end{fig}
\begin{fig}
  % Equivalent nodal loads (solid) and concentrated nodal loads (dashed)
\begin{tikzpicture}[scale=0.62,font=\small]
  \draw[member] (0,0) -- (20,0);
  \foreach \x/\n in {0/1,6/2,10/3,20/4}{\node[below,inner sep=3pt] at (\x,0) {\n};}
  % node 1
  \draw[load] (0,0.2) -- (0,1.6) node[above] {$-63$};
  \draw[->,thick] (0.6,0.4) arc (30:180:0.6); \node[left] at (-0.5,0.4) {$-54$};
  % node 2
  \draw[load] (6,0.2) -- (6,1.6) node[above right] {$-27$};
  \draw[load,dashed] (6,1.7) -- (6,3.2) node[above] {$-200$};
  \draw[->,thick] (6.6,0.4) arc (30:180:0.6); \node[left] at (5.5,0.6) {36};
  % node 3
  \draw[load] (10,0.2) -- (10,1.6) node[above] {$-75$};
  \draw[->,thick] (10.6,0.4) arc (30:180:0.6); \node[right] at (10.6,0.7) {$-187.5$};
  \draw[->,thick,dashed] (10.9,0.9) arc (20:180:1.0); \node[right] at (11.0,1.4) {$-90$};
  % node 4
  \draw[load] (20,0.2) -- (20,1.6) node[above] {$-75$};
  \draw[->,thick] (20.6,0.4) arc (30:180:0.6); \node[right] at (20.6,0.7) {187.5};
  % legend
  \draw (22.5,1.6) -- (24,1.6) node[right] {$\mat{f}_{eq}$};
  \draw[dashed] (22.5,0.6) -- (24,0.6) node[right] {$\Pvec_c$};
\end{tikzpicture}

\end{fig}
Note that in the program, the fixed end forces \texttt{FEF} shall be input as follows:
\begin{lstlisting}[language={},numbers=none,basicstyle=\small\ttfamily]
*FEF
1   63    54    27   -36
2    0     0     0     0
3   75   187.5  75  -187.5
\end{lstlisting}
After assembling, we obtain the final nodal loads:
\begin{fig}
  % Final nodal loads after assembling
\begin{tikzpicture}[scale=0.7,font=\small]
  \draw[member] (0,0) -- (20,0);
  \pic at (0,0) {wallL}; \pic at (10,0) {roller}; \pic at (20,0) {roller};
  \foreach \x/\n in {0/1,6/2,10/3,20/4}{\node[below right,inner sep=3pt] at (\x,0) {\n};}
  \draw[load] (6,0.2) -- (6,2.2) node[above] {$-227$\,kN};
  \draw[->,thick] (6.6,0.5) arc (30:180:0.7); \node[left] at (5.6,0.9) {36\,kNm};
  \draw[->,thick] (10.8,0.5) arc (30:180:0.9); \node[right] at (10.6,1.1) {$-277.5$\,kNm};
  \draw[->,thick] (20.8,0.5) arc (30:180:0.9); \node[above] at (20,1.3) {187.5\,kNm};
\end{tikzpicture}

\end{fig}
\begin{equation*}
  \Longrightarrow\quad \Pvec = \Pvec_c + \Pvec_{eq} = \Bmat{-200\\ 0\\ -90\\ 0} + \Bmat{-27\\ 36\\ -187.5\\ 187.5}
  = \Bmat{-227\text{ kN}\\ 36\text{ kN}\cdot\text{m}\\ -277.5\text{ kN}\cdot\text{m}\\ 187.5\text{ kN}\cdot\text{m}}
\end{equation*}
\begin{equation*}
  \Kmat = \bmat{43983 & 20300 & 60900 & 0\\ 20300 & 324800 & 81200 & 0\\ 60900 & 81200 & 227360 & 32480\\ 0 & 0 & 32480 & 64960}\quad\text{(with appropriate units)}
\end{equation*}
Then,
\begin{equation*}
  \bs{\Delta} = \Bmat{-4.4729\text{ m}\\ 0.56143\text{ rad}\\ -0.68415\text{ rad}\\ 3.2285\text{ rad}}\times10^{-3}
\end{equation*}

\topic{Force recovery (element one)}
\begin{fig}
  % Element 1: fixed-end forces and local displacements
\begin{tikzpicture}[scale=0.8,font=\small]
  \begin{scope}
    \draw[member] (0,0) -- (6,0);
    \node[below right] at (0.2,0) {1}; \node[below left] at (5.8,0) {2};
    \draw[->,thick] (-0.15,-0.6) arc (-100:100:0.6); \node[above left] at (-0.1,0.5) {54};
    \draw[->,thick] (6.15,-0.6) arc (280:80:0.6); \node[above right] at (6.1,0.5) {$-36$};
    \draw[load] (0,-1.2) -- (0,-0.1) node[pos=0,below] {63};
    \draw[load] (6,-1.2) -- (6,-0.1) node[pos=0,below] {27};
  \end{scope}
  \begin{scope}[xshift=9.5cm]
    \draw[member] (0,0) -- (6,0);
    \node[below right] at (0.2,0) {1}; \node[below left] at (5.8,0) {2};
    \draw[->,thick] (-0.15,-0.6) arc (-100:100:0.6); \node[above left] at (-0.1,0.5) {0};
    \draw[->,thick] (6.15,-0.6) arc (280:80:0.6); \node[above right] at (6.1,0.5) {$\Delta_2$};
    \draw[load] (0,-1.2) -- (0,-0.1) node[pos=0,below] {0};
    \draw[load] (6,-1.2) -- (6,-0.1) node[pos=0,below] {$\Delta_1$};
  \end{scope}
\end{tikzpicture}

\end{fig}
Fixed end forces $\mat{f}_{fef} = \Bmat{63\\ 54\\ 27\\ -36}$, and local displacement vector $\bs{\Delta}_l = \Bmat{0\\ 0\\ \Delta_1\\ \Delta_2}$. Then from Eq.~\eqref{eq10:a},
\begin{equation*}
  \mat{f}_1 = \kmat\bs{\Delta}_l + \mat{f}_{fef} = \Bmat{146.3\text{ kN}\\ 281.19\text{ kN}\cdot\text{m}\\ -56.33\text{ kN}\\ 236.78\text{ kN}\cdot\text{m}}.
\end{equation*}
Similarly, we can obtain the member forces of elements 2 and 3 as follows:
\begin{equation*}
  \mat{f}_2 = \Bmat{-143.67\text{ kN}\\ -236.78\text{ kN}\cdot\text{m}\\ 143.67\text{ kN}\\ -337.92\text{ kN}\cdot\text{m}}
  \quad\text{and}\quad
  \mat{f}_3 = \Bmat{99.79\text{ kN}\\ 247.92\text{ kN}\cdot\text{m}\\ 50.21\text{ kN}\\ 0\text{ kN}\cdot\text{m}}
\end{equation*}
\end{example}

\section{Self-straining problems}
Self-straining problems include those dealing with temperature changes, fabrication errors (misfit), and support settlement. For these kinds of problems, the member stiffness equations can also be written as
\begin{equation}
  \kmat\bs{\Delta}_l + \mat{f}_{fef} = \mat{f}\qquad\text{or}\qquad \kmat\bs{\Delta}_l = \mat{f} + \mat{f}_{eq}
\end{equation}
where $\mat{f}_{eq} = -\mat{f}_{fef}$ are defined as the forces that need to be applied to the member so that all nodal displacements are equal to zero, or equivalently defined as the forces that counterbalance the initial strains. For example, if an axial member is subjected to a temperature rise $\Delta T$, then obviously a compressive force with the magnitude of $AE\alpha\Delta T$ needs to be applied to the ends of the member in order not to induce deformations or forces for the other members, where $\alpha =$ thermal expansion coefficient, a material constant. Actually, this compressive force produces an axial strain $= -\alpha\Delta T$, which counterbalances the initial strain due to the temperature rise $= \alpha\Delta T$. In summary, if $\eps_i = \alpha\Delta T$ represents the initial strain, then the fixed-force $f_{fef} = -EA\eps_i$ and therefore the equivalent force $f_{eq} = -f_{fef} = EA\eps_i$.

Similarly, if an axial bar has the fabrication error $\Delta L$, then it has the initial strain $\eps_i = \Delta L/L$. The fixed-end force $f_{fef} = -EA\eps_i$, and the equivalent force $f_{eq} = EA\eps_i = EA\Delta L/L$.

In general, from the principle of virtual work, we have, for an axial bar, that
\begin{equation}
  f_{eq} = EA\frac1L\int_0^L \eps_i(x)\,dx
\end{equation}
For beams, the fixed-end forces corresponding to support settlement or temperature changes can be determined similarly as we discuss in class.
\begin{fig}
  % Self-straining: temperature rise in two bars of a three-bar truss = locked + unlocked
\begin{tikzpicture}[scale=0.8,font=\small]
  \newcommand{\truss}{%
    \draw[member] (-2,0) -- (0,3) -- (2,0); \draw[member] (0,0) -- (0,3);
    \pic at (-2,0) {pin}; \pic at (0,0) {pin}; \pic at (2,0) {pin}; \node[node dot] at (0,3) {};}
  \begin{scope}
    \truss
    \node[above left] at (-1.2,1.8) {$\uparrow\Delta T_1$}; \node[above right] at (1.2,1.8) {$\downarrow\Delta T_2$};
    \node at (-1.2,3.9) {$\Delta T_1 > 0$}; \node at (1.2,3.9) {$\Delta T_2 > 0$};
  \end{scope}
  \node at (3.6,1.5) {$=$};
  \begin{scope}[xshift=7.2cm]
    \truss \node at (0,4.6) {\itshape locked};
    \draw[load] ($(0,3)+(-1.4,2.1)$) -- ($(0,3)+(-0.15,0.22)$) node[pos=0,above left] {$EA\alpha\Delta T_2$};
    \draw[load] ($(0,3)+(1.4,2.1)$) -- ($(0,3)+(0.15,0.22)$) node[pos=0,above right] {$EA\alpha\Delta T_1$};
  \end{scope}
  \node at (10.8,1.5) {$+$};
  \begin{scope}[xshift=14.4cm]
    \truss \node at (0,4.6) {\itshape unlocked};
    \draw[load] ($(0,3)+(-0.15,0.22)$) -- ($(0,3)+(-1.4,2.1)$) node[above left] {$EA\alpha\Delta T_2$};
    \draw[load] ($(0,3)+(0.15,0.22)$) -- ($(0,3)+(1.4,2.1)$) node[above right] {$EA\alpha\Delta T_1$};
  \end{scope}
\end{tikzpicture}

\end{fig}

\begin{example}[Example 5.10 in the textbook]
Member $ad$ of the truss is 3~mm fabricated too short and $E = 200{,}000$~MPa.
\begin{fig}
  % Example 5.10: truss with member ad fabricated 3 mm too short
\begin{tikzpicture}[scale=0.75,font=\scriptsize]
  \coordinate (c) at (0,0); \coordinate (b) at (10.928,0); \coordinate (a) at (6.928,4); \coordinate (d) at (0,4);
  \draw[member] (c) -- (a) -- (b) -- (c); \draw[member] (d) -- (a);
  \pic[rotate=-90] at (d) {pin}; \pic at (c) {pin}; \pic at (b) {roller};
  \node[above right] at (a) {$a$}; \node[below right] at (b) {$b$}; \node[below left] at (c) {$c$}; \node[above right] at (d) {$d$};
  \node[above,sloped] at ($(c)!0.5!(a)$) {$A = 15\times10^3$ mm$^2$};
  \node[above,sloped] at ($(a)!0.5!(b)$) {$A = 20\times10^3$ mm$^2$};
  \node[above] at (3.5,4) {$A = 20\times10^3$ mm$^2$};
  \node[below] at (5.5,0) {$A = 18\times10^3$ mm$^2$};
  \draw[->] (1.6,0) arc (0:30:1.6); \node at (2.0,0.4) {$30^\circ$};
  \draw[->] (9.4,0) arc (180:135:1.5); \node at (8.9,0.5) {$45^\circ$};
  \draw[dim] (0,4.8) -- (6.928,4.8) node[midway,fill=white,inner sep=1pt] {6.928 m};
  \draw[dim] (0,-1.2) -- (10.928,-1.2) node[midway,fill=white,inner sep=1pt] {10.928 m};
  \draw[dim] (11.8,0) -- (11.8,4) node[midway,right] {4 m};
  \draw[->] (-0.6,0.3) -- (-0.6,1.3) node[above] {$y$}; \draw[->] (11.2,-0.5) -- (12.2,-0.5) node[right] {$x$};
\end{tikzpicture}

\end{fig}
DOFs and fixed end forces:
\begin{fig}
  % Example 5.10: DOFs and reactions
\begin{tikzpicture}[scale=0.55,font=\small]
  \coordinate (c) at (0,0); \coordinate (b) at (10.928,0); \coordinate (a) at (6.928,4); \coordinate (d) at (0,4);
  \draw[thinmember] (c) -- (a) -- (b) -- (c); \draw[thinmember] (d) -- (a);
  \foreach \p in {a,b,c,d}{\fill (\p) circle (0.07);}
  \node[above left] at (a) {$a$}; \node[above] at (b) {$b$}; \node[above right] at (c) {$c$}; \node[below right] at (d) {$d$};
  \draw[load] (a) -- ++(1.6,0) node[right] {$P_1,\Delta_1$};
  \draw[load] (a) -- ++(0,1.8) node[right] {$P_2,\Delta_2$};
  \draw[load] (b) -- ++(1.6,0) node[right] {$P_3,\Delta_3$};
  \draw[load] ($(d)+(-1.8,0)$) -- (d) node[pos=0,below] {$R_{xd}$};
  \draw[load] ($(c)+(-1.8,0)$) -- (c) node[pos=0,below] {$R_{xc}$};
  \draw[load] ($(c)+(0,-1.8)$) -- (c) node[pos=0,right] {$R_{yc}$};
  \draw[load] ($(b)+(0,-1.8)$) -- (b) node[pos=0,right] {$R_{yb}$};
\end{tikzpicture}

\end{fig}
\begin{fig}
  % Example 5.10: fixed-end forces of member ad (3 mm too short)
\begin{tikzpicture}[scale=0.8,font=\small]
  \draw[member] (0,0) -- (6.928,0);
  \draw (6.928,-0.4) -- (6.928,0.9); \draw (7.3,-0.4) -- (7.3,0.9);
  \draw[dim] (7.3,0.7) -- (8.6,0.7); \node[right] at (8.6,0.7) {3 mm};
  \draw[<-] (6.928,0.7) -- (5.8,0.7);
  \node[below right] at (0,0) {$d$}; \node[above left] at (6.928,0) {$a$};
  \draw[load] (-1.6,0) -- (-0.1,0) node[pos=0,below] {$F^F_{xd}$};
  \draw[load] (8.6,0) -- (7.4,0) node[pos=0,right] {$F^F_{xa}$};
  \draw[dim] (0,-0.8) -- (6.928,-0.8) node[midway,fill=white,inner sep=1pt] {6.928 m};
\end{tikzpicture}

\end{fig}
\begin{equation*}
  F^F_{xa} = -F^F_{xd} = \Bigl(\frac{EA}{L}\Bigr)_{ad}\cdot\Delta L_{ad} = 577.37\times3 = 1732.1\text{ kN}
\end{equation*}
Obviously, the fixed end force is zero except member $ad$.

Q: $\mat{f}^{ad}_{fef} = \langle -1732.1, 1732.1\rangle$~kN or $\mat{f}^{ad}_{fef} = \langle 1732.1, -1732.1\rangle$~kN? Why?

Load vector:
\begin{equation*}
  \Pvec = \Pvec_c + \Pvec_{eq} = \Bmat{0\\ 0\\ 0} + \Bmat{-1732.1\\ 0\\ 0}\text{ kN}
\end{equation*}
Stiffness equations:
\begin{equation*}
  10^2\bmat{12.122 & -1.912 & -3.536\\ & 4.473 & 3.536\\ & & 6.830}\Bmat{\Delta_1\\ \Delta_2\\ \Delta_3} = \Bmat{-1732.1\\ 0\\ 0}
\end{equation*}
Displacement: $\langle\Delta_1, \Delta_2, \Delta_3\rangle = \langle -1.683, -0.0158, -8.448\rangle$~mm.
\end{example}

\begin{example}[Example 5.11 in the textbook]
The support at $b$ settles 15~mm; analyze the problem.
\begin{fig}
  % Example 5.11: two-span beam, support b settles 15 mm
\begin{tikzpicture}[scale=0.8,font=\small]
  \draw[member] (0,0) -- (13,0);
  \pic at (0,0) {pin}; \pic at (8,0) {roller}; \pic at (13,0) {wallR};
  \node[above] at (0,0.1) {$a$}; \node[above] at (8,0.1) {$b$}; \node[above left] at (13,0.1) {$c$};
  \draw[dim] (0,-1.0) -- (8,-1.0) node[midway,fill=white,inner sep=1pt] {8 m};
  \draw[dim] (8,-1.0) -- (13,-1.0) node[midway,fill=white,inner sep=1pt] {5 m};
  \foreach \x in {0,8,13}{\draw (\x,-0.8) -- (\x,-1.2);}
\end{tikzpicture}

\end{fig}
Fixed end forces (note that $\theta_{zb} = 0$ in the figure below):
\begin{fig}
  % Example 5.11: fixed-end state with support b settled 15 mm (theta_zb = 0)
\begin{tikzpicture}[scale=0.8,font=\small]
  \pic at (0,0) {wallL}; \pic at (13,0) {wallR};
  \draw[dashed] (0,0) -- (13,0);
  \draw[member] (0,0) .. controls (3,0) and (5,-0.9) .. (8,-0.9) .. controls (10,-0.9) and (11.5,0) .. (13,0);
  \pic at (8,-0.9) {roller};
  \draw[dim] (13.6,0) -- (13.6,-0.9) node[midway,right] {15 mm};
  \draw[dim] (8.4,0.6) -- (8.4,-0.9); \node[right] at (8.4,0.5) {15 mm};
  \node[above] at (0,0.1) {$a$}; \node[above right] at (8,-0.8) {$b$}; \node[above left] at (13,0.1) {$c$};
\end{tikzpicture}

\end{fig}
\begin{fig}
  % Example 5.11: fixed-end forces of members ab and bc
\begin{tikzpicture}[scale=0.8,font=\small]
  \begin{scope}
    \draw[member] (0,0) -- (7,0); \fill (0,0) circle (0.06); \fill (7,0) circle (0.06);
    \node[above right] at (0,0) {$a$}; \node[above left] at (7,0) {$b$};
    \draw[->,thick] (-0.1,-0.6) arc (-100:100:0.6); \node[above left] at (-0.1,0.5) {$M^F_{za}$};
    \draw[->,thick] (7.1,-0.6) arc (280:80:0.6); \node[above right] at (7.1,0.5) {$M^F_{zb}$};
    \draw[load] (0,-1.3) -- (0,-0.1) node[pos=0,below] {$F^F_{ya}$};
    \draw[load] (7,-1.3) -- (7,-0.1) node[pos=0,below] {$F^F_{yb}$};
  \end{scope}
  \begin{scope}[xshift=10cm]
    \draw[member] (0,0) -- (6,0); \fill (0,0) circle (0.06); \fill (6,0) circle (0.06);
    \node[above right] at (0,0) {$b$}; \node[above left] at (6,0) {$c$};
    \draw[->,thick] (-0.1,-0.6) arc (-100:100:0.6); \node[above left] at (-0.1,0.5) {$M^F_{zb}$};
    \draw[->,thick] (6.1,-0.6) arc (280:80:0.6); \node[above right] at (6.1,0.5) {$M^F_{zc}$};
    \draw[load] (0,-1.3) -- (0,-0.1) node[pos=0,below] {$F^F_{yb}$};
    \draw[load] (6,-1.3) -- (6,-0.1) node[pos=0,below] {$F^F_{yc}$};
  \end{scope}
\end{tikzpicture}

\end{fig}
\begin{minipage}{0.5\linewidth}
\begin{align*}
  M^F_{za} = M^F_{zb} &= 200(-18.75)(-15)\\ &= 56.25\times10^3\text{ kNmm}\\
  F^F_{ya} = -F^F_{yb} &= 200(-0.00469)(-15)\\ &= 14.07\text{ kN}
\end{align*}
\end{minipage}\hfill
\begin{minipage}{0.48\linewidth}
\begin{align*}
  M^F_{zb} = M^F_{zc} &= 200(12.00)(-15)\\ &= -36.0\times10^3\text{ kNmm}\\
  F^F_{yb} = -F^F_{yc} &= 200(0.00480)(-15)\\ &= -14.40\text{ kN}
\end{align*}
\end{minipage}
\begin{equation*}
  \mat{f}^{ab}_{fef} = \langle 14.07, 56250, -14.07, 56250\rangle^\tr;\qquad
  \mat{f}^{bc}_{fef} = \langle -14.4, -36000, 14.4, -36000\rangle^\tr
\end{equation*}
Load vector:
\begin{equation*}
  \Pvec = \Pvec_c + \Pvec_{eq} = \Bmat{0\\ 0} + \Bmat{-56.25\times10^3\\ -56.25\times10^3 + 36\times10^3}\text{ kN}\cdot\text{mm}
  = \Bmat{-56.25\times10^3\\ -20.25\times10^3}\text{ kN}\cdot\text{mm}
\end{equation*}
Stiffness equations:
\begin{equation*}
  200\times10^5\bmat{1 & 0.5\\ 0.5 & 1.4}\Bmat{\theta_{za}\\ \theta_{zb}} = \Bmat{-56.25\times10^3\\ -20.25\times10^3}
\end{equation*}
Displacement: $\langle\theta_{za}, \theta_{zb}\rangle = \langle -0.002984, 0.000342\rangle$~rad.

\topic{Comparison: the slope-deflection method}
\begin{fig}
  % Comparison with the slope-deflection method
\begin{tikzpicture}[scale=0.8,font=\small]
  \begin{scope}
    \draw[member] (0,0) -- (10,0); \pic at (0,0) {pin}; \pic at (4,0) {roller}; \pic at (10,0) {wallR};
    \draw[->,thick] (0.6,0.5) arc (20:170:0.6); \node[above] at (-0.1,0.9) {$\theta_a$};
    \draw[->,thick] (4.6,0.5) arc (20:170:0.6); \node[above] at (3.9,0.9) {$\theta_b$};
    \node[below] at (2,-0.6) {$EI_1, L_1$}; \node[below] at (7,-0.6) {$EI_2, L_2$};
  \end{scope}
  \begin{scope}[yshift=-3.6cm]
    \draw[member] (0,0) -- (10,0); \pic at (0,0) {pin}; \pic at (4,0) {roller}; \pic at (10,0) {wallR};
    \node[above] at (0,0.6) {$M_a = 0$};
    \draw[->,thick] (4.6,0.5) arc (20:170:0.6); \node[above] at (3.9,0.9) {$\theta_b$};
    \node[below] at (2,-0.6) {$EI_1, L_1$}; \node[below] at (7,-0.6) {$EI_2, L_2$};
  \end{scope}
\end{tikzpicture}

\end{fig}
\begin{equation*}
  \Kmat = \begin{bmatrix}\dfrac{4E_1I_1}{L_1} & \dfrac{2E_1I_1}{L_1}\\[8pt] \dfrac{2E_1I_1}{L_1} & \dfrac{4E_1I_1}{L_1}+\dfrac{4E_2I_2}{L_2}\end{bmatrix}
  \qquad\qquad
  k = \frac{3E_1I_1}{L_1} + \frac{4E_2I_2}{L_2},\quad M_b = k\cdot\theta_b
\end{equation*}
\end{example}

\section{Thermal load problem of beam}
\begin{fig}
  % Thermal load problem of a beam (after the textbook)
\begin{tikzpicture}[scale=0.8,font=\scriptsize]
  % ---- (a) beam with temperature gradient ----
  \begin{scope}
    \draw (0,0) rectangle (7,0.9);
    \pic at (0,0) {pin}; \pic at (7,0) {roller};
    \node[left] at (0,0.45) {1}; \node[right] at (7,0.45) {2};
    \draw[->] (7.4,0.45) -- (8.8,0.45) node[right] {$x$};
    \draw[->] (0,1.2) -- (0,2.4) node[above] {$y$};
    \draw (3.0,0) -- (3.0,0.9); \draw (3.2,0) -- (3.2,0.9); \node[above] at (3.1,0.95) {$dx$};
    \draw[dim] (0.4,0.45) -- (0.4,0.9); \node[right] at (0.45,0.6) {$h$};
    \node[below right] at (0.3,-0.05) {$EI_z = $ constant};
    % temperature profile
    \draw (5.0,0.9) -- (5.0,0); \draw (5.0,0.9) -- (6.0,0.9) -- (5.0,0) -- cycle;
    \node[above] at (6.0,0.95) {$T_u$}; \node[below] at (5.5,-0.05) {$T_l$};
    \draw[dim] (0,-0.9) -- (7,-0.9) node[midway,fill=white,inner sep=1pt] {$L$};
    \node at (3.5,-1.5) {(a)};
    % I section
    \draw (8.9,1.6) -- (9.5,1.6) -- (9.5,1.5) -- (9.25,1.5) -- (9.25,0.9) -- (9.5,0.9) -- (9.5,0.8) -- (8.9,0.8) -- (8.9,0.9) -- (9.15,0.9) -- (9.15,1.5) -- (8.9,1.5) -- cycle;
  \end{scope}
  % ---- (b) deformed beam ----
  \begin{scope}[yshift=-4.6cm]
    \draw (0,0.9) .. controls (2.3,0.0) and (4.7,0.0) .. (7,0.9);
    \draw (0,1.7) .. controls (2.3,0.9) and (4.7,0.9) .. (7,1.7);
    \draw (0,0.9) -- (0,1.7); \draw (7,0.9) -- (7,1.7);
    \fill[hatch] (3.35,0.32) rectangle (3.65,1.12);
    \pic at (0,0.9) {pin}; \pic at (7,0.9) {roller};
    \draw[dim] (3.35,2.1) -- (3.65,2.1); \node[right] at (3.7,2.1) {$dx+\alpha T_u dx$};
    \draw[dim] (3.35,-0.2) -- (3.65,-0.2); \node[right] at (3.7,-0.2) {$dx+\alpha T_l dx$};
    \node at (3.5,-0.9) {(b)};
  \end{scope}
  % ---- (d) fixed-end moments ----
  \begin{scope}[yshift=-8.4cm]
    \draw (0,0) rectangle (7,0.7);
    \pic at (0,0) {pin}; \pic at (7,0) {roller};
    \node[below left] at (0,0) {1}; \node[below right] at (7,0) {2};
    \draw[->,thick] (-0.3,-0.4) arc (-100:120:0.55); \node[above left,align=right] at (-0.2,0.9) {$M^F_{z1} = \dfrac{EI_z\alpha(T_l-T_u)}{h}$};
    \draw[->,thick] (7.3,-0.4) arc (280:60:0.55); \node[above right,align=left] at (7.2,0.9) {$M^F_{z2} = \dfrac{EI_z\alpha(T_l-T_u)}{h}$};
    \node at (3.5,-1.2) {(d)};
  \end{scope}
  % ---- (c) differential element ----
  \begin{scope}[xshift=11.5cm,yshift=-5.5cm]
    \coordinate (O) at (0,5.2);
    \draw (O) -- (-1.0,0); \draw (O) -- (1.0,0);
    \draw[thick] (-0.35,1.6) -- (0.35,1.6) -- (0.6,0.3) -- (-0.6,0.3) -- cycle;
    \draw (-0.2,1.6) rectangle (0.2,1.9);
    \draw[dim] (-0.8,1.6) -- (-0.8,0.3) node[midway,left] {$h$};
    \draw[dim] (-1.4,5.2) -- (-1.4,0.3); \node[left] at (-1.4,2.8) {$\rho$};
    \draw[->] (-0.35,4.3) arc (250:290:1.0); \node[below] at (0,3.9) {$d\theta$};
    \node[right] at (0.4,1.75) {$(1+\alpha T_u)dx$};
    \draw[dim] (-0.6,-0.3) -- (0.6,-0.3); \node[right] at (0.7,-0.3) {$(1+\alpha T_l)dx$};
    \node at (0,-0.9) {(c)};
  \end{scope}
\end{tikzpicture}

\end{fig}
Referring to the above figure, the fixed end moment can be written as $EI\dfrac{\alpha\Delta T}{h}$.

\chapter{Principle of Virtual Displacements}
\label{ch:pvd1}

\section{Virtual work principles}
The virtual work principles include the principle of virtual displacements (PVD) and the principle of virtual forces (PVF). It is often that the principle of virtual displacements is referred to as the principle of virtual work and the principle of virtual forces is referred to as the principle of complementary virtual work. In this chapter, we will discuss the principle of virtual displacements or the principle of virtual work, which represents equivalently the \emph{equilibrium equations} of a system. In Ch.~\ref{ch:pvf}, we shall then discuss the principle of virtual forces or the principle of complementary virtual work, which represents equivalently the \emph{compatibility equations} of a system.

\section{The principle of virtual displacements (PVD)}
\topic{Particles}
\begin{equation}
  \dW = 0 \quad\Longleftrightarrow\quad \text{Equilibrium equations: } \sum F_x = \sum F_y = \sum F_z = 0
\end{equation}
where $\dW$ is the external virtual work defined as
\begin{equation}
  \dW = \Bigl(\sum\F\Bigr)\bullet\delta\mat{u} = \Bigl(\sum F_x\Bigr)\delta u_x + \Bigl(\sum F_y\Bigr)\delta u_y + \Bigl(\sum F_z\Bigr)\delta u_z
\end{equation}
with $\delta u_x$, $\delta u_y$, and $\delta u_z$ being the virtual displacements along the $x$-, $y$-, and $z$-axis.

\topic{Rigid bodies}
\begin{equation}
  \begin{aligned}
    \dW = 0 \quad\Longleftrightarrow\quad &\text{Equilibrium equations: }\\
    &\sum F_x = \sum F_y = \sum F_z = 0\ \text{ and }\ \sum M_x = \sum M_y = \sum M_z = 0
  \end{aligned}
\end{equation}
where
\begin{equation}
  \begin{aligned}
    \dW = \Bigl(\sum\F\Bigr)\bullet\delta\mat{u} + \Bigl(\sum\mat{M}\Bigr)\bullet\delta\bs{\theta}
    &= \Bigl(\sum F_x\Bigr)\delta u_x + \Bigl(\sum F_y\Bigr)\delta u_y + \Bigl(\sum F_z\Bigr)\delta u_z\\
    &\quad + \Bigl(\sum M_x\Bigr)\delta\theta_x + \Bigl(\sum M_y\Bigr)\delta\theta_y + \Bigl(\sum M_z\Bigr)\delta\theta_z
  \end{aligned}
\end{equation}
For both particles and rigid bodies, if the equilibrium equations are satisfied, then $\dW = 0$ is readily seen. Conversely, if $\dW = 0$, it is easy to conclude that the equilibrium equations are satisfied because the virtual displacements $\delta u_x$, $\delta u_y$, and $\delta u_z$, and virtual rotations $\delta\theta_x$, $\delta\theta_y$, and $\delta\theta_z$ are independent and arbitrary. Therefore, the principle is proved.

\topic{Deformable bodies}
\begin{keypoint}
With \emph{kinematically admissible} virtual displacement fields $\delta\mat{u}(\mat{x})$,
\begin{equation*}
  \dW = \dU_i \quad\Longleftrightarrow\quad \text{Equilibrium equations}
\end{equation*}
\end{keypoint}
Virtual displacement fields are called kinematically admissible if the following two conditions are satisfied:
\begin{enumerate}[label=(\arabic*)]
  \item $\delta u_i(\mat{x}) = 0$ for $\mat{x}\in S_u$, where $S_u =$ the surfaces where displacements are specified.
  \item $\delta\eps_{ij} = \frac12(\delta u_{i,j}+\delta u_{j,i})$ for $\mat{x}\in V$.
\end{enumerate}
Note that for structural systems, virtual displacements should satisfy the \emph{homogeneous form} of the essential (or geometric) boundary conditions.

The external virtual work $\dW$ is defined as
\begin{equation}
  \dW = \int_{S_t}\bar{\mat{t}}\bullet\delta\mat{u}\,dS + \int_V\bar{\mat{f}}\bullet\delta\mat{u}\,dV
      = \int_{S_t}\bar t_i\,\delta u_i\,dS + \int_V\bar f_i\,\delta u_i\,dV
\end{equation}
where $S_t =$ the surfaces where tractions are specified, and $\bar{\mat{t}}$ and $\bar{\mat{f}} =$ the specified tractions and body forces, respectively.

The virtual internal energy $\dU_i$ is defined as
\begin{equation}
  \dU_i = \int_V\sigma_{ij}\,\delta\eps_{ij}\,dV
\end{equation}
\begin{fig}
  % Boundary element of a deformable body: stresses and surface tractions
\begin{tikzpicture}[scale=1.0,font=\small]
  % body
  \draw[thick] plot[smooth cycle,tension=0.8] coordinates {(0,0) (1.6,-0.6) (3.2,0.2) (3.0,1.6) (1.4,2.2) (-0.2,1.4)};
  \fill[hatch] (2.55,1.55) -- (3.05,1.35) -- (2.85,1.05) -- cycle;
  \foreach \x in {1.9,2.2,...,3.1}{\draw (\x,-0.35) -- (\x-0.2,-0.6);}
  \draw[->] (3.1,1.5) .. controls (4.2,2.2) and (5.4,2.0) .. (6.0,1.4);
  % wedge element
  \begin{scope}[shift={(6.6,-0.2)}]
    \draw[thick] (0,0) -- (0,1.8) -- (1.4,0) -- cycle;
    \draw[load] (0,0.9) -- (-1.2,0.9) node[left] {$\sigma_{xx}$};
    \draw[load] (0,0.75) -- (0,-0.5) node[below left] {$\sigma_{xy}$};
    \draw[load] (0.7,0) -- (0.7,-1.2) node[below] {$\sigma_{yy}$};
    \draw[load] (0.55,0) -- (-0.5,0) node[below left] {$\sigma_{xy}$};
    \draw[->,thick] (0.7,0.9) -- (1.7,1.7) node[above right] {$\mat{n}$};
    \draw[load] (0.7,0.9) -- (0.7,2.2) node[above] {$\bar t_y$};
    \draw[load] (0.7,0.9) -- (2.0,0.9) node[right] {$\bar t_x$};
  \end{scope}
\end{tikzpicture}

\end{fig}
\begin{align*}
  \textstyle\sum F_x = 0:&\qquad \sigma_{xx}\cdot n_x + \sigma_{xy}\cdot n_y = \bar t_x\\
  \textstyle\sum F_y = 0:&\qquad \sigma_{xy}\cdot n_x + \sigma_{yy}\cdot n_y = \bar t_y
\end{align*}

\begin{proof}
\textbf{Part I: Equilibrium equations $\Rightarrow$ $\dW = \dU_i$.}\\
Equilibrium equations:
\begin{equation*}
  \sigma_{ij,j} + \bar f_i = 0\ \text{ in } V \qquad\text{and}\qquad \sigma_{ij}n_j = \bar t_i\ \text{ on } S_t
\end{equation*}
\begin{equation*}
  \begin{aligned}
    \dU_i &= \int_V\sigma_{ij}\delta\eps_{ij}\,dV = \int_V\sigma_{ij}\tfrac12(\delta u_{i,j}+\delta u_{j,i})\,dV = \int_V\sigma_{ij}\delta u_{i,j}\,dV\quad(\because\sigma_{ij}\text{ is symmetric})\\
          &= \int_V(\sigma_{ij}\delta u_i)_{,j}\,dV - \int_V\sigma_{ij,j}\delta u_i\,dV\\
          &= \int_S\sigma_{ij}n_j\,\delta u_i\,dS - \int_V\sigma_{ij,j}\delta u_i\,dV = \int_{S_t}\sigma_{ij}n_j\,\delta u_i\,dS - \int_V\sigma_{ij,j}\delta u_i\,dV\quad(\because\delta u_i = 0\text{ on }S_u)
  \end{aligned}
\end{equation*}
Making use of the two equilibrium equations yields:
\begin{equation*}
  \dU_i = \int_{S_t}\bar t_i\,\delta u_i\,dS + \int_V\bar f_i\,\delta u_i\,dV
\end{equation*}
Therefore $\dU_i = \dW$.

\textbf{Part II: $\dW = \dU_i$ $\Rightarrow$ Equilibrium equations.}\\
From the expression for $\dU_i$ given in Part I, $\dW = \dU_i$ can be written as
\begin{equation*}
  \int_{S_t}\bar t_i\delta u_i\,dS + \int_V\bar f_i\delta u_i\,dV = \int_{S_t}\sigma_{ij}n_j\delta u_i\,dS + \int_V(-\sigma_{ij,j})\delta u_i\,dV
\end{equation*}
or
\begin{equation*}
  \int_{S_t}(\bar t_i - \sigma_{ij}n_j)\delta u_i\,dS + \int_V(\bar f_i + \sigma_{ij,j})\delta u_i\,dV = 0
\end{equation*}
Because the virtual displacements are arbitrary, the above identity results in (choose $\delta u_i$ as a delta function on $S_t$ or in $V$):
\begin{equation*}
  \sigma_{ij}n_j = \bar t_i\ \text{ on } S_t,\qquad \sigma_{ij,j}+\bar f_i = 0\ \text{ in } V
\end{equation*}
which are the equilibrium equations (one is on the boundary and the other is in the body).
\end{proof}

We shall notice that the principle of virtual displacements has nothing to do with constitutive laws and therefore can be applied to all kinds of materials, elastic, inelastic, etc. For an inelastic body,
\begin{equation}
  \dU_i = \int_V\sigma_{ij}\delta\eps_{ij}\,dV = \int_V\sigma_{ij}(\delta\eps^e_{ij}+\delta\eps^i_{ij})\,dV = \dU + \dU_d
\end{equation}
where $\delta\eps^e_{ij} =$ virtual elastic strains and $\delta\eps^i_{ij} =$ virtual inelastic strains, $\dU =$ virtual strain energy, and $\dU_d =$ virtual dissipation. For example, we often employ the PVD to calculate the limit (collapse) load of structures in conjunction with the upper bound theorem. In this case, we often assume that the rigid plastic law is followed. This means that we neglect $\dU$ and therefore $\dU_i = \dU_d$ and the PVD becomes:
\begin{equation}
  \dW = \dU_d
\end{equation}
\begin{center}
  \begin{minipage}{0.35\linewidth}\centering% Moment-rotation of a rigid-plastic hinge: dissipation (shaded) >> strain energy
\begin{tikzpicture}[scale=1.0,font=\small]
  \draw[->] (0,0) -- (3.4,0) node[right] {$\delta\theta$};
  \draw[->] (0,0) -- (0,2.4) node[above] {$M$};
  \fill[black!15] (0.35,0) -- (0.35,1.7) -- (2.8,1.7) -- (2.8,0) -- cycle;
  \draw[thick] (0,0) -- (0.35,1.7) -- (2.8,1.7);
  \draw[dashed] (0,1.7) -- (0.35,1.7); \node[left] at (0,1.7) {$M_P$};
\end{tikzpicture}
\end{minipage}\hfill
  \begin{minipage}{0.6\linewidth}The energy dissipation due to plastic work (shaded area) is much larger than the strain energy.\end{minipage}
\end{center}

\section{The principle of stationary (minimum) total potential energy and Castigliano's first theorem}
\topic{Elastic bodies}
There exists a strain energy function $U$ such that $\dU = \dU_i$ for elastic bodies. The strain energy is a function of displacements or positions only and is independent of how the body is deformed to this position. Notice that the existence of $U$ does not require the body to be linear. For the special case of linearly elastic materials, $U$ is a quadratic function of displacements or strains.

\topic{Strain energy for structural members}
\begin{align*}
  \text{Axial deformation:}& \quad U = \tfrac12\int EA\,u'^2\,dx\\
  \text{Flexural deformation:}& \quad U = \tfrac12\int EI\,v''^2\,dx\\
  \text{Torsional deformation:}& \quad U = \tfrac12\int GJ\,\theta_x'^2\,dx
\end{align*}

\topic{Conservative forces}
A conservative force is defined as a force that the work done by the force is independent of the path along which it travels. Therefore, the work done is only dependent on the initial and final positions. In this case, there exists a load potential $V$, a function of positions only, such that $\dW = -\delta V$.

\topic{Conservative systems and total potential energy}
If a body is elastic (conservative internally) and the external forces are conservative (conservative externally), the system (internal plus external) is called conservative. For a conservative system, the PVD then becomes
\begin{equation}
  \dU + \delta V = 0
\end{equation}
Define the total potential energy as $\Pi = U + V$; the PVD becomes
\begin{equation}
  \delta\Pi = 0
\end{equation}
This is often called the \emph{principle of stationary total potential energy} or the \emph{principle of minimum total potential energy}.

\topic{Examples of total potential energy}
Total potential energy $\Pi = U + V$.
\begin{fig}
  % Linear spring subjected to P
\begin{tikzpicture}[scale=1.0,font=\small]
  \fill[left color=black!60,right color=black!15] (-0.3,-0.4) rectangle (0,0.4);
  \draw[thick,decorate,decoration={zigzag,segment length=6pt,amplitude=4pt,pre length=4pt,post length=4pt}] (0,0) -- (3,0);
  \node[above] at (1.5,0.3) {$K$};
  \draw[load] (3.2,0) -- (4.2,0) node[right] {$P(\Delta)$};
\end{tikzpicture}

\end{fig}
\begin{gather*}
  \Pi = U + V = \tfrac12 K\Delta^2 - P\Delta\\
  \delta\Pi = K\Delta\cdot\delta\Delta - P\cdot\delta\Delta = (K\Delta - P)\delta\Delta = 0\\
  \Rightarrow K\Delta = P
\end{gather*}
\begin{fig}
  % Beam with distributed load, point load, end moment, inclined bar and spring
\begin{tikzpicture}[scale=0.85,font=\small]
  \def\L{3.2}
  \pic at (0,0) {wallL};
  \draw[member] (0,0) -- (3*\L,0);
  \node[above] at (1.6*\L,0.1) {$EI$};
  % distributed load on first span
  \foreach \x in {0.2,0.55,...,3.0}{\draw[->] (\x,0.9) -- (\x,0.12);} \draw (0.2,0.9) -- (3.0,0.9);
  \node[above] at (1.6,0.95) {$w$};
  \draw[->] (0,1.2) -- (0,2.4) node[right] {$v(x)$};
  % point load at 2L
  \draw[load] (2*\L,1.8) -- (2*\L,0.15) node[pos=0,above] {$\overline{P}$};
  % end moment at 3L
  \draw[->,thick] (3*\L+0.5,0.55) arc (40:200:0.5); \node[above left] at (3*\L-0.3,0.7) {$\overline{M}$};
  % inclined bar at 30 deg from vertical, up-right from the end
  \coordinate (E) at (3*\L,0);
  \coordinate (T) at ($(E)+(60:2.8)$);
  \draw[thinmember] (E) -- (T);
  \pic[rotate=60] at (T) {wallR};
  \node[left] at ($(E)!0.5!(T)+(-0.1,0.1)$) {$EA$};
  \draw[dashed] (E) -- ($(E)+(0,2.2)$);
  \draw[->] ($(E)+(0,1.4)$) arc (90:60:1.4); \node at ($(E)+(0.45,1.75)$) {$30^\circ$};
  \draw[dim] ($(E)+(0.9,0.2)$) -- ($(T)+(0.9,0.2)$) node[midway,right] {$L$};
  % spring below the end
  \draw[thick,decorate,decoration={zigzag,segment length=6pt,amplitude=4pt,pre length=6pt,post length=6pt}] (E) -- ($(E)+(0,-2.6)$);
  \node[right] at ($(E)+(0.2,-1.6)$) {$K_0$};
  \pic at ($(E)+(0,-2.6)$) {fixed};
  % dims
  \draw[dim] (0,-0.9) -- (\L,-0.9) node[midway,fill=white,inner sep=1pt] {$L$};
  \draw[dim] (\L,-0.9) -- (2*\L,-0.9) node[midway,fill=white,inner sep=1pt] {$L$};
  \draw[dim] (2*\L,-0.9) -- (3*\L,-0.9) node[midway,fill=white,inner sep=1pt] {$L$};
  \foreach \x in {0,\L,2*\L,3*\L}{\draw (\x,-0.7) -- (\x,-1.1);}
\end{tikzpicture}

\end{fig}
\begin{gather*}
  U = \frac12\int_0^{3L}EIv''^2\,dx + \frac12\frac{EA}{L}\Delta^2 + \frac12K_0\bigl[v(3L)\bigr]^2\\
  \Delta = -\frac{\sqrt3}{2}v(3L)\quad(\text{``$-$'' means shortening})\\
  V = \int_0^L -w\cdot v(x)\,dx + \overline{P}\cdot v(2L) - \overline{M}\cdot v'(3L)
\end{gather*}

\topic{Rayleigh--Ritz method}
Let the displacement field be constructed through the use of some functions, $\phi_i(x)$ ($i = 1, n$), e.g., $v(x) = \sum_{i=1}^n c_i\phi_i(x)$, where $c_i =$ generalized coefficients. Substituting $v(x)$ into $\Pi$ results in $\Pi(\mat{c}) = U(\mat{c}) + V(\mat{c})$. Then according to the principle of stationary total potential energy,
\begin{equation}
  \delta\Pi = \sum_{i=1}^n\pd{\Pi}{c_i}\delta c_i = 0 \quad\Rightarrow\quad \pd{\Pi}{c_i} = 0,\ i = 1, n \quad\because\delta c_i\text{ are independent}
\end{equation}
The load potential can be expressed as a linear function of $\mat{c}$, $V(\mat{c}) = -\mat{c}^\tr\Pvec_c$, where $\Pvec_c$ is called the generalized load vector. When the material is linearly elastic, the strain energy is a quadratic function of $\mat{c}$, $U(\mat{c}) = \frac12\mat{c}^\tr\Kmat_c\mat{c}$, where $\Kmat_c$ is called the generalized stiffness matrix. With these definitions, the expression $\pd{\Pi}{\mat{c}} = \mat{0}$ becomes
\begin{equation}
  \Kmat_c\mat{c} = \Pvec_c
\end{equation}

\topic{Castigliano's first theorem}
When the external forces $\Pvec$ acting on a structural system are point loads or concentrated loads (including couples), then the strain energy and the load potential can be expressed as functions of the corresponding displacements (including rotations) $\bs{\Delta}$ as
\begin{equation}
  U = U(\Delta_1, \Delta_2, \dots, \Delta_n)\quad\text{and}\quad V = -\sum_{i=1}^n P_i\Delta_i
\end{equation}
According to the principle of stationary total potential energy, we obtain
\begin{equation}
  \delta\Pi = \sum_{i=1}^n\pd{U}{\Delta_i}\delta\Delta_i - \sum_{i=1}^n P_i\delta\Delta_i = \sum_{i=1}^n\Bigl(\pd{U}{\Delta_i} - P_i\Bigr)\delta\Delta_i = 0
\end{equation}
As $\delta\bs{\Delta}$ are independent and arbitrary, we conclude from the above equation that
\begin{equation}
  \pd{U}{\Delta_i} = P_i,\quad i = 1, n
\end{equation}
This is often referred to as \emph{Castigliano's first theorem}. Note that the theorem can be applied to nonlinearly elastic material as well. In this case, $U$ is not a quadratic function of $\bs{\Delta}$.

The stiffness equations for a linear structural system can be derived from the principle of stationary total potential energy or Castigliano's first theorem. The strain energy and load potential are of the form:
\begin{equation}
  U = \tfrac12\bs{\Delta}^\tr\Kmat\bs{\Delta}\quad\text{and}\quad V = -\bs{\Delta}^\tr\Pvec
\end{equation}
Therefore, from $\delta\Pi = 0$ (or $\partial\Pi/\partial\Delta_i = 0$ or $\partial\Pi/\partial\bs{\Delta} = \mat{0}$) or $\partial U/\partial\Delta_i = P_i$ ($\partial U/\partial\bs{\Delta} = \Pvec$), we have
\begin{equation}
  \Kmat\bs{\Delta} = \Pvec
\end{equation}

\topic{Example of Castigliano's first theorem (nonlinear elastic case)}
\begin{fig}
  % Nonlinear elastic bar subjected to F
\begin{tikzpicture}[scale=1.0,font=\small]
  \fill[left color=black!60,right color=black!15] (-0.3,-0.4) rectangle (0,0.4);
  \draw[thick] (0,-0.15) rectangle (4,0.15);
  \draw[load] (4.1,0) -- (5.0,0) node[right] {$F(\delta)$};
\end{tikzpicture}

\end{fig}
\begin{gather*}
  F = c\cdot\delta^{\frac12}\\
  U = \int_0^\delta F\,d\delta = \int_0^\delta c\cdot\delta^{\frac12}\,d\delta = \frac23 c\cdot\delta^{\frac32}\\
  \pd{U}{\delta} = c\cdot\delta^{\frac12} = F
\end{gather*}

\begin{example}[Plastic collapse load or limit load of beams]
Determine this structure's plastic collapse load by P.V.W.
\begin{fig}
  % Propped cantilever with two point loads W at the third points
\begin{tikzpicture}[scale=1.0,font=\small]
  \def\L{9}
  \pic at (0,0) {wallL}; \pic at (\L,0) {pin};
  \draw[member] (0,0) -- (\L,0);
  \node[below right] at (0,0) {$A$}; \node[below] at (3,-0.05) {$B$}; \node[below] at (6,-0.05) {$C$}; \node[below left] at (\L,0) {$D$};
  \draw[load] (3,1.6) -- (3,0.15) node[pos=0,above] {$W$};
  \draw[load] (6,1.6) -- (6,0.15) node[pos=0,above] {$W$};
  \draw[dim] (0,-0.9) -- (3,-0.9) node[midway,fill=white,inner sep=1pt] {$\frac L3$};
  \draw[dim] (3,-0.9) -- (6,-0.9) node[midway,fill=white,inner sep=1pt] {$\frac L3$};
  \draw[dim] (6,-0.9) -- (\L,-0.9) node[midway,fill=white,inner sep=1pt] {$\frac L3$};
  \foreach \x in {0,3,6,\L}{\draw (\x,-0.7) -- (\x,-1.1);}
\end{tikzpicture}

\end{fig}
\underline{Mechanism 1:} nodes $A$ and $C$ become plastic hinges.
\begin{fig}
  % Mechanism 1: plastic hinges at A and C
\begin{tikzpicture}[scale=1.0,font=\small]
  \pic at (0,0) {wallL}; \pic at (9,0) {pin};
  \draw[thick] (0,0) -- (9,0);
  \draw[dashed,thick] (0,0) -- (6,-1.2) -- (9,0);
  \node[node dot] at (0,0) {}; \node[node dot] at (6,-1.2) {};
  \draw[->] (2.2,0) arc (0:-11.3:2.2); \node[right] at (2.3,-0.3) {$\delta\theta$};
\end{tikzpicture}

\end{fig}
\begin{equation*}
  \dW = W\Bigl(\frac L3\delta\theta\Bigr) + W\Bigl(\frac{2L}{3}\delta\theta\Bigr) = WL\delta\theta
\end{equation*}
and virtual dissipation
\begin{equation*}
  \dU_d = M_p\delta\theta + M_p(3\delta\theta) = 4M_p\delta\theta
\end{equation*}
where $M_p =$ sectional plastic moment or moment capacity. By P.V.W., we have
\begin{equation*}
  WL\delta\theta = 4M_p\delta\theta \Rightarrow W = \frac{4M_p}{L}
\end{equation*}
\underline{Mechanism 2:} nodes $A$ and $B$ become plastic hinges.
\begin{fig}
  % Mechanism 2: plastic hinges at A and B
\begin{tikzpicture}[scale=1.0,font=\small]
  \pic at (0,0) {wallL}; \pic at (9,0) {pin};
  \draw[thick] (0,0) -- (9,0);
  \draw[dashed,thick] (0,0) -- (3,-1.2) -- (9,0);
  \node[node dot] at (0,0) {}; \node[node dot] at (3,-1.2) {};
  \draw[->] (6.8,0) arc (180:191.3:2.2); \node[left] at (6.7,-0.3) {$\delta\theta$};
\end{tikzpicture}

\end{fig}
\begin{equation*}
  \dW = W\Bigl(\frac L3 2\delta\theta\Bigr) + W\Bigl(\frac L3\delta\theta\Bigr) = WL\delta\theta
\end{equation*}
and virtual dissipation
\begin{equation*}
  \dU_d = M_p(2\delta\theta) + M_p(3\delta\theta) = 5M_p\delta\theta
\end{equation*}
By P.V.W., we have
\begin{equation*}
  WL\delta\theta = 5M_p\delta\theta \Rightarrow W = \frac{5M_p}{L}
\end{equation*}
\underline{Mechanism 3:} nodes $B$ and $C$ become plastic hinges.
\begin{fig}
  % Mechanism 3: plastic hinges at B and C
\begin{tikzpicture}[scale=1.0,font=\small]
  \pic at (0,0) {wallL}; \pic at (9,0) {pin};
  \draw[thick] (0,0) -- (9,0);
  \draw[dashed,thick] (3,0) -- (6,-1.2) -- (9,0);
  \node[node dot] at (3,0) {}; \node[node dot] at (6,-1.2) {};
  \draw[->] (4.6,0) arc (0:-21.8:1.6); \node[right] at (4.7,-0.35) {$\delta\theta$};
\end{tikzpicture}

\end{fig}
\begin{equation*}
  \dW = W\Bigl(\frac L3\delta\theta\Bigr) = \frac{WL}{3}\delta\theta
\end{equation*}
and virtual dissipation
\begin{equation*}
  \dU_d = M_p\delta\theta + M_p(2\delta\theta) = 3M_p\delta\theta
\end{equation*}
By P.V.W., we have
\begin{equation*}
  \frac{WL}{3}\delta\theta = 3M_p\delta\theta \Rightarrow W = \frac{9M_p}{L}
\end{equation*}
According to the upper bound theory, the collapse load shall be the minimum of these three values. However, it doesn't mean that the real collapse load is the minimum one according to the upper bound theorem. We can plot the moment diagram for the mechanism with the minimum failure load. If there is no section that violates the yield condition, then the upper bound solution is also the exact one; try to draw the moment diagrams for the above three mechanisms. You will find that the internal moments for mechanism 1 are all less than or equal to $M_p$. Therefore, $W = \dfrac{4M_p}{L}$ is the exact failure load. However, some of the internal moments for mechanisms 2 and 3 are greater than $M_p$.
\end{example}

\begin{example}[Braced frame]
Find the plastic collapse load in terms of $M_p$ and $L$. Note that $ac$ is an axial bar with the axial capacity of $A\sigma_y$, which is equal to $\dfrac{M_p}{L}$.
\begin{fig}
  % Braced frame: geometry (left) and collapse mechanism (right)
\begin{tikzpicture}[scale=0.8,font=\small]
  % ---- geometry ----
  \begin{scope}
    \coordinate (a) at (0,0); \coordinate (b) at (0,3); \coordinate (c) at (4,3); \coordinate (d) at (4,-2);
    \draw[member] (a) -- (b) -- (c) -- (d); \draw[thinmember] (a) -- (c);
    \pic at (a) {pin}; \pic at (d) {pin};
    \node[left] at (a) {a}; \node[above left] at (b) {b}; \node[above right] at (c) {c}; \node[right] at (d) {d};
    \draw[load] (c) -- ++(1.4,0) node[right] {P};
    \node[above] at (2,3.05) {$2M_p$}; \node[below] at (2,2.95) {$4L$};
    \node[left,align=right] at (-0.1,1.5) {$3M_p$\\$3L$};
    \node[right,align=left] at (4.1,0.5) {$M_p$\\$5L$};
    \node[below right] at (1.6,1.2) {$\sigma_y A = \dfrac{M_p}{L}$};
  \end{scope}
  % ---- mechanism ----
  \begin{scope}[xshift=9cm]
    \coordinate (a) at (0,0); \coordinate (b) at (-0.6,3); \coordinate (c) at (3.4,3.2); \coordinate (d) at (4,-2);
    \draw[thinmember] (a) -- (b) -- (c) -- (d);
    \pic at (a) {pin}; \pic at (d) {pin};
    \draw[dashed] (a) -- (0,3.2); \draw[dashed] (d) -- (4,3.4);
    \draw[->] (0,1.2) arc (90:101.3:1.2); \node[left] at (-0.3,1.5) {$\delta\theta$};
    \draw[->] (4,1.0) arc (90:96.8:3.0); \node[left] at (3.85,0.9) {$\frac35\delta\theta$};
    \draw[load] (c) -- ++(1.4,0) node[right] {P};
    \draw[->] (b) -- ++(1.0,0) node[above] {$3L\delta\theta$}; \draw[->] (c) -- ++(1.0,0.2) ; \node[above] at (3.8,3.5) {$3L\delta\theta$};
    \draw[load] ($(a)!0.55!(c)$) -- ++(-0.6,-0.5); \node[right] at ($(a)!0.55!(c)+(0.1,0.2)$) {$\frac{M_p}{L}$};
    \draw[load] ($(a)!0.1!(c)$) -- ++(0.6,0.5);
  \end{scope}
\end{tikzpicture}

\end{fig}
\begin{align*}
  \dU_d &= 2M_p(\delta\theta) + M_p\Bigl(\frac35\delta\theta\Bigr) &&\text{or}& \dU_d &= 2M_p(\delta\theta) + M_p\Bigl(\frac35\delta\theta\Bigr) + \frac{M_p}{L}\Bigl(\frac45 3L\delta\theta\Bigr)\\
  \dW &= P(3L\delta\theta) - \frac{M_p}{L}\Bigl(\frac45 3L\delta\theta\Bigr) &&\text{or}& \dW &= P(3L\delta\theta)
\end{align*}
Note the former case considers the axial bar effect as applied loads while the latter treats it as internal forces.
\begin{equation*}
  \dU_d = \dW \;\rightarrow\; P = \frac{5M_p}{3L}.\qquad\text{Without the brace, } P = \frac{2.6M_p}{3L} < \frac{5M_p}{3L},\text{ which is expected.}
\end{equation*}
\end{example}

\begin{example}[Portal frame]
\leavevmode
\begin{fig}
  % Portal frame: (a) loading and (b)-(d) collapse mechanisms
\begin{tikzpicture}[scale=0.75,font=\footnotesize]
  % ---- (a) ----
  \begin{scope}
    \draw[member] (0,0) -- (0,3) -- (5,3) -- (5,0);
    \pic at (0,0) {fixed}; \pic at (5,0) {fixed};
    \node[below] at (0,-0.3) {$A$}; \node[below] at (5,-0.3) {$D$};
    \node[above left] at (0,3) {$B$}; \node[above right] at (5,3) {$C$}; \node[below] at (2.5,2.9) {$E$};
    \draw[load] (-1.2,3) -- (-0.1,3) node[pos=0,above] {$P$};
    \draw[load] (2.5,4.3) -- (2.5,3.1) node[pos=0,right] {$2P$};
    \node[above] at (1.4,3.05) {$2M_p$}; \node[right] at (0.05,1.5) {$M_p$}; \node[left] at (4.95,1.5) {$M_p$};
    \draw[dim] (0,4.6) -- (2.5,4.6) node[midway,fill=white,inner sep=1pt] {$l/2$};
    \draw[dim] (2.5,4.6) -- (5,4.6) node[midway,fill=white,inner sep=1pt] {$l/2$};
    \draw[dim] (-1.6,0) -- (-1.6,3) node[midway,left] {$0.6l$};
    \node at (2.5,-1.1) {(a)};
  \end{scope}
  % ---- (b) beam mechanism ----
  \begin{scope}[xshift=8.5cm]
    \draw[member] (0,0) -- (0,3) -- (2.5,2.2) -- (5,3) -- (5,0);
    \pic at (0,0) {fixed}; \pic at (5,0) {fixed};
    \node[node dot] at (0,3) {}; \node[node dot] at (2.5,2.2) {}; \node[node dot] at (5,3) {};
    \node[below] at (0,-0.3) {$A$}; \node[below] at (5,-0.3) {$D$};
    \node[above left] at (0,3) {$B$}; \node[above right] at (5,3) {$C$}; \node[below] at (2.5,2.1) {$E$};
    \draw[load] (-1.2,3) -- (-0.1,3) node[pos=0,above] {$P_{c1}$};
    \draw[load] (2.5,3.5) -- (2.5,2.3) node[pos=0,right] {$2P_{c1}$};
    \draw[->] (1.2,3) arc (0:-17.7:1.2); \node[above] at (1.0,2.95) {$\theta$};
    \draw[->] (3.8,3) arc (180:197.7:1.2); \node[above] at (4.0,2.95) {$\theta$};
    \node[below] at (2.5,1.5) {$2\theta$};
    \draw[dim] (5.6,2.2) -- (5.6,3) node[midway,right] {$\frac{l\theta}{2}$};
    \node at (2.5,-1.1) {(b)};
  \end{scope}
  % ---- (c) sway mechanism ----
  \begin{scope}[yshift=-7.5cm]
    \draw[member] (0,0) -- (0.6,3) -- (5.6,3) -- (5,0);
    \pic at (0,0) {fixed}; \pic at (5,0) {fixed};
    \node[node dot] at (0,0) {}; \node[node dot] at (0.6,3) {}; \node[node dot] at (5.6,3) {}; \node[node dot] at (5,0) {};
    \node[below] at (0,-0.3) {$A$}; \node[below] at (5,-0.3) {$D$};
    \node[above left] at (0.6,3) {$B$}; \node[above right] at (5.6,3) {$C$}; \node[below] at (3.1,2.9) {$E$};
    \draw[load] (-0.7,3) -- (0.5,3) node[pos=0,above] {$P_{c2}$};
    \draw[load] (3.1,4.3) -- (3.1,3.1) node[pos=0,right] {$2P_c$};
    \draw[dashed] (0,0) -- (0,3.3); \draw[dashed] (5,0) -- (5,3.3);
    \draw[dim] (0,3.8) -- (0.6,3.8) node[midway,above] {$0.6l\theta$};
    \node[right] at (0.4,1.6) {$\theta$}; \node[left] at (4.6,1.6) {$\theta$};
    \node at (2.5,-1.1) {(c)};
  \end{scope}
  % ---- (d) combined mechanism ----
  \begin{scope}[xshift=8.5cm,yshift=-7.5cm]
    \draw[member] (0,0) -- (0.6,3) -- (3.1,2.2) -- (5.6,3) -- (5,0);
    \pic at (0,0) {fixed}; \pic at (5,0) {fixed};
    \node[node dot] at (0,0) {}; \node[node dot] at (3.1,2.2) {}; \node[node dot] at (5.6,3) {}; \node[node dot] at (5,0) {};
    \node[below] at (0,-0.3) {$A$}; \node[below] at (5,-0.3) {$D$};
    \node[above left] at (0.6,3) {$B$}; \node[above right] at (5.6,3) {$C$}; \node[below] at (3.1,2.1) {$E$};
    \draw[load] (-0.7,3) -- (0.5,3) node[pos=0,above] {$P_{c3}$};
    \draw[load] (3.1,3.5) -- (3.1,2.3) node[pos=0,right] {$2P_c$};
    \draw[dashed] (0,0) -- (0,3.3); \draw[dashed] (5,0) -- (5,3.3);
    \draw[dim] (0,3.8) -- (0.6,3.8) node[midway,above] {$0.6l\theta$};
    \draw[dim] (-0.8,2.2) -- (-0.8,3) node[midway,left] {$\frac{l\theta}{2}$};
    \node[right] at (0.4,1.6) {$\theta$}; \node[left] at (4.6,1.2) {$\theta$};
    \node[below] at (3.1,1.5) {$2\theta$}; \node[above] at (4.6,2.7) {$2\theta$};
    \node at (2.5,-1.1) {(d)};
  \end{scope}
\end{tikzpicture}

\end{fig}
\begin{enumerate}[label=\alph*.,start=2,leftmargin=*]
  \item $\dW = 2P\bigl(\frac l2\delta\theta\bigr) = Pl\delta\theta$,\quad $\dU_d = M_P(\delta\theta)\times2 + 2M_P(2\delta\theta) = 6M_P\delta\theta$
    \begin{equation*}
      Pl\delta\theta = 6M_P\delta\theta \Rightarrow P = \frac{6M_P}{l}
    \end{equation*}
  \item $\dW = P\bigl(\frac{3l}{5}\delta\theta\bigr)$,\quad $\dU_d = M_P(\delta\theta)\times4 = 4M_P\delta\theta$
    \begin{equation*}
      \frac{3Pl}{5}\delta\theta = 4M_P\delta\theta \Rightarrow P = \frac{20M_P}{3l}
    \end{equation*}
  \item $\dW = P\bigl(\frac{3l}{5}\delta\theta\bigr) + 2P\bigl(\frac l2\delta\theta\bigr) = \frac{8Pl}{5}\delta\theta$,\quad $\dU_d = M_P(\delta\theta)\times2 + M_P(2\delta\theta) + 2M_P(2\delta\theta) = 8M_P\delta\theta$
    \begin{equation*}
      \frac{8Pl}{5}\delta\theta = 8M_P\delta\theta \Rightarrow P = \frac{5M_P}{l}
    \end{equation*}
\end{enumerate}
\end{example}

\begin{example}[Truss]
This example is adopted from the textbook, p.~141, Example~6.1.
\end{example}

\begin{example}[P.V.D.]
\leavevmode
\begin{fig}
  % Cantilever beam AB with a tip load P
\begin{tikzpicture}[scale=1.0,font=\small]
  \pic at (0,0) {wallL};
  \draw[thick] (0,-0.15) rectangle (5,0.15);
  \node[above right] at (0,0.15) {A}; \node[above left] at (5,0.15) {B};
  \draw[->] (0.1,0) -- (1.4,0) node[right] {$x$};
  \draw[->] (0,0.5) -- (0,2.0) node[above] {$Y$};
  \draw[load] (5,0.2) -- (5,1.5) node[above] {$P$};
\end{tikzpicture}

\end{fig}
The exact displacement field is known as $v(x) = \dfrac{P}{EI}\Bigl(L\dfrac{x^2}{2} - \dfrac{x^3}{6}\Bigr)$, which can be obtained by integrating $EIv'' = M(x) = P(L-x)$ twice. Assume the virtual displacement field is $u(x)$, which can be any function of $x$ but shall be kinematically admissible, i.e.\ $u(0) = 0$ and $u'(0) = 0$.

Since
\begin{equation*}
  \dU = \int_0^L EIv''u''\,dx\qquad\text{and}\qquad \dW = Pu_B = Pu(L)
\end{equation*}
By P.V.D., we have
\begin{equation*}
  Pu(L) = \int_0^L EIv''u''\,dx
\end{equation*}
By making use of the kinematically admissible conditions, $u(0) = 0$ and $u'(0) = 0$, and $EIv'' = P(L-x)$, we have
\begin{equation*}
  Pu(L) = \int_0^L P(L-x)u''\,dx
\end{equation*}
We will choose different $u(x)$ to verify the above identity.
\begin{description}[leftmargin=2em]
  \item[Case 1:] $u(x) = \dfrac{u_B}{L^2}x^2$
    \begin{align*}
      \text{LHS} &= Pu_B\\
      \text{RHS} &= \int_0^L P(L-x)\frac{2u_B}{L^2}\,dx = Pu_B
    \end{align*}
  \item[Case 2:] $u(x) = u_B\Bigl(-\cos\dfrac{\pi x}{2L}+1\Bigr)$\\
    That RHS $=$ LHS can be shown easily.
\end{description}
\emph{As the real displacement field is exact, the P.V.D. relationship would be satisfied regardless of the chosen virtual displacement fields. However, if $v(x)$ is not exact, there will be some differences as discussed below.}

Assume $v(x) = \dfrac{v_B}{L^2}x^2$; we would like to find $v_B$.
\begin{description}[leftmargin=2em]
  \item[Case 1:] $u(x) = \dfrac{u_B}{L^2}x^2$
    \begin{align*}
      \text{LHS} &= Pu_B\\
      \text{RHS} &= \frac{4EI}{L^3}u_Bv_B\\
      \because\text{RHS} = \text{LHS} &\Rightarrow v_B = \frac{PL^3}{4EI}\qquad\Bigl(\text{The exact one is } v_B = \frac{PL^3}{3EI}\Bigr)
    \end{align*}
  \item[Case 2:] $u(x) = \dfrac{u_B}{L^4}x^4$
    \begin{align*}
      \text{LHS} &= Pu_B\\
      \text{RHS} &= \frac{8EI}{L^3}u_Bv_B\\
      \because\text{RHS} = \text{LHS} &\Rightarrow v_B = \frac{PL^3}{8EI}
    \end{align*}
\end{description}
\begin{note}
You may prefer to assume $v(x) = cx^2$, where $c$ is to be determined by the PVD. For case 1, you assume $u(x) = x^2$ and for case 2 $u(x) = x^4$. You will get the same solutions; namely $c = \dfrac{PL}{4EI}$ for case 1 and $c = \dfrac{PL}{8EI}$ for case 2. Please check.

When you assume $v(x) = c\Bigl(L\dfrac{x^2}{2} - \dfrac{x^3}{6}\Bigr)$ you will get the same $c = \dfrac{P}{EI}$ for cases 1 and 2, which is the exact solution. Of course, in reality, it is impossible to know the exact function form. We will employ the Rayleigh--Ritz method to solve it by assuming $v(x) = c_1x^2 + c_2x^3$.
\end{note}

\topic{Conclusion}
\begin{enumerate}
  \item The real and virtual displacement fields are independent, but the virtual displacement field should be kinematically admissible.
  \item When the real displacement field is exact, whatever you choose for the virtual displacement field, exact solutions will be obtained from P.V.D.
  \item When the real displacement field is not exact, the solutions (approximate) obtained using P.V.D. depend on the chosen virtual displacement fields. Also, these approximate displacements will be less than the exact ones.
\end{enumerate}
\end{example}

\begin{example}[Rayleigh--Ritz method and PVD]
Solve the above example using the Rayleigh--Ritz method. Assume $v(x) = c_1x^2 + c_2x^3$, which already satisfies the geometric boundary conditions $v(0) = 0$ and $v'(0) = 0$. Then the total potential energy $\Pi$ can be obtained as follows:
\begin{gather*}
  \Pi = U + V\\
  U = \frac12\int EIv''^2\,dx = \frac12EI(4c_1^2L + 12c_1c_2L^2 + 12c_2^2L^3) = \frac12\mat{c}^\tr\Kmat_c\mat{c}\\
  V = -Pv(L) = -P(c_1L^2 + c_2L^3) = -\mat{c}^\tr\Pvec_c\\
  \Pi = \frac12\mat{c}^\tr\Kmat_c\mat{c} - \mat{c}^\tr\Pvec_c
\end{gather*}
$\pd{\Pi}{c_1} = \pd{\Pi}{c_2} = 0$, we have $\Kmat_c\mat{c} = \Pvec_c$ or
\begin{equation*}
  EI\bmat{4L & 6L^2\\ 6L^2 & 12L^3}\Bmat{c_1\\ c_2} = \Bmat{PL^2\\ PL^3}
  \;\Rightarrow\; \Bmat{c_1\\ c_2} = \frac{P}{EI}\Bmat{L/2\\ -1/6}
  \;\Rightarrow\; v(x) = \frac{P}{EI}\Bigl(\frac L2x^2 - \frac16x^3\Bigr)
\end{equation*}
\begin{note}
In terms of computational efforts, the PVD
\begin{equation*}
  \dW = \dU\qquad \dW = Pu(L)\qquad \dU = \int_0^L EIv''u''\,dx
\end{equation*}
is more effective to use. Let $v(x) = c_1x^2 + c_2x^3$ and

Case 1: $u(x) = x^2$:\quad $PL^2 = \int_0^L EI(2c_1+6c_2x)(2)\,dx \Rightarrow PL^2 = EI(4c_1L + 6c_2L^2)$\\
\phantom{Case 1: }$u(x) = x^3$:\quad $PL^3 = \int_0^L EI(2c_1+6c_2x)(6x)\,dx \Rightarrow PL^3 = EI(6c_1L^2 + 12c_2L^3)$\\
the same system as above.

Case 2: $u(x) = x^4$ and $x^5$
\begin{equation*}
  EI\bmat{8L^3 & 18L^4\\ 10L^4 & 24L^5}\Bmat{c_1\\ c_2} = \Bmat{PL^4\\ PL^5}
  \;\Rightarrow\; \Bmat{c_1\\ c_2} = \frac{P}{EI}\Bmat{L/2\\ -1/6}
  \;\Rightarrow\; v(x) = \frac{P}{EI}\Bigl(\frac L2x^2 - \frac16x^3\Bigr)
\end{equation*}
You shall be aware that if $u(x)$ (virtual displacement) has the same form as $v(x)$ the system matrix is symmetric; otherwise, it is unsymmetric. The same solutions ($c_1$ and $c_2$) are obtained even though the systems are different, which is because the form of $v(x)$ is exact. If an inexact form, say $v(x) = c_1x^3 + c_2x^4$, is assumed, then you will have different solutions. Try it.
\end{note}
\end{example}

\chapter[PVD in Framework Analysis]{Principle of Virtual Displacements in Framework Analysis}
\label{ch:pvd2}

\section{Shape functions and stiffness matrix}
\subsection{Two-noded (linear) axial element}
\begin{equation}
  u(x) = \Bigl(1-\frac xL\Bigr)u_1 + \frac xL u_2 = N_1(x)u_1 + N_2(x)u_2 = \langle N_1, N_2\rangle\Bmat{u_1\\ u_2} = \NN\bs{\Delta}
\end{equation}
where the shape functions are
\begin{equation}
  N_1(x) = 1-\frac xL\quad\text{and}\quad N_2(x) = \frac xL
\end{equation}
Note that shape function $N_i(x)$ represents the displacement field associated with $\Delta_i = 1$ and other $\Delta_i$'s $= 0$.

The strain field is then given by
\begin{equation}
  \eps(x) = u'(x) = \NN'\bs{\Delta}
\end{equation}
where $\NN' = \langle N_1', N_2'\rangle = \langle -1/L, 1/L\rangle$.

The virtual strain energy is
\begin{equation}
  \dU = \int_0^L\delta\eps\,E\eps\,A\,dx = \int_0^L(\delta\eps)^\tr E\eps\,A\,dx = \delta\bs{\Delta}^\tr\Bigl(\int_0^L EA\NN'^{\tr}\NN'\,dx\Bigr)\bs{\Delta}
\end{equation}
Noting that $\dU = \delta\bs{\Delta}^\tr\kmat\bs{\Delta}$, therefore,
\begin{gather}
  \kmat = \int_0^L EA\NN'^{\tr}\NN'\,dx = EA\int_0^L\NN'^{\tr}\NN'\,dx\quad\text{if } EA = \text{constant}\\
  \Rightarrow\quad \kmat = \frac{EA}{L}\bmat{1 & -1\\ -1 & 1}\quad\text{if } EA = \text{constant.}
\end{gather}

\subsection{Flexural element: bent about the \texorpdfstring{$z$}{z}-axis}
The displacement is
\begin{equation}
  v(x) = N_1(x)v_1 + N_2(x)\theta_{z1} + N_3(x)v_2 + N_4(x)\theta_{z2} = \sum_{i=1}^4 N_i(x)\Delta_i = \NN\bs{\Delta}
\end{equation}
where the shape functions are
\begin{align*}
  N_1(x) &= 2\Bigl(\frac xL\Bigr)^3 - 3\Bigl(\frac xL\Bigr)^2 + 1\\
  N_2(x) &= \frac{x^3}{L^2} - 2\frac{x^2}{L} + x\\
  N_3(x) &= -2\Bigl(\frac xL\Bigr)^3 + 3\Bigl(\frac xL\Bigr)^2\\
  N_4(x) &= \frac{x^3}{L^2} - \frac{x^2}{L}
\end{align*}
Similarly, shape function $N_i(x)$ is the displacement field due to $\Delta_i = 1$ and the other $\Delta_i$'s $= 0$. For example, $N_2(x)$ satisfies the following four conditions: $N_2(0) = 0$, $N_2'(0) = 1$, $N_2(L) = 0$, $N_2'(L) = 0$; these correspond to a hinged support with a unit rotation for the left end and a fixed support for the right end.

The above shape functions can be derived as follows. First, as there are four degrees of freedom, a cubic displacement field can be assumed; namely
\begin{equation}
  v(x) = a_0 + a_1x + a_2x^2 + a_3x^3
\end{equation}
Then use of the four boundary conditions: $v(0) = v_1$, $v(L) = v_2$, $v'(0) = \theta_{z1}$, and $v'(L) = \theta_{z2}$ results in
\begin{equation}
  \bmat{1 & 0 & 0 & 0\\ 1 & L & L^2 & L^3\\ 0 & 1 & 0 & 0\\ 0 & 1 & 2L & 3L^2}\Bmat{a_0\\ a_1\\ a_2\\ a_3} = \Bmat{v_1\\ v_2\\ \theta_{z1}\\ \theta_{z2}}
  \quad\Rightarrow\quad
  \begin{aligned}
    a_0 &= v_1\\ a_1 &= \theta_{z1}\\ a_2 &= \frac{3}{L^2}(-v_1+v_2) - \frac1L(2\theta_{z1}+\theta_{z2})\\ a_3 &= \frac{2}{L^3}(v_1-v_2) + \frac{1}{L^2}(\theta_{z1}+\theta_{z2})
  \end{aligned}
\end{equation}
Substituting $a_0, a_1, a_2$, and $a_3$ into $v(x)$ and rearranging the terms results in the shape functions immediately.

The axial strain along the depth of a cross section due to bending is given by
\begin{equation}
  \eps(y) = -yv'' = -y\NN''\bs{\Delta}
\end{equation}
The virtual strain energy is
\begin{equation}
  \begin{aligned}
    \dU &= \int_V\delta\eps\,\sigma\,dV = \int_V\delta\eps\,E\eps\,dV\\
        &= \int_V\delta v''Ey^2v''\,dV = \int_0^L EI_z\,\delta v''^{\tr}v''\,dx\\
        &= \delta\bs{\Delta}^\tr\Bigl(\int_0^L EI_z\NN''^{\tr}\NN''\,dx\Bigr)\bs{\Delta} = \delta\bs{\Delta}^\tr\kmat\bs{\Delta}
  \end{aligned}
\end{equation}
Therefore,
\begin{equation}
  \kmat = \int_0^L EI_z\NN''^{\tr}\NN''\,dx\qquad\text{or}\qquad k_{ij} = \int_0^L EI_zN_i''(x)N_j''(x)\,dx
\end{equation}
If $EI_z =$ constant, it is easy to obtain, on the basis of this equation and the shape functions, that
\begin{equation}
  \kmat = \frac{EI_z}{L}
  \begin{bmatrix}
    \dfrac{12}{L^2} & \dfrac6L & -\dfrac{12}{L^2} & \dfrac6L\\[8pt]
    & 4 & -\dfrac6L & 2\\[8pt]
    & & \dfrac{12}{L^2} & -\dfrac6L\\[8pt]
    \multicolumn{2}{l}{\textit{Symm.}} & & 4
  \end{bmatrix}
\end{equation}

\subsection{Flexural element: bent about the \texorpdfstring{$y$}{y}-axis}
The displacement field is:
\begin{equation}
  w(x) = N_1(x)w_1 + N_2(x)\theta_{y1} + N_3(x)w_2 + N_4(x)\theta_{y2} = \sum_{i=1}^4 N_i(x)\Delta_i = \NN\bs{\Delta}
\end{equation}
where the shape functions are given by
\begin{align*}
  N_1(x) &= 2\Bigl(\frac xL\Bigr)^3 - 3\Bigl(\frac xL\Bigr)^2 + 1\\
  N_2(x) &= -\Bigl(\frac{x^3}{L^2} - 2\frac{x^2}{L} + x\Bigr)\\
  N_3(x) &= -2\Bigl(\frac xL\Bigr)^3 + 3\Bigl(\frac xL\Bigr)^2\\
  N_4(x) &= -\Bigl(\frac{x^3}{L^2} - \frac{x^2}{L}\Bigr)
\end{align*}
Note that shape functions $N_1(x)$ and $N_3(x)$ are the same as those for the $z$-axis bending but shape functions $N_2(x)$ and $N_4(x)$ have a sign change compared with those for the $z$-axis bending. The sign change is due to the fact that $\theta_y(x) = -w'(x)$ but $\theta_z(x) = v'(x)$.

\subsection{Torsional element}
The shape functions for $\theta_x$ are the same as those for the axial displacement. Therefore, the stiffness matrix for a torsional element takes the form:
\begin{equation}
  \kmat = \int_0^L GJ\NN'^{\tr}\NN'\,dx \quad\Rightarrow\quad \kmat = \frac{GJ}{L}\bmat{1 & -1\\ -1 & 1}\ \text{if } GJ = \text{constant.}
\end{equation}

\section{Special purpose elements}
\subsection{Three-noded (quadratic) axial element}
The displacement field is given by
\begin{equation}
  u(x) = \sum_{i=1}^3 N_i(x)u_i = \NN(x)\bs{\Delta}
\end{equation}
where $\NN(x) = \langle N_1(x), N_2(x), N_3(x)\rangle$ and $\bs{\Delta} = \langle u_1, u_2, u_3\rangle^\tr$. The shape functions are given by
\begin{fig}
  % Shape functions of the three-noded (quadratic) axial element
\begin{tikzpicture}[scale=1.0,font=\small,
  declare function={Na(\x)=1-3*\x+2*\x*\x; Nb(\x)=4*\x-4*\x*\x; Nc(\x)=-\x+2*\x*\x;}]
  \def\L{6}
  % N1
  \begin{scope}
    \draw (0,0) -- (\L,0);
    \draw[member,domain=0:1,samples=60,variable=\t] plot ({\t*\L},{Na(\t)*1.2});
    \draw[dim] (-0.3,0) -- (-0.3,1.2) node[midway,left] {1};
    \node[below] at (0,-0.05) {0}; \node[below] at (\L/2,-0.05) {$\frac L2$}; \node[below] at (\L,-0.05) {$L$};
    \node[right] at (\L+0.4,0.6) {$N_1(x) = \dfrac{2}{L^2}\Bigl(x-\dfrac L2\Bigr)(x-L) = 1-3\dfrac xL+2\dfrac{x^2}{L^2}$};
  \end{scope}
  % N2
  \begin{scope}[yshift=-2.6cm]
    \draw (0,0) -- (\L,0);
    \draw[member,domain=0:1,samples=60,variable=\t] plot ({\t*\L},{Nb(\t)*1.2});
    \draw[dim] (\L/2,0) -- (\L/2,1.2) node[midway,right] {1};
    \node[below] at (0,-0.05) {0}; \node[below] at (\L/2,-0.05) {$\frac L2$}; \node[below] at (\L,-0.05) {$L$};
    \node[right] at (\L+0.4,0.6) {$N_2(x) = \dfrac{-4}{L^2}x(x-L) = 4\dfrac xL-4\dfrac{x^2}{L^2}$};
  \end{scope}
  % N3
  \begin{scope}[yshift=-5.2cm]
    \draw (0,0) -- (\L,0);
    \draw[member,domain=0:1,samples=60,variable=\t] plot ({\t*\L},{Nc(\t)*1.2});
    \draw[dim] (\L+0.3,0) -- (\L+0.3,1.2) node[midway,right] {1};
    \node[below] at (0,-0.05) {0}; \node[below] at (\L/2,-0.05) {$\frac L2$}; \node[below] at (\L,-0.05) {$L$};
    \node[right] at (\L+0.8,0.6) {$N_3(x) = \dfrac{2}{L^2}x\Bigl(x-\dfrac L2\Bigr) = -\dfrac xL+2\dfrac{x^2}{L^2}$};
  \end{scope}
\end{tikzpicture}

\end{fig}
Let the cross-sectional area be of the form $A(x) = A_0(1-mx/L)$, where $m =$ constant. The stiffness matrix for this linearly tapered axial bar is:
\begin{fig}
  % Linearly tapered axial bar with three nodes
\begin{tikzpicture}[scale=1.0,font=\small]
  \draw[thick] (0,-1.0) -- (6,-0.35) -- (6,0.35) -- (0,1.0) -- cycle;
  \draw[dashed] (-0.8,0) -- (7.2,0);
  \draw (3,-0.68) -- (3,0.68);
  \draw[load] (0,0) -- (0.9,0) node[below] {$u_1$};
  \draw[load] (3,0) -- (3.9,0) node[below] {$u_2$};
  \draw[load] (6,0) -- (6.9,0) node[below] {$u_3$};
  \node[above left] at (0,1.0) {$A_0$};
  \node[above] at (3,0.7) {$A_0\bigl(1-m\frac12\bigr)$};
  \node[above right] at (6,0.35) {$A_0(1-m)$};
  \node[below] at (0,-1.1) {$x=0$}; \node[below] at (3,-0.8) {$x=0.5L$}; \node[below] at (6,-0.45) {$x=L$};
\end{tikzpicture}

\end{fig}
\begin{equation}
  \kmat = \int_0^L EA\NN'^{\tr}\NN'\,dx = \frac{EA_0}{6L}\bmat{-3m+14 & 4m-16 & -m+2\\ 4m-16 & -16m+32 & 12m-16\\ -m+2 & 12m-16 & -11m+14}
\end{equation}
It is easy to verify that each column of $\kmat$ is a set of forces in equilibrium (summation of each column is zero). A special case is that of $A(x) = A_0$, i.e.\ $m = 0$, for which
\begin{equation}
  \kmat = \frac{EA_0}{3L}\bmat{7 & -8 & 1\\ -8 & 16 & -8\\ 1 & -8 & 7}
\end{equation}

\subsection{Exact shape functions and exact stiffness matrix for a tapered bar}
The governing differential equation for a tapered bar is given by
\begin{equation}
  \od{}{x}\Bigl(EA(x)\od{u}{x}\Bigr) = 0
\end{equation}
where $A(x) = A_0(1-mx/L)$ is assumed. Integrating the differential equation twice and using the boundary conditions $u(0) = u_1$ and $u(L) = u_2$ yields
\begin{align*}
  &\od{}{x}\Bigl(EA_0(1-mx/L)\od{u}{x}\Bigr) = 0\\
  &\Rightarrow (1-mx/L)\od{u}{x} = c_1 \quad\Rightarrow\quad \od{u}{x} = \frac{c_1}{(1-mx/L)}\\
  &\Rightarrow u(x) = c_1\Bigl(\frac{-L}{m}\Bigr)\ln(1-mx/L) + c_2\\
  &u(0) = u_1 \quad\Rightarrow\quad c_2 = u_1\\
  &u(L) = u_2 \quad\Rightarrow\quad c_1\Bigl(\frac{-L}{m}\Bigr)\ln(1-m) + u_1 = u_2\\
  &\Rightarrow c_1 = (u_2-u_1)\Bigl(\frac{-m}{L}\Bigr)\frac{1}{\ln(1-m)}\\
  &\Rightarrow u(x) = (u_2-u_1)\frac{\ln(1-mx/L)}{\ln(1-m)} + u_1 = N_1(x)u_1 + N_2(x)u_2
\end{align*}
where
\begin{equation}
  N_1(x) = 1 - \frac{\ln(1-mx/L)}{\ln(1-m)},\qquad N_2(x) = \frac{\ln(1-mx/L)}{\ln(1-m)}
\end{equation}
The above shape functions are exact in the sense that they satisfy the differential equation. The ``exact'' stiffness matrix, calculated by $k_{ij} = E\int_0^L A(x)N_i'(x)N_j'(x)\,dx$, is of the form:
\begin{equation}
  \kmat = -\frac{EA_0}{L}\frac{m}{\ln(1-m)}\bmat{1 & -1\\ -1 & 1}
\end{equation}
If a two-noded linear element is used, then $u(x) = (1-\frac xL)u_1 + \frac xL u_2 \Rightarrow \od{u}{x} = \dfrac{u_2-u_1}{L}$
\begin{equation*}
  \od{}{x}\Bigl[EA_0(1-mx/L)\Bigl(\frac{u_2-u_1}{L}\Bigr)\Bigr] \neq 0
\end{equation*}

\begin{example}
Consider a linearly tapered bar fixed at the left end and a force $P$ is applied at the right end. Assuming $m = 2/3$, the exact stiffness equations are:
\begin{equation*}
  -\frac{EA_0}{L}\frac{2/3}{\ln(1/3)}\bmat{1 & -1\\ -1 & 1}\Bmat{u_1\\ u_2} = \Bmat{R\\ P}
\end{equation*}
where $R$ is the reaction at the fixed end. Using the boundary condition $u_1 = 0$, we obtain from the above stiffness equations that
\begin{equation*}
  u_2 = 1.6479\frac{PL}{EA_0}
\end{equation*}
which is, of course, exact.

If a quadratic axial element is used to solve the problem, the stiffness equations can be written as
\begin{equation*}
  \frac{2EA_0}{9L}\bmat{9 & -10 & 1\\ -10 & 16 & -6\\ 1 & -6 & 5}\Bmat{u_1\\ u_2\\ u_3} = \Bmat{R\\ 0\\ P}
\end{equation*}
Knowing that $u_1 = 0$, we can solve for $u_2$ and $u_3$ from the above equations that
\begin{equation*}
  \Bmat{u_2\\ u_3} = \Bmat{0.6136\\ 1.6364}\frac{PL}{EA_0}
\end{equation*}
The error for $u_3$, compared with the exact value, is only 0.7\%.

If a two-noded linear element is used and the cross-sectional area is taken as the mean cross-sectional area, i.e.\ $\overline{A} = \frac23A_0$, then we have
\begin{equation*}
  \frac{2EA_0}{3L}\bmat{1 & -1\\ -1 & 1}\Bmat{u_1\\ u_2} = \Bmat{R\\ P} \Rightarrow u_2 = 1.5\frac{PL}{EA_0}
\end{equation*}
The error is greater than that with the quadratic element. If the bar is divided into three or more elements, the error will become very small.
\end{example}

\subsection{Exact stiffness matrix of a tapered beam}
\begin{minipage}{0.45\linewidth}
Assume a tapered beam with width $b =$ constant and depth $h(x) = h_1\bigl(1+\frac xL\bigr)$. Therefore
\begin{equation*}
  I(x) = \frac{bh^3}{12} = \frac{bh_1^3}{12}\Bigl(1+\frac xL\Bigr)^3
\end{equation*}
\end{minipage}\hfill
\begin{minipage}{0.5\linewidth}\centering
  % Tapered beam: depth h(x) = h1(1 + x/L)
\begin{tikzpicture}[scale=1.0,font=\small]
  \draw[thick] (0,-0.5) -- (6,-1.0) -- (6,1.0) -- (0,0.5) -- cycle;
  \draw[dim] (3,-0.75) -- (3,0.75); \node[right] at (3.1,0.3) {$h(x)$};
  \node[left] at (-0.1,0) {$h_1$}; \node[right] at (6.1,0) {$2h_1$};
  \node[below] at (0,-0.6) {$x=0$}; \node[below] at (6,-1.1) {$x=L$};
\end{tikzpicture}

\end{minipage}

Governing equation (G.E.):
\begin{equation}
  \frac{d^2}{dx^2}\bigl[EI(x)v''\bigr] = 0
\end{equation}
By integrating G.E. twice, we have
\begin{equation}
  \Bigl(1+\frac xL\Bigr)^3v'' = c_1 + c_2x
\end{equation}
By integrating the above equation twice, we obtain
\begin{equation}
  v(x) = \frac{L^2}{2}\Bigl(\frac{c_1}{1+\frac xL}\Bigr) - L^3\Bigl[\frac{1}{2(1+\frac xL)} + \ln\Bigl(1+\frac xL\Bigr)\Bigr]c_2 + c_3x + c_4
\end{equation}
Boundary conditions (BCs):
\begin{equation*}
  v(0) = v_1,\quad v'(0) = \theta_{z1},\quad v(L) = v_2,\quad v'(L) = \theta_{z2}
\end{equation*}
Substituting BCs into $v(x)$, we can obtain $c_i$ ($i = 1, 4$) in terms of $v_1, \theta_{z1}, v_2$, and $\theta_{z2}$. Finally, rearrange the terms to result in the form below:
\begin{equation}
  v(x) = N_1(x)v_1 + N_2(x)\theta_{z1} + N_3(x)v_2 + N_4(x)\theta_{z2}
\end{equation}
The shape functions $N_i(x)$ ($i = 1, 4$) are therefore obtained; these shape functions are exact since they satisfy the governing equation.

Stiffness matrix:
\begin{equation}
  \kmat = \int_0^L EI(x)\NN''^{\tr}\NN''\,dx = \text{exact stiffness matrix.}
\end{equation}

\section{Use of PVD to calculate the equivalent nodal forces or fixed-end forces}
In general, we can determine $\mat{f}_{eq}$ using the principle of virtual displacements. In this way, we let the virtual work done by $\mat{f}_{eq}$ be equal to that done by the intermediate concentrated loads (can be expressed as delta functions) or distributed loads. That is,
\begin{fig}
  % Distributed load q(x) on a beam -> equivalent nodal loads
\begin{tikzpicture}[scale=0.9,font=\small]
  \begin{scope}
    \draw[thick] (0,0) rectangle (5,0.4);
    \draw[thick] (0,0.4) .. controls (0.8,2.2) and (1.6,2.0) .. (2.2,1.6) .. controls (2.8,1.2) and (3.6,1.4) .. (4.2,1.9) .. controls (4.6,2.2) .. (5,0.4);
    \foreach \x in {0.4,0.8,...,4.7}{\draw (\x,0.42) -- (\x,{0.4+1.3});}
    \node[above] at (2.6,1.9) {$q(x)$};
  \end{scope}
  \node at (6.2,0.6) {$\Longrightarrow$};
  \begin{scope}[xshift=7.5cm]
    \draw[thick] (0,0) rectangle (5,0.4);
    \draw[load] (0,0.5) -- (0,1.6); \draw[load] (5,0.5) -- (5,1.6);
    \draw[->,thick] (-0.3,-0.2) arc (-90:90:0.4); \draw[->,thick] (5.3,-0.2) arc (270:90:0.4);
  \end{scope}
\end{tikzpicture}

\end{fig}
\begin{gather}
  \delta\bs{\Delta}^\tr\mat{f}_{eq} = \int_0^L(\NN\delta\bs{\Delta})^\tr q(x)\,dx
  \;\rightarrow\; \delta\bs{\Delta}^\tr\mat{f}_{eq} = \delta\bs{\Delta}^\tr\int_0^L\NN^\tr q(x)\,dx\\
  \rightarrow\quad \mat{f}_{eq} = \int_0^L\NN^\tr q(x)\,dx
\end{gather}
where $\NN = \langle N_1, N_2, \dots\rangle$ are shape functions. In component form,
\begin{equation}
  (\mat{f}_{eq})_i = \int_0^L N_i(x)q(x)\,dx
\end{equation}
The corresponding fixed-end forces are $\mat{f}_{fef} = -\mat{f}_{eq}$.

\begin{example}
\leavevmode
\begin{fig}
  % Uniform load q -> equivalent nodal loads f_eq1..f_eq4
\begin{tikzpicture}[scale=0.9,font=\small]
  \begin{scope}
    \draw[thick] (0,0) rectangle (5,0.4);
    \foreach \x in {0.2,0.7,...,4.9}{\draw[->] (\x,0.45) -- (\x,1.3);}
    \node[above] at (2.5,1.3) {$q$};
  \end{scope}
  \node at (6.2,0.4) {$\Longrightarrow$};
  \begin{scope}[xshift=7.5cm]
    \draw[thick] (0,0) rectangle (5,0.4);
    \draw[load] (0,0.5) -- (0,1.6) node[above] {$f_{eq1}$}; \draw[load] (5,0.5) -- (5,1.6) node[above] {$f_{eq3}$};
    \draw[->,thick] (-0.3,-0.2) arc (-90:90:0.4); \node[below left] at (-0.1,-0.1) {$f_{eq2}$};
    \draw[->,thick] (5.3,-0.2) arc (270:90:0.4); \node[below right] at (5.1,-0.1) {$f_{eq4}$};
  \end{scope}
\end{tikzpicture}

\end{fig}
Using the equations given above, we obtain
\begin{alignat*}{2}
  f_{eq1} &= q\int_0^L N_1(x)\,dx = \frac{qL}{2};&\qquad f_{eq3} &= q\int_0^L N_3(x)\,dx = \frac{qL}{2}\\
  f_{eq2} &= q\int_0^L N_2(x)\,dx = \frac{qL^2}{12};&\qquad f_{eq4} &= q\int_0^L N_4(x)\,dx = -\frac{qL^2}{12}
\end{alignat*}
\end{example}

\begin{example}
\leavevmode
\begin{fig}
  % Point load P at x = a -> equivalent nodal loads
\begin{tikzpicture}[scale=0.9,font=\small]
  \begin{scope}
    \draw[thick] (0,0) rectangle (5,0.4);
    \draw[load] (1.5,0.45) -- (1.5,1.8) node[above] {$P$};
    \draw[dim] (0,-0.5) -- (1.5,-0.5) node[midway,below] {$a$};
    \draw[dim] (1.5,-0.5) -- (5,-0.5) node[midway,below] {$b$};
    \draw (0,-0.7) -- (0,-0.3); \draw (1.5,-0.7) -- (1.5,-0.3); \draw (5,-0.7) -- (5,-0.3);
  \end{scope}
  \node at (6.2,0.4) {$\Longrightarrow$};
  \begin{scope}[xshift=7.5cm]
    \draw[thick] (0,0) rectangle (5,0.4);
    \draw[load] (0,0.5) -- (0,1.6) node[above] {$f_{eq1}$}; \draw[load] (5,0.5) -- (5,1.6) node[above] {$f_{eq3}$};
    \draw[->,thick] (-0.3,-0.2) arc (-90:90:0.4); \node[below left] at (-0.1,-0.1) {$f_{eq2}$};
    \draw[->,thick] (5.3,-0.2) arc (270:90:0.4); \node[below right] at (5.1,-0.1) {$f_{eq4}$};
  \end{scope}
\end{tikzpicture}

\end{fig}
We can use the same equations to calculate the equivalent nodal loads. In this case, $q(x) = P\delta(x-a)$. By using the property $\int_0^L N_i(x)\delta(x-a)\,dx = N_i(a)$, one can obtain that $f_{eqi} = PN_i(a)$ and therefore
\begin{equation*}
  \mat{f}_{eq} = \Bmat{\dfrac{Pb^2}{L^3}(3a+b)\\[6pt] \dfrac{ab^2P}{L^2}\\[6pt] \dfrac{Pa^2}{L^3}(a+3b)\\[6pt] -\dfrac{ba^2P}{L^2}}
  \qquad\text{Check special case } a = b = \frac L2:\quad
  \mat{f}_{eq} = \Bmat{\dfrac P2\\[6pt] \dfrac{PL}{8}\\[6pt] \dfrac P2\\[6pt] -\dfrac{PL}{8}}
\end{equation*}
or
\begin{equation*}
  \dW = P\delta v(x)\big|_{x=a} = P\bigl[N_1(x)\big|_{x=a}\delta v_1 + \dots + N_4(x)\big|_{x=a}\delta\theta_2\bigr] = f_{eqi}\delta\Delta_i \Rightarrow f_{eq_i} = PN_i(a)
\end{equation*}
Similarly, for an applied moment, $M$, at $x = a$, we have
\begin{equation*}
  \dW = M\delta v'(x)\big|_{x=a} = M\bigl[N_1'(x)\big|_{x=a}\delta v_1 + \dots + N_4'(x)\big|_{x=a}\delta\theta_2\bigr] = f_{eqi}\delta\Delta_i \Rightarrow f_{eq_i} = MN_i'(a)
\end{equation*}
(see Fig.~\ref{fig:fef-table}).
\end{example}

\begin{example}[Axial element]
\leavevmode
\begin{fig}
  % Triangular axial load q(x) -> equivalent nodal loads
\begin{tikzpicture}[scale=0.9,font=\small]
  \begin{scope}
    \draw[thick] (0,0) rectangle (5,0.4);
    \draw[thick] (0,0.4) -- (5,1.9) -- (5,0.4);
    \node[above] at (5,1.9) {$q$}; \node[right] at (5.05,1.1) {$q(x)$};
    \foreach \x in {0.3,1.3,...,4.4}{\draw[->] (\x,0.2) -- (\x+0.7,0.2);}
  \end{scope}
  \node at (6.2,0.4) {$\Longrightarrow$};
  \begin{scope}[xshift=7.5cm]
    \draw[thick] (0,0) rectangle (5,0.4);
    \draw[load] (0.1,0.2) -- (0.9,0.2); \node[above] at (0.4,0.5) {$f_{eq1}$};
    \draw[load] (5.05,0.2) -- (5.9,0.2); \node[above] at (5.3,0.5) {$f_{eq2}$};
  \end{scope}
\end{tikzpicture}

\end{fig}
\begin{align*}
  f_{eq1} &= \int_0^L q(x)N_1(x)\,dx = \int_0^L q\frac xL\Bigl(1-\frac xL\Bigr)dx = \frac{qL}{6}\\
  f_{eq2} &= \int_0^L q(x)N_2(x)\,dx = \int_0^L q\frac xL\Bigl(\frac xL\Bigr)dx = \frac{qL}{3}
\end{align*}
Consider the axial bar below discretized into 2 elements. We can use the above formulas to obtain nodal loads as
\begin{equation*}
  \Pvec = \Pvec_c + \Pvec_{eq} = \Bmat{0\\ P} + \Bmat{\frac12qL\bigl(\frac23\bigr)+\frac12qL\bigl(\frac23\bigr)\\[4pt] \frac12qL\bigl(\frac13\bigr)} = \Bmat{\frac23qL\\[4pt] P+\frac16qL}
\end{equation*}
By $\Kmat\bs{\Delta} = \Pvec$, $\bs{\Delta}$ can be obtained.
\begin{fig}
  % Axial bar discretized into two elements with triangular load and a tip load P
\begin{tikzpicture}[scale=1.0,font=\small]
  \fill[left color=black!60,right color=black!15] (-0.3,-0.45) rectangle (0,0.45);
  \draw[thick] (0,-0.2) rectangle (6,0.2);
  \draw (3,-0.2) -- (3,0.2);
  \draw[thick] (0,0.2) -- (3,1.4) -- (6,0.2);
  \node[above] at (3,1.4) {$q$};
  \foreach \x in {0.3,1.2,...,5.4}{\draw[->] (\x,0) -- (\x+0.6,0);}
  \draw[load] (6.1,0) -- (7.0,0) node[right] {$P$};
  \node[above left] at (0.15,0.25) {1}; \node[above] at (3,0.25) {2}; \node[above right] at (5.85,0.25) {3};
  \draw[dim] (0,-0.8) -- (3,-0.8) node[midway,fill=white,inner sep=1pt] {$L$};
  \draw[dim] (3,-0.8) -- (6,-0.8) node[midway,fill=white,inner sep=1pt] {$L$};
  \foreach \x in {0,3,6}{\draw (\x,-0.6) -- (\x,-1.0);}
\end{tikzpicture}

\end{fig}
\end{example}

\chapter{Special Analysis Procedures}
\label{ch:special}

\section{Static condensation}
A simple example: cantilever beam.
\begin{fig}
  % Cantilever beam with tip load P (static condensation example)
\begin{tikzpicture}[scale=1.0,font=\small]
  \pic at (0,0) {wallL};
  \draw[thick] (0,-0.2) rectangle (6,0.2);
  \draw[load] (6,0.3) -- (6,1.6) node[right] {$P$};
  \draw[dim] (0,-0.8) -- (6,-0.8) node[midway,fill=white,inner sep=1pt] {$L$};
  \draw (0,-0.6) -- (0,-1.0); \draw (6,-0.6) -- (6,-1.0);
\end{tikzpicture}

\end{fig}
\begin{equation*}
  \begin{bmatrix}\dfrac{12EI}{L^3} & \dfrac{-6EI}{L^2}\\[8pt] \dfrac{-6EI}{L^2} & \dfrac{4EI}{L}\end{bmatrix}\Bmat{v_2\\ \theta_2} = \Bmat{P\\ 0}
\end{equation*}
2nd Eq.: $\ -\dfrac{6EI}{L^2}v_2 + \dfrac{4EI}{L}\theta_2 = 0 \;\rightarrow\; \theta_2 = \dfrac{3v_2}{2L}$ \hfill (a)

Substitute Eq.~(a) into the 1st Eq.
\begin{gather*}
  \frac{12EI}{L^3}v_2 - \frac{6EI}{L^2}\Bigl(\frac{3v_2}{2L}\Bigr) = P \;\rightarrow\; \frac{3EI}{L^3}v_2 = P\quad\text{(Condensed system!)}\\
  \rightarrow\; v_2 = \frac{PL^3}{3EI} \;\rightarrow\; \text{(a)} \;\rightarrow\; \theta_2 = \frac{3v_2}{2L} = \frac{PL^2}{2EI}
\end{gather*}
The condensed system can be derived systematically as follows. The equations of equilibrium are
\begin{equation}
  \Kmat\bs{\Delta} = \Pvec\qquad\text{or}\qquad
  \bmat{\Kmat_{bb} & \Kmat_{bc}\\ \Kmat_{cb} & \Kmat_{cc}}\Bmat{\bs{\Delta}_b\\ \bs{\Delta}_c} = \Bmat{\Pvec_b\\ \Pvec_c}
\end{equation}
where $\bs{\Delta}_b$ are the degrees of freedom to be eliminated (condensed), $\bs{\Delta}_c$ are the remaining degrees of freedom, and $\Kmat_{cb} = \Kmat_{bc}^\tr$.

We can derive the condensed system using the following two methods.

\topic{Method 1: variable substitution}
Expand the upper partition to obtain
\begin{equation*}
  \Kmat_{bb}\bs{\Delta}_b + \Kmat_{bc}\bs{\Delta}_c = \Pvec_b
\end{equation*}
Solving for $\bs{\Delta}_b$ from the above equation, we have
\begin{equation*}
  \bs{\Delta}_b = \Kmat_{bb}^{-1}(-\Kmat_{bc}\bs{\Delta}_c + \Pvec_b)
\end{equation*}
Then by substituting $\bs{\Delta}_b$ into the lower partition, we obtain the condensed system:
\begin{equation}
  \widehat{\Kmat}_{cc}\bs{\Delta}_c = \widehat{\Pvec}_c
\end{equation}
where
\begin{equation}
  \widehat{\Kmat}_{cc} = \Kmat_{cc} - \Kmat_{cb}\Kmat_{bb}^{-1}\Kmat_{bc},\qquad
  \widehat{\Pvec}_c = \Pvec_c - \Kmat_{cb}\Kmat_{bb}^{-1}\Pvec_b
\end{equation}
Note that $\widehat{\Kmat}_{cc}$ is symmetric. Also, $\widehat{\Pvec}_c = \Pvec_c$ when $\Pvec_b = \mat{0}$.

\topic{Method 2: stationary of total potential energy}
\begin{equation}
  \bs{\Delta} = \Bmat{\bs{\Delta}_b\\ \bs{\Delta}_c} = \bmat{-\Kmat_{bb}^{-1}\Kmat_{bc}\\ \mat{I}}\bs{\Delta}_c + \Bmat{\Kmat_{bb}^{-1}\Pvec_b\\ \mat{0}} = \Tmat\bs{\Delta}_c + \mat{R}
\end{equation}
where
\begin{equation*}
  \Tmat = \bmat{-\Kmat_{bb}^{-1}\Kmat_{bc}\\ \mat{I}};\qquad \mat{R} = \Bmat{\Kmat_{bb}^{-1}\Pvec_b\\ \mat{0}}
\end{equation*}
By making use of the above equation, the strain energy in terms of $\bs{\Delta}_c$ is:
\begin{equation*}
  \begin{aligned}
    U(\bs{\Delta}) &= \tfrac12\bs{\Delta}^\tr\Kmat\bs{\Delta}\\
    &= \tfrac12(\mat{R}^\tr + \bs{\Delta}_c^\tr\Tmat^\tr)\Kmat(\Tmat\bs{\Delta}_c + \mat{R})\\
    &= \tfrac12\bs{\Delta}_c^\tr\Tmat^\tr\Kmat\Tmat\bs{\Delta}_c + \bs{\Delta}_c^\tr\Tmat^\tr\Kmat\mat{R} + \tfrac12\mat{R}^\tr\Kmat\mat{R} \equiv \widehat U(\bs{\Delta}_c)
  \end{aligned}
\end{equation*}
By definition, the load potential is
\begin{equation*}
  V(\bs{\Delta}) = -\bs{\Delta}^\tr\Pvec = -(\Tmat\bs{\Delta}_c + \mat{R})^\tr\Pvec = -\bs{\Delta}_c^\tr\Tmat^\tr\Pvec - \mat{R}^\tr\Pvec \equiv \widehat V(\bs{\Delta}_c)
\end{equation*}
where $\Pvec = \Bmat{\Pvec_b\\ \Pvec_c}$. Knowing the total potential energy $\widehat\Pi(\bs{\Delta}_c) = \widehat U + \widehat V$, the equilibrium equations can be obtained from the stationary condition $\pd{\widehat\Pi}{\bs{\Delta}_c} = \mat{0}$, which results in
\begin{equation}
  \Tmat^\tr\Kmat\Tmat\bs{\Delta}_c = \Tmat^\tr\Pvec - \Tmat^\tr\Kmat\mat{R}
\end{equation}
By making use of the definitions of $\Tmat$ and $\mat{R}$ given above, it can be verified that
\begin{equation*}
  \Tmat^\tr\Kmat\Tmat = \widehat{\Kmat}_{cc}\quad\text{and}\quad \Tmat^\tr\Pvec - \Tmat^\tr\Kmat\mat{R} = \widehat{\Pvec}_c.
\end{equation*}

\section{Application of static condensation: beam element with an end hinge}
\label{sec:end-hinge}
Calculate the stiffness matrix for a beam element with a hinge at one end (\emph{moment release}).
\begin{fig}
  % Beam element (left) and beam element with a hinge at the left end (right)
\begin{tikzpicture}[scale=0.9,font=\small]
  \begin{scope}
    \draw[thinmember] (0,0) -- (5,0);
    \draw[load] (0,0.1) -- (0,1.0) node[above] {$F_1$}; \draw[load] (5,0.1) -- (5,1.0) node[above] {$F_2$};
    \draw[->,thick] (-0.05,-0.4) arc (-100:100:0.4); \node[below right] at (0.1,-0.3) {$M_1$};
    \draw[->,thick] (5.05,-0.4) arc (280:80:0.4); \node[below right] at (5.1,-0.3) {$M_2$};
  \end{scope}
  \draw[->,thick] (6.3,0) -- (7.3,0);
  \begin{scope}[xshift=8.5cm]
    \draw[thinmember] (0,0) -- (5,0); \node[node dot] at (0,0) {};
    \draw[load] (0,0.15) -- (0,1.0) node[above] {$F_1$}; \draw[load] (5,0.1) -- (5,1.0) node[above] {$F_2$};
    \draw[->,thick] (5.05,-0.4) arc (280:80:0.4); \node[below right] at (5.1,-0.3) {$M_2$};
  \end{scope}
\end{tikzpicture}

\end{fig}
The stiffness matrix for the left element is familiar to us. We can use the equation $M_1 = 0$ (the second Eq.) to express $\theta_1$ as a function of the other three degrees of freedom; that is
\begin{equation}
  \frac{EI}{L}\Bigl(\frac6L v_1 + 4\theta_1 - \frac6L v_2 + 2\theta_2\Bigr) = 0 \;\rightarrow\; \theta_1 = -\frac12\theta_2 + \frac{3}{2L}(v_2 - v_1)\quad\text{(constraint Eq.)}
\end{equation}
Substituting this equation into the other three stiffness equations, we will obtain the condensed stiffness matrix as follows.
\begin{equation}
  \kmat = \frac{EI}{L}
  \begin{array}{c}
    \begin{array}{ccc} \;\;v_1\;\; & \;\;v_2\;\; & \;\;\theta_2\;\; \end{array}\\
    \left[\begin{array}{ccc}
      \dfrac{3}{L^2} & -\dfrac{3}{L^2} & \dfrac3L\\[8pt]
      -\dfrac{3}{L^2} & \dfrac{3}{L^2} & -\dfrac3L\\[8pt]
      \dfrac3L & -\dfrac3L & 3
    \end{array}\right]
  \end{array}
\end{equation}
3rd eq.\ $\rightarrow$
\begin{equation*}
  M_2 = \frac{3EI}{L}\Bigl(\frac{v_1}{L} - \frac{v_2}{L} + \theta_2\Bigr) = \frac{3EI}{L}(\theta_2 - R),\quad
  \text{where } R = \frac{v_2-v_1}{L}\ \text{(positive if counterclockwise)}
\end{equation*}

\topic{(Modified) slope-deflection equation:}
\begin{gather*}
  M_{12} = \frac{2EI}{L}(2\theta_1 + \theta_2 - 3R);\qquad M_{21} = \frac{2EI}{L}(2\theta_2 + \theta_1 - 3R)\quad\text{with } R = \frac{\Delta}{L} = \frac{v_2-v_1}{L}\\
  \text{Let } M_{12} = 0 \;\rightarrow\; \theta_1 = \frac{-\theta_2}{2} + \frac{3R}{2} \;\rightarrow\; M_{21} = \frac{2EI}{L}(1.5\theta_2 - 1.5R) = \frac{3EI}{L}(\theta_2 - R)
\end{gather*}
Note that the degrees of freedom are ordered in the sequence of $v_1$, $v_2$, and $\theta_2$. For the convenience of computer implementation, it is usual to expand the matrix into
\begin{equation}
  \kmat = \frac{EI}{L}
  \begin{bmatrix}
    \dfrac{3}{L^2} & 0 & -\dfrac{3}{L^2} & \dfrac3L\\[8pt]
    0 & 0 & 0 & 0\\[4pt]
    -\dfrac{3}{L^2} & 0 & \dfrac{3}{L^2} & -\dfrac3L\\[8pt]
    \dfrac3L & 0 & -\dfrac3L & 3
  \end{bmatrix}
\end{equation}
As for the element with a hinge at the right end, the constraint equation (from $M_2 = 0$) and stiffness matrix can be derived similarly, which are given by
\begin{equation}
  \theta_2 = -\frac12\theta_1 + \frac{3}{2L}(v_2 - v_1);\qquad
  \kmat = \frac{EI}{L}
  \begin{array}{c}
    \begin{array}{ccc} \;\;v_1\;\; & \;\;\theta_1\;\; & \;\;v_2\;\; \end{array}\\
    \left[\begin{array}{ccc}
      \dfrac{3}{L^2} & \dfrac3L & -\dfrac{3}{L^2}\\[8pt]
      \dfrac3L & 3 & -\dfrac3L\\[8pt]
      -\dfrac{3}{L^2} & -\dfrac3L & \dfrac{3}{L^2}
    \end{array}\right]
  \end{array}
\end{equation}
$\rightarrow$ 2nd Eq.
\begin{equation*}
  M_2 = \frac{3EI}{L}\Bigl(\frac{v_1}{L} - \frac{v_2}{L} + \theta_1\Bigr) = \frac{3EI}{L}(\theta_1 - R)\quad\text{where } R = \frac{v_2-v_1}{L}
\end{equation*}
Note that the degrees of freedom are ordered in the sequence of $v_1$, $\theta_1$ and $v_2$. It can also be expanded into a 4 by 4 matrix as we did for moment release at the left end.

\begin{example}
Three approaches can be used to solve the problem shown below.
\begin{fig}
  % Fixed-fixed beam with an internal hinge at B and load P; DOFs v, theta_l, theta_r
\begin{tikzpicture}[scale=0.9,font=\small]
  \begin{scope}
    \pic at (0,0) {wallL}; \pic at (7.5,0) {wallR};
    \draw[member] (0,0) -- (7.5,0); \node[node dot] at (2.5,0) {};
    \node[above right] at (0,0.05) {$A$}; \node[above right] at (2.5,0.05) {$B$}; \node[above left] at (7.5,0.05) {$C$};
    \draw[load] (2.5,0.2) -- (2.5,1.4) node[above] {$P$};
    \draw[dim] (0,-0.7) -- (2.5,-0.7) node[midway,fill=white,inner sep=1pt] {$L$};
    \draw[dim] (2.5,-0.7) -- (7.5,-0.7) node[midway,fill=white,inner sep=1pt] {$2L$};
  \end{scope}
  \begin{scope}[xshift=10cm]
    \pic at (0,0) {wallL}; \pic at (5.5,0) {wallR};
    \draw[member] (0,0) -- (5.5,0); \node[node dot] at (1.8,0) {};
    \draw[load] (1.8,0.2) -- (1.8,1.4) node[above] {$v$};
    \draw[->,thick] (1.55,-0.15) arc (-160:-30:0.45); \node[below] at (1.3,-0.5) {$\theta_l$};
    \draw[->,thick] (2.5,0.15) arc (30:150:0.45); \node[below] at (2.4,-0.5) {$\theta_r$};
  \end{scope}
\end{tikzpicture}

\end{fig}
\textbf{Method 1:} Conventional elements are used for $AB$ and $BC$.

The DOFs include $v$, $\theta_l$, and $\theta_r$. The stiffness equations can be obtained by assembling two conventional beam elements, which are
\begin{equation*}
  \begin{bmatrix}
    \dfrac{13.5EI}{L^3} & -\dfrac{6EI}{L^2} & \dfrac{1.5EI}{L^2}\\[8pt]
    & \dfrac{4EI}{L} & 0\\[8pt]
    \textit{Sym.} & & \dfrac{2EI}{L}
  \end{bmatrix}\Bmat{v\\ \theta_l\\ \theta_r} = \Bmat{P\\ 0\\ 0}
\end{equation*}
The solution is: $v = \dfrac{PL^3}{3.375EI}$, $\theta_l = \dfrac49\dfrac{PL^2}{EI}$, and $\theta_r = -\dfrac29\dfrac{PL^2}{EI}$.

The advantages of this method are that all the variables can be obtained simultaneously and only conventional elements are used. The method is implemented in FRAME20 by making use of \emph{the double node technique} as discussed in Ch.~\ref{ch:frame21}.

\textbf{Method 2:} Pinned-end (condensed) elements are used for $AB$ and $BC$.
\begin{equation*}
  \Bigl(\frac{3EI}{L^3} + \frac{3EI}{(2L)^3}\Bigr)v = P
\end{equation*}
This is like two cantilevers arranged in parallel. From the above equation, we obtain $v = \dfrac{PL^3}{3.375EI}$. We can use the constraint equations to obtain $\theta_l$ and $\theta_r$ as follows.
\begin{align*}
  \theta_l = \theta_2 &= -\frac12\theta_1 + \frac{3}{2L}(v_2 - v_1) = \frac{3}{2L}(v - 0) = \frac49\frac{PL^2}{EI}\\
  \theta_r = \theta_1 &= -\frac12\theta_2 + \frac{3}{2(2L)}(v_2 - v_1) = \frac{3}{2L}(0 - v) = -\frac29\frac{PL^2}{EI}
\end{align*}

\textbf{Method 3:} Pinned-end element is used for $AB$ and conventional element is used for $BC$.
\begin{equation*}
  \begin{bmatrix}\dfrac{4.5EI}{L^3} & \dfrac{1.5EI}{L^2}\\[8pt] \dfrac{1.5EI}{L^2} & \dfrac{2EI}{L}\end{bmatrix}\Bmat{v\\ \theta_r} = \Bmat{P\\ 0}
\end{equation*}
Solving the above system, we obtain $v = \dfrac{PL^3}{3.375EI}$ and $\theta_r = -\dfrac29\dfrac{PL^2}{EI}$.

Similarly, if a pinned-end element is used for $BC$ and a conventional element for $AB$, we will obtain $v = \dfrac{PL^3}{3.375EI}$ and $\theta_l = \dfrac49\dfrac{PL^2}{EI}$.
\end{example}

\begin{example}[Stiffness matrix for a beam with an internal hinge]
The stiffness matrix for a beam with an internal hinge has been given in the textbook by McGuire et al.\ (Example~4.14 on page 86), in which the condensation method is employed. As you can see, the procedure involves lots of computations.

We will use a very simple method to obtain the same stiffness matrix. First, we support the structure so that it becomes \emph{statically determinate}. One possibility is shown below; you can support it in other ways. Then we apply \emph{the unit load method} to find the flexibility coefficient corresponding to the free DOF.
\begin{fig}
  % Statically determinate element with an internal hinge and its moment diagram
\begin{tikzpicture}[scale=0.9,font=\small]
  \begin{scope}
    \pic at (0,0) {wallL}; \draw[member] (0,0) -- (6,0); \node[node dot] at (2.2,0) {};
    \pic at (6,0) {roller}; \node[node dot] at (6,0) {};
    \draw[->,thick] (6.3,0.6) arc (60:200:0.5); \node[right] at (6.4,0.6) {1};
    \draw[dim] (0,-0.9) -- (2.2,-0.9) node[midway,fill=white,inner sep=1pt] {$a$};
    \draw[dim] (2.2,-0.9) -- (6,-0.9) node[midway,fill=white,inner sep=1pt] {$b$};
    \node[below] at (3,-1.4) {(a) statically determinate element};
  \end{scope}
  \begin{scope}[xshift=9cm]
    \draw (0,0) -- (6,0);
    \draw[thick] (0,0) -- (0,-0.9) -- (6,1.5) -- (6,0);
    \node[left] at (0,-0.6) {$-\dfrac ab$}; \node[right] at (6,1.5) {1};
    \node[below] at (3,-1.4) {(b) moment diagram};
  \end{scope}
\end{tikzpicture}

\end{fig}
\begin{equation*}
  1\cdot\theta_2 = \int\frac{Mm}{EI}\,dx = \int_0^{a+b}\frac{1}{EIb^2}(x-a)^2\,dx = \frac{1}{3EI}\frac{1}{b^2}(a^3 + b^3)
\end{equation*}
Therefore, the flexibility coefficient $d = \theta_2 = \dfrac{1}{3EIb^2}(a^3+b^3)$ and
\begin{equation*}
  k_{44} = d^{-1} = \frac{3EI}{a^3+b^3}b^2
\end{equation*}
From $k_{44}$, the stiffness matrix without support can be obtained by $\kmat = \Tmat^\tr k_{44}\Tmat$.
\begin{fig}
  % Free bodies of the two segments of the hinged beam element
\begin{tikzpicture}[scale=0.9,font=\small]
  \begin{scope}
    \draw[member] (0,0) -- (3,0); \node[node dot] at (3,0) {};
    \draw[->,thick] (-0.5,0.7) arc (120:-60:0.5); \node[above left] at (-0.4,0.6) {$M_1$};
    \draw[load] (0,-1.1) -- (0,-0.1) node[pos=0,below] {$F_1$};
    \draw[load] (3,-0.1) -- (3,-1.1) node[below] {$F_1$};
    \draw[dim] (0,-0.5) -- (3,-0.5) node[midway,fill=white,inner sep=1pt] {$a$};
  \end{scope}
  \begin{scope}[xshift=4.5cm]
    \draw[member] (0,0) -- (5,0); \node[node dot] at (0,0) {};
    \draw[->,thick] (5.5,0.7) arc (60:240:0.5); \node[above right] at (5.4,0.6) {$M_2$};
    \draw[load] (0,-0.1) -- (0,-1.1) node[below] {$F_2$};
    \draw[load] (5,-1.1) -- (5,-0.1) node[pos=0,below] {$F_2$};
    \draw[dim] (0,-0.5) -- (5,-0.5) node[midway,fill=white,inner sep=1pt] {$b$};
  \end{scope}
\end{tikzpicture}

\end{fig}
\begin{gather*}
  M_2 = -F_2\cdot b \Rightarrow F_2 = \frac{-M_2}{b}\\
  F_1 = -F_2 = \frac{M_2}{b}\\
  M_1 = F_1\cdot a = -F_2\cdot a = \frac{M_2\cdot a}{b}
\end{gather*}
\begin{equation*}
  \F = \Tmat^\tr\F' \Rightarrow \Bmat{F_1\\ M_1\\ F_2\\ M_2} = \Bigl\{\frac1b\quad\frac ab\quad\frac{-1}{b}\quad 1\Bigr\}^\tr\cdot M_2
\end{equation*}
\begin{equation*}
  \Tmat_{1\times4} = \Bigl\{\frac1b\quad\frac ab\quad\frac{-1}{b}\quad 1\Bigr\}\qquad
  \kmat = \Tmat^\tr k_{44}\Tmat = \frac{3EI}{a^3+b^3}
  \bmat{1 & a & -1 & b\\ & a^2 & -a & ab\\ & & 1 & -b\\ \textit{Sym.} & & & b^2}
\end{equation*}
\end{example}

\section{Modification of fixed-end forces considering moment release}
\begin{fig}
  % Beam element with moment release at the left end (M1 = 0)
\begin{tikzpicture}[scale=0.9,font=\small]
  \draw[thinmember] (0,0) -- (6,0); \node[node dot] at (0,0) {};
  \draw[load] (0,0.15) -- (0,1.0); \draw[load] (6,0.1) -- (6,1.0);
  \draw[->,thick] (6.05,-0.4) arc (280:80:0.4);
  \node[below] at (0,-0.2) {$M_1 = 0$};
\end{tikzpicture}

\end{fig}
\begin{gather*}
  M_1 = 0 \Rightarrow \frac{6EI}{L^2}v_1 + \frac{4EI}{L}\theta_1 - \frac{6EI}{L^2}v_2 + \frac{2EI}{L}\theta_2 + f^2_{fef} = 0\\
  \Rightarrow \theta_1 = -\frac12\theta_2 + \frac{3}{2L}(v_2 - v_1) - \frac{f^2_{fef}L}{4EI}\quad\Rightarrow\quad\text{Eqs.\ 1, 3, 4}\;\rightarrow\;\text{Modified Eqs.}
\end{gather*}
Equivalently, (row 2 multiplied by $-\frac{3}{2L}$, $\frac{3}{2L}$, $-\frac12$ is added to rows 1, 3, 4)
{\small\begin{equation*}
  \begin{bmatrix}
    \dfrac{12EI}{L^3} & \dfrac{6EI}{L^2} & -\dfrac{12EI}{L^3} & \dfrac{6EI}{L^2}\\[8pt]
    & \dfrac{4EI}{L} & -\dfrac{6EI}{L^2} & \dfrac{2EI}{L}\\[8pt]
    & & \dfrac{12EI}{L^3} & -\dfrac{6EI}{L^2}\\[8pt]
    \textit{sym.} & & & \dfrac{4EI}{L}
  \end{bmatrix}\Bmat{v_1\\ \theta_1\\ v_2\\ \theta_2} + \Bmat{f^1_{fef}\\ f^2_{fef}\\ f^3_{fef}\\ f^4_{fef}} = \Bmat{F_1\\ M_1\\ F_2\\ M_2}
  \;\Rightarrow\; \mat{f}_{fef}' = \Bmat{f^1_{fef} - \dfrac{3}{2L}f^2_{fef}\\[6pt] f^2_{fef}\\[4pt] f^3_{fef} + \dfrac{3}{2L}f^2_{fef}\\[6pt] \boxed{f^4_{fef} - \dfrac12 f^2_{fef}}}
\end{equation*}
}
\textless EX\textgreater
\begin{fig}
  % Fixed-end forces for propped cantilevers (uniform load and centre point load)
\begin{tikzpicture}[scale=0.8,font=\small]
  % pin left, fixed right, UDL
  \begin{scope}
    \draw[member] (0,0) -- (7,0); \pic at (0,0) {pin}; \pic at (7,0) {wallR};
    \foreach \x in {0.3,0.8,...,6.8}{\draw[->] (\x,0.9) -- (\x,0.1);} \draw (0.3,0.9) -- (6.8,0.9);
    \node[above] at (3.5,0.95) {$w$}; \node[below] at (3.5,-0.1) {$l$};
    \draw[load] (0,-1.5) -- (0,-0.6) node[pos=0,below] {$\frac38wl$};
    \draw[load] (7,-1.5) -- (7,-0.6) node[pos=0,below] {$\frac58wl$};
    \draw[->,thick] (7.4,-0.9) arc (-70:110:0.4); \node[left] at (6.9,-0.9) {$\frac18wl^2$};
  \end{scope}
  % fixed left, pin right, UDL
  \begin{scope}[yshift=-4.2cm]
    \draw[member] (0,0) -- (7,0); \pic at (0,0) {wallL}; \pic at (7,0) {pin};
    \foreach \x in {0.3,0.8,...,6.8}{\draw[->] (\x,0.9) -- (\x,0.1);} \draw (0.3,0.9) -- (6.8,0.9);
    \node[above] at (3.5,0.95) {$w$}; \node[below] at (3.5,-0.1) {$l$};
    \draw[load] (0,-1.5) -- (0,-0.6) node[pos=0,below] {$\frac58wl$};
    \draw[load] (7,-1.5) -- (7,-0.6) node[pos=0,below] {$\frac38wl$};
    \draw[->,thick] (-0.4,-0.9) arc (250:70:0.4); \node[right] at (0.1,-0.9) {$\frac18wl^2$};
  \end{scope}
  % fixed left, pin right, centre point load
  \begin{scope}[yshift=-8.4cm]
    \draw[member] (0,0) -- (7,0); \pic at (0,0) {wallL}; \pic at (7,0) {pin};
    \draw[load] (3.5,1.2) -- (3.5,0.1) node[pos=0,above] {$P$}; \node[below] at (3.5,-0.1) {$l$};
    \draw[load] (0,-1.5) -- (0,-0.6) node[pos=0,below] {$\frac{11}{16}P$};
    \draw[load] (7,-1.5) -- (7,-0.6) node[pos=0,below] {$\frac{5}{16}P$};
    \draw[->,thick] (-0.4,-0.9) arc (250:70:0.4); \node[right] at (0.1,-0.9) {$\frac{3}{16}Pl$};
  \end{scope}
\end{tikzpicture}

\end{fig}
\noindent$***$ You shall always remember that the fixed end moment will be increased once the other node becomes a hinge, which is very obvious. Also, the amount of increase will be half of the fixed end moment of the original case. The end shear forces are just calculated from the equilibrium equations.

\section{Substructures}
Example~13.2 in the textbook is a good example to understand this method.

\section{Systems with constraints}
The equations of equilibrium for a structural system:
\begin{equation}
  \Kmat\bs{\Delta} = \Pvec\qquad\text{or}\qquad \bmat{\Kmat_{aa} & \Kmat_{ab}\\ \Kmat_{ba} & \Kmat_{bb}}\Bmat{\bs{\Delta}_a\\ \bs{\Delta}_b} = \Bmat{\Pvec_a\\ \Pvec_b}
\end{equation}
where $\Kmat_{ab} = \Kmat_{ab}^\tr$. The system is subjected to $r$ constraint equations of the form
\begin{equation}
  \mat{A}_{r\times(n-r)}(\bs{\Delta}_a)_{(n-r)\times1} + \mat{B}_{r\times r}(\bs{\Delta}_b)_{r\times1} = \mat{H}_{r\times1}
\end{equation}
where $\mat{A}$ and $\mat{B}$ are given matrices and $\mat{H}$ is a given vector.

Rewrite the constraint equations to express $\bs{\Delta}_b$ in terms of $\bs{\Delta}_a$ and $\mat{H}$ as follows.
\begin{equation}
  \bs{\Delta}_b = \mat{B}^{-1}(-\mat{A}\bs{\Delta}_a + \mat{H}) = \mat{G}_1\bs{\Delta}_a + \mat{G}_2
\end{equation}
where $\mat{G}_1 = -\mat{B}^{-1}\mat{A}$ and $\mat{G}_2 = \mat{B}^{-1}\mat{H}$.

\topic{Method 1: Transformation method (DOFs are reduced)}
\begin{equation}
  \bs{\Delta} = \Bmat{\bs{\Delta}_a\\ \bs{\Delta}_b} = \bmat{\mat{I}\\ \mat{G}_1}\bs{\Delta}_a + \Bmat{\mat{0}\\ \mat{G}_2} = \Tmat\bs{\Delta}_a + \mat{R}
\end{equation}
The total potential energy
\begin{equation*}
  \Pi = \tfrac12\bs{\Delta}^\tr\Kmat\bs{\Delta} - \bs{\Delta}^\tr\Pvec
\end{equation*}
Rewrite it considering the constraints so that it becomes a function of only $\bs{\Delta}_a$.
{\small\begin{equation*}
  \begin{aligned}
    \widehat\Pi(\bs{\Delta}_a) &= \Pi(\bs{\Delta})\big|_{\bs{\Delta} = \Tmat\bs{\Delta}_a + \mat{R}}\\
    &= \tfrac12(\bs{\Delta}_a^\tr\Tmat^\tr + \mat{R}^\tr)\Kmat(\Tmat\bs{\Delta}_a + \mat{R}) - (\bs{\Delta}_a^\tr\Tmat^\tr + \mat{R}^\tr)\Pvec\\
    &= \tfrac12\bs{\Delta}_a^\tr\Tmat^\tr\Kmat\Tmat\bs{\Delta}_a + \tfrac12\bs{\Delta}_a^\tr\Tmat^\tr\Kmat\mat{R} + \tfrac12\mat{R}^\tr\Kmat\Tmat\bs{\Delta}_a + \tfrac12\mat{R}^\tr\Kmat\mat{R} - \bs{\Delta}_a^\tr\Tmat^\tr\Pvec - \mat{R}^\tr\Pvec\\
    &= \tfrac12\bs{\Delta}_a^\tr\Tmat^\tr\Kmat\Tmat\bs{\Delta}_a + \bs{\Delta}_a^\tr\Tmat^\tr\Kmat\mat{R} - \bs{\Delta}_a^\tr\Tmat^\tr\Pvec + \tfrac12\mat{R}^\tr\Kmat\mat{R} - \mat{R}^\tr\Pvec
  \end{aligned}
\end{equation*}}
Equations of equilibrium:
\begin{equation}
  \pd{\widehat\Pi}{\bs{\Delta}_a} = \mat{0}\quad\rightarrow\quad \boxed{\Tmat^\tr\Kmat\Tmat\bs{\Delta}_a = \Tmat^\tr\Pvec - \Tmat^\tr\Kmat\mat{R}}
\end{equation}
As can be seen from the above equation, the stiffness matrix corresponding to $\bs{\Delta}_a$ can be obtained using the congruent transformation but the load vector cannot be obtained using the contragredient principle. However, when the constraint equations are homogeneous, i.e.\ $\mat{H} = \mat{0}$, then $\mat{R} = \mat{0}$ and the contragredient principle works.

By substituting $\Tmat = \bmat{\mat{I}\\ \mat{G}_1}$ and $\mat{R} = \Bmat{\mat{0}\\ \mat{G}_2}$ into the above equation, we obtain
\begin{equation}
  \widehat{\Kmat}_{aa}\bs{\Delta}_a = \widehat{\Pvec}_a,\quad\text{where}\quad
  \begin{aligned}
    \widehat{\Kmat}_{aa} &= \Kmat_{aa} + \Kmat_{ab}\mat{G}_1 + \mat{G}_1^\tr\Kmat_{ba} + \mat{G}_1^\tr\Kmat_{bb}\mat{G}_1\\
    \widehat{\Pvec}_a &= \Pvec_a + \mat{G}_1\Pvec_b - \Kmat_{ab}\mat{G}_2 - \mat{G}_1\Kmat_{bb}\mat{G}_2
  \end{aligned}
\end{equation}
Note that $\widehat{\Kmat}_{aa}$ is symmetric.

\emph{Special case:} If the constraints are homogeneous, i.e.\ $\mat{H} = \mat{0}$, then $\mat{G}_2 = \mat{0}$, and
\begin{equation*}
  \widehat{\Pvec}_a = \Pvec_a + \mat{G}_1\Pvec_b
\end{equation*}

\topic{Method 2: Lagrange multiplier method (DOFs are increased)}
The modified total potential energy:
\begin{equation}
  \widehat\Pi(\bs{\Delta}_a, \bs{\Delta}_b, \bs{\lambda}) = \Pi(\bs{\Delta}_a, \bs{\Delta}_b) + \bs{\lambda}^\tr(\mat{A}\bs{\Delta}_a + \mat{B}\bs{\Delta}_b - \mat{H})
  = \tfrac12\bs{\Delta}^\tr\Kmat\bs{\Delta} - \bs{\Delta}^\tr\Pvec + \bs{\lambda}^\tr\bs{\delta}
\end{equation}
where $\bs{\lambda}_{r\times1}$ are the Lagrange multipliers (can be considered as generalized forces that are work conjugate to $-\bs{\delta}$) and $\bs{\delta} = \mat{A}\bs{\Delta}_a + \mat{B}\bs{\Delta}_b - \mat{H}$ (can be considered as generalized displacements).

The equations of equilibrium can be obtained by
\begin{equation*}
  \pd{\widehat\Pi}{\bs{\Delta}_a} = \mat{0},\qquad \pd{\widehat\Pi}{\bs{\Delta}_b} = \mat{0},\qquad \pd{\widehat\Pi}{\bs{\lambda}} = \mat{0}
\end{equation*}
The resulting expressions are:
\begin{equation}
  \boxed{\bmat{\Kmat_{aa} & \Kmat_{ab} & \mat{A}^\tr\\ \Kmat_{ba} & \Kmat_{bb} & \mat{B}^\tr\\ \mat{A} & \mat{B} & \mat{0}}\Bmat{\bs{\Delta}_a\\ \bs{\Delta}_b\\ \bs{\lambda}} = \Bmat{\Pvec_a\\ \Pvec_b\\ \mat{H}}}
\end{equation}

\topic{Method 3: Penalty function (substituting member) method (DOFs remain unchanged)}
\begin{fig}
  % Penalty (substituting member) method for a truss with an inclined support
\begin{tikzpicture}[scale=0.8,font=\small]
  \begin{scope}
    \coordinate (c) at (0,0); \coordinate (b) at (6,0); \coordinate (a) at (3.5,2.2);
    \draw[member] (c) -- (a) -- (b) -- (c);
    \pic at (c) {pin}; \pic[rotate=30] at (b) {roller};
    \foreach \p in {a,b,c}{\node[node dot] at (\p) {};}
  \end{scope}
  \node at (7.5,1.0) {$\Longrightarrow$};
  \begin{scope}[xshift=9cm]
    \coordinate (c) at (0,0); \coordinate (b) at (6,0); \coordinate (a) at (3.5,2.2);
    \draw[member] (c) -- (a) -- (b) -- (c);
    \pic at (c) {pin};
    \foreach \p in {a,b,c}{\node[node dot] at (\p) {};}
    \node[above left] at ($(c)!0.5!(a)$) {$k_3$}; \node[above right] at ($(a)!0.5!(b)$) {$k_1$}; \node[below] at (3,-0.1) {$k_2$};
    % substituting member perpendicular to the incline
    \coordinate (s) at ($(b)+(-60:1.6)$);
    \draw[line width=3pt,black] (b) -- (s); \pic[rotate=-60] at (s) {pin};
    \node[right] at ($(b)+(0.9,-0.6)$) {$\alpha \gg k_1\sim k_3$};
  \end{scope}
\end{tikzpicture}

\end{fig}
Let $\bs{\delta}(\bs{\Delta}_a, \bs{\Delta}_b) = \mat{A}\bs{\Delta}_a + \mat{B}\bs{\Delta}_b - \mat{H}$ and the modified total potential energy is defined as
\begin{equation}
  \widehat\Pi(\bs{\Delta}_a, \bs{\Delta}_b) = \Pi + U_s = \Pi + \tfrac12\bs{\delta}^\tr\bs{\alpha}\bs{\delta}
\end{equation}
where $U_s$ represents the strain energy of the substituting members and the penalty coefficient (the stiffness of the substituting members) matrix is given by
\begin{equation*}
  \bs{\alpha} = \bmat{\alpha_1 & \cdots & 0\\ \cdots & \cdots & \cdots\\ 0 & \cdots & \alpha_r}
\end{equation*}
Therefore,
\begin{equation*}
  \tfrac12\bs{\delta}^\tr\bs{\alpha}\bs{\delta} = \tfrac12\delta_1^2\alpha_1 + \tfrac12\delta_2^2\alpha_2 + \dots + \tfrac12\delta_r^2\alpha_r.
\end{equation*}
In general, the penalty coefficients shall be large enough to let the constraints be satisfied almost.

Rewrite the constraint equations as
\begin{equation*}
  \mat{C}\bs{\Delta} = \mat{H}
\end{equation*}
where $\mat{C} = \langle\mat{A}, \mat{B}\rangle$ and $\bs{\Delta} = \langle\bs{\Delta}_a, \bs{\Delta}_b\rangle^\tr$. Then the generalized displacements can be rewritten as
\begin{equation*}
  \bs{\delta} = \mat{C}\bs{\Delta} - \mat{H}
\end{equation*}
The modified total potential then becomes
\begin{equation*}
  \begin{aligned}
    \widehat\Pi(\bs{\Delta}_a, \bs{\Delta}_b) &= \Pi + \tfrac12\bs{\delta}^\tr\bs{\alpha}\bs{\delta}\\
    &= \tfrac12\bs{\Delta}^\tr\Kmat\bs{\Delta} - \bs{\Delta}^\tr\Pvec + \tfrac12(\mat{C}\bs{\Delta} - \mat{H})^\tr\bs{\alpha}(\mat{C}\bs{\Delta} - \mat{H})\\
    &= \tfrac12\bs{\Delta}^\tr\Kmat\bs{\Delta} - \bs{\Delta}^\tr\Pvec + \tfrac12\bs{\Delta}^\tr\mat{C}^\tr\bs{\alpha}\mat{C}\bs{\Delta} - \tfrac12\bs{\Delta}^\tr\mat{C}^\tr\bs{\alpha}\mat{H} - \tfrac12\mat{H}^\tr\bs{\alpha}\mat{C}\bs{\Delta} + \tfrac12\mat{H}^\tr\bs{\alpha}\mat{H}\\
    &= \tfrac12\bs{\Delta}^\tr\Kmat\bs{\Delta} - \bs{\Delta}^\tr\Pvec + \tfrac12\bs{\Delta}^\tr\mat{C}^\tr\bs{\alpha}\mat{C}\bs{\Delta} - \bs{\Delta}^\tr\mat{C}^\tr\bs{\alpha}\mat{H} + \tfrac12\mat{H}^\tr\bs{\alpha}\mat{H}
  \end{aligned}
\end{equation*}
Equilibrium equations:
\begin{equation}
  \pd{\widehat\Pi}{\bs{\Delta}} = \mat{0}\;\rightarrow\;
  \Kmat\bs{\Delta} - \Pvec + \mat{C}^\tr\bs{\alpha}\mat{C}\bs{\Delta} - \mat{C}^\tr\bs{\alpha}\mat{H} = \mat{0}\quad\text{or}\quad
  \boxed{(\Kmat + \mat{C}^\tr\bs{\alpha}\mat{C})\bs{\Delta} = \Pvec + \mat{C}^\tr\bs{\alpha}\mat{H}}
\end{equation}

\begin{example}[Truss with an inclined support; McGuire et al.\ (2000), page 390]
\leavevmode
\begin{fig}
  % Truss with an inclined support (after McGuire et al., 2000, p. 390)
\begin{tikzpicture}[scale=0.8,font=\small]
  \coordinate (c) at (0,0); \coordinate (b) at (8,0); \coordinate (a) at (5,2.6);
  \draw[thinmember] (c) -- (a) -- (b) -- (c);
  \pic at (c) {pin}; \node[node dot] at (b) {};
  \node[left] at (c) {c}; \node[above left] at (a) {a}; \node[above] at ($(b)+(0.15,0.1)$) {b};
  \draw[load] (a) -- ++(0,1.4) node[right] {$P_{ya} = 321.4$ kN};
  \draw[load] (a) -- ++(1.6,0) node[right] {$P_{xa} = 383$ kN};
  % inclined support at b (30 deg)
  \draw[thick] ($(b)+(-0.9,-0.52)$) -- ($(b)+(0.9,0.52)$);
  \foreach \t in {-0.8,-0.5,...,0.8}{\draw ($(b)+(\t,{\t*0.577})+(0,-0.05)$) -- ($(b)+(\t,{\t*0.577})+(-0.3,-0.45)$);}
  \draw[->,thick] (b) -- ++(30:1.8) node[right] {$x''$};
  \draw[->,thick] (b) -- ++(120:1.4) node[above] {$y''$};
  \draw[dashed] (b) -- ++(2.4,0);
  \draw[->] ($(b)+(2.0,0)$) arc (0:30:2.0); \node[right] at ($(b)+(2.05,0.55)$) {$30^\circ$};
  \draw[->,thick] ($(b)+(2.4,0)$) -- ++(1.2,0) node[below] {$x$};
  \draw[->,thick] (c) ++(0,0.3) -- ++(0,1.6) node[above] {$y$};
\end{tikzpicture}

\end{fig}
\topic{Constraint equations}
\begin{equation*}
  v_b'' = -0.5u_b + \sqrt3/2\,v_b = 0\qquad\text{or}\qquad \mat{C}\bs{\Delta} = 0,
\end{equation*}
where $\mat{C} = \langle 0, 0, -0.5, \sqrt3/2\rangle$, $\bs{\Delta}^\tr = \langle u_a, v_a, u_b, v_b\rangle$.

\emph{In the following calculations, all displacements are in mm and all forces are in kN.}

\subsection*{Lagrange multiplier method}
Total potential:
\begin{equation*}
  \widehat\Pi = \tfrac12\bs{\Delta}^\tr\Kmat\bs{\Delta} - u_aP_{xa} - v_aP_{ya} + \lambda v_b''
  = \tfrac12\bs{\Delta}^\tr\Kmat\bs{\Delta} - u_aP_{xa} - v_aP_{ya} + \lambda(-0.5u_b + \sqrt3/2\,v_b)
\end{equation*}
Note that $(-\lambda)$ and $v_b''$ are a work conjugate pair.

$\pd{\widehat\Pi}{\bs{\Delta}} = \mat{0}$ and $\pd{\widehat\Pi}{\lambda} = 0$ result in the system of equations:
{\small\begin{equation*}
  \bmat{634.8 & -191.2 & -353.6 & 353.6 & 0\\ -191.2 & 447.3 & 353.6 & -353.6 & 0\\ -353.6 & 353.6 & 683.0 & -353.6 & -.5\\ 353.6 & -353.6 & -353.6 & 353.6 & .5\sqrt3\\ 0 & 0 & -.5 & .5\sqrt3 & 0}
  \Bmat{u_a\\ v_a\\ u_b\\ v_b\\ \lambda} = \Bmat{383\\ 321.4\\ 0\\ 0\\ 0},\qquad
  \text{Solution is: }\Bmat{.92828\\ 1.1426\\ -8.1588\times10^{-2}\\ -4.7105\times10^{-2}\\ 73.425}
\end{equation*}}
Check satisfaction of constraint:
\begin{equation*}
  v_b'' = -0.5(-8.1588\times10^{-2}) + \sqrt3/2(-4.7105\times10^{-2}) = -1.2666\times10^{-7}\simeq 0
\end{equation*}
Displacement along the incline:
\begin{equation*}
  u_b'' = \sqrt3/2(-8.1588\times10^{-2}) + 0.5(-4.7105\times10^{-2}) = -.09421
\end{equation*}

\subsection*{Penalty function method}
Total potential:
\begin{equation*}
  \widehat\Pi = \tfrac12\bs{\Delta}^\tr\Kmat\bs{\Delta} - u_aP_{xa} - v_aP_{ya} + \tfrac12\alpha(v_b'')^2 = \dots + \tfrac12\alpha(\mat{C}\bs{\Delta})^\tr(\mat{C}\bs{\Delta}) = \dots + \tfrac12\alpha\bs{\Delta}^\tr\mat{C}^\tr\mat{C}\bs{\Delta}
\end{equation*}
Note that $\alpha$ represents the axial stiffness, $EA/L$, of the axial bar perpendicular to the incline. Theoretically, $\alpha$ should be infinite but this will make the system become ill-conditioned. Generally, $\alpha = 1000(k_{ii})_{\max}$ is a good choice.

$\pd{\widehat\Pi}{\bs{\Delta}} = \mat{0}$ results in the system of equations:
\begin{equation*}
  (\Kmat + \alpha\mat{C}^\tr\mat{C})\bs{\Delta} = \Pvec,
\end{equation*}
where
{\small\begin{equation*}
  \Pvec = \Bmat{P_{xa}\\ P_{ya}\\ 0\\ 0} = \Bmat{383\\ 321.4\\ 0\\ 0},\qquad
  \mat{C}^\tr\mat{C} = \Bmat{0\\ 0\\ -0.5\\ \sqrt3/2}\langle 0\quad 0\quad -0.5\quad \sqrt3/2\rangle
  = \bmat{0 & 0 & 0 & 0\\ 0 & 0 & 0 & 0\\ 0 & 0 & .25 & -.43301\\ 0 & 0 & -.43301 & .75}
\end{equation*}}
{\footnotesize
\begin{enumerate}[label=(\arabic*),leftmargin=*]
  \item $\alpha = 1000(k_{ii})_{\max} = 683000$
    \begin{gather*}
      \Kmat = 683000\bmat{0 & 0 & 0 & 0\\ 0 & 0 & 0 & 0\\ 0 & 0 & .25 & -.43301\\ 0 & 0 & -.43301 & .75} + \bmat{634.8 & -191.2 & -353.6 & 353.6\\ -191.2 & 447.3 & 353.6 & -353.6\\ -353.6 & 353.6 & 683.0 & -353.6\\ 353.6 & -353.6 & -353.6 & 353.6},\\
      \bmat{634.8 & -191.2 & -353.6 & 353.6\\ -191.2 & 447.3 & 353.6 & -353.6\\ -353.6 & 353.6 & 1.7143\times10^5 & -2.961\times10^5\\ 353.6 & -353.6 & -2.961\times10^5 & 5.126\times10^5}\Bmat{u_a\\ v_a\\ u_b\\ v_b} = \Bmat{383\\ 321.4\\ 0\\ 0},\qquad
      \text{Solution is: }\Bmat{.92813\\ 1.1428\\ -8.2269\times10^{-2}\\ -4.7374\times10^{-2}}
    \end{gather*}
    Check satisfaction of constraint: $v_b'' = -0.5(-8.2269\times10^{-2}) + \sqrt3/2(-4.7374\times10^{-2}) = 1.0741\times10^{-4}$\\
    Displacement along the incline: $u_b'' = \sqrt3/2(-8.2269\times10^{-2}) + 0.5(-4.7374\times10^{-2}) = -9.4934\times10^{-2}$
  \item $\alpha = 100(k_{ii})_{\max} = 68300$
    \begin{gather*}
      \Kmat = 68300\bmat{0 & 0 & 0 & 0\\ 0 & 0 & 0 & 0\\ 0 & 0 & .25 & -.43301\\ 0 & 0 & -.43301 & .75} + \bmat{634.8 & -191.2 & -353.6 & 353.6\\ -191.2 & 447.3 & 353.6 & -353.6\\ -353.6 & 353.6 & 683.0 & -353.6\\ 353.6 & -353.6 & -353.6 & 353.6},\\
      \bmat{634.8 & -191.2 & -353.6 & 353.6\\ -191.2 & 447.3 & 353.6 & -353.6\\ -353.6 & 353.6 & 17758. & -29928.\\ 353.6 & -353.6 & -29928. & 51579.}\Bmat{u_a\\ v_a\\ u_b\\ v_b} = \Bmat{383\\ 321.4\\ 0\\ 0},\qquad
      \text{Solution is: }\Bmat{.92784\\ 1.1433\\ -8.1453\times10^{-2}\\ -4.5785\times10^{-2}}
    \end{gather*}
    Check satisfaction of constraint: $v_b'' = -0.5(-8.1453\times10^{-2}) + \sqrt3/2(-4.5785\times10^{-2}) = 1.0755\times10^{-3}$\\
    Displacement along the incline: $u_b'' = \sqrt3/2(-8.1453\times10^{-2}) + 0.5(-4.5785\times10^{-2}) = -9.3433\times10^{-2}$
  \item $\alpha = 10(k_{ii})_{\max} = 6830$
    \begin{gather*}
      \Kmat = 6830\bmat{0 & 0 & 0 & 0\\ 0 & 0 & 0 & 0\\ 0 & 0 & .25 & -.43301\\ 0 & 0 & -.43301 & .75} + \bmat{634.8 & -191.2 & -353.6 & 353.6\\ -191.2 & 447.3 & 353.6 & -353.6\\ -353.6 & 353.6 & 683.0 & -353.6\\ 353.6 & -353.6 & -353.6 & 353.6},\\
      \bmat{634.8 & -191.2 & -353.6 & 353.6\\ -191.2 & 447.3 & 353.6 & -353.6\\ -353.6 & 353.6 & 2390.5 & -3311.1\\ 353.6 & -353.6 & -3311.1 & 5476.1}\Bmat{u_a\\ v_a\\ u_b\\ v_b} = \Bmat{383\\ 321.4\\ 0\\ 0},\qquad
      \text{Solution is: }\Bmat{.92373\\ 1.1505\\ -8.1595\times10^{-2}\\ -3.4696\times10^{-2}}
    \end{gather*}
    Check satisfaction of constraint: $v_b'' = -0.5(-8.1595\times10^{-2}) + \sqrt3/2(-3.4696\times10^{-2}) = .01075$\\
    Displacement along the incline: $u_b'' = \sqrt3/2(-8.1595\times10^{-2}) + 0.5(-3.4696\times10^{-2}) = -8.8011\times10^{-2}$
\end{enumerate}

}
\subsection*{Transformation method}
{\small
\topic{With symmetric system}
\begin{enumerate}[label=(\arabic*),leftmargin=*]
  \item Independent DOFs: $u_a$, $v_a$, and $u_b$ ($v_b = u_b/\sqrt3$)
    \begin{equation*}
      \Bmat{u_a\\ v_a\\ u_b\\ v_b} = \bmat{1 & 0 & 0\\ 0 & 1 & 0\\ 0 & 0 & 1\\ 0 & 0 & 1/\sqrt3}\Bmat{u_a\\ v_a\\ u_b}\qquad\text{or}\qquad \bs{\Delta} = \Tmat\bs{\Delta}'
    \end{equation*}
    \begin{align*}
      \Kmat' = \Tmat^\tr\Kmat\Tmat &= \bmat{1 & 0 & 0 & 0\\ 0 & 1 & 0 & 0\\ 0 & 0 & 1 & 1/\sqrt3}100\bmat{6.348 & -1.912 & -3.536 & 3.536\\ -1.912 & 4.473 & 3.536 & -3.536\\ -3.536 & 3.536 & 6.830 & -3.536\\ 3.536 & -3.536 & -3.536 & 3.536}\bmat{1 & 0 & 0\\ 0 & 1 & 0\\ 0 & 0 & 1\\ 0 & 0 & 1/\sqrt3}\\
      &= \bmat{634.8 & -191.2 & -149.45\\ -191.2 & 447.3 & 149.45\\ -149.45 & 149.45 & 392.56}\\
      \Pvec' = \Tmat^\tr\Pvec &= \bmat{1 & 0 & 0 & 0\\ 0 & 1 & 0 & 0\\ 0 & 0 & 1 & 1/\sqrt3}\Bmat{383\\ 321.4\\ -0.5r\\ 0.5r\sqrt3} = \Bmat{383.0\\ 321.4\\ 0}
    \end{align*}
    Note that $r$ is the reaction of the inclined support, being positive along the $y''$ axis. Namely, $r = -\lambda$.
    \begin{equation*}
      \bmat{634.8 & -191.2 & -149.45\\ -191.2 & 447.3 & 149.45\\ -149.45 & 149.45 & 392.56}\Bmat{u_a\\ v_a\\ u_b} = \Bmat{383\\ 321.4\\ 0},\quad\text{Solution is: }\Bmat{.92828\\ 1.1426\\ -.08159}
    \end{equation*}
    From the constraint equation, $v_b = u_b/\sqrt3 = -.08159/\sqrt3 = -4.7106\times10^{-2}$.
  \item Independent DOFs: $u_a$, $v_a$, and $v_b$ ($u_b = \sqrt3v_b$)
    \begin{equation*}
      \Bmat{u_a\\ v_a\\ v_b\\ u_b} = \bmat{1 & 0 & 0\\ 0 & 1 & 0\\ 0 & 0 & 1\\ 0 & 0 & \sqrt3}\Bmat{u_a\\ v_a\\ v_b}\qquad\text{or}\qquad \bs{\Delta} = \Tmat\bs{\Delta}'
    \end{equation*}
    \begin{align*}
      \Kmat' = \Tmat^\tr\Kmat\Tmat &= \bmat{1 & 0 & 0 & 0\\ 0 & 1 & 0 & 0\\ 0 & 0 & 1 & \sqrt3}100\bmat{6.348 & -1.912 & 3.536 & -3.536\\ -1.912 & 4.473 & -3.536 & 3.536\\ 3.536 & -3.536 & 3.536 & -3.536\\ -3.536 & 3.536 & -3.536 & 6.830}\bmat{1 & 0 & 0\\ 0 & 1 & 0\\ 0 & 0 & 1\\ 0 & 0 & \sqrt3}\\
      &= \bmat{634.8 & -191.2 & -258.85\\ -191.2 & 447.3 & 258.85\\ -258.85 & 258.85 & 1177.7}\\
      \Pvec' = \Tmat^\tr\Pvec &= \bmat{1 & 0 & 0 & 0\\ 0 & 1 & 0 & 0\\ 0 & 0 & 1 & \sqrt3}\Bmat{383\\ 321.4\\ 0.5r\sqrt3\\ -0.5r} = \Bmat{383.0\\ 321.4\\ 0}
    \end{align*}
    \begin{equation*}
      \bmat{634.8 & -191.2 & -258.85\\ -191.2 & 447.3 & 258.85\\ -258.85 & 258.85 & 1177.7}\Bmat{u_a\\ v_a\\ v_b} = \Bmat{383.0\\ 321.4\\ 0},\quad\text{Solution is: }\Bmat{.92828\\ 1.1426\\ -4.7104\times10^{-2}}
    \end{equation*}
    From the constraint equation, $u_b = \sqrt3v_b = \sqrt3(-4.7104\times10^{-2}) = -8.1587\times10^{-2}$.
\end{enumerate}

}
\topic{With unsymmetric system}
\begin{equation*}
  \bmat{634.8 & -191.2 & -353.6 & 353.6\\ -191.2 & 447.3 & 353.6 & -353.6\\ -353.6 & 353.6 & 683.0 & -353.6\\ 353.6 & -353.6 & -353.6 & 353.6}\Bmat{u_a\\ v_a\\ u_b\\ v_b} = \Bmat{383\\ 321.4\\ -0.5r\\ 0.5r\sqrt3}
\end{equation*}
Making use of the constraint, $v_b = u_b/\sqrt3$, the above system becomes
\begin{equation*}
  \bmat{634.8 & -191.2 & -149.45\\ -191.2 & 447.3 & 149.45\\ -353.6 & 353.6 & 478.85\\ 353.6 & -353.6 & -149.45}\Bmat{u_a\\ v_a\\ u_b} = \Bmat{383\\ 321.4\\ -0.5r\\ 0.5r\sqrt3}
\end{equation*}
This can also be rewritten as
\begin{equation*}
  \bmat{634.8 & -191.2 & -149.45 & 0\\ -191.2 & 447.3 & 149.45 & 0\\ -353.6 & 353.6 & 478.85 & .5\\ 353.6 & -353.6 & -149.45 & -.5\sqrt3}\Bmat{u_a\\ v_a\\ u_b\\ r} = \Bmat{383\\ 321.4\\ 0\\ 0},\quad
  \text{Solution is: }\Bmat{.92828\\ 1.1426\\ -8.1588\times10^{-2}\\ -73.425}
\end{equation*}
From the constraint equation, $v_b = u_b/\sqrt3 = (-8.1588\times10^{-2})/\sqrt3 = -4.7105\times10^{-2}$.

If the system is formulated corresponding to $u_b = \sqrt3v_b$ (take $v_b$ as the independent variable), the same results will be obtained.
\end{example}

\section{Skew support (joint coordinate): global DOFs + local DOFs}
Use Example~13.5 in the textbook (on pages 390--391) to describe how this method works.
\begin{fig}
  % Skew support with joint (local) coordinates at b (Example 13.5 of the textbook)
\begin{tikzpicture}[scale=0.8,font=\small]
  \coordinate (c) at (0,0); \coordinate (b) at (8,0); \coordinate (a) at (5,2.6);
  \draw[member] (c) -- (a) -- (b) -- (c);
  \pic at (c) {pin}; \pic[rotate=30] at (b) {roller};
  \foreach \p in {a,b,c}{\node[node dot] at (\p) {};}
  \node[left] at (c) {$c$}; \node[above left] at (a) {$a$}; \node[below right] at ($(b)+(0.05,-0.05)$) {$b$};
  \draw[load] (a) -- ++(0,1.5) node[right] {$v_a$}; \draw[load] (a) -- ++(1.5,0) node[right] {$u_a$};
  \draw[load] (b) -- ++(0,1.5) node[right] {$v_b$}; \draw[load] (b) -- ++(1.5,0) node[right] {$u_b$};
  \draw[load] (b) -- ++(30:1.5) node[right] {$u_b''$}; \draw[load] (b) -- ++(120:1.5) node[left] {$v_b''$};
  \draw[->] ($(b)+(1.1,0)$) arc (0:30:1.1); \node[right] at ($(b)+(1.2,0.35)$) {$30^\circ$};
\end{tikzpicture}

\end{fig}
At joint $b$
\begin{equation*}
  \Bmat{u_b\\ v_b} = \bmat{\cos\phi & -\sin\phi\\ \sin\phi & \cos\phi}\Bmat{u_b''\\ v_b''} = \bmat{0.866 & -0.5\\ 0.5 & 0.866}\Bmat{u_b''\\ v_b''}
\end{equation*}
Therefore, we have
\begin{equation*}
  \bs{\Delta} = \Bmat{u_a\\ v_a\\ u_b\\ v_b} = \bmat{1 & 0 & 0 & 0\\ 0 & 1 & 0 & 0\\ 0 & 0 & 0.866 & -0.5\\ 0 & 0 & 0.5 & 0.866}\Bmat{u_a\\ v_a\\ u_b''\\ v_b''}\qquad\text{or}\qquad \bs{\Delta} = \Tmat\bs{\Delta}''
\end{equation*}
Then the stiffness matrix $\kmat'' = \Tmat^\tr\kmat\Tmat$ and load vector $\Pvec'' = \Tmat^\tr\Pvec$ can be obtained as follows.
\begin{equation*}
  \kmat'' = 10^2\bmat{6.348 & -1.192 & -1.294 & 4.83\\ & 4.473 & 1.294 & -4.83\\ & & 2.944 & -3.194\\ \textit{sym.} & & & 7.421}
\end{equation*}
\begin{equation*}
  \Pvec'' = \Tmat^\tr\Pvec = \langle P_{xa}, P_{ya}, 0.866P_{xb} + 0.5P_{yb}, -0.5P_{xb} + 0.866P_{yb}\rangle^\tr = \langle P_{xa}, P_{ya}, 0, r\rangle^\tr
\end{equation*}
where $r$ is the reaction, $P_{xb} = -r/2$ and $P_{yb} = r\sqrt3/2$.

Although $r$ is unknown, considering that $v_b'' = 0$, we only need to solve the first three equations; the solution is:
\begin{equation*}
  \bs{\Delta}'' = \langle 0.928\quad 1.143\quad -0.0947\quad 0\rangle\ \text{mm}
\end{equation*}

\section{Flexible connection (semi-rigid joint)}
\begin{fig}
  % Idealized rigid / pinned connection systems and the real system
\begin{tikzpicture}[scale=0.8,font=\small]
  \newcommand{\frameshape}{\draw[thinmember] (0,0) -- (0,4) -- (4,4) -- (4,0); \pic at (0,0) {fixed}; \pic at (4,0) {fixed};}
  \begin{scope}
    \frameshape
    \node[below,align=center] at (2,-0.5) {idealized rigid connection\\system};
  \end{scope}
  \begin{scope}[xshift=6.5cm]
    \frameshape \node[node dot] at (0,4) {}; \node[node dot] at (4,4) {};
    \node[below,align=center] at (2,-0.5) {idealized pinned connection\\system};
  \end{scope}
  \begin{scope}[xshift=13cm]
    \frameshape
    \foreach \x in {0,4}{\draw[thick] (\x,4) circle (0.18); \draw[thick] (\x,4) circle (0.09);}
    \node[below] at (2,-0.5) {real system};
  \end{scope}
\end{tikzpicture}

\end{fig}
The real one shall be discussed by regarding the joint as an internal rotational spring. We will derive the stiffness matrix for a beam element with rotational springs at its ends.

\begin{example}
Determine the stiffness matrix with rotational springs at both ends for the simple beam shown below.
\begin{fig}
  % Simple beam with rotational springs at both ends
\begin{tikzpicture}[scale=0.9,font=\small]
  \draw[thinmember] (0,0) -- (7,0);
  \pic at (0,0) {roller}; \pic at (7,0) {pin};
  \foreach \x in {0,7}{\draw[thick] (\x,0) circle (0.2); \draw[thick] (\x,0) circle (0.1);}
  \draw[->,thick] (-0.4,0.7) arc (110:-40:0.5); \node[above left] at (-0.2,0.8) {$M_1(\alpha_1)$};
  \draw[->,thick] (7.4,0.7) arc (70:220:0.5); \node[above right] at (7.2,0.8) {$M_2(\alpha_2)$};
  \node[below right] at (0.2,-0.2) {$k_{\theta1}$}; \node[below left] at (6.8,-0.2) {$k_{\theta2}$};
  \node[below] at (3.5,-0.2) {$EI, L$};
\end{tikzpicture}

\end{fig}
The relationship between the end rotations $\bs{\alpha}$ and end moments $\mat{M}$ is given by
\begin{fig}
  % Elastic beam + rigid bar with end rotational springs (springs in series with the beam)
\begin{tikzpicture}[scale=0.85,font=\small]
  \begin{scope}
    \draw[thinmember] (0,0) -- (6,0); \pic at (0,0) {roller}; \pic at (6,0) {pin};
    \node[node dot] at (0,0) {}; \node[node dot] at (6,0) {};
    \node[below] at (3,-0.2) {$EI, L$};
  \end{scope}
  \node at (7.3,0.3) {\large $+$};
  \begin{scope}[xshift=8.6cm]
    \draw[member] (0,0) -- (6,0); \pic at (0,0) {roller}; \pic at (6,0) {pin};
    \foreach \x in {0,6}{\draw[thick] (\x,0) circle (0.2); \draw[thick] (\x,0) circle (0.1);}
    \node[below right] at (0.2,-0.2) {$k_{\theta1}$}; \node[below left] at (5.8,-0.2) {$k_{\theta2}$};
    \node[below] at (3,-0.2) {\itshape rigid};
  \end{scope}
\end{tikzpicture}

\end{fig}
\begin{equation*}
  \begin{aligned}
    \bs{\alpha} = \Bmat{\alpha_1\\ \alpha_2} &= \begin{bmatrix}\dfrac{L}{3EI} & -\dfrac{L}{6EI}\\[8pt] -\dfrac{L}{6EI} & \dfrac{L}{3EI}\end{bmatrix}\Bmat{M_1\\ M_2} + \Bmat{\dfrac{M_1}{k_{\theta1}}\\[8pt] \dfrac{M_2}{k_{\theta2}}}\\
    &= \begin{bmatrix}\dfrac{L}{3EI} + \dfrac{1}{k_{\theta1}} & -\dfrac{L}{6EI}\\[8pt] -\dfrac{L}{6EI} & \dfrac{L}{3EI} + \dfrac{1}{k_{\theta2}}\end{bmatrix}\Bmat{M_1\\ M_2} = \mat{d}\mat{M}
  \end{aligned}
\end{equation*}
where $\mat{d}$ is the flexibility matrix of this beam. Note that the two rotational springs act in series with the elastic beam.

By $\kmat = \Tmat^\tr\kmat_n\Tmat = \Tmat^\tr\mat{d}^{-1}\Tmat$, the stiffness matrix is
\begin{equation*}
  \kmat = \frac{\alpha EI}{L}
  \begin{bmatrix}
    \dfrac{12}{L^2}\Bigl(1+\dfrac{\alpha_1+\alpha_2}{\alpha_1\alpha_2}\Bigr) & \dfrac6L\Bigl(1+\dfrac{2}{\alpha_2}\Bigr) & -\dfrac{12}{L^2}\Bigl(1+\dfrac{\alpha_1+\alpha_2}{\alpha_1\alpha_2}\Bigr) & \dfrac6L\Bigl(1+\dfrac{2}{\alpha_1}\Bigr)\\[10pt]
    & 4\Bigl(1+\dfrac{3}{\alpha_2}\Bigr) & -\dfrac6L\Bigl(1+\dfrac{2}{\alpha_2}\Bigr) & 2\\[10pt]
    & & \dfrac{12}{L^2}\Bigl(1+\dfrac{\alpha_1+\alpha_2}{\alpha_1\alpha_2}\Bigr) & -\dfrac6L\Bigl(1+\dfrac{2}{\alpha_1}\Bigr)\\[10pt]
    \textit{sym.} & & & 4\Bigl(1+\dfrac{3}{\alpha_1}\Bigr)
  \end{bmatrix}
\end{equation*}
where $\alpha = \dfrac{\alpha_1\alpha_2}{\alpha_1\alpha_2 + 4\alpha_1 + 4\alpha_2 + 12}$, $k_{\theta1} = \alpha_1\dfrac{EI}{L}$ and $k_{\theta2} = \alpha_2\dfrac{EI}{L}$.

\topic{Special cases:}
\begin{enumerate}[leftmargin=*]
  \item $\alpha_1\to\infty$, $\alpha_2\to\infty$ (rigid connections at two ends) $\Rightarrow\alpha = 1$
    \begin{equation*}
      \kmat = \frac{EI}{L}\begin{bmatrix}\dfrac{12}{L^2} & \dfrac6L & -\dfrac{12}{L^2} & \dfrac6L\\[8pt] & 4 & -\dfrac6L & 2\\[8pt] & & \dfrac{12}{L^2} & -\dfrac6L\\[8pt] \textit{sym.} & & & 4\end{bmatrix}
    \end{equation*}
  \item $\alpha_1\to\infty$, $\alpha_2 =$ finite (rigid connection at left end) $\Rightarrow\alpha = \dfrac{\alpha_2}{\alpha_2+4}$
    \begin{equation*}
      \kmat = \frac{EI}{L}
      \begin{bmatrix}
        \dfrac{12}{L^2}\dfrac{\alpha_2+1}{\alpha_2+4} & \dfrac6L\dfrac{\alpha_2+2}{\alpha_2+4} & -\dfrac{12}{L^2}\dfrac{\alpha_2+1}{\alpha_2+4} & \dfrac6L\dfrac{\alpha_2}{\alpha_2+4}\\[10pt]
        & 4\dfrac{\alpha_2+3}{\alpha_2+4} & -\dfrac6L\dfrac{\alpha_2+2}{\alpha_2+4} & 2\dfrac{\alpha_2}{\alpha_2+4}\\[10pt]
        & & \dfrac{12}{L^2}\dfrac{\alpha_2+1}{\alpha_2+4} & -\dfrac6L\dfrac{\alpha_2}{\alpha_2+4}\\[10pt]
        \textit{sym.} & & & 4\dfrac{\alpha_2}{\alpha_2+4}
      \end{bmatrix}
    \end{equation*}
  \item $\alpha_1\to\infty$, $\alpha_2 = 0$ (rigid connection at left end and pin connection at right end)
    \begin{equation*}
      \kmat = \frac{EI}{L}\begin{bmatrix}\dfrac{3}{L^2} & \dfrac3L & -\dfrac{3}{L^2} & 0\\[8pt] & 3 & -\dfrac3L & 0\\[8pt] & & \dfrac{3}{L^2} & 0\\[8pt] \textit{sym.} & & & 0\end{bmatrix}\qquad\text{(same as the element with a hinge at the right end)}
    \end{equation*}
\end{enumerate}
\end{example}

\begin{example}
The example is used to show how rigid a joint is. If the rotational spring with the stiffness $k_\theta = \alpha\dfrac{EI}{L}$, what is the displacement at the loaded point?
\begin{fig}
  % Cantilever with a rotational spring at the fixed end, loaded by P at the tip
\begin{tikzpicture}[scale=0.9,font=\small]
  \pic at (0,0) {wallL};
  \draw[thick] (0.2,0) circle (0.2); \draw[thick] (0.2,0) circle (0.1);
  \draw[thinmember] (0.4,0) -- (6,0);
  \draw[load] (6,1.2) -- (6,0.1) node[pos=0,above left] {$P$};
  \node[below] at (0.4,-0.3) {$k_\theta$}; \node[below] at (3.2,-0.3) {$EI, L$};
\end{tikzpicture}

\end{fig}
The vertical displacement of the free end comes from two parts. One is the elastic deformation of the beam with the left end fixed, and the other is the deformation due to a rigid bar with a rotational spring at the left end.
\begin{fig}
  % Elastic cantilever + rigid bar with a rotational spring
\begin{tikzpicture}[scale=0.85,font=\small]
  \begin{scope}
    \pic at (0,0) {wallL}; \draw[thinmember] (0,0) -- (6,0);
    \draw[load] (6,1.2) -- (6,0.1) node[pos=0,above left] {$P$};
    \node[below] at (3,-0.3) {$EI, L$};
  \end{scope}
  \node at (7.2,0.3) {\large $+$};
  \begin{scope}[xshift=8.4cm]
    \pic at (0,0) {wallL}; \draw[thick] (0.2,0) circle (0.2); \draw[thick] (0.2,0) circle (0.1);
    \draw[member] (0.4,0) -- (6,0);
    \draw[load] (6,1.2) -- (6,0.1) node[pos=0,above left] {$P$};
    \node[below] at (0.4,-0.3) {$k_\theta$}; \node[below] at (3.2,-0.3) {\itshape rigid};
  \end{scope}
\end{tikzpicture}

\end{fig}
\begin{equation*}
  \Delta_v = \Delta_e + \Delta_s = \frac{PL^3}{3EI} + \frac{PL^3}{\alpha EI} = \frac{PL^3}{EI}\Bigl(\frac13 + \frac1\alpha\Bigr)
\end{equation*}
If $\alpha > 100$, the displacement contributed from the rotational spring is very small.

Other derivations of the stiffness matrix of a beam element with rotational springs can be found in the textbook (see Examples~13.6 and 13.7).

Note that if the end rotational springs are attached to the beam externally, then
\begin{equation*}
  \kmat_n = \frac{2EI}{L}\bmat{2 & 1\\ 1 & 2} + \bmat{k_{\theta1} & \\ & k_{\theta2}}.
\end{equation*}
The stiffness matrix for the unsupported element can be derived similarly using the congruent transformation. When do we need this kind of external rotational springs?
\begin{fig}
  % Beam with external rotational springs at both ends
\begin{tikzpicture}[scale=0.9,font=\small]
  \draw[thinmember] (0,0) -- (7,0); \node[node dot] at (0,0) {}; \node[node dot] at (7,0) {};
  \pic at (0,0) {roller}; \pic at (7,0) {pin};
  % spiral springs to inclined walls
  \draw[thick] (0,0) .. controls (-0.2,0.6) and (0.5,0.9) .. (0.3,0.4) .. controls (0.1,0.1) and (-0.4,0.5) .. (-0.6,0.8);
  \pic[rotate=-40] at (-0.7,0.9) {wallR};
  \draw[thick] (7,0) .. controls (7.2,0.6) and (6.5,0.9) .. (6.7,0.4) .. controls (6.9,0.1) and (7.4,0.5) .. (7.6,0.8);
  \pic[rotate=40] at (7.7,0.9) {wallL};
  \node[above] at (-0.4,1.1) {$k_{\theta1}$}; \node[above] at (7.4,1.1) {$k_{\theta2}$};
\end{tikzpicture}

\end{fig}
\end{example}

\subsection*{Chen and Toma (1994)}
\begin{figure}[H]
  \centering
  % Typical beam-to-column connections (schematic, after Chen & Toma 1994, Fig. 1.7)
\begin{tikzpicture}[scale=0.55,font=\scriptsize]
  % column flange (vertical) + beam (horizontal) template
  \newcommand{\colbeam}{%
    \draw[thick] (0,-2.2) -- (0,2.2); \draw[thick] (-0.5,-2.2) -- (-0.5,2.2);   % column flanges
    \draw (0.6,0.9) -- (3.4,0.9); \draw (0.6,-0.9) -- (3.4,-0.9);            % beam flanges
    \draw (0.6,0.9) -- (0.6,-0.9); \draw[dashed] (3.4,0.9) -- (3.4,-0.9);}
  \newcommand{\bolt}[1]{\fill #1 circle (0.08);}
  % 1 single/double web angle
  \begin{scope}
    \colbeam
    \draw[thick] (0,0.6) -- (0.6,0.6) -- (0.6,-0.6) -- (0,-0.6); \bolt{(0.25,0.3)} \bolt{(0.25,-0.3)} \bolt{(0.9,0.3)} \bolt{(0.9,-0.3)}
    \node[below,align=center] at (1.7,-2.3) {Single web angle\\or double web angle};
  \end{scope}
  % 2 top and seat angle
  \begin{scope}[xshift=6cm]
    \colbeam
    \draw[thick] (0,1.6) -- (0,0.9) -- (1.4,0.9); \draw[thick] (0,-1.6) -- (0,-0.9) -- (1.4,-0.9);
    \bolt{(0.3,1.3)} \bolt{(1.0,1.1)} \bolt{(0.3,-1.3)} \bolt{(1.0,-1.1)}
    \node[below] at (1.7,-2.3) {Top and seat angle};
  \end{scope}
  % 3 header plate
  \begin{scope}[xshift=12cm]
    \colbeam
    \draw[thick] (0,0.7) rectangle (0.25,-0.7); \bolt{(0.12,0.4)} \bolt{(0.12,-0.4)} \draw (0.25,0.7) -- (0.6,0.7); \draw (0.25,-0.7) -- (0.6,-0.7);
    \node[below] at (1.7,-2.3) {Header plate};
  \end{scope}
  % 4 bottom flange and web angle
  \begin{scope}[xshift=18cm]
    \colbeam
    \draw[thick] (0,0.6) -- (0.6,0.6) -- (0.6,-0.6) -- (0,-0.6); \bolt{(0.25,0.2)} \bolt{(0.9,0.2)}
    \draw[thick] (0,-1.6) -- (0,-0.9) -- (1.4,-0.9); \bolt{(0.3,-1.3)} \bolt{(1.0,-1.1)}
    \node[below,align=center] at (1.7,-2.3) {Bottom flange and\\web angle};
  \end{scope}
  % 5 T-stub
  \begin{scope}[yshift=-6.5cm]
    \colbeam
    \draw[thick] (0,1.7) -- (0.25,1.7) -- (0.25,0.9) -- (1.4,0.9); \draw[thick] (0,0.2) -- (0.25,0.2); \draw[thick] (0.25,0.2) -- (0.25,0.9);
    \draw[thick] (0,-1.7) -- (0.25,-1.7) -- (0.25,-0.9) -- (1.4,-0.9); \draw[thick] (0,-0.2) -- (0.25,-0.2); \draw[thick] (0.25,-0.2) -- (0.25,-0.9);
    \bolt{(0.12,1.4)} \bolt{(0.12,0.45)} \bolt{(0.12,-1.4)} \bolt{(0.12,-0.45)} \bolt{(1.0,1.1)} \bolt{(1.0,-1.1)}
    \node[below] at (1.7,-2.3) {T-stub};
  \end{scope}
  % 6 top and seat angle with double web angle
  \begin{scope}[xshift=6cm,yshift=-6.5cm]
    \colbeam
    \draw[thick] (0,1.6) -- (0,0.9) -- (1.4,0.9); \draw[thick] (0,-1.6) -- (0,-0.9) -- (1.4,-0.9);
    \draw[thick] (0,0.6) -- (0.6,0.6) -- (0.6,-0.6) -- (0,-0.6);
    \bolt{(0.3,1.3)} \bolt{(1.0,1.1)} \bolt{(0.3,-1.3)} \bolt{(1.0,-1.1)} \bolt{(0.25,0.25)} \bolt{(0.25,-0.25)} \bolt{(0.9,0.25)} \bolt{(0.9,-0.25)}
    \node[below,align=center] at (1.7,-2.3) {Top and seat angle\\with double web angle};
  \end{scope}
  % 7 flush end plate
  \begin{scope}[xshift=12cm,yshift=-6.5cm]
    \colbeam
    \draw[thick] (0,1.0) rectangle (0.25,-1.0); \draw (0.25,0.9) -- (0.6,0.9); \draw (0.25,-0.9) -- (0.6,-0.9);
    \bolt{(0.12,0.6)} \bolt{(0.12,0.2)} \bolt{(0.12,-0.2)} \bolt{(0.12,-0.6)}
    \node[below] at (1.7,-2.3) {Flush end plate};
  \end{scope}
  % 8 extended end plate
  \begin{scope}[xshift=18cm,yshift=-6.5cm]
    \colbeam
    \draw[thick] (0,1.6) rectangle (0.25,-1.6); \draw (0.25,0.9) -- (0.6,0.9); \draw (0.25,-0.9) -- (0.6,-0.9);
    \bolt{(0.12,1.3)} \bolt{(0.12,0.55)} \bolt{(0.12,-0.55)} \bolt{(0.12,-1.3)}
    \fill[hatch] (-0.5,0.85) rectangle (0,1.0); \fill[hatch] (-0.5,-1.0) rectangle (0,-0.85);
    \draw (-0.25,-1.0) -- (-0.9,-1.9) node[left] {Stiffeners};
    \node[below] at (1.7,-2.3) {Extended end plate};
  \end{scope}
\end{tikzpicture}

  \caption{Typical beam-to-column connections (after Chen and Toma, 1994, Fig.~1.7).}
\end{figure}
\begin{enumerate}[label=(\arabic*),start=4]
  \item Flange angle.
  \item Header plate.
  \item Flush/extended end plate.
  \item Combined web and flange angles.
  \item T-stubs.
  \item Top and seat angle.
  \item T-stubs and web angles.
\end{enumerate}
Some of the most commonly used connection types reported in Goverdhan (1984) and Nethercot (1985) are shown in the figure above. Typical moment-rotation curves for these connections are illustrated in the figure below. It should be noted that for any connection type, the stiffness and strength are dependent upon the geometric parameters such as plate or angle thickness, bolt size, method of tightening, and connection depth. The more flexible ones, such as single and double web angle connections, may be approximated as pin connections represented by the horizontal axis in the figure. These types of connections might be classified as ``flexible.'' The stiffest ones, such as the extended end plate connection, might be classified as ``rigid,'' represented by the vertical axis. Those with intermediate stiffness might be classified as ``semi-rigid.'' All welded connections, which are not shown in the figure, tend to be extremely stiff and they are usually designed to transmit full moment capacity of the adjoining members. For the purpose of design, they may be assumed as perfectly rigid.

\paragraph{Kishi and Chen data base.} In 1986, Kishi and Chen conducted a comprehensive search on beam-column connection data. Characteristic results on moment-rotation response and the corresponding parameters of beam-to-column connections used frequently in steel construction were collected and stored in a data base (Kishi and Chen, 1986, Kishi et al., 1992). The data base consisted of experimental data on rivet, bolted, and welded connections that were published from 1936 to 1986. The experimental data collected were compared with some prediction equations in order to develop a rational method of analysis for semi-rigid frame design.
\begin{figure}[H]
  \centering
  % Typical moment-rotation curves of beam-to-column connections (after Chen & Toma 1994, Fig. 1.8)
\begin{tikzpicture}
  \begin{axis}[
    width=13cm,height=9cm,
    xlabel={Rotation (radians)},ylabel={Moment (kNm)},
    xmin=0,xmax=0.036,ymin=0,ymax=400,
    xtick={0,0.004,...,0.036},xticklabel style={/pgf/number format/fixed,/pgf/number format/precision=3},
    ytick={0,40,...,400},
    clip=false,samples=80,domain=0:0.024,
    every axis plot/.append style={thick,black},
    cycle list={{black},{black},{black},{black},{black},{black},{black}},
    scaled x ticks=false,
  ]
    % curves of the form M = M0 * (1 - exp(-x/c))
    \addplot {390*(1-exp(-x/0.006))} node[right,font=\scriptsize] {Extended end plate};
    \addplot {350*(1-exp(-x/0.007))} node[right,font=\scriptsize] {Flush end plate};
    \addplot {265*(1-exp(-x/0.008))} node[right,align=left,font=\scriptsize] {Top and seat angle\\with web angles};
    \addplot {210*(1-exp(-x/0.008))} node[right,font=\scriptsize] {Top and seat angle};
    \addplot {105*(1-exp(-x/0.009))} node[right,font=\scriptsize] {Header plate};
    \addplot {70*(1-exp(-x/0.010))} node[right,font=\scriptsize] {Double web angle};
    \addplot {40*(1-exp(-x/0.010))} node[right,font=\scriptsize] {Single web angle};
    \draw[thick] (axis cs:0,0) -- (axis cs:0.0005,290); \node[right,font=\scriptsize] at (axis cs:0.0006,265) {``Rigid''};
    \draw[thick] (axis cs:0,0) -- (axis cs:0.024,12); \node[right,font=\scriptsize] at (axis cs:0.0245,12) {``Pinned''};
  \end{axis}
\end{tikzpicture}

  \caption{Typical moment-rotation curves (after Chen and Toma, 1994, Fig.~1.8).}
\end{figure}
In particular, three prediction equations were discussed in detail. The first was the analytical polynomial equation proposed by Frye and Morris (1976). The second was a curve-fitting equation using a modified exponential model proposed by Kishi and Chen (1986). The third model utilizes a three-parameter power model proposed by Kishi and Chen (1990), which is a simplified analytical model suitable for use in the analysis and design of steel frames.

The Kishi and Chen connection data base is made available to the reader in Chapter~3 of Chen and Toma. This data base includes connection types with semi-rigid moment-rotation characteristics similar to those reported in Goverdhan's data bank. Available test data on steel beam-to-column connections and comparison between experiment and theoretical results calculated using the three prediction models are available there.

\section{Symmetry and anti-symmetry}
The loads can be separated into two parts, one is the symmetric one, and the other is the anti-symmetric one. The DOFs for these two problems are listed in the table below.
\begin{center}
  \begin{tabular}{|c|c|c|}\hline
    & symmetric part & anti-symmetric part\\\hline
    DOFs & $u_b, v_b, \theta_b, v_c$ & $u_b, v_b, \theta_b, u_c, \theta_c$\\\hline
  \end{tabular}
\end{center}
\begin{fig}
  % Gable frame: loads, symmetric part, anti-symmetric part, and the half models
\begin{tikzpicture}[scale=0.62,font=\small]
  \newcommand{\gable}{\draw[thinmember] (0,0) -- (0,3) -- (2.5,4.2) -- (5,3) -- (5,0); \pic at (0,0) {fixed}; \pic at (5,0) {fixed};
    \node[left] at (0,0.3) {a}; \node[left] at (0,2.8) {b}; \node[above right] at (2.5,4.2) {c}; \node[right] at (5,2.8) {d}; \node[right] at (5,0.3) {e};}
  \newcommand{\halfgable}{\draw[thinmember] (0,0) -- (0,3) -- (2.5,4.2); \pic at (0,0) {fixed};
    \node[left] at (0,0.3) {a}; \node[left] at (0,2.8) {b};}
  % ---- original loading ----
  \begin{scope}[xshift=13cm,yshift=1cm]
    \gable
    \draw[load] (-1.5,3) -- (-0.1,3) node[pos=0,above] {8};
    \draw[load] (1.25,5.2) -- (1.25,3.7) node[pos=0,above] {12};
    \draw[load] (2.5,5.8) -- (2.5,4.3) node[pos=0,above] {10};
  \end{scope}
  % ---- symmetric ----
  \begin{scope}[yshift=-7cm]
    \gable
    \draw[load] (-1.5,3) -- (-0.1,3) node[pos=0,above] {4}; \draw[load] (6.5,3) -- (5.1,3) node[pos=0,above] {4};
    \draw[load] (1.25,5.2) -- (1.25,3.7) node[pos=0,above] {6}; \draw[load] (3.75,5.2) -- (3.75,3.7) node[pos=0,above] {6};
    \draw[load] (2.5,5.8) -- (2.5,4.3) node[pos=0,above] {5+5};
    \node[below] at (2.5,-0.6) {symmetric};
    \node at (7.5,2) {$\Longrightarrow$};
    \begin{scope}[xshift=9cm]
      \halfgable
      \draw[load] (-1.5,3) -- (-0.1,3) node[pos=0,above] {4};
      \draw[load] (1.25,5.2) -- (1.25,3.7) node[pos=0,above] {6};
      \draw[load] (2.5,5.8) -- (2.5,4.3) node[pos=0,above] {5};
      \node[below right] at (2.5,4.1) {c};
      \pic[rotate=-90] at (2.5,4.2) {roller}; \fill[hatch] (3.1,3.5) rectangle (3.3,4.9); \draw (3.1,3.5) -- (3.1,4.9);
    \end{scope}
  \end{scope}
  % ---- anti-symmetric ----
  \begin{scope}[yshift=-15cm]
    \gable
    \draw[load] (-1.5,3) -- (-0.1,3) node[pos=0,above] {4}; \draw[load] (5.1,3) -- (6.5,3) node[above] {4};
    \draw[load] (1.25,5.2) -- (1.25,3.7) node[pos=0,above] {6}; \draw[load] (3.75,3.7) -- (3.75,5.2) node[above] {6};
    \draw[load] (2.35,5.8) -- (2.35,4.3); \draw[load] (2.65,4.3) -- (2.65,5.8); \node[above] at (2.5,5.8) {5\ \,5};
    \node[below] at (2.5,-0.6) {anti-symmetric};
    \node at (7.5,2) {$\Longrightarrow$};
    \begin{scope}[xshift=9cm]
      \halfgable
      \draw[load] (-1.5,3) -- (-0.1,3) node[pos=0,above] {4};
      \draw[load] (1.25,5.2) -- (1.25,3.7) node[pos=0,above] {6};
      \node[above right] at (2.5,4.2) {c};
      \pic[rotate=-90] at (2.5,4.2) {roller}; % roller allowing horizontal movement, v_c = 0
    \end{scope}
  \end{scope}
\end{tikzpicture}

\end{fig}
\begin{center}
  \renewcommand{\arraystretch}{1.2}
  \begin{tabular}{|c|c|c|}\hline
    \multirow{6}{*}{The other DOFs} & \multirow{6}{*}{Sym.\ $\begin{aligned}u_d &= -u_b\\ v_d &= v_b\\ \theta_d &= -\theta_b\\ u_c &= 0\\ \theta_c &= 0\\ M_c &\neq 0\end{aligned}$}
    & \multirow{6}{*}{Anti-sym.\ $\begin{aligned}u_d &= u_b\\ v_d &= -v_b\\ \theta_d &= \theta_b\\ v_c &= 0\\ u_c &\neq 0\\ M_c &= 0\end{aligned}$}\\
    &&\\ &&\\ &&\\ &&\\ &&\\\hline
  \end{tabular}
\end{center}

\section{Effect of axial deformation (portal frame)}
\label{sec:axial-portal}
\begin{fig}
  % Portal frame with six DOFs (axial deformation included)
\begin{tikzpicture}[scale=0.9,font=\small]
  \draw[member] (0,0) -- (0,4) -- (4,4) -- (4,0);
  \pic at (0,0) {fixed}; \pic at (4,0) {fixed};
  \node at (2,2) {$EA, I, L$};
  \draw[load] (-1.6,4) -- (-0.9,4) node[pos=0,below] {$P$};
  \draw[load] (-0.8,4) -- (-0.1,4) node[pos=0,below] {$u_1$};
  \draw[load] (0,4.1) -- (0,5.0) node[above] {$u_2$};
  \draw[->,thick] (0.5,4.4) arc (30:170:0.5); \node[above left] at (-0.1,4.5) {$u_3$};
  \draw[load] (4.1,4) -- (5.0,4) node[right] {$u_4$};
  \draw[load] (4,4.1) -- (4,5.0) node[above] {$u_5$};
  \draw[->,thick] (4.5,4.4) arc (30:170:0.5); \node[above left] at (3.9,4.5) {$u_6$};
\end{tikzpicture}

\end{fig}
{\small
\begin{equation*}
  \Kmat = \begin{bmatrix}
    12E\frac{I}{L^3} + E\frac AL & 0 & 6E\frac{I}{L^2} & -E\frac AL & 0 & 0\\[3pt]
    0 & 12E\frac{I}{L^3} + E\frac AL & 6E\frac{I}{L^2} & 0 & -12E\frac{I}{L^3} & 6E\frac{I}{L^2}\\[3pt]
    6E\frac{I}{L^2} & 6E\frac{I}{L^2} & 8E\frac IL & 0 & -6E\frac{I}{L^2} & 2E\frac IL\\[3pt]
    -E\frac AL & 0 & 0 & 12E\frac{I}{L^3} + E\frac AL & 0 & 6E\frac{I}{L^2}\\[3pt]
    0 & -12E\frac{I}{L^3} & -6E\frac{I}{L^2} & 0 & 12E\frac{I}{L^3} + E\frac AL & -6E\frac{I}{L^2}\\[3pt]
    0 & 6E\frac{I}{L^2} & 2E\frac IL & 6E\frac{I}{L^2} & -6E\frac{I}{L^2} & 8E\frac IL
  \end{bmatrix},\qquad
  \Pvec = \Bmat{P\\ 0\\ 0\\ 0\\ 0\\ 0}
\end{equation*}
\begin{equation*}
  \bs{\Delta} = \Kmat^{-1}\Pvec = \Bmat{
    \frac{1}{12EI}L^3\frac{216I^2 + 84IAL^2 + 5A^2L^4}{72I^2 + 45IAL^2 + 7A^2L^4}P\\[4pt]
    3\frac{L^3}{(24I + 7AL^2)E}P\\[4pt]
    -\frac{1}{4EI}L^2\frac{96I^2 + 34IAL^2 + A^2L^4}{72I^2 + 45IAL^2 + 7A^2L^4}P\\[4pt]
    \frac{1}{12EI}L^3\frac{72I^2 + 42IAL^2 + 5A^2L^4}{72I^2 + 45IAL^2 + 7A^2L^4}P\\[4pt]
    -3\frac{L^3}{(24I + 7AL^2)E}P\\[4pt]
    -\frac{1}{4EI}L^2\frac{48I^2 + A^2L^4 + 20IAL^2}{72I^2 + 45IAL^2 + 7A^2L^4}P}
  = \Bmat{\frac{5}{84}\frac{PL^3}{EI}\\[4pt] 0\\[4pt] -\frac{1}{28}\frac{PL^2}{EI}\\[4pt] \frac{5}{84}\frac{PL^3}{EI}\\[4pt] 0\\[4pt] -\frac{1}{28}\frac{PL^2}{EI}}
  \quad\text{when }\frac{I}{AL^2}\to 0
\end{equation*}}
\begin{minipage}{0.4\linewidth}
\begin{align*}
  E &= 1.0\\ A &= 40(40)\ \mathrm{cm^2}\\ I &= \tfrac{1}{12}(40)^4\ \mathrm{cm^4}\\ L &= 300\ \mathrm{cm}
\end{align*}
\end{minipage}
\begin{minipage}{0.5\linewidth}\centering
\begin{tikzpicture}[scale=0.8]
  \fill[hatch] (0,0) rectangle (1.2,1.2); \draw (0,0) rectangle (1.2,1.2);
  \node[above] at (0.6,1.2) {40 cm}; \node[right] at (1.2,0.6) {40 cm};
\end{tikzpicture}
\end{minipage}
{\small
\begin{equation*}
  \begin{array}{cccc}
    A & 0.1A & 10A & 1000A\\[4pt]
    \Bmat{7.6487P\\ 7.9951\times10^{-2}P\\ -1.5679\times10^{-2}P\\ 7.5553P\\ -7.9951\times10^{-2}P\\ -1.5368\times10^{-2}P} &
    \Bmat{8.6378P\\ .76473P\\ -2.0933\times10^{-2}P\\ 7.7402P\\ -.76473P\\ -1.7941\times10^{-2}P} &
    \Bmat{7.5451P\\ 8.0316\times10^{-3}P\\ -1.5128\times10^{-2}P\\ 7.5357P\\ -8.0316\times10^{-3}P\\ -1.5097\times10^{-2}P} &
    \Bmat{7.5336P\\ 8.0357\times10^{-5}P\\ -1.5068\times10^{-2}P\\ 7.5335P\\ -8.0357\times10^{-5}P\\ -1.5067\times10^{-2}P}
  \end{array}
\end{equation*}}

\topic{Neglecting axial deformation ($u_2 = 0$, $u_5 = 0$; $u_1 = u_4$):}
\begin{equation*}
  \Kmat' = \begin{bmatrix}
    12E\frac{I}{L^3} + E\frac AL & 6E\frac{I}{L^2} & -E\frac AL & 0\\[3pt]
    6E\frac{I}{L^2} & 8E\frac IL & 0 & 2E\frac IL\\[3pt]
    -E\frac AL & 0 & 12E\frac{I}{L^3} + E\frac AL & 6E\frac{I}{L^2}\\[3pt]
    0 & 2E\frac IL & 6E\frac{I}{L^2} & 8E\frac IL
  \end{bmatrix}\quad\text{(DOFs } u_1, u_3, u_4, u_6\text{)}
\end{equation*}
\begin{equation*}
  \bs{\Delta}' = \Bmat{u_1\\ u_3\\ u_4\\ u_6} = \bmat{1 & 0 & 0\\ 0 & 1 & 0\\ 1 & 0 & 0\\ 0 & 0 & 1}\Bmat{u_1\\ u_3\\ u_6} = \Tmat\bs{\Delta}
\end{equation*}
\begin{equation*}
  \Kmat = \Tmat^\tr\Kmat'\Tmat = \bmat{1 & 0 & 1 & 0\\ 0 & 1 & 0 & 0\\ 0 & 0 & 0 & 1}
  \begin{bmatrix}
    12E\frac{I}{L^3} + E\frac AL & 6E\frac{I}{L^2} & -E\frac AL & 0\\[3pt]
    6E\frac{I}{L^2} & 8E\frac IL & 0 & 2E\frac IL\\[3pt]
    -E\frac AL & 0 & 12E\frac{I}{L^3} + E\frac AL & 6E\frac{I}{L^2}\\[3pt]
    0 & 2E\frac IL & 6E\frac{I}{L^2} & 8E\frac IL
  \end{bmatrix}\bmat{1 & 0 & 0\\ 0 & 1 & 0\\ 1 & 0 & 0\\ 0 & 0 & 1}
\end{equation*}
\begin{equation*}
  \therefore\Kmat = \bmat{24E\frac{I}{L^3} & 6E\frac{I}{L^2} & 6E\frac{I}{L^2}\\[3pt] 6E\frac{I}{L^2} & 8E\frac IL & 2E\frac IL\\[3pt] 6E\frac{I}{L^2} & 2E\frac IL & 8E\frac IL}
  \Rightarrow \Kmat^{-1} = \bmat{\frac{5}{84EI}L^3 & -\frac{1}{28EI}L^2 & -\frac{1}{28EI}L^2\\[3pt] -\frac{1}{28EI}L^2 & \frac{13}{84EI}L & -\frac{1}{84EI}L\\[3pt] -\frac{1}{28EI}L^2 & -\frac{1}{84EI}L & \frac{13}{84EI}L}
\end{equation*}
$\Pvec = \Tmat^\tr\Pvec'$, where $\Pvec'^\tr = \langle P, 0, 0, 0\rangle$
\begin{equation*}
  \bs{\Delta} = \Bmat{u_1\\ u_3\\ u_6} = \bmat{\frac{5}{84EI}L^3 & -\frac{1}{28EI}L^2 & -\frac{1}{28EI}L^2\\[3pt] -\frac{1}{28EI}L^2 & \frac{13}{84EI}L & -\frac{1}{84EI}L\\[3pt] -\frac{1}{28EI}L^2 & -\frac{1}{84EI}L & \frac{13}{84EI}L}\Bmat{P\\ 0\\ 0} = \Bmat{\frac{5}{84EI}PL^3\\[3pt] -\frac{1}{28EI}PL^2\\[3pt] -\frac{1}{28EI}PL^2}
\end{equation*}
With $E = 1.0$, $A = 40(40)$~cm$^2$, $I = \frac{1}{12}(40)^4$~cm$^4$, $L = 300$~cm:
\begin{equation*}
  \bs{\Delta} = \Bmat{\frac{3375}{448}P\\[3pt] -\frac{27}{1792}P\\[3pt] -\frac{27}{1792}P} = \Bmat{7.5335P\\ -1.5067\times10^{-2}P\\ -1.5067\times10^{-2}P}
\end{equation*}
\begin{center}
\begin{minipage}{0.4\linewidth}\centering% Portal frame with three DOFs (axial deformation neglected)
\begin{tikzpicture}[scale=0.9,font=\small]
  \draw[member] (0,0) -- (0,4) -- (4,4) -- (4,0);
  \pic at (0,0) {fixed}; \pic at (4,0) {fixed};
  \node at (2,2) {$EI, L$};
  \node[circle,draw,inner sep=1pt,fill=white] at (0,2) {\scriptsize 1};
  \node[circle,draw,inner sep=1pt,fill=white] at (2,4) {\scriptsize 2};
  \node[circle,draw,inner sep=1pt,fill=white] at (4,2) {\scriptsize 3};
  \draw[load] (-1.4,4) -- (-0.1,4) node[pos=0,left] {$v$};
  \draw[->,thick] (0.5,4.4) arc (30:170:0.5); \node[above left] at (-0.1,4.5) {$\theta_1$};
  \draw[->,thick] (4.5,4.4) arc (30:170:0.5); \node[above left] at (3.9,4.5) {$\theta_2$};
\end{tikzpicture}
\end{minipage}\hfill
\begin{minipage}{0.55\linewidth}
\begin{equation*}
  \Kmat = \begin{bmatrix}
    \dfrac{12EI}{L^3} + \dfrac{12EI}{L^3} & \dfrac{6EI}{L^2} & \dfrac{6EI}{L^2}\\[8pt]
    \dfrac{6EI}{L^2} & \dfrac{4EI}{L} + \dfrac{4EI}{L} & \dfrac{2EI}{L}\\[8pt]
    \dfrac{6EI}{L^2} & \dfrac{2EI}{L} & \dfrac{4EI}{L} + \dfrac{4EI}{L}
  \end{bmatrix}
\end{equation*}
\end{minipage}
\end{center}

\section{Beam elements in a rectangular frame with neglect of axial deformation}
No vertical displacements at both ends if the axial deformation of columns are very small and can be neglected.
\begin{center}
  \renewcommand{\arraystretch}{1.3}
  \begin{tabular}{|p{0.22\linewidth}|c|c|}\hline
    B1 & \begin{tikzpicture}[scale=0.7,font=\small,baseline=(current bounding box.center)]
  \draw[thinmember] (0,0) -- (5,0);
  \draw[->,thick] (0.6,0.45) arc (30:170:0.5); \node[above] at (0.1,0.7) {$\theta_1$};
  \draw[->,thick] (5.0,0.45) arc (30:170:0.5); \node[above] at (4.5,0.7) {$\theta_2$};
\end{tikzpicture}
 & $[k] = \dfrac{2EI}{L}\bmat{2 & 1\\ 1 & 2}$\\\hline
    B2\newline Pin at node 1\newline ($\theta_1$ is condensed) & \begin{tikzpicture}[scale=0.7,font=\small,baseline=(current bounding box.center)]
  \draw[thinmember] (0,0) -- (5,0); \node[node dot] at (0,0) {};
  \draw[->,thick] (5.0,0.45) arc (30:170:0.5); \node[above] at (4.5,0.7) {$\theta_2$};
\end{tikzpicture}
 & $k = \dfrac{3EI}{L}$\\\hline
    B3\newline Pin at node 2\newline ($\theta_2$ is condensed) & \begin{tikzpicture}[scale=0.7,font=\small,baseline=(current bounding box.center)]
  \draw[thinmember] (0,0) -- (5,0); \node[node dot] at (5,0) {};
  \draw[->,thick] (0.6,0.45) arc (30:170:0.5); \node[above] at (0.1,0.7) {$\theta_1$};
\end{tikzpicture}
 & $k = \dfrac{3EI}{L}$\\\hline
  \end{tabular}
\end{center}

\section{Column elements in a rectangular frame with neglect of axial deformation}
\begin{center}
  \renewcommand{\arraystretch}{1.3}
  \begin{tabular}{|c|c|c|}\hline
    C1 & \begin{tikzpicture}[scale=0.7,font=\small,baseline=(current bounding box.center)]
  \draw[thinmember] (0,0) -- (0,4);
  \draw[load] (0.1,4) -- (1.1,4) node[right] {$u_2$}; \draw[->,thick] (0.5,4.4) arc (30:170:0.5); \node[left] at (-0.4,4.5) {$\theta_2$};
  \draw[load] (0.1,0) -- (1.1,0) node[right] {$u_1$}; \draw[->,thick] (0.5,0.4) arc (30:170:0.5); \node[left] at (-0.4,0.5) {$\theta_1$};
\end{tikzpicture}
 &
      $[k] = \begin{array}{c}\begin{array}{cccc}\ \ u_1\ \ & \ \ \theta_1\ \ & \ \ u_2\ \ & \ \ \theta_2\ \ \end{array}\\
      \left[\begin{array}{cccc} 12EI/L^3 & -6EI/L^2 & -12EI/L^3 & -6EI/L^2\\ -6EI/L^2 & 4EI/L & 6EI/L^2 & 2EI/L\\ -12EI/L^3 & 6EI/L^2 & 12EI/L^3 & 6EI/L^2\\ -6EI/L^2 & 2EI/L & 6EI/L^2 & 4EI/L\end{array}\right]\end{array}$\\\hline
    C2 & \begin{tikzpicture}[scale=0.7,font=\small,baseline=(current bounding box.center)]
  \draw[thinmember] (0,0) -- (0,4); \node[node dot] at (0,4) {};
  \draw[load] (0.15,4) -- (1.1,4) node[right] {$u_2$};
  \draw[load] (0.1,0) -- (1.1,0) node[right] {$u_1$}; \draw[->,thick] (0.5,0.4) arc (30:170:0.5); \node[left] at (-0.4,0.5) {$\theta_1$};
\end{tikzpicture}
 &
      $[k] = \begin{array}{c}\begin{array}{ccc}\ \ u_1\ \ & \ \ \theta_1\ \ & \ \ u_2\ \ \end{array}\\
      \left[\begin{array}{ccc} 3EI/L^3 & -3EI/L^2 & -3EI/L^3\\ -3EI/L^2 & 3EI/L & 3EI/L^2\\ -3EI/L^3 & 3EI/L^2 & 3EI/L^3\end{array}\right]\end{array}$\\\hline
    C3 & \begin{tikzpicture}[scale=0.7,font=\small,baseline=(current bounding box.center)]
  \draw[thinmember] (0,0) -- (0,4); \node[node dot] at (0,0) {};
  \draw[load] (0.1,4) -- (1.1,4) node[right] {$u_2$}; \draw[->,thick] (0.5,4.4) arc (30:170:0.5); \node[left] at (-0.4,4.5) {$\theta_2$};
  \draw[load] (0.15,0) -- (1.1,0) node[right] {$u_1$};
\end{tikzpicture}
 &
      $[k] = \begin{array}{c}\begin{array}{ccc}\ \ u_1\ \ & \ \ u_2\ \ & \ \ \theta_2\ \ \end{array}\\
      \left[\begin{array}{ccc} 3EI/L^3 & -3EI/L^3 & -3EI/L^2\\ -3EI/L^3 & 3EI/L^3 & 3EI/L^2\\ -3EI/L^2 & 3EI/L^2 & 3EI/L\end{array}\right]\end{array}$\\\hline
  \end{tabular}
\end{center}

\begin{example}
\leavevmode
\begin{fig}
  %% Frame decomposition: whole frame = column 1 + beam 2 + column 3
\begin{tikzpicture}[scale=0.62,font=\small]
  \begin{scope}
    \draw[member] (0,0) -- (0,4) -- (4,4) -- (4,0); \pic at (0,0) {fixed}; \pic at (4,0) {fixed};
    \node[circle,draw,fill=white,inner sep=1pt,font=\scriptsize] at (0,2) {1}; \node[circle,draw,fill=white,inner sep=1pt,font=\scriptsize] at (2,4) {2}; \node[circle,draw,fill=white,inner sep=1pt,font=\scriptsize] at (4,2) {3};
    \draw[load] (4.1,4) -- (5.3,4) node[right] {$u_1$};
    \draw[->,thick] (0.45,4.35) arc (30:170:0.45); \node[above left,inner sep=1pt] at (-0.1,4.4) {$\theta_1$};
    \draw[->,thick] (4.45,4.35) arc (30:170:0.45); \node[above left,inner sep=1pt] at (3.9,4.4) {$\theta_2$};
  \end{scope}
  \node at (6.6,2) {\large $=$};
  \begin{scope}[xshift=8cm]
    \draw[member] (0,0) -- (0,4); \pic at (0,0) {fixed}; \node[circle,draw,fill=white,inner sep=1pt,font=\scriptsize] at (0,2) {1};
    \draw[load] (0.1,4) -- (1.3,4) node[right] {$u_1$};
    \draw[->,thick] (0.45,4.35) arc (30:170:0.45); \node[above left,inner sep=1pt] at (-0.1,4.4) {$\theta_1$};
    \node[below] at (0,-0.9) {C1};
  \end{scope}
  \node at (10,2) {\large $+$};
  \begin{scope}[xshift=11.5cm]
    \draw[member] (0,4) -- (4,4); \node[circle,draw,fill=white,inner sep=1pt,font=\scriptsize] at (2,4) {2};
    \draw[->,thick] (0.45,4.35) arc (30:170:0.45); \node[above left,inner sep=1pt] at (-0.1,4.4) {$\theta_1$};
    \draw[->,thick] (4.45,4.35) arc (30:170:0.45); \node[above left,inner sep=1pt] at (3.9,4.4) {$\theta_2$};
    \node[below] at (2,-0.9) {B1};
  \end{scope}
  \node at (17,2) {\large $+$};
  \begin{scope}[xshift=18.5cm]
    \draw[member] (0,0) -- (0,4); \node[circle,draw,fill=white,inner sep=1pt,font=\scriptsize] at (0,2) {3};
    \draw[load] (0.1,4) -- (1.3,4) node[right] {$u_1$};
    \draw[->,thick] (0.45,4.35) arc (30:170:0.45); \node[above left,inner sep=1pt] at (-0.1,4.4) {$\theta_2$};
    \node[below] at (0,-0.9) {C1};
  \end{scope}
\end{tikzpicture}

\end{fig}
{\small
\begin{align*}
  [k] &= \begin{array}{c}\begin{array}{ccc}\ u_1\ & \ \theta_1\ & \ \theta_2\ \end{array}\\ \left[\begin{array}{ccc}12EI/L^3 & 6EI/L^2 & 0\\ 6EI/L^2 & 4EI/L & 0\\ 0 & 0 & 0\end{array}\right]\end{array}
  + \begin{array}{c}\begin{array}{ccc}\ u_1\ & \ \theta_1\ & \ \theta_2\ \end{array}\\ \left[\begin{array}{ccc}0 & 0 & 0\\ 0 & 4EI/L & 2EI/L\\ 0 & 2EI/L & 4EI/L\end{array}\right]\end{array}
  + \begin{array}{c}\begin{array}{ccc}\ u_1\ & \ \theta_1\ & \ \theta_2\ \end{array}\\ \left[\begin{array}{ccc}12EI/L^3 & 0 & 6EI/L^2\\ 0 & 0 & 0\\ 6EI/L^2 & 0 & 4EI/L\end{array}\right]\end{array}\\
  &= \begin{array}{c}\begin{array}{ccc}\ u_1\ & \ \theta_1\ & \ \theta_2\ \end{array}\\ \left[\begin{array}{ccc}24EI/L^3 & 6EI/L^2 & 6EI/L^2\\ 6EI/L^2 & 8EI/L & 2EI/L\\ 6EI/L^2 & 2EI/L & 8EI/L\end{array}\right]\end{array}\qquad\text{(same as that given in the previous section)}
\end{align*}}
\end{example}

\begin{example}[Treatment of the internal hinge is the same as that given in Sec.~\ref{sec:end-hinge}]
\leavevmode

\textbf{Case 1:}
\begin{fig}
  %% Frame decomposition: whole frame = column 1 + beam 2 + column 3
\begin{tikzpicture}[scale=0.62,font=\small]
  \begin{scope}
    \draw[member] (0,0) -- (0,4) -- (4,4) -- (4,0); \pic at (0,0) {fixed}; \pic at (4,0) {fixed};
    \node[circle,draw,fill=white,inner sep=1pt,font=\scriptsize] at (0,2) {1}; \node[circle,draw,fill=white,inner sep=1pt,font=\scriptsize] at (2,4) {2}; \node[circle,draw,fill=white,inner sep=1pt,font=\scriptsize] at (4,2) {3};
    \node[node dot] at (0,4) {};
    \draw[load] (4.1,4) -- (5.3,4) node[right] {$u_1$};
    \draw[->,thick] (0.45,4.35) arc (30:170:0.45); \node[above left,inner sep=1pt] at (-0.1,4.4) {$\theta_1$};
    \draw[->,thick] (4.45,4.35) arc (30:170:0.45); \node[above left,inner sep=1pt] at (3.9,4.4) {$\theta_2$};
  \end{scope}
  \node at (6.6,2) {\large $=$};
  \begin{scope}[xshift=8cm]
    \draw[member] (0,0) -- (0,4); \pic at (0,0) {fixed}; \node[circle,draw,fill=white,inner sep=1pt,font=\scriptsize] at (0,2) {1};
    \draw[load] (0.1,4) -- (1.3,4) node[right] {$u_1$};
    \draw[->,thick] (0.45,4.35) arc (30:170:0.45); \node[above left,inner sep=1pt] at (-0.1,4.4) {$\theta_1$};
    \node[below] at (0,-0.9) {C1};
  \end{scope}
  \node at (10,2) {\large $+$};
  \begin{scope}[xshift=11.5cm]
    \draw[member] (0,4) -- (4,4); \node[circle,draw,fill=white,inner sep=1pt,font=\scriptsize] at (2,4) {2};
    \node[node dot] at (0,4) {};
    \draw[->,thick] (4.45,4.35) arc (30:170:0.45); \node[above left,inner sep=1pt] at (3.9,4.4) {$\theta_2$};
    \node[below] at (2,-0.9) {B2};
  \end{scope}
  \node at (17,2) {\large $+$};
  \begin{scope}[xshift=18.5cm]
    \draw[member] (0,0) -- (0,4); \node[circle,draw,fill=white,inner sep=1pt,font=\scriptsize] at (0,2) {3};
    \draw[load] (0.1,4) -- (1.3,4) node[right] {$u_1$};
    \draw[->,thick] (0.45,4.35) arc (30:170:0.45); \node[above left,inner sep=1pt] at (-0.1,4.4) {$\theta_2$};
    \node[below] at (0,-0.9) {C1};
  \end{scope}
\end{tikzpicture}

\end{fig}
{\small
\begin{align*}
  [k] &= \begin{array}{c}\begin{array}{ccc}\ u_1\ & \ \theta_1\ & \ \theta_2\ \end{array}\\ \left[\begin{array}{ccc}12EI/L^3 & 6EI/L^2 & 0\\ 6EI/L^2 & 4EI/L & 0\\ 0 & 0 & 0\end{array}\right]\end{array}
  + \begin{array}{c}\begin{array}{ccc}u_1 & \theta_1 & \ \theta_2\ \end{array}\\ \left[\begin{array}{ccc}0 & 0 & 0\\ 0 & 0 & 0\\ 0 & 0 & 3EI/L\end{array}\right]\end{array}
  + \begin{array}{c}\begin{array}{ccc}\ u_1\ & \ \theta_1\ & \ \theta_2\ \end{array}\\ \left[\begin{array}{ccc}12EI/L^3 & 0 & 6EI/L^2\\ 0 & 0 & 0\\ 6EI/L^2 & 0 & 4EI/L\end{array}\right]\end{array}\\
  &= \begin{array}{c}\begin{array}{ccc}\ u_1\ & \ \theta_1\ & \ \theta_2\ \end{array}\\ \left[\begin{array}{ccc}24EI/L^3 & 6EI/L^2 & 6EI/L^2\\ 6EI/L^2 & 4EI/L & 0\\ 6EI/L^2 & 0 & 7EI/L\end{array}\right]\end{array}
\end{align*}}
\textbf{Case 2:}
\begin{fig}
  %% Frame decomposition: whole frame = column 1 + beam 2 + column 3
\begin{tikzpicture}[scale=0.62,font=\small]
  \begin{scope}
    \draw[member] (0,0) -- (0,4) -- (4,4) -- (4,0); \pic at (0,0) {fixed}; \pic at (4,0) {fixed};
    \node[circle,draw,fill=white,inner sep=1pt,font=\scriptsize] at (0,2) {1}; \node[circle,draw,fill=white,inner sep=1pt,font=\scriptsize] at (2,4) {2}; \node[circle,draw,fill=white,inner sep=1pt,font=\scriptsize] at (4,2) {3};
    \node[node dot] at (0,4) {};
    \draw[load] (4.1,4) -- (5.3,4) node[right] {$u_1$};
    \draw[->,thick] (4.45,4.35) arc (30:170:0.45); \node[above left,inner sep=1pt] at (3.9,4.4) {$\theta_2$};
    \draw[->,thick] (0.45,4.35) arc (30:170:0.45); \node[above right,inner sep=1pt] at (0.1,4.4) {$\theta_1$};
  \end{scope}
  \node at (6.6,2) {\large $=$};
  \begin{scope}[xshift=8cm]
    \draw[member] (0,0) -- (0,4); \pic at (0,0) {fixed}; \node[circle,draw,fill=white,inner sep=1pt,font=\scriptsize] at (0,2) {1};
    \node[node dot] at (0,4) {};
    \draw[load] (0.1,4) -- (1.3,4) node[right] {$u_1$};
    \node[below] at (0,-0.9) {C2};
  \end{scope}
  \node at (10,2) {\large $+$};
  \begin{scope}[xshift=11.5cm]
    \draw[member] (0,4) -- (4,4); \node[circle,draw,fill=white,inner sep=1pt,font=\scriptsize] at (2,4) {2};
    \draw[->,thick] (0.45,4.35) arc (30:170:0.45); \node[above left,inner sep=1pt] at (-0.1,4.4) {$\theta_1$};
    \draw[->,thick] (4.45,4.35) arc (30:170:0.45); \node[above left,inner sep=1pt] at (3.9,4.4) {$\theta_2$};
    \node[below] at (2,-0.9) {B1};
  \end{scope}
  \node at (17,2) {\large $+$};
  \begin{scope}[xshift=18.5cm]
    \draw[member] (0,0) -- (0,4); \node[circle,draw,fill=white,inner sep=1pt,font=\scriptsize] at (0,2) {3};
    \draw[load] (0.1,4) -- (1.3,4) node[right] {$u_1$};
    \draw[->,thick] (0.45,4.35) arc (30:170:0.45); \node[above left,inner sep=1pt] at (-0.1,4.4) {$\theta_2$};
    \node[below] at (0,-0.9) {C1};
  \end{scope}
\end{tikzpicture}

\end{fig}
{\small
\begin{align*}
  [k] &= \begin{array}{c}\begin{array}{ccc}\ u_1\ & \theta_1 & \theta_2\end{array}\\ \left[\begin{array}{ccc}3EI/L^3 & 0 & 0\\ 0 & 0 & 0\\ 0 & 0 & 0\end{array}\right]\end{array}
  + \begin{array}{c}\begin{array}{ccc}\ u_1\ & \ \theta_1\ & \ \theta_2\ \end{array}\\ \left[\begin{array}{ccc}0 & 0 & 0\\ 0 & 4EI/L & 2EI/L\\ 0 & 2EI/L & 4EI/L\end{array}\right]\end{array}
  + \begin{array}{c}\begin{array}{ccc}\ u_1\ & \ \theta_1\ & \ \theta_2\ \end{array}\\ \left[\begin{array}{ccc}12EI/L^3 & 0 & 6EI/L^2\\ 0 & 0 & 0\\ 6EI/L^2 & 0 & 4EI/L\end{array}\right]\end{array}\\
  &= \begin{array}{c}\begin{array}{ccc}\ u_1\ & \ \theta_1\ & \ \theta_2\ \end{array}\\ \left[\begin{array}{ccc}15EI/L^3 & 0 & 6EI/L^2\\ 0 & 4EI/L & 2EI/L\\ 6EI/L^2 & 2EI/L & 8EI/L\end{array}\right]\end{array}
\end{align*}}
\textbf{Case 3:}
\begin{fig}
  %% Frame decomposition: whole frame = column 1 + beam 2 + column 3
\begin{tikzpicture}[scale=0.62,font=\small]
  \begin{scope}
    \draw[member] (0,0) -- (0,4) -- (4,4) -- (4,0); \pic at (0,0) {fixed}; \pic at (4,0) {fixed};
    \node[circle,draw,fill=white,inner sep=1pt,font=\scriptsize] at (0,2) {1}; \node[circle,draw,fill=white,inner sep=1pt,font=\scriptsize] at (2,4) {2}; \node[circle,draw,fill=white,inner sep=1pt,font=\scriptsize] at (4,2) {3};
    \node[node dot] at (0,4) {};
    \draw[load] (4.1,4) -- (5.3,4) node[right] {$u_1$};
    \draw[->,thick] (4.45,4.35) arc (30:170:0.45); \node[above left,inner sep=1pt] at (3.9,4.4) {$\theta_2$};
  \end{scope}
  \node at (6.6,2) {\large $=$};
  \begin{scope}[xshift=8cm]
    \draw[member] (0,0) -- (0,4); \pic at (0,0) {fixed}; \node[circle,draw,fill=white,inner sep=1pt,font=\scriptsize] at (0,2) {1};
    \node[node dot] at (0,4) {};
    \draw[load] (0.1,4) -- (1.3,4) node[right] {$u_1$};
    \node[below] at (0,-0.9) {C2};
  \end{scope}
  \node at (10,2) {\large $+$};
  \begin{scope}[xshift=11.5cm]
    \draw[member] (0,4) -- (4,4); \node[circle,draw,fill=white,inner sep=1pt,font=\scriptsize] at (2,4) {2};
    \node[node dot] at (0,4) {};
    \draw[->,thick] (4.45,4.35) arc (30:170:0.45); \node[above left,inner sep=1pt] at (3.9,4.4) {$\theta_2$};
    \node[below] at (2,-0.9) {B2};
  \end{scope}
  \node at (17,2) {\large $+$};
  \begin{scope}[xshift=18.5cm]
    \draw[member] (0,0) -- (0,4); \node[circle,draw,fill=white,inner sep=1pt,font=\scriptsize] at (0,2) {3};
    \draw[load] (0.1,4) -- (1.3,4) node[right] {$u_1$};
    \draw[->,thick] (0.45,4.35) arc (30:170:0.45); \node[above left,inner sep=1pt] at (-0.1,4.4) {$\theta_2$};
    \node[below] at (0,-0.9) {C1};
  \end{scope}
\end{tikzpicture}

\end{fig}
{\small
\begin{align*}
  [k] &= \begin{array}{c}\begin{array}{cc}\ u_1\ & \theta_2\end{array}\\ \left[\begin{array}{cc}3EI/L^3 & 0\\ 0 & 0\end{array}\right]\end{array}
  + \begin{array}{c}\begin{array}{cc}u_1 & \ \theta_2\ \end{array}\\ \left[\begin{array}{cc}0 & 0\\ 0 & 3EI/L\end{array}\right]\end{array}
  + \begin{array}{c}\begin{array}{cc}\ u_1\ & \ \theta_2\ \end{array}\\ \left[\begin{array}{cc}12EI/L^3 & 6EI/L^2\\ 6EI/L^2 & 4EI/L\end{array}\right]\end{array}
  = \begin{array}{c}\begin{array}{cc}\ u_1\ & \ \theta_2\ \end{array}\\ \left[\begin{array}{cc}15EI/L^3 & 6EI/L^2\\ 6EI/L^2 & 7EI/L\end{array}\right]\end{array}
\end{align*}}
\textbf{Case 4:}
\begin{fig}
  %% Frame decomposition: whole frame = column 1 + beam 2 + column 3
\begin{tikzpicture}[scale=0.62,font=\small]
  \begin{scope}
    \draw[member] (0,0) -- (0,4) -- (4,4) -- (4,0); \pic at (0,0) {fixed}; \pic at (4,0) {fixed};
    \node[circle,draw,fill=white,inner sep=1pt,font=\scriptsize] at (0,2) {1}; \node[circle,draw,fill=white,inner sep=1pt,font=\scriptsize] at (2,4) {2}; \node[circle,draw,fill=white,inner sep=1pt,font=\scriptsize] at (4,2) {3};
    \node[node dot] at (0,4) {};
    \draw[load] (4.1,4) -- (5.3,4) node[right] {$u_1$};
    \draw[->,thick] (0.45,4.35) arc (30:170:0.45); \node[above left,inner sep=1pt] at (-0.1,4.4) {$\theta_{1l}$};
    \draw[->,thick] (4.45,4.35) arc (30:170:0.45); \node[above left,inner sep=1pt] at (3.9,4.4) {$\theta_2$};
    \draw[->,thick] (0.45,4.35) arc (30:170:0.45); \node[above right,inner sep=1pt] at (0.1,4.4) {$\theta_{1r}$};
  \end{scope}
  \node at (6.6,2) {\large $=$};
  \begin{scope}[xshift=8cm]
    \draw[member] (0,0) -- (0,4); \pic at (0,0) {fixed}; \node[circle,draw,fill=white,inner sep=1pt,font=\scriptsize] at (0,2) {1};
    \draw[load] (0.1,4) -- (1.3,4) node[right] {$u_1$};
    \draw[->,thick] (0.45,4.35) arc (30:170:0.45); \node[above left,inner sep=1pt] at (-0.1,4.4) {$\theta_{1l}$};
    \node[below] at (0,-0.9) {C1};
  \end{scope}
  \node at (10,2) {\large $+$};
  \begin{scope}[xshift=11.5cm]
    \draw[member] (0,4) -- (4,4); \node[circle,draw,fill=white,inner sep=1pt,font=\scriptsize] at (2,4) {2};
    \node[node dot] at (0,4) {};
    \draw[->,thick] (0.45,4.35) arc (30:170:0.45); \node[above left,inner sep=1pt] at (-0.1,4.4) {$\theta_{1r}$};
    \draw[->,thick] (4.45,4.35) arc (30:170:0.45); \node[above left,inner sep=1pt] at (3.9,4.4) {$\theta_2$};
    \node[below] at (2,-0.9) {B1};
  \end{scope}
  \node at (17,2) {\large $+$};
  \begin{scope}[xshift=18.5cm]
    \draw[member] (0,0) -- (0,4); \node[circle,draw,fill=white,inner sep=1pt,font=\scriptsize] at (0,2) {3};
    \draw[load] (0.1,4) -- (1.3,4) node[right] {$u_1$};
    \draw[->,thick] (0.45,4.35) arc (30:170:0.45); \node[above left,inner sep=1pt] at (-0.1,4.4) {$\theta_2$};
    \node[below] at (0,-0.9) {C1};
  \end{scope}
\end{tikzpicture}

\end{fig}
{\footnotesize\setlength{\arraycolsep}{3pt}
\begin{align*}
  [k] &= \begin{array}{c}\begin{array}{cccc}\ u_1\ & \ \theta_{1l}\ & \theta_{1r} & \theta_2\end{array}\\ \left[\begin{array}{cccc}12EI/L^3 & 6EI/L^2 & 0 & 0\\ 6EI/L^2 & 4EI/L & 0 & 0\\ 0 & 0 & 0 & 0\\ 0 & 0 & 0 & 0\end{array}\right]\end{array}
  + \begin{array}{c}\begin{array}{cccc}u_1 & \theta_{1l} & \ \theta_{1r}\ & \ \theta_2\ \end{array}\\ \left[\begin{array}{cccc}0 & 0 & 0 & 0\\ 0 & 0 & 0 & 0\\ 0 & 0 & 4EI/L & 2EI/L\\ 0 & 0 & 2EI/L & 4EI/L\end{array}\right]\end{array}
  + \begin{array}{c}\begin{array}{cccc}\ u_1\ & \theta_{1l} & \theta_{1r} & \ \theta_2\ \end{array}\\ \left[\begin{array}{cccc}12EI/L^3 & 0 & 0 & 6EI/L^2\\ 0 & 0 & 0 & 0\\ 0 & 0 & 0 & 0\\ 6EI/L^2 & 0 & 0 & 4EI/L\end{array}\right]\end{array}\\
  &= \begin{array}{c}\begin{array}{cccc}\ u_1\ & \ \theta_{1l}\ & \ \theta_{1r}\ & \ \theta_2\ \end{array}\\ \left[\begin{array}{cccc}24EI/L^3 & 6EI/L^2 & 0 & 6EI/L^2\\ 6EI/L^2 & 4EI/L & 0 & 0\\ 0 & 0 & 4EI/L & 2EI/L\\ 6EI/L^2 & 0 & 2EI/L & 8EI/L\end{array}\right]\end{array}
\end{align*}}
Discussion: To reduce the number of DOFs, condensed elements are usually used for both beams and columns, i.e.\ case 3.
\end{example}

\begin{example}[Note that the column is continuous at node B]
\leavevmode

\begin{minipage}{0.55\linewidth}\centering% Example 3: two-storey frame with hinges at A and B (column continuous at B)
\begin{tikzpicture}[scale=0.85,font=\small]
  \draw[member] (0,0) -- (0,5) -- (3.5,5) -- (3.5,0); \draw[member] (0,2.5) -- (6,2.5) -- (6,0);
  \pic at (0,0) {pin}; \pic at (3.5,0) {fixed}; \pic at (6,0) {fixed};
  \node[node dot] at (3.5,5) {}; \node[node dot] at (3.5,2.5) {};
  \node[above] at (3.5,5.1) {$A$}; \node[above right] at (3.55,2.55) {$B$};
  \newcommand{\mem}[2]{\node[circle,draw,fill=white,inner sep=1pt,font=\scriptsize] at #1 {#2};}
  \mem{(0,1.2)}{1} \mem{(3.5,1.2)}{2} \mem{(6,1.2)}{3} \mem{(0,3.8)}{4} \mem{(3.5,3.8)}{5}
  \mem{(1.75,2.5)}{6} \mem{(4.75,2.5)}{7} \mem{(1.75,5)}{8}
  \draw[load] (3.6,5) -- (4.8,5) node[right] {$u_2$};
  \draw[load] (6.1,2.5) -- (7.3,2.5) node[right] {$u_1$};
  \draw[->,thick] (0.5,5.35) arc (30:170:0.5); \node[left] at (-0.4,5.5) {$\theta_4$};
  \draw[->,thick] (0.5,2.85) arc (30:170:0.5); \node[left] at (-0.4,3.0) {$\theta_1$};
  \draw[->,thick] (3.3,2.85) arc (30:170:0.5); \node[left] at (2.4,3.0) {$\theta_2$};
  \draw[->,thick] (6.5,2.85) arc (30:170:0.5); \node[left] at (5.6,3.0) {$\theta_3$};
\end{tikzpicture}
\end{minipage}\hfill
\begin{minipage}{0.42\linewidth}
  \begin{tabular}{lll}
    \textbf{Element} & \textbf{Type} & \textbf{DOFs}\\
    1 & C3 & $u_1, \theta_1$\\
    2 & C1 & $u_1, \theta_2$\\
    3 & C1 & $u_1, \theta_3$\\
    4 & C1 & $u_1, \theta_1, u_2, \theta_4$\\
    5 & C2 & $u_1, \theta_2, u_2$\\
    6 & B1 & $\theta_1, \theta_2$\\
    7 & B2 & $\theta_3$\\
    8 & B3 & $\theta_4$
  \end{tabular}
\end{minipage}
\end{example}

\section{Equivalent forces for uniformly distributed loads and centre point load}
\begin{figure}[H]
  \centering
  % Equivalent forces for uniformly distributed loads and centre point load (B1-B3, C1-C3)
\begin{tikzpicture}[scale=0.72,font=\scriptsize]
  \newcommand{\udl}{\foreach \x in {0.3,0.8,...,4.8}{\draw[->] (\x,1.0) -- (\x,0.15);} \draw (0.3,1.0) -- (4.8,1.0);}
  \newcommand{\hudl}{\foreach \y in {0.3,0.8,...,3.8}{\draw[->] (-1.0,\y) -- (-0.15,\y);} \draw (-1.0,0.3) -- (-1.0,3.8);}
  % ================= beams =================
  % B1
  \begin{scope}
    \node[left] at (-0.5,0) {B1};
    \draw[thinmember] (0,0) -- (5,0); \udl
    \draw[->,thick] (0.6,-0.1) arc (30:170:0.5); \node[left] at (-0.2,-0.2) {$\frac{wL^2}{12}$};
    \draw[->,thick] (4.4,-0.4) arc (200:20:0.5); \node[right] at (5.2,-0.2) {$\frac{wL^2}{12}$};
  \end{scope}
  \begin{scope}[xshift=9cm]
    \draw[thinmember] (0,0) -- (5,0); \draw[load] (2.5,1.2) -- (2.5,0.15) node[pos=0,above] {$P$};
    \draw[->,thick] (0.6,-0.1) arc (30:170:0.5); \node[left] at (-0.2,-0.2) {$\frac{PL}{8}$};
    \draw[->,thick] (4.4,-0.4) arc (200:20:0.5); \node[right] at (5.2,-0.2) {$\frac{PL}{8}$};
  \end{scope}
  % B2
  \begin{scope}[yshift=-2.8cm]
    \node[left] at (-0.5,0) {B2};
    \draw[thinmember] (0,0) -- (5,0); \node[node dot] at (0,0) {}; \udl
    \draw[->,thick] (4.4,-0.4) arc (200:20:0.5); \node[right] at (5.2,-0.2) {$\frac{wL^2}{8}$};
  \end{scope}
  \begin{scope}[xshift=9cm,yshift=-2.8cm]
    \draw[thinmember] (0,0) -- (5,0); \node[node dot] at (0,0) {}; \draw[load] (2.5,1.2) -- (2.5,0.15) node[pos=0,above] {$P$};
    \draw[->,thick] (4.4,-0.4) arc (200:20:0.5); \node[right] at (5.2,-0.2) {$\frac{3PL}{16}$};
  \end{scope}
  % B3
  \begin{scope}[yshift=-5.6cm]
    \node[left] at (-0.5,0) {B3};
    \draw[thinmember] (0,0) -- (5,0); \node[node dot] at (5,0) {}; \udl
    \draw[->,thick] (0.6,-0.1) arc (30:170:0.5); \node[left] at (-0.2,-0.2) {$\frac{wL^2}{8}$};
  \end{scope}
  \begin{scope}[xshift=9cm,yshift=-5.6cm]
    \draw[thinmember] (0,0) -- (5,0); \node[node dot] at (5,0) {}; \draw[load] (2.5,1.2) -- (2.5,0.15) node[pos=0,above] {$P$};
    \draw[->,thick] (0.6,-0.1) arc (30:170:0.5); \node[left] at (-0.2,-0.2) {$\frac{3PL}{16}$};
  \end{scope}
  % ================= columns =================
  % C1
  \begin{scope}[yshift=-12.5cm]
    \node[left] at (-1.6,2) {C1};
    \draw[thinmember] (0,0) -- (0,4); \hudl
    \draw[load] (0.1,4) -- (1.0,4) node[right] {$\frac{wL}{2}$}; \draw[->,thick] (0.5,4.35) arc (30:170:0.45); \node[above right] at (0.5,4.4) {$\frac{wL^2}{12}$};
    \draw[load] (0.1,0) -- (1.0,0) node[right] {$\frac{wL}{2}$}; \draw[->,thick] (0.5,0.35) arc (30:170:0.45); \node[above right] at (0.5,0.4) {$\frac{wL^2}{12}$};
  \end{scope}
  \begin{scope}[xshift=9cm,yshift=-12.5cm]
    \draw[thinmember] (0,0) -- (0,4); \draw[load] (-1.2,2) -- (-0.15,2) node[pos=0,above] {$P$};
    \draw[load] (0.1,4) -- (1.0,4) node[right] {$\frac P2$}; \draw[->,thick] (0.5,4.35) arc (30:170:0.45); \node[above right] at (0.5,4.4) {$\frac{PL}{8}$};
    \draw[load] (0.1,0) -- (1.0,0) node[right] {$\frac P2$}; \draw[->,thick] (0.5,0.35) arc (30:170:0.45); \node[above right] at (0.5,0.4) {$\frac{PL}{8}$};
  \end{scope}
  % C2 (hinge at top)
  \begin{scope}[yshift=-18.5cm]
    \node[left] at (-1.6,2) {C2};
    \draw[thinmember] (0,0) -- (0,4); \node[node dot] at (0,4) {}; \hudl
    \draw[load] (0.15,4) -- (1.0,4) node[right] {$\frac38wL$};
    \draw[load] (0.1,0) -- (1.0,0) node[right] {$\frac58wL$}; \draw[->,thick] (0.5,0.35) arc (30:170:0.45); \node[above right] at (0.5,0.4) {$\frac{wL^2}{8}$};
  \end{scope}
  \begin{scope}[xshift=9cm,yshift=-18.5cm]
    \draw[thinmember] (0,0) -- (0,4); \node[node dot] at (0,4) {}; \draw[load] (-1.2,2) -- (-0.15,2) node[pos=0,above] {$P$};
    \draw[load] (0.15,4) -- (1.0,4) node[right] {$\frac{5}{16}P$};
    \draw[load] (0.1,0) -- (1.0,0) node[right] {$\frac{11}{16}P$}; \draw[->,thick] (0.5,0.35) arc (30:170:0.45); \node[above right] at (0.5,0.4) {$\frac{3}{16}PL$};
  \end{scope}
  % C3 (hinge at bottom)
  \begin{scope}[yshift=-24.5cm]
    \node[left] at (-1.6,2) {C3};
    \draw[thinmember] (0,0) -- (0,4); \node[node dot] at (0,0) {}; \hudl
    \draw[load] (0.1,4) -- (1.0,4) node[right] {$\frac58wL$}; \draw[->,thick] (0.5,4.35) arc (30:170:0.45); \node[above right] at (0.5,4.4) {$\frac{wL^2}{8}$};
    \draw[load] (0.15,0) -- (1.0,0) node[right] {$\frac38wL$};
  \end{scope}
  \begin{scope}[xshift=9cm,yshift=-24.5cm]
    \draw[thinmember] (0,0) -- (0,4); \node[node dot] at (0,0) {}; \draw[load] (-1.2,2) -- (-0.15,2) node[pos=0,above] {$P$};
    \draw[load] (0.1,4) -- (1.0,4) node[right] {$\frac{11}{16}P$}; \draw[->,thick] (0.5,4.35) arc (30:170:0.45); \node[above right] at (0.5,4.4) {$\frac{3}{16}PL$};
    \draw[load] (0.15,0) -- (1.0,0) node[right] {$\frac{5}{16}P$};
  \end{scope}
\end{tikzpicture}

\end{figure}

\begin{example}[Frame with a hinge at the top-left corner under three loadings]
\leavevmode
\begin{fig}
  % Example 2: the three loading cases on the frame with a hinge at the top-left corner
\begin{tikzpicture}[scale=0.6,font=\small]
  \newcommand{\hframe}{\draw[member] (0,0) -- (0,4) -- (4,4) -- (4,0); \pic at (0,0) {fixed}; \pic at (4,0) {fixed}; \node[node dot] at (0,4) {};}
  \begin{scope} \hframe \draw[load] (-1.4,4) -- (-0.15,4) node[pos=0,below] {$P$}; \end{scope}
  \begin{scope}[xshift=7cm] \hframe \draw[load] (-1.4,2) -- (-0.1,2) node[pos=0,below] {$P$}; \end{scope}
  \begin{scope}[xshift=14cm] \hframe \foreach \x in {0.3,0.8,...,3.8}{\draw[->] (\x,5.0) -- (\x,4.15);} \draw (0.3,5.0) -- (3.8,5.0); \node[above] at (2,5.0) {$w$}; \end{scope}
\end{tikzpicture}

\end{fig}
{\scriptsize
\renewcommand{\arraystretch}{1.0}
\begin{longtable}{|c|c|c|c|}\hline
  \textbf{Case} & Load 1 ($P$ at the corner) & Load 2 ($P$ at mid-height) & Load 3 ($w$ on the beam)\\\hline
  \shortstack{Case 1:\\[2pt] \begin{tikzpicture}[scale=0.5,font=\scriptsize,baseline=(current bounding box.center)]
  \draw[member] (0,0) -- (0,4) -- (4,4) -- (4,0); \pic at (0,0) {fixed}; \pic at (4,0) {fixed};
  \node[node dot] at (0,4) {};
  \draw[load] (4.1,4) -- (5.3,4) node[right] {$u_1$};
  \draw[->,thick] (4.45,4.35) arc (30:170:0.45); \node[above left,inner sep=1pt] at (3.9,4.4) {$\theta_2$};
  \draw[->,thick] (0.45,4.35) arc (30:170:0.45); \node[above left,inner sep=1pt] at (-0.1,4.4) {$\theta_1$};
\end{tikzpicture}
} &
  $\begin{aligned}\Pvec &= \Bmat{P\\ 0\\ 0}\\ \Bmat{u_1\\ \theta_1(\theta_{1l})\\ \theta_2} &= \Bmat{\frac{7}{69}\frac{PL^3}{EI}\\[2pt] -\frac{7}{46}\frac{PL^2}{EI}\\[2pt] -\frac{2}{23}\frac{PL^2}{EI}}\end{aligned}$ &
  $\begin{aligned}\Pvec &= \Bmat{\frac P2\\[2pt] \frac{PL}{8}\\[2pt] 0}\\ \Bmat{u_1\\ \theta_1(\theta_{1l})\\ \theta_2} &= \Bmat{\frac{35}{1104}\frac{PL^3}{EI}\\[2pt] -\frac{3}{184}\frac{PL^2}{EI}\\[2pt] -\frac{5}{184}\frac{PL^2}{EI}}\end{aligned}$ &
  $\begin{aligned}\Pvec &= \Bmat{0\\ 0\\ \frac{wL^2}{8}}\\ \Bmat{u_1\\ \theta_1(\theta_{1l})\\ \theta_2} &= \Bmat{-\frac{1}{92}\frac{wL^4}{EI}\\[2pt] \frac{3}{184}\frac{wL^3}{EI}\\[2pt] \frac{5}{184}\frac{wL^3}{EI}}\end{aligned}$\\\hline
  \shortstack{Case 2:\\[2pt] \begin{tikzpicture}[scale=0.5,font=\scriptsize,baseline=(current bounding box.center)]
  \draw[member] (0,0) -- (0,4) -- (4,4) -- (4,0); \pic at (0,0) {fixed}; \pic at (4,0) {fixed};
  \node[node dot] at (0,4) {};
  \draw[load] (4.1,4) -- (5.3,4) node[right] {$u_1$};
  \draw[->,thick] (4.45,4.35) arc (30:170:0.45); \node[above left,inner sep=1pt] at (3.9,4.4) {$\theta_2$};
  \draw[->,thick] (0.9,4.35) arc (30:170:0.45); \node[above right,inner sep=1pt] at (0.6,4.4) {$\theta_1$};
\end{tikzpicture}
} &
  $\begin{aligned}\Pvec &= \Bmat{P\\ 0\\ 0}\\ \Bmat{u_1\\ \theta_1(\theta_{1r})\\ \theta_2} &= \Bmat{\frac{7}{69}\frac{PL^3}{EI}\\[2pt] \frac{1}{23}\frac{PL^2}{EI}\\[2pt] -\frac{2}{23}\frac{PL^2}{EI}}\end{aligned}$ &
  $\begin{aligned}\Pvec &= \Bmat{\frac{5P}{16}\\[2pt] 0\\ 0}\\ \Bmat{u_1\\ \theta_1(\theta_{1r})\\ \theta_2} &= \Bmat{\frac{35}{1104}\frac{PL^3}{EI}\\[2pt] \frac{5}{368}\frac{PL^2}{EI}\\[2pt] -\frac{5}{184}\frac{PL^2}{EI}}\end{aligned}$ &
  $\begin{aligned}\Pvec &= \Bmat{0\\ -\frac{wL^2}{12}\\[2pt] \frac{wL^2}{12}}\\ \Bmat{u_1\\ \theta_1(\theta_{1r})\\ \theta_2} &= \Bmat{-\frac{1}{92}\frac{wL^4}{EI}\\[2pt] -\frac{19}{552}\frac{wL^3}{EI}\\[2pt] \frac{5}{184}\frac{wL^3}{EI}}\end{aligned}$\\\hline
  \shortstack{Case 3:\\[2pt] \begin{tikzpicture}[scale=0.5,font=\scriptsize,baseline=(current bounding box.center)]
  \draw[member] (0,0) -- (0,4) -- (4,4) -- (4,0); \pic at (0,0) {fixed}; \pic at (4,0) {fixed};
  \node[node dot] at (0,4) {};
  \draw[load] (4.1,4) -- (5.3,4) node[right] {$u_1$};
  \draw[->,thick] (4.45,4.35) arc (30:170:0.45); \node[above left,inner sep=1pt] at (3.9,4.4) {$\theta_2$};
\end{tikzpicture}
} &
  $\begin{aligned}\Pvec &= \Bmat{P\\ 0}\\ \Bmat{u_1\\ \theta_2} &= \Bmat{\frac{7}{69}\frac{PL^3}{EI}\\[2pt] -\frac{2}{23}\frac{PL^2}{EI}}\end{aligned}$ &
  $\begin{aligned}\Pvec &= \Bmat{\frac{5P}{16}\\[2pt] 0}\\ \Bmat{u_1\\ \theta_2} &= \Bmat{\frac{35}{1104}\frac{PL^3}{EI}\\[2pt] -\frac{5}{184}\frac{PL^2}{EI}}\end{aligned}$ &
  $\begin{aligned}\Pvec &= \Bmat{0\\ \frac{wL^2}{8}}\\ \Bmat{u_1\\ \theta_2} &= \Bmat{-\frac{1}{92}\frac{wL^4}{EI}\\[2pt] \frac{5}{184}\frac{wL^3}{EI}}\end{aligned}$\\\hline
  \shortstack{Case 4:\\[2pt] \begin{tikzpicture}[scale=0.5,font=\scriptsize,baseline=(current bounding box.center)]
  \draw[member] (0,0) -- (0,4) -- (4,4) -- (4,0); \pic at (0,0) {fixed}; \pic at (4,0) {fixed};
  \node[node dot] at (0,4) {};
  \draw[load] (4.1,4) -- (5.3,4) node[right] {$u_1$};
  \draw[->,thick] (4.45,4.35) arc (30:170:0.45); \node[above left,inner sep=1pt] at (3.9,4.4) {$\theta_2$};
  \draw[->,thick] (0.45,4.35) arc (30:170:0.45); \node[above left,inner sep=1pt] at (-0.1,4.4) {$\theta_{1l}$};
  \draw[->,thick] (0.9,4.35) arc (30:170:0.45); \node[above right,inner sep=1pt] at (0.6,4.4) {$\theta_{1r}$};
\end{tikzpicture}
} &
  $\begin{aligned}\Pvec &= \Bmat{P\\ 0\\ 0\\ 0}\\ \Bmat{u_1\\ \theta_{1l}\\ \theta_{1r}\\ \theta_2} &= \Bmat{\frac{7}{69}\frac{PL^3}{EI}\\[2pt] -\frac{7}{46}\frac{PL^2}{EI}\\[2pt] \frac{1}{23}\frac{PL^2}{EI}\\[2pt] -\frac{2}{23}\frac{PL^2}{EI}}\end{aligned}$ &
  $\begin{aligned}\Pvec &= \Bmat{\frac P2\\[2pt] \frac{PL}{8}\\[2pt] 0\\ 0}\\ \Bmat{u_1\\ \theta_{1l}\\ \theta_{1r}\\ \theta_2} &= \Bmat{\frac{35}{1104}\frac{PL^3}{EI}\\[2pt] -\frac{3}{184}\frac{PL^2}{EI}\\[2pt] \frac{5}{368}\frac{PL^2}{EI}\\[2pt] -\frac{5}{184}\frac{PL^2}{EI}}\end{aligned}$ &
  $\begin{aligned}\Pvec &= \Bmat{0\\ 0\\ -\frac{wL^2}{12}\\[2pt] \frac{wL^2}{12}}\\ \Bmat{u_1\\ \theta_{1l}\\ \theta_{1r}\\ \theta_2} &= \Bmat{-\frac{1}{92}\frac{wL^4}{EI}\\[2pt] \frac{3}{184}\frac{wL^3}{EI}\\[2pt] -\frac{19}{552}\frac{wL^3}{EI}\\[2pt] \frac{5}{184}\frac{wL^3}{EI}}\end{aligned}$\\\hline
\end{longtable}}
\end{example}

\chapter{Element Flexibility Matrix}
\label{ch:flex}

\section{Element flexibility matrix}
To define the flexibility coefficients for an element (member), the member shall be well supported, i.e.\ stable. Let the flexibility equations of a member be written as
\begin{equation}
  \mat{r} = \mat{d}\F_f
\end{equation}
where $\mat{r}$ indicates the (relative) deformation vector associated with the independent force vector $\F_f$ and $\mat{d}$ is the flexibility matrix. We emphasize that the components of $\F_f$ are independent. According to the definition, the flexibility coefficient $d_{ij}$ is the deformation $r_i$ due to $F_{fj} = 1$ and the other components of $\F_f$ equal zero.

\topic{Truss element:}
\begin{equation}
  d = L/EA,\quad\text{i.e.}\quad \delta = F_x\frac{L}{EA}
\end{equation}

\topic{Beam element: three different supported elements can be distinguished}
\begin{enumerate}[label=(\alph*),leftmargin=*]
  \item Simply supported at both ends
    \begin{fig}
      % Simply supported beam: end moments and deformed shape due to M1
\begin{tikzpicture}[scale=0.85,font=\small]
  \begin{scope}
    \draw[member] (0,0) -- (6,0); \pic at (0,0) {pin}; \pic at (6,0) {roller};
    \draw[->,thick] (-0.35,0.7) arc (110:-40:0.5); \node[above left] at (-0.2,0.75) {$M_1(\bar\theta_1)$};
    \draw[->,thick] (6.35,0.7) arc (70:220:0.5); \node[above right] at (6.2,0.75) {$M_2(\bar\theta_2)$};
  \end{scope}
  \begin{scope}[xshift=8.5cm]
    \pic at (0,0) {pin}; \pic at (6,0) {roller};
    \draw[dashed] (0,0) -- (6,0);
    \draw[thinmember] (0,0) .. controls (2,1.4) and (4.5,0.9) .. (6,0);
    \draw[thin] (0,0) -- (1.6,1.1); \draw[thin] (6,0) -- (4.6,0.55);
    \draw[->] (0.9,0) arc (0:34:0.9); \node[below right] at (0.9,0.15) {$\bar\theta_1 = \dfrac{M_1L}{3EI}$};
    \draw[->] (5.1,0) arc (180:158:0.9); \node[above left] at (4.9,0.2) {$\bar\theta_2 = \dfrac{M_1L}{6EI}$};
    \draw[->,thick] (-0.35,0.7) arc (110:-40:0.5); \node[above left] at (-0.2,0.75) {$M_1$};
  \end{scope}
\end{tikzpicture}

    \end{fig}
    \begin{equation}
      \Bmat{\bar\theta_1\\ \bar\theta_2} = \frac{L}{6EI}\bmat{2 & -1\\ -1 & 2}\Bmat{M_1\\ M_2},\quad\text{i.e.}\quad \mat{d}_a = \frac{L}{6EI}\bmat{2 & -1\\ -1 & 2}
    \end{equation}
  \item Cantilever (fixed at $x = 0$ and free at $x = L$)
    \begin{fig}
      % Cantilever: F2 and M2 at the free end
\begin{tikzpicture}[scale=0.85,font=\small]
  \pic at (0,0) {wallL}; \draw[member] (0,0) -- (6,0);
  \draw[load] (6,0.15) -- (6,1.3) node[above] {$F_2(\bar v_2)$};
  \draw[->,thick] (6.5,-0.35) arc (-45:170:0.4); \node[right] at (6.4,-0.2) {$M_2(\bar\theta_2)$};
\end{tikzpicture}

    \end{fig}
    \begin{equation}
      \Bmat{\bar v_2\\ \bar\theta_2} = \frac{L}{6EI}\bmat{2L^2 & 3L\\ 3L & 6}\Bmat{F_2\\ M_2},\quad\text{i.e.}\quad \mat{d}_b = \frac{L}{6EI}\bmat{2L^2 & 3L\\ 3L & 6}
    \end{equation}
  \item Slide at $x = 0$ and simple support at $x = L$
    \begin{fig}
      % Slide at x = 0 and simple support at x = L: F1 and M2
\begin{tikzpicture}[scale=0.85,font=\small]
  % slide support at the left: wall with rollers
  \fill[hatch] (-0.6,-0.6) rectangle (-0.4,0.6); \draw (-0.4,-0.6) -- (-0.4,0.6);
  \draw[fill=white] (-0.25,0.25) circle (0.1); \draw[fill=white] (-0.25,-0.25) circle (0.1);
  \draw (-0.12,-0.5) -- (-0.12,0.5);
  \draw[member] (-0.1,0) -- (6,0); \pic at (6,0) {pin};
  \draw[load] (0,0.15) -- (0,1.3) node[above] {$F_1(\bar v_1)$};
  \draw[->,thick] (6.5,0.5) arc (45:230:0.4); \node[above right] at (6.3,0.5) {$M_2(\bar\theta_2)$};
\end{tikzpicture}

    \end{fig}
    \begin{equation}
      \Bmat{\bar v_1\\ \bar\theta_2} = \frac{L}{6EI}\bmat{2L^2 & -3L\\ -3L & 6}\Bmat{F_1\\ M_2},\quad\text{i.e.}\quad \mat{d}_c = \frac{L}{6EI}\bmat{2L^2 & -3L\\ -3L & 6}
    \end{equation}
\end{enumerate}

\section{Complementary strain energy}
The complementary strain energy of an element is defined as follows.

Truss element:
\begin{equation}
  U^* = \int_0^L\frac{F_x^2}{2EA}\,dx;\quad\text{if } EA = \text{constant},\ U^* = \tfrac12F_x^\tr dF_x,\ \text{where } d = \frac{L}{EA}.
\end{equation}
Beam element:
\begin{equation}
  U^* = \int_0^L\frac{M^2}{2EI}\,dx = \tfrac12\mat{r}^\tr\F_f = \tfrac12\F_f^\tr\mat{d}\F_f
\end{equation}
In general, we have the following relation:
\begin{equation}
  \boxed{U^* = \tfrac12\F_f^\tr\mat{d}\F_f}
\end{equation}
Therefore,
\begin{equation}
  \mat{d} = \frac{\partial^2U^*}{\partial^2\F_f}\qquad\text{or}\qquad d_{ij} = \frac{\partial^2U^*}{\partial F_{fi}\,\partial F_{fj}}
\end{equation}
With the above relation, once the complementary strain energy $U^*$ can be formulated as a function of independent member forces $\F_f$, the flexibility coefficients can be readily obtained.

For example, for the simply supported beam, the complementary strain energy indicated as $U^*_a$ is:
\begin{equation*}
  U^*_a = \int_0^L\frac{M^2}{2EI}\,dx
\end{equation*}
where $M(x) = -M_1\bigl(1-\frac xL\bigr) + M_2\frac xL = \bigl\langle -1+\frac xL, \frac xL\bigr\rangle\Bmat{M_1\\ M_2} = \NN(\F_f)_a$. Therefore,
\begin{equation*}
  U^*_a = \int_0^L\frac{M^2}{2EI}\,dx = \frac{1}{2EI}\int_0^L M^\tr M\,dx = \frac{1}{2EI}\int_0^L(\F_f)_a^\tr\NN^\tr\NN(\F_f)_a\,dx = \tfrac12(\F_f)_a^\tr\mat{d}_a(\F_f)_a
\end{equation*}
where
\begin{equation*}
  \mat{d}_a = \frac{1}{EI}\int_0^L\NN^\tr\NN\,dx = \frac{L}{6EI}\bmat{2 & -1\\ -1 & 2}
\end{equation*}
Similarly, for case (b): since $M(x) = F_2(L-x) + M_2 = \langle L-x, 1\rangle\Bmat{F_2\\ M_2} = \NN(\F_f)_b$, we have
\begin{gather*}
  \mat{d}_b = \frac{1}{EI}\int_0^L\NN^\tr\NN\,dx = \frac{L}{6EI}\bmat{2L^2 & 3L\\ 3L & 6}\\
  U^*_b = \tfrac12\langle F_2, M_2\rangle\frac{L}{6EI}\bmat{2L^2 & 3L\\ 3L & 6}\Bmat{F_2\\ M_2} = \tfrac12(\F_f)_b^\tr\mat{d}_b(\F_f)_b
\end{gather*}
Therefore, with different choices of $\F_f$, we seem to get different values of $U^*$. Actually, this is not true at all. From the equilibrium equations, we know that
\begin{equation*}
  \Bmat{F_2\\ M_2} = \bmat{-1/L & -1/L\\ 0 & 1}\Bmat{M_1\\ M_2}\qquad\text{or}\qquad\boxed{(\F_f)_b = \Tmat_{ba}(\F_f)_a}
\end{equation*}
Substituting $(\F_f)_b = \Tmat_{ba}(\F_f)_a$ (or $\F_f' = \Tmat\F_f$) into $U^*_b$ yields
\begin{equation*}
  U^*_b = \tfrac12(\F_f)_b^\tr\mat{d}_b(\F_f)_b = \tfrac12(\F_f)_a^\tr\Tmat_{ba}^\tr\mat{d}_b\Tmat_{ba}(\F_f)_a = \tfrac12(\F_f)_a^\tr\mat{d}_a(\F_f)_a = U^*_a
\end{equation*}
where
\begin{equation*}
  \boxed{\mat{d}_a = \Tmat_{ba}^\tr\mat{d}_b\Tmat_{ba}}\qquad(\text{or } \mat{d} = \Tmat^\tr\mat{d}'\Tmat)
\end{equation*}
Substituting $\Tmat_{ba}$ and $\mat{d}_b$ into this equation gives
\begin{equation*}
  \mat{d}_a = \frac{L}{6EI}\bmat{2 & -1\\ -1 & 2}
\end{equation*}
With the above derivations, we have two important conclusions. First, the same value of $U^*$ is obtained for different choices of $\F_f$. Second, the transformation between two flexibility matrices follows the following rule:
\begin{equation}
  \boxed{(\F_f)_b = \Tmat_{ba}(\F_f)_a\quad\Longleftrightarrow\quad\mat{d}_a = \Tmat_{ba}^\tr\mat{d}_b\Tmat_{ba}}
\end{equation}
or
\begin{equation}
  \F_f' = \Tmat\F_f\quad\Longleftrightarrow\quad\mat{d} = \Tmat^\tr\mat{d}'\Tmat
\end{equation}

\section{Flexibility--stiffness transformation}
For a linearly elastic body, the strain energy for a supported element is
\begin{equation}
  U = \tfrac12\mat{r}^\tr\F_f = \tfrac12\mat{r}^\tr\kmat_n\mat{r}
\end{equation}
where $\kmat_n = \mat{d}^{-1}$ is the stiffness matrix corresponding to the supported element.

The stiffness matrix for the corresponding unsupported element can be obtained from $\kmat_n$ by the congruent transformation. To do so, we need to know the relation:
\begin{equation*}
  \mat{r} = \Tmat\bs{\Delta}
\end{equation*}
or equivalently
\begin{equation*}
  \F = \Tmat^\tr\F_f
\end{equation*}
It is usually convenient to obtain $\Tmat$ from the relation between forces. According to the equilibrium conditions, we have
\begin{equation}
  \F_s = \bs{\Phi}\F_f
\end{equation}
where $\bs{\Phi}$ is called the \emph{equilibrium matrix}, which depends on how the element is supported. For example, for the simply supported case
\begin{fig}
  % Unsupported beam element: F1, M1, F2, M2
\begin{tikzpicture}[scale=0.85,font=\small]
  \draw[member] (0,0) -- (6,0);
  \draw[load] (0,0.1) -- (0,1.3) node[above] {$F_1$}; \draw[load] (6,0.1) -- (6,1.3) node[above] {$F_2$};
  \draw[->,thick] (-0.05,-0.4) arc (-100:100:0.4); \node[left] at (-0.4,0) {$M_1$};
  \draw[->,thick] (6.05,-0.4) arc (280:80:0.4); \node[right] at (6.4,0) {$M_2$};
\end{tikzpicture}

\end{fig}
\begin{equation*}
  \Bmat{F_1\\ F_2} = \bmat{1/L & 1/L\\ -1/L & -1/L}\Bmat{M_1\\ M_2},\quad\text{i.e.}\quad \bs{\Phi} = \bmat{1/L & 1/L\\ -1/L & -1/L}
\end{equation*}
Therefore,
\begin{equation*}
  \Bmat{\F_f\\ \F_s} = \Bmat{\mat{I}\\ \bs{\Phi}}\F_f\qquad\text{or}\qquad \F = \Tmat^\tr\F_f\quad\rightarrow\quad \Tmat = \langle\mat{I}, \bs{\Phi}^\tr\rangle
\end{equation*}
According to the congruent transformation, the stiffness matrix for the unsupported element
\begin{equation}
  \kmat = \Tmat^\tr\kmat_n\Tmat = \Bmat{\mat{I}\\ \bs{\Phi}}\kmat_n\langle\mat{I}, \bs{\Phi}^\tr\rangle = \bmat{\kmat_n & \kmat_n\bs{\Phi}^\tr\\ \bs{\Phi}\kmat_n & \bs{\Phi}\kmat_n\bs{\Phi}^\tr}
\end{equation}
If expressed in terms of $\mat{d}$,
\begin{equation}
  \kmat = \bmat{\mat{d}^{-1} & \mat{d}^{-1}\bs{\Phi}^\tr\\ \bs{\Phi}\mat{d}^{-1} & \bs{\Phi}\mat{d}^{-1}\bs{\Phi}^\tr}
\end{equation}
This is the so-called \emph{stiffness--flexibility transformation}.

Substituting the expressions for $\mat{d}$ and $\bs{\Phi}$ of the simply supported beam into the above transformation results in
\begin{equation*}
  \kmat = \begin{array}{c}
    \begin{array}{cccc}\ \ \theta_1\ \ & \ \ \theta_2\ \ & \ \ v_1\ \ & \ \ v_2\ \ \end{array}\\
    \left[\begin{array}{cccc}
      \dfrac{4EI}{L} & \dfrac{2EI}{L} & \dfrac{6EI}{L^2} & -\dfrac{6EI}{L^2}\\[8pt]
      \dfrac{2EI}{L} & \dfrac{4EI}{L} & \dfrac{6EI}{L^2} & -\dfrac{6EI}{L^2}\\[8pt]
      \dfrac{6EI}{L^2} & \dfrac{6EI}{L^2} & \dfrac{12EI}{L^3} & -\dfrac{12EI}{L^3}\\[8pt]
      -\dfrac{6EI}{L^2} & -\dfrac{6EI}{L^2} & -\dfrac{12EI}{L^3} & \dfrac{12EI}{L^3}
    \end{array}\right]
  \end{array}
\end{equation*}
It should be noted that the ordering of degrees of freedom is different from the conventional one. To get the conventional stiffness matrix needs reordering the stiffness coefficients. Such a reordering can be avoided if we order the components of $\F$ and redefine $\Tmat$ as follows:
\begin{equation*}
  \Bmat{F_1\\ M_1\\ F_2\\ M_2} = \bmat{1/L & 1/L\\ 1 & 0\\ -1/L & -1/L\\ 0 & 1}\Bmat{M_1\\ M_2}\qquad\text{or}\qquad \F = \Tmat^\tr\F_f,\quad\text{where } \Tmat^\tr = \bmat{1/L & 1/L\\ 1 & 0\\ -1/L & -1/L\\ 0 & 1}
\end{equation*}
Then $\kmat = \Tmat^\tr\kmat_n\Tmat$ will result in the conventional stiffness matrix.

\section{Deformation vector of supported element versus displacement vector of unsupported element}
As discussed above, we can establish, from the equilibrium equations, that
\begin{equation*}
  \F = \Tmat^\tr\F_f
\end{equation*}
According to the contragredient principle, the deformation vector $\mat{r}$ conjugate to $\F_f$ and the displacement vector $\bs{\Delta}$ conjugate to $\F$ ($= \langle F_1, M_1, F_2, M_2\rangle^\tr$) are related by
\begin{equation*}
  \mat{r} = \Tmat\bs{\Delta}
\end{equation*}
From the equilibrium equations, we can obtain $\Tmat$ for the three types of supported beam and the corresponding $\mat{r}$ as follows.
\begin{enumerate}[label=(\alph*),leftmargin=*]
  \item Simply supported at both ends
    \begin{fig}
      % Deformation of the simply supported element vs. displacements of the unsupported element
\begin{tikzpicture}[scale=1.0,font=\small]
  \coordinate (A) at (0,0); \coordinate (B) at (9,3);
  \draw[very thick] (A) -- (B);
  \draw[dashed,thick] (A) .. controls (2.5,2.6) and (6.0,1.2) .. (B);
  \draw[thin] (A) -- ++(52:1.8); \draw[thin] (B) -- ++(232:2.0); \draw[thin] (B) -- ++(18.4:0.9) ; \draw[thin] (B) -- ++(198.4:2.2);
  \draw[thin] (A) -- ++(2.4,0);
  \draw[->] ($(A)+(1.7,0)$) arc (0:18.4:1.7); \node at (2.05,0.3) {$\theta_r$};
  \draw[->] ($(A)+(1.1,0)$) arc (0:52:1.1); \node at (1.05,0.95) {$\theta_1$};
  \draw[->] ($(A)+(18.4:1.45)$) arc (18.4:52:1.45); \node at (1.15,1.4) {$\bar\theta_1$};
  \draw[thin] (B) -- ++(-2.4,0);
  \draw[->] ($(B)+(180:1.7)$) arc (180:198.4:1.7); \node at ($(B)+(-2.0,-0.3)$) {$\theta_r$};
  \draw[->] ($(B)+(180:1.1)$) arc (180:232:1.1); \node at ($(B)+(-1.0,-0.95)$) {$\theta_2$};
  \draw[->] ($(B)+(198.4:1.45)$) arc (198.4:232:1.45); \node at ($(B)+(-1.65,-0.85)$) {$\bar\theta_2$};
  \draw[dim] (-0.3,-1.5) -- (-0.3,0) node[midway,left] {$v_1$};
  \draw[dim] (9.3,-1.5) -- (9.3,3) node[midway,right] {$v_2$};
  \draw[thin] (0,-1.5) -- (-0.5,-1.5); \draw[thin] (9,-1.5) -- (9.5,-1.5);
  \draw[dim] (0,-1.5) -- (9,-1.5) node[midway,fill=white,inner sep=1pt] {$L$};
\end{tikzpicture}

    \end{fig}
    \begin{gather*}
      \Tmat = \bmat{1/L & 1 & -1/L & 0\\ 1/L & 0 & -1/L & 1}\\
      \bar\theta_1 = \theta_1 - (v_2-v_1)/L = \theta_1 - \theta_r\\
      \bar\theta_2 = \theta_2 - (v_2-v_1)/L = \theta_2 - \theta_r\\
      \text{where } \theta_r = (v_2-v_1)/L
    \end{gather*}
  \item Cantilever (fixed at $x = 0$ and free at $x = L$)
    \begin{fig}
      % Deformation of the cantilever element vs. displacements of the unsupported element
\begin{tikzpicture}[scale=1.0,font=\small]
  \coordinate (A) at (0,0); \coordinate (B) at (9,3);
  % rigid-body line: rotation theta_1 about node 1
  \coordinate (Br) at (9,1.6);
  \draw[very thick] (A) -- (Br);
  \draw[dashed,thick] (A) .. controls (3,0.8) and (6.5,1.6) .. (B);
  \draw[thin] (A) -- ++(2.6,0);
  \draw[->] ($(A)+(2.2,0)$) arc (0:10.1:2.2); \node[right] at (2.3,0.25) {$\theta_1$};
  \draw[thin] (B) -- ++(210.1:2.4); \draw[thin] (B) -- ++(190.1:2.4);
  \draw[->] ($(B)+(190.1:1.8)$) arc (190.1:210.1:1.8); \node at ($(B)+(-2.2,-0.85)$) {$\bar\theta_2$};
  \draw[->] ($(B)+(180:1.2)$) arc (180:210.1:1.2); \node at ($(B)+(-1.1,-0.75)$) {$\theta_2$};
  \draw[thin] (B) -- ++(-2.4,0); \node[above left] at ($(B)+(-1.8,0)$) {$\theta_1$};
  \draw[dim] (9.3,1.6) -- (9.3,3) node[midway,right] {$\bar v_2$};
  \draw[dim] (9.3,-1.5) -- (9.3,1.6) node[midway,right] {$v_2$};
  \node[left] at (8.9,0.9) {$L\theta_1$};
  \draw[dim] (-0.3,-1.5) -- (-0.3,0) node[midway,left] {$v_1$};
  \draw[thin] (0,-1.5) -- (-0.5,-1.5); \draw[thin] (9,-1.5) -- (9.5,-1.5);
  \draw[dim] (0,-1.5) -- (9,-1.5) node[midway,fill=white,inner sep=1pt] {$L$};
\end{tikzpicture}

    \end{fig}
    \begin{gather*}
      \Tmat = \bmat{-1 & -L & 1 & 0\\ 0 & -1 & 0 & 1}\\
      \bar v_2 = v_2 - v_1 - L\theta_1\\
      \bar\theta_2 = \theta_2 - \theta_1
    \end{gather*}
  \item Slide at $x = 0$ and simple support at $x = L$
    \begin{gather*}
      \Tmat = \bmat{1 & L & -1 & 0\\ 0 & -1 & 0 & 1}\\
      \bar v_1 = v_1 - v_2 + L\theta_1\\
      \bar\theta_2 = \theta_2 - \theta_1
    \end{gather*}
\end{enumerate}
The deformation vectors $\mat{r}$ can be obtained by excluding from $\bs{\Delta}$ the rigid body modes; the associated rigid body modes can be constructed as follows.
\begin{enumerate}[label=(\alph*)]
  \item $v_1\rightarrow\theta_r$ w.r.t.\ node 1
  \item $v_1\rightarrow\theta_1$ w.r.t.\ node 1
  \item $v_2\rightarrow\theta_1$ w.r.t.\ node 2
\end{enumerate}

\section{Flexibility matrices for some special elements}
\begin{enumerate}[label=(\alph*)]
  \item Beam with an internal hinge
  \item Beam with internal rotational spring(s)
  \item Beam with an end rotational spring
\end{enumerate}
\begin{fig}
  % Special elements: (a) internal hinge, (b) internal rotational springs, (c) end rotational spring
\begin{tikzpicture}[scale=0.8,font=\small]
  \newcommand{\rspring}[1]{\draw[thick] #1 circle (0.18); \draw[thick] #1 circle (0.09);}
  % (a)
  \begin{scope}
    \node[left] at (-0.6,0) {(a)};
    \pic at (0,0) {wallL}; \draw[member] (0,0) -- (6,0); \node[node dot] at (2.5,0) {}; \node[node dot] at (6,0) {}; \pic at (6,0) {roller};
  \end{scope}
  % (b)
  \begin{scope}[yshift=-2.2cm]
    \node[left] at (-0.6,0) {(b)};
    \pic at (0,0) {wallL}; \draw[member] (0,0) -- (6,0); \rspring{(2.5,0)}
    \begin{scope}[xshift=8.5cm]
      \draw[member] (0,0) -- (6,0); \node[node dot] at (0,0) {}; \node[node dot] at (6,0) {}; \pic at (0,0) {roller}; \pic at (6,0) {pin}; \rspring{(1.8,0)}
    \end{scope}
  \end{scope}
  % (c)
  \begin{scope}[yshift=-4.4cm]
    \node[left] at (-0.6,0) {(c)};
    \draw[member] (0,0) -- (6,0); \node[node dot] at (6,0) {}; \pic at (0,0) {roller}; \pic at (6,0) {pin}; \rspring{(0.25,0.1)}
  \end{scope}
\end{tikzpicture}

\end{fig}
Consider the case of (b)-2:
\begin{fig}
  % Simple beam with an internal rotational spring at x = a, and its moment diagram
\begin{tikzpicture}[scale=0.85,font=\small]
  \begin{scope}
    \draw[member] (0,0) -- (6,0); \node[node dot] at (0,0) {}; \node[node dot] at (6,0) {}; \pic at (0,0) {roller}; \pic at (6,0) {pin};
    \draw[thick] (1.8,0) circle (0.18); \draw[thick] (1.8,0) circle (0.09);
  \end{scope}
  \begin{scope}[xshift=8.5cm]
    \draw (0,0) -- (6,0);
    \draw[thick] (0,0) -- (0,-1.0) -- (6,1.2) -- (6,0);
    \node[left] at (0,-0.6) {$-M_1$}; \node[right] at (6,0.7) {$M_2$};
    \draw[dotted] (1.8,0) -- (1.8,-0.34); \node[below] at (1.8,-0.4) {$M_s$};
    \draw[->] (0,0.4) -- (1.8,0.4) node[above,pos=0.6] {$x = a$};
  \end{scope}
\end{tikzpicture}

\end{fig}
\begin{gather*}
  U^* = U^*_b + U^*_s\\
  M(x) = -M_1 + \frac{M_1+M_2}{L}x\\
  M_s = M(x)\big|_{x=a} = \Bigl\langle -1+\frac aL, \frac aL\Bigr\rangle\Bmat{M_1\\ M_2} = \NN\F_f
\end{gather*}
where $\NN = \bigl\langle\frac{-b}{L}, \frac aL\bigr\rangle$ and $b = L-a$.
\begin{gather*}
  U^*_b = \tfrac12\F_f^\tr\mat{d}_b\F_f\\
  U^*_s = \tfrac12M_s\Delta\theta = \tfrac12M_s\frac{M_s}{K_\theta} = \tfrac12\frac{M_s^\tr M_s}{K_\theta} = \tfrac12\frac{\F_f^\tr\NN^\tr\NN\F_f}{K_\theta} = \tfrac12\F_f^\tr\mat{d}_s\F_f
\end{gather*}
where
\begin{equation*}
  \mat{d}_s = \frac{\NN^\tr\NN}{K_\theta} = \frac{1}{K_\theta L^2}\bmat{b^2 & -ab\\ -ab & a^2}
\end{equation*}
\begin{gather*}
  U^* = U^*_b + U^*_s = \tfrac12\F_f^\tr\mat{d}_b\F_f + \tfrac12\F_f^\tr\mat{d}_s\F_f = \tfrac12\F_f^\tr\mat{d}\F_f\\
  \Rightarrow\quad \mat{d} = \mat{d}_b + \mat{d}_s
\end{gather*}

% =====================================================================
\clearpage
\section[Appendix: Conjugate beam method for beams with springs (Leu, Sung \& Tsai, 2000)]{Appendix: Conjugate beam method for beams with springs}
\label{sec:conjugate-paper}
\begin{center}
  \emph{L.~J.~Leu, Y.~C.~Sung, and M.~H.~Tsai}\\
  Department of Civil Engineering, National Taiwan University, Taipei 10617, Taiwan\\[2pt]
  {\small pp.~49--54, Proceedings of the Thirteenth KKNN Symposium on Civil Engineering, December 7--8, 2000, Taipei, Taiwan}
\end{center}

\subsection*{Abstract}
The conjugate beam method is a commonly used and convenient method for calculating the elastic deformations of a beam. However, it seems that this method is not applicable to the cases in which the beam is supported by linear elastic translational/rotational springs, or has internal rotational springs. This paper has made an extension of the conjugate method to deal with those cases. It is found that the conjugate beam also contains springs for a beam with springs but the types of springs are different for the real and conjugate beams. However, it is not necessary to know the stiffness of the springs for the conjugate beam; instead a force condition at the location of each spring can be stated. Several examples on the calculation of beam deflections validate the present extension.

\subsection*{Introduction}
The conjugate beam method was first presented by Otto Mohr in 1860. Since then, it has become a powerful method for the calculation of slope and deflection of beams. However, to the best knowledge of the authors, how to apply the method to beams with translational/rotational springs has not been investigated so far. This paper aims at bridging this gap.

\subsection*{Conjugate beams for beams with spring}
\paragraph{Sign convention.}
The main ideas of the conjugate beam method can be summarized as follows. The slopes and deflections of a beam subjected to applied loads can be evaluated by analyzing a conjugate beam having an applied distributed load that varies as the curvature of the real beam. The shear and moment at any point in the conjugate beam are then equal to the slope and vertical displacement, respectively, of the real beam at that point. Therefore, the calculation of slopes and deflections of a beam is reduced to obtaining the variation of the shear and moment in the corresponding conjugate beam.

The sign convention for slopes and deflections in a real beam and that for shear and moment in a conjugate beam are shown in Fig.~\ref{fig:cb1}. Note that the sign convention for shear and moment in the real beam is the same as that for the conjugate beam and that the distributed load acting on the conjugate beam is positive or upward if the curvature of the real beam is positive.
\begin{figure}[H]
  \centering
  % Sign convention for the conjugate beam method (paper Fig. 1)
\begin{tikzpicture}[scale=0.9,font=\scriptsize]
  \begin{scope}
    \draw[->] (0,0) -- (5.2,0) node[right] {x}; \draw[->] (0,0) -- (0,2.6) node[above] {y};
    \node[above] at (0.9,2.2) {Real Beam}; \node[above] at (3.4,2.2) {Conjugate Beam};
    \draw[thick] (0.5,1.0) .. controls (0.8,1.6) .. (1.3,1.9); \draw[dim] (0.7,0) -- (0.7,1.0) node[midway,right] {$v$};
    \draw[->] (1.9,1.4) -- (2.2,1.4); \draw[->] (2.2,1.4) -- (1.9,1.4);
    \draw[thick] (2.7,0.9) .. controls (3.4,1.1) .. (4.1,0.9) -- (4.3,1.9) .. controls (3.4,1.7) .. (2.5,1.9) -- cycle;
    \draw[->,thick] (2.5,0.7) arc (-120:-240:0.5); \node[below] at (2.4,0.5) {$\overline M$};
    \draw[->,thick] (4.3,2.1) arc (60:-60:0.5); \node[below] at (4.4,1.0) {$\overline M$};
    \node[below,align=center] at (0.9,-0.1) {Positive\\displacement}; \node[below] at (3.4,-0.1) {Positive moment};
  \end{scope}
  \begin{scope}[xshift=7cm]
    \draw[->] (0,0) -- (5.2,0) node[right] {x}; \draw[->] (0,0) -- (0,2.6) node[above] {y};
    \node[above] at (0.9,2.2) {Real Beam}; \node[above] at (3.4,2.2) {Conjugate Beam};
    \draw[thick] (0.4,1.0) .. controls (0.8,1.5) .. (1.3,1.7); \draw[dashed] (0.6,1.15) -- (1.5,1.15); \draw[->] (1.3,1.15) arc (0:28:0.7); \node[right] at (1.2,1.45) {$\theta$};
    \draw[->] (1.9,1.4) -- (2.2,1.4); \draw[->] (2.2,1.4) -- (1.9,1.4);
    \draw[thick] (2.7,0.9) rectangle (4.1,1.9);
    \draw[->,thick] (2.5,1.9) -- (2.5,0.9) node[left,pos=0.5] {$\overline V$};
    \draw[->,thick] (4.3,0.9) -- (4.3,1.9) node[right,pos=0.5] {$\overline V$};
    \node[below] at (0.9,-0.1) {Positive slope}; \node[below] at (3.4,-0.1) {Positive shear};
  \end{scope}
\end{tikzpicture}

  \caption{Sign convention.}
  \label{fig:cb1}
\end{figure}

\paragraph{Beam with translational spring (Cases 1 and 2).}
Consider a real beam with translational springs. The shear force at the location of a typical translational spring, $x$, will have a jump $\Delta V = V(x^+) - V(x^-)$. From the equilibrium equation at that location, we obtain $\Delta V = -F$, where $F$ is the force in the translational spring and is positive if tensile. Hence, the deflection at $x$ in the real beam is
\begin{equation}
  \delta = \frac Fk = -\frac{\Delta V}{k} \tag{1}\label{cb:1}
\end{equation}
Since there is a jump on the shear force at the location of a translational spring for the real beam, there must exist an internal rotational spring for the conjugate beam at the same location as shown in case 1 of Fig.~\ref{fig:cb2}. There comes a question, what is the stiffness of the internal rotational spring? It turns out that its value is not needed. According to the conjugate beam theory, the deflection of the real beam given in Eq.~\eqref{cb:1} shall be equal to the moment at the same location in the conjugate beam. Hence, we have
\begin{equation}
  \overline M = -\frac{\Delta V}{k} \tag{2}\label{cb:2}
\end{equation}
This means that the moment $\overline M$ in the internal rotational spring in the conjugate beam can be determined from the jump on the shear force and the stiffness of the translational spring in the real beam.

Referring to case 2 of Fig.~\ref{fig:cb2}, note that if the translational spring is located at the right end of a beam, then $\Delta V = 0 - V_2 = -V_2$ in Eqs.~\eqref{cb:1} and \eqref{cb:2}, where $V_2$ is the shear force at the right end. Also, the right end of the conjugate beam is a hinged support with an external rotational spring or equivalently a fixed support with an internal rotational spring. However, if the spring is located at the left end, $\Delta V = V_1 - 0 = V_1$ in Eqs.~\eqref{cb:1} and \eqref{cb:2}, where $V_1$ is the shear force at the left end.

\paragraph{Beam with internal rotational spring (Cases 3 and 4).}
Consider a real beam with internal rotational springs. The rotation at the location of a typical internal rotational spring, $x$, will have a jump $\Delta\theta$ of the following value
\begin{equation}
  \Delta\theta = \theta(x^+) - \theta(x^-) = \frac{M(x)}{k_\theta} \tag{3}\label{cb:3}
\end{equation}
In Eq.~\eqref{cb:3}, $M(x)$ is the moment at $x$ in the real beam and $k_\theta$ is the stiffness of the internal rotational spring.

To produce such a rotation jump, there must be a shear force jump that equals the rotation jump in the conjugate beam at the same location according to the conjugate beam theory. Hence,
\begin{equation}
  \Delta\overline V = \frac{M(x)}{k_\theta} \tag{4}\label{cb:4}
\end{equation}
where $\Delta\overline V = \overline V(x^+) - \overline V(x^-)$. To introduce $\Delta\overline V$ in the conjugate beam, a translational spring is placed at $x$ as shown in case 3 of Fig.~\ref{fig:cb2}. However, its stiffness is not required since the force acting on the conjugate beam from the translational spring shall equal $\Delta\overline V$ according to the conjugate beam theory.

Note that if the rotational spring is located at the right end of a beam as shown in case 4 of Fig.~\ref{fig:cb2}, then $\Delta\overline V$ in Eq.~\eqref{cb:4} should be replaced by $-\overline V_2$, where $\overline V_2$ is the shear force at the right end. On the other hand, if the rotational spring is located at the left end of a beam, then $\Delta\overline V$ should be replaced by $\overline V_1$, the shear force at the left end.

\paragraph{Beam with external rotational spring (Cases 5 and 6).}
Consider a real beam with external rotational springs. The rotation at the location of a typical external rotational spring, $x$, can be expressed by
\begin{equation}
  \theta = \frac{\Delta M}{k_\theta} \tag{5}\label{cb:5}
\end{equation}
where $\Delta M = M(x^+) - M(x^-)$, the moment jump induced by the external rotational spring. The external rotational spring in the real beam corresponds to an ``elastic slide'' in the conjugate beam as shown in case 5 of Fig.~\ref{fig:cb2}. However, the stiffness of the slide is not required. According to the conjugate beam theory, the shear force it carries shall be equal to the rotation in the real beam; by making use of Eq.~\eqref{cb:5}, this can be stated as
\begin{equation}
  \overline V = \frac{\Delta M}{k_\theta} \tag{6}\label{cb:6}
\end{equation}
Note that if the external rotational spring is located at the right end of a beam, then $\Delta M$ should be replaced by $-M_2$ in Eq.~\eqref{cb:6}, where $M_2$ is the moment at the right end. However, if the external rotational spring is located at the left end of a beam, then $\Delta M$ should be replaced by $M_1$, where $M_1$ is the moment at the left end. These two special cases are shown in case 6 of Fig.~\ref{fig:cb2}.
\begin{figure}[H]
  \centering
  % Conjugate beams for beams with springs and the corresponding force conditions (paper Fig. 2)
\begin{tikzpicture}[scale=0.68,font=\scriptsize,
  tspring/.pic={\draw[thick,decorate,decoration={zigzag,segment length=4pt,amplitude=2.5pt,pre length=2pt,post length=2pt}] (0,0) -- (0,-1.0); \pic at (0,-1.0) {fixed};},
  rspring/.pic={\draw[thick] (0,0) circle (0.16); \draw[thick] (0,0) circle (0.08);},
  cut/.pic={\draw (0,-0.25) -- (0,0.25);},
  slide/.pic={\draw[fill=white] (-0.12,-0.35) rectangle (0.12,0.35); \draw (-0.12,-0.2) -- (0.12,-0.2); \draw (-0.12,0) -- (0.12,0); \draw (-0.12,0.2) -- (0.12,0.2);}]
  \newcommand{\seg}{\draw[member] (0,0) -- (3,0); \pic at (0,0) {cut}; \pic at (3,0) {cut};}
  % header
  \node at (-0.8,1.5) {Case}; \node[align=center] at (1.5,1.5) {Actual Beam\\$(V,M)^*$}; \node at (7.0,1.5) {Conjugate Beam}; \node at (12.5,1.5) {Conjugate Beam};
  \draw (-1.4,0.9) -- (14.8,0.9);
  % ---------- case 1 ----------
  \begin{scope}[yshift=-1.4cm]
    \node at (-0.8,0) {1};
    \seg \pic at (1.5,0) {tspring}; \node[right] at (1.6,-0.5) {$k$};
    \draw[<->] (4.0,0) -- (4.8,0);
    \begin{scope}[xshift=5.5cm] \seg \pic at (1.5,0) {rspring}; \end{scope}
    \draw[->] (9.5,0) -- (10.3,0);
    \begin{scope}[xshift=11cm] \draw[member] (0,0) -- (1.3,0); \draw[member] (1.7,0) -- (3,0); \pic at (0,0) {cut}; \pic at (3,0) {cut};
      \draw[->,thick] (1.05,-0.4) arc (-100:100:0.4); \draw[->,thick] (1.95,-0.4) arc (280:80:0.4);
      \node[above] at (1.5,0.6) {$\overline M = -\dfrac{\Delta V}{k}$}; \end{scope}
  \end{scope}
  \draw (-1.4,-3.1) -- (14.8,-3.1);
  % ---------- case 2 ----------
  \begin{scope}[yshift=-4.9cm]
    \node at (-0.8,0) {2};
    \draw[member] (0,0) -- (3,0); \pic at (1.5,0) {cut}; \pic at (0,0) {tspring}; \pic at (3,0) {tspring}; \node[right] at (0.1,-0.5) {$k_1$}; \node[right] at (3.1,-0.5) {$k_2$};
    \draw[<->] (4.0,0) -- (4.8,0);
    \begin{scope}[xshift=5.5cm,yshift=0.5cm] \draw[member] (0,0) -- (3,0); \pic at (1.5,0) {cut}; \pic at (0,0) {pin}; \pic at (3,0) {pin}; \pic at (0,0.1) {rspring}; \pic at (3,0.1) {rspring}; \end{scope}
    \begin{scope}[xshift=5.5cm,yshift=-0.7cm] \draw[member] (0,0) -- (3,0); \pic at (1.5,0) {cut}; \pic at (0,0) {wallL}; \pic at (3,0) {wallR}; \pic at (0.25,0) {rspring}; \pic at (2.75,0) {rspring}; \end{scope}
    \draw[->] (9.5,0) -- (10.3,0);
    \begin{scope}[xshift=11cm] \draw[member] (0,0) -- (3,0); \pic at (1.5,0) {cut}; \pic at (0,0) {pin}; \pic at (3,0) {pin};
      \draw[->,thick] (-0.4,0.6) arc (110:-40:0.5); \draw[->,thick] (3.4,0.6) arc (70:220:0.5);
      \node[above] at (0.4,0.7) {$\overline M_1 = -\dfrac{V_1}{k_1}$}; \node[above] at (2.7,0.7) {$\overline M_2 = \dfrac{V_2}{k_2}$}; \end{scope}
  \end{scope}
  \draw (-1.4,-6.6) -- (14.8,-6.6);
  % ---------- case 3 ----------
  \begin{scope}[yshift=-8.3cm]
    \node at (-0.8,0) {3};
    \seg \pic at (1.5,0) {rspring}; \node[below] at (1.5,-0.3) {$k_\theta$};
    \draw[<->] (4.0,0) -- (4.8,0);
    \begin{scope}[xshift=5.5cm] \seg \pic at (1.5,0) {tspring}; \end{scope}
    \draw[->] (9.5,0) -- (10.3,0);
    \begin{scope}[xshift=11cm] \seg \draw[load] (1.5,0.1) -- (1.5,1.0); \node[right] at (1.6,0.7) {$\Delta\overline V = \dfrac{M}{k_\theta}$}; \end{scope}
  \end{scope}
  \draw (-1.4,-9.9) -- (14.8,-9.9);
  % ---------- case 4 ----------
  \begin{scope}[yshift=-11.8cm]
    \node at (-0.8,0) {4};
    \begin{scope}[yshift=0.5cm] \draw[member] (0,0) -- (3,0); \pic at (1.5,0) {cut}; \pic at (0,0) {pin}; \pic at (3,0) {pin}; \pic at (0,0.1) {rspring}; \pic at (3,0.1) {rspring}; \node[above] at (0.6,0.2) {$k_{\theta1}$}; \node[above] at (2.4,0.2) {$k_{\theta2}$}; \end{scope}
    \begin{scope}[yshift=-1.4cm] \draw[member] (0,0) -- (3,0); \pic at (1.5,0) {cut}; \pic at (0,0) {wallL}; \pic at (3,0) {wallR}; \pic at (0.25,0) {rspring}; \pic at (2.75,0) {rspring}; \node[below] at (0.6,-0.2) {$k_{\theta1}$}; \node[below] at (2.4,-0.2) {$k_{\theta2}$}; \end{scope}
    \draw[<->] (4.0,0) -- (4.8,0);
    \begin{scope}[xshift=5.5cm] \draw[member] (0,0) -- (3,0); \pic at (1.5,0) {cut}; \pic at (0,0) {tspring}; \pic at (3,0) {tspring}; \end{scope}
    \draw[->] (9.5,0) -- (10.3,0);
    \begin{scope}[xshift=11cm] \draw[member] (0,0) -- (3,0); \pic at (1.5,0) {cut};
      \draw[load] (0,-0.9) -- (0,-0.1) node[pos=0,below] {$\overline V_1 = \dfrac{M_1}{k_{\theta1}}$};
      \draw[load] (3,0.9) -- (3,0.1) node[pos=0,above] {$\overline V_2 = -\dfrac{M_2}{k_{\theta2}}$}; \end{scope}
  \end{scope}
  \draw (-1.4,-13.7) -- (14.8,-13.7);
  % ---------- case 5 ----------
  \begin{scope}[yshift=-15.4cm]
    \node at (-0.8,0) {5};
    \seg \pic at (1.5,0) {rspring}; \draw[thin] (1.5,-0.16) -- (1.5,-0.6); \pic at (1.5,-0.6) {fixed}; \node[left] at (1.3,-0.4) {$k_\theta$};
    \draw[<->] (4.0,0) -- (4.8,0);
    \begin{scope}[xshift=5.5cm] \seg \pic at (1.5,0) {slide}; \end{scope}
    \draw[->] (9.5,0) -- (10.3,0);
    \begin{scope}[xshift=11cm] \draw[member] (0,0) -- (1.3,0); \draw[member] (1.7,0) -- (3,0); \pic at (0,0) {cut}; \pic at (3,0) {cut};
      \draw[load] (1.3,0.7) -- (1.3,-0.1); \draw[load] (1.7,-0.7) -- (1.7,0.1);
      \node[right] at (3.2,-0.3) {$\overline V = \dfrac{\Delta M}{k_\theta}$}; \end{scope}
  \end{scope}
  \draw (-1.4,-17.0) -- (14.8,-17.0);
  % ---------- case 6 ----------
  \begin{scope}[yshift=-18.6cm]
    \node at (-1.1,0) {6};
    \draw[member] (0,0) -- (3,0); \pic at (1.5,0) {cut}; \pic at (0,0) {rspring}; \pic at (3,0) {rspring};
    \draw[thin] (-0.16,0) -- (-0.6,0.35); \pic[rotate=-40] at (-0.7,0.42) {wallR}; \draw[thin] (3.16,0) -- (3.6,0.35); \pic[rotate=40] at (3.7,0.42) {wallL};
    \node[above] at (0.6,0.2) {$k_{\theta1}$}; \node[above] at (2.4,0.2) {$k_{\theta2}$};
    \draw[<->] (4.0,0) -- (4.8,0);
    \begin{scope}[xshift=5.5cm] \draw[member] (0,0) -- (3,0); \pic at (1.5,0) {cut}; \pic at (0,0) {slide}; \pic at (3,0) {slide}; \pic at (0,-0.35) {fixed}; \pic at (3,-0.35) {fixed}; \end{scope}
    \draw[->] (9.5,0) -- (10.3,0);
    \begin{scope}[xshift=11cm] \draw[member] (0,0) -- (3,0); \pic at (1.5,0) {cut};
      \draw[load] (0,-0.9) -- (0,-0.1) node[pos=0,below] {$\overline V_1 = \dfrac{M_1}{k_{\theta1}}$};
      \draw[load] (3,0.9) -- (3,0.1) node[pos=0,above] {$\overline V_2 = -\dfrac{M_2}{k_{\theta2}}$}; \end{scope}
  \end{scope}
  \draw (-1.4,-20.4) -- (14.8,-20.4);
  \node[below] at (6.7,-20.5) {$*$ Internal forces of beam at the location of spring};
\end{tikzpicture}

  \caption{Conjugate beams for beams with springs and the corresponding force conditions.}
  \label{fig:cb2}
\end{figure}

\subsection*{Examples}
\paragraph{Example 1: Beam with translational springs.}
The beam considered and its conjugate beam are shown, respectively, in Figs.~3(a) and 3(b). We like to calculate $v_c$ and $\theta_c$. The shear force at point $a$ and the shear force jump at point $b$ are $V_a = -P/2$ and $\Delta V_b = 3P/2$, respectively. Therefore, $\overline M_a = -V_a/k_a = PL^3/6EI$ and $\overline M_b = -\Delta V_b/k_b = -PL^3/32EI$ according to case 2 of Fig.~\ref{fig:cb2} and Eq.~\eqref{cb:2}. Then, two equations of $\overline M_c$ and $\overline V_c$ need to be formulated to solve them. Referring to Fig.~3(b), equilibrium of moment with respect to point $a$ results in one equation; referring to Fig.~3(c), taking moment about point $b$ results in the other equation. Solving these two equations and according to the conjugate beam theory, we obtain $v_c = \overline M_c = -49PL^3/192EI$ and $\theta_c = \overline V_c = -94PL^2/192EI$.
\begin{figure}[H]
  \centering
  % Paper Fig. 3: beam with translational springs, its conjugate beam, and free body of bc
\begin{tikzpicture}[scale=0.8,font=\scriptsize,
  tspring/.pic={\draw[thick,decorate,decoration={zigzag,segment length=4pt,amplitude=2.5pt,pre length=2pt,post length=2pt}] (0,0) -- (0,-1.0); \pic at (0,-1.0) {fixed};}]
  % (a)
  \begin{scope}
    \draw[member] (0,0) -- (7.5,0); \pic at (0,0) {tspring}; \pic at (5,0) {tspring};
    \node[above] at (0,0.1) {a}; \node[above] at (5,0.1) {b}; \node[above] at (7.5,0.1) {c};
    \node[right] at (0.1,-0.5) {$k_a = \frac{3EI}{L^3}$}; \node[right] at (5.1,-0.5) {$k_b = \frac{48EI}{L^3}$};
    \draw[load] (7.5,1.2) -- (7.5,0.1) node[pos=0,right] {$P$};
    \draw[dim] (0,-1.6) -- (5,-1.6) node[midway,fill=white,inner sep=1pt] {$L$}; \draw[dim] (5,-1.6) -- (7.5,-1.6) node[midway,fill=white,inner sep=1pt] {$L/2$};
    \node[below] at (3.75,-2.2) {Fig.~3(a) Beam with translational springs};
  \end{scope}
  % (b)
  \begin{scope}[xshift=10cm]
    \draw[member] (0,0) -- (7.5,0); \pic at (0,0) {pin}; \draw[thick] (0,0.1) circle (0.14); \draw[thick] (5,0) circle (0.14); \draw[thick] (5,0) circle (0.07);
    \pic at (7.5,0) {wallR};
    \node[above] at (0,0.3) {a}; \node[above] at (5,0.2) {b}; \node[above] at (7.5,0.2) {c};
    \draw (0,0) -- (5,-1.6) -- (7.5,0);
    \foreach \x in {0.7,1.4,...,4.9}{\draw[->] (\x,-0.05) -- (\x,{-1.6*\x/5+0.05});}
    \foreach \x in {5.6,6.2,6.8}{\draw[->] (\x,-0.05) -- (\x,{-1.6*(7.5-\x)/2.5+0.05});}
    \node[below] at (5,-1.7) {$\frac{PL}{2EI}$};
    \draw[load] (7.5,-0.1) -- (7.5,-1.1) node[below] {$\overline V_c$};
    \draw[->,thick] (7.9,0.5) arc (60:-90:0.5); \node[right] at (7.9,0.6) {$\overline M_c$};
    \node[below] at (3.75,-2.2) {Fig.~3(b) Conjugate beam of Fig.~3(a)};
  \end{scope}
  % (c)
  \begin{scope}[yshift=-5.5cm,xshift=2cm]
    \draw[member] (0,0) -- (3.5,0); \node[above] at (0,0.1) {b}; \node[above] at (3.5,0.1) {c};
    \draw (0,0) -- (0,-1.6) -- (3.5,0); \foreach \x in {0.5,1.0,...,3.2}{\draw[->] (\x,-0.05) -- (\x,{-1.6*(3.5-\x)/3.5+0.05});}
    \draw[load] (-0.4,-1.2) -- (-0.4,-0.1) node[pos=0,below] {$\overline V_b$}; \draw[->,thick] (-0.6,0.8) arc (100:260:0.7); \node[above] at (-0.6,0.8) {$\overline M_b$};
    \draw[load] (3.9,-0.1) -- (3.9,-1.2) node[below] {$\overline V_c$}; \draw[->,thick] (4.1,0.8) arc (80:-80:0.7); \node[above] at (4.1,0.8) {$\overline M_c$};
    \node[below] at (1.75,-2.2) {Fig.~3(c) Free body of element $bc$};
  \end{scope}
\end{tikzpicture}

\end{figure}

\paragraph{Example 2: Beam with internal rotational spring.}
Consider the beam with an internal rotational spring as depicted in Fig.~4(a); we like to calculate $\theta_c$. The corresponding conjugate beam is shown in Fig.~4(b), in which $P = M_b/k_\theta = M_cL_1/Lk_\theta$ according to Eq.~\eqref{cb:4}; note that $L = L_1 + L_2$. Referring to Fig.~4(b), we can obtain $\overline V_c = [PL_1 + (M_cL/2EI)(2L/3)]/L = M_cL(L_1^2EI/k_\theta L^3 + 1/3)/EI$, which equals $\theta_c$ in accordance with the conjugate beam theory.
\begin{figure}[H]
  \centering
  % Paper Fig. 4: beam with an internal rotational spring and its conjugate beam
\begin{tikzpicture}[scale=0.8,font=\scriptsize]
  \begin{scope}
    \draw[member] (0,0) -- (7.5,0); \pic at (0,0) {pin}; \pic at (7.5,0) {roller};
    \draw[thick] (2.5,0) circle (0.14); \draw[thick] (2.5,0) circle (0.07);
    \node[above] at (0,0.1) {a}; \node[above left] at (2.4,0.15) {b}; \node[above right] at (2.6,0.15) {$k_\theta$}; \node[above] at (7.5,0.1) {c};
    \draw[->,thick] (8.0,0.5) arc (60:200:0.5); \node[right] at (8.0,0.5) {$M_c$};
    \draw[dim] (0,-1.0) -- (2.5,-1.0) node[midway,fill=white,inner sep=1pt] {$L_1$}; \draw[dim] (2.5,-1.0) -- (7.5,-1.0) node[midway,fill=white,inner sep=1pt] {$L_2$};
    \node[below] at (3.75,-1.5) {Fig.~4(a) Beam with internal rotational spring};
  \end{scope}
  \begin{scope}[xshift=10cm]
    \draw[member] (0,0) -- (7.5,0); \pic at (0,0) {pin}; \pic at (7.5,0) {roller};
    \draw (0,0) -- (7.5,1.4); \foreach \x in {0.8,1.6,...,7.2}{\draw[->] (\x,0.05) -- (\x,{1.4*\x/7.5-0.05});}
    \node[right] at (7.5,1.4) {$\frac{M_c}{EI}$};
    \draw[load] (2.5,0.1) -- (2.5,1.6) node[right] {$P$};
    \node[above] at (0,0.1) {a}; \node[below] at (2.5,-0.05) {b}; \node[above] at (7.5,0.1) {c};
    \draw[dim] (0,-1.0) -- (2.5,-1.0) node[midway,fill=white,inner sep=1pt] {$L_1$}; \draw[dim] (2.5,-1.0) -- (7.5,-1.0) node[midway,fill=white,inner sep=1pt] {$L_2$};
    \node[below] at (3.75,-1.5) {Fig.~4(b) Conjugate beam of Fig.~4(a)};
  \end{scope}
\end{tikzpicture}

\end{figure}

\paragraph{Example 3: Beam with external rotational spring.}
Consider the beam with an external rotational spring as depicted in Fig.~5(a); we like to calculate $v_c$ and $\theta_c$. Clearly there is a moment jump $\Delta M_b = 2PL$ induced by the external rotational spring at point $b$. The corresponding conjugate beam is shown in Fig.~5(b), where the shear force at point $b$ shall equal $\overline V_b = \Delta M_b/k_\theta = 2PL^2/EI$ from Eq.~\eqref{cb:6}. From the force equilibrium equation for the free body shown in Fig.~5(c), we can obtain $\overline V_a = PL^2/2EI + \overline V_b = 5PL^2/2EI$. On the basis of the force and moment equilibrium equations for the whole structure of Fig.~5(a), we can obtain $\overline V_c = \overline V_a = 5PL^2/2EI$ and $\overline M_c = (PL^2/2EI)2L/3 - (PL^2/2EI)(4L/3) + 2\overline V_aL = 14PL^3/3EI$. Finally, $v_c = \overline M_c$ and $\theta_c = \overline V_c$.
\begin{figure}[H]
  \centering
  % Paper Fig. 5: beam with an external rotational spring, its conjugate beam, and free body of ab
\begin{tikzpicture}[scale=0.8,font=\scriptsize]
  \begin{scope}
    \draw[member] (0,0) -- (7.5,0); \pic at (0,0) {pin};
    \draw[thick] (3.75,0) circle (0.14); \draw[thick] (3.75,0) circle (0.07); \draw[thin] (3.75,0.14) -- (3.5,0.6); \pic[rotate=-30] at (3.45,0.68) {wallR};
    \node[above] at (0,0.1) {a}; \node[below] at (3.75,-0.1) {b}; \node[right] at (3.95,0.35) {$k_\theta = \frac{EI}{L}$}; \node[above] at (7.5,0.1) {c};
    \draw[load] (7.5,0.1) -- (7.5,1.2) node[right] {$P$};
    \draw[dim] (0,-0.9) -- (3.75,-0.9) node[midway,fill=white,inner sep=1pt] {$L$}; \draw[dim] (3.75,-0.9) -- (7.5,-0.9) node[midway,fill=white,inner sep=1pt] {$L$};
    \node[below] at (3.75,-1.5) {Fig.~5(a) Beam with external rotational spring};
  \end{scope}
  \begin{scope}[xshift=10cm]
    \draw[member] (0,0) -- (7.5,0); \pic at (0,0) {pin};
    \draw[fill=white] (3.63,-0.35) rectangle (3.87,0.35); \draw (3.63,-0.15) -- (3.87,-0.15); \draw (3.63,0.15) -- (3.87,0.15);
    \pic at (7.5,0) {wallR};
    \draw (0,0) -- (3.75,-1.2); \foreach \x in {0.6,1.2,...,3.5}{\draw[->] (\x,-0.05) -- (\x,{-1.2*\x/3.75+0.05});} \node[below] at (3.5,-1.2) {$\frac{PL}{EI}$};
    \draw (3.75,1.2) -- (7.5,0); \foreach \x in {4.2,4.8,...,7.3}{\draw[->] (\x,0.05) -- (\x,{1.2*(7.5-\x)/3.75-0.05});} \node[above] at (4.0,1.2) {$\frac{PL}{EI}$};
    \draw[load] (0,-1.0) -- (0,-0.1) node[pos=0,below] {$\overline V_a$};
    \draw[load] (7.5,-0.1) -- (7.5,-1.0) node[below] {$\overline V_c$};
    \draw[->,thick] (7.9,0.5) arc (60:-90:0.5); \node[right] at (7.9,0.6) {$\overline M_c$};
    \node[above] at (0,0.1) {a}; \node[above left] at (3.6,0.1) {b}; \node[above] at (7.5,0.3) {c};
    \node[below] at (3.75,-1.8) {Fig.~5(b) Conjugate beam of Fig.~5(a)};
  \end{scope}
  \begin{scope}[yshift=-5cm,xshift=2cm]
    \draw[member] (0,0) -- (3.75,0); \node[above] at (0,0.1) {a}; \node[above] at (3.75,0.1) {b};
    \draw (0,0) -- (3.75,-1.2); \foreach \x in {0.6,1.2,...,3.5}{\draw[->] (\x,-0.05) -- (\x,{-1.2*\x/3.75+0.05});}
    \draw[load] (-0.4,-1.2) -- (-0.4,-0.1) node[pos=0,below] {$\overline V_a$};
    \draw[load] (4.1,-0.1) -- (4.1,-1.2) node[below] {$\overline V_b$}; \draw[->,thick] (4.3,0.9) arc (80:-80:0.8); \node[above] at (4.3,0.9) {$\overline M_b$};
    \node[below] at (1.9,-2.0) {Fig.~5(c) Free body of element $ab$};
  \end{scope}
\end{tikzpicture}

\end{figure}

\subsection*{Conclusions and discussions}
This paper has extended the conjugate beam method to beams with translational and/or rotational springs. The solutions to all the above examples have been verified by the direct stiffness method; the stiffness matrix for a beam with internal or end rotational springs can be found, for example, in McGuire et al.\ (2000). Although the above examples are all statically determinate, the proposed method is applicable to indeterminate beams as well; in particular it has been used conveniently to obtain the fixed-end forces for beams with internal or end rotational springs (Tsai 2000).

\subsection*{Acknowledgements}
This research was sponsored by the National Science Council of the Republic of China under Grant No.\ NSC 89-2211-E-002-124.

\subsection*{References}
\begin{itemize}[leftmargin=*]
  \item McGuire, W., Gallagher, R.H., and Ziemian, R.D. (2000), \emph{Matrix Structural Analysis}, 2nd Edition, John Wiley \& Sons, New York, USA.
  \item Tsai, M.~H. (2000), \emph{Nonlinear Analysis of Frames and Continuous Beams with Softening Moment-Curvature Relationship}, Master Thesis, Dept.\ of Civil Engineering, National Taiwan University, Taipei, Taiwan, (in Chinese).
\end{itemize}

\chapter{Principle of Virtual Forces}
\label{ch:pvf}

\section{Principle of virtual forces (PVF) and related principles or theorems}
The virtual force system should be ``statically admissible'' (the virtual displacements should be ``kinematically admissible''); this means that the external virtual forces and the internal forces are in equilibrium.

For an admissible virtual force system, the principle of virtual forces is stated as
\begin{equation}
  \dW^* = \dU_i^* \quad\Longleftrightarrow\quad \text{Compatibility Equations}
\end{equation}
where $\dW^* =$ the complementary virtual work done by the external virtual forces through the ``real displacements'' (or the external complementary virtual work), and $\dU_i^* =$ the complementary virtual internal energy (or the internal complementary virtual work). They are defined, respectively, by
\begin{align}
  \dW^* &= \int_S\delta t_iu_i\,dS + \int_V\delta f_iu_i\,dV\\
  \dU_i^* &= \int_V\delta\sigma_{ij}\eps_{ij}\,dV
\end{align}
The statically admissible conditions consist of the following two equilibrium conditions.
\begin{gather*}
  \delta\sigma_{ij,j} + \delta f_i = 0\ \text{ in } V\\
  \delta\sigma_{ij}n_j = \delta t_i\ \text{ on } S\ (\text{not } S_t)
\end{gather*}
For an (linearly or nonlinearly) elastic body, there exists a function of stresses (internal forces), called complementary strain energy $U^*$, such that $\dU^* = \dU_i^*$. If the external forces are conservative, then there exists a complementary load potential $V^*$ such that
\begin{equation*}
  \delta V^* = -\dW^*
\end{equation*}
Therefore, for a conservative system because the material is elastic and the external forces are conservative, there exists the complementary total potential energy defined by
\begin{equation*}
  \Pi^* = U^* + V^*
\end{equation*}
The PVF then becomes
\begin{equation}
  \delta\Pi^* = 0 \quad\Longleftrightarrow\quad \text{Compatibility Equations}
\end{equation}
which is often referred to as the \emph{principle of stationary (minimum) total complementary potential energy}.

If the external forces acting on a conservative system are point (concentrated) loads, then
\begin{equation*}
  \Pi^* = U^* + V^* = U^*(\Pvec) - \Pvec^\tr\bs{\Delta} = U^*(\Pvec) - \sum_{i=1}^nP_i\Delta_i
\end{equation*}
Therefore,
\begin{equation*}
  \delta\Pi^* = 0 \;\Rightarrow\; \sum_{i=1}^n\Bigl(\pd{U^*}{P_i} - \Delta_i\Bigr)\delta P_i = 0
\end{equation*}
As $\delta P_i$'s are independent, we can conclude that
\begin{equation}
  \pd{U^*}{P_i} = \Delta_i,\quad i = 1, n
\end{equation}
This is often referred to as \emph{Engesser's first theorem}. Note that for \emph{linearly} elastic bodies, $U^*$ is a \emph{quadratic} function of $\Pvec$ and its value is the same as $U$. Therefore, in this case we often use the notation $U(\Pvec)$, i.e.\ the strain energy, instead of $U^*(\Pvec)$. The above relation then becomes
\begin{equation}
  \pd{U}{P_i} = \Delta_i,\quad i = 1, n
\end{equation}
This is often referred to as \emph{Castigliano's second theorem}.

For a statically indeterminate conservative structural system, the total complementary potential energy can be written as
\begin{equation*}
  \Pi^* = U^* + V^* = U^*(\Pvec, \mat{R}) - \Pvec^\tr\bs{\Delta} - \mat{R}^\tr\mat{r} = U^*(\Pvec, \mat{R}) - \sum_{i=1}^nP_i\Delta_i - \sum_{i=1}^rR_ir_i
\end{equation*}
The principle of stationary (minimum) total complementary potential energy states that
\begin{equation}
  \pd{U^*}{P_i} = \Delta_i,\quad i = 1, n
\end{equation}
and
\begin{equation}
  \pd{U^*}{R_i} = r_i = 0,\quad i = 1, r
\end{equation}
The condition $r_i = 0$ is very obvious if the associated $R_i$ is a reaction because no movement is allowed at the corresponding support. If $R_i$ is an internal force, the condition means consistency of deformation (no gap, continuous slope, etc.). \emph{The second set of equations is often referred to as the theorem of least work or Engesser's second theorem.}

In matrix form, the above two sets of equations can be written as
\begin{equation}
  \Bmat{\dfrac{\partial U^*}{\partial\Pvec}\\[8pt] \dfrac{\partial U^*}{\partial\mat{R}}} = \Bmat{\bs{\Delta}\\ \mat{0}}
\end{equation}
Having obtained the above equations, we can solve for $\mat{R}$ from the second set of equations. Finally, substitute $\mat{R}$ into the first set of equations to obtain $\bs{\Delta}$.

\section{Complementary strain energy and the unit-load method}
The complementary strain energy for a \emph{linearly} elastic structural member:
\begin{enumerate}[label=(\alph*)]
  \item due to axial deformation,
    \begin{equation*}
      U^* = \int\frac{N^2}{2EA}\,dx = \frac{N^2L}{2EA}\ \text{ if } \frac{N}{EA} = \text{constant, where } N = \text{axial force.}
    \end{equation*}
  \item due to flexural deformation,
    \begin{equation*}
      U^* = \int\frac{M^2}{2EI}\,dx,\ \text{ where } M = \text{bending moment}
    \end{equation*}
  \item due to torsional deformation,
    \begin{equation*}
      U^* = \int\frac{T^2}{2GJ}\,dx = \frac{T^2L}{2GJ}\ \text{ if } \frac{T}{GJ} = \text{constant, where } T = \text{torsion.}
    \end{equation*}
\end{enumerate}
The unit-load method is an application of the principle of virtual forces. Consider a structural system consisting of \emph{linearly elastic} members, for which we like to find the displacement $\Delta_i$. Define a virtual force system that consists of only one virtual force $\delta P_i$ acting on the structural system and along the direction of $\Delta_i$. According to the PVF,
\begin{equation}
  \delta P_i\Delta_i = \int\frac{Nn}{EA}\,dx + \int\frac{Mm}{EI}\,dx + \int\frac{Tt}{GJ}\,dx
\end{equation}
where $N$, $M$, and $T =$ axial force, bending moment, and torsion due to the real external loads, and $n$, $m$, and $t =$ axial force, bending moment, and torsion due to the virtual force $\delta P_i$. As the virtual internal forces $n$, $m$, and $t$ depend on $\delta P_i$ linearly, without loss of generality we could let $\delta P_i = 1$; this results in
\begin{equation}
  \Delta_i = \int\frac{Nn}{EA}\,dx + \int\frac{Mm}{EI}\,dx + \int\frac{Tt}{GJ}\,dx
\end{equation}
where $n$, $m$, and $t =$ axial force, bending moment, and torsion due to $\delta P_i = 1$. This is the so-called \emph{unit load method}.
\begin{fig}
  % Real system (N, M, T) and virtual system (n, m, t); support conditions can be different
\begin{tikzpicture}[scale=0.9,font=\small]
  \begin{scope}
    \draw[thick] plot[smooth cycle,tension=0.9] coordinates {(0,0) (1.5,-0.9) (3.2,-0.3) (3.4,1.2) (2.0,2.0) (0.4,1.4)};
    \draw[load] (-1.1,1.2) -- (0.4,0.6) node[pos=0,above] {$P_1$};
    \draw[load] (1.6,0.4) -- (1.6,2.4) node[above] {$P_2$};
    \draw[load] (2.2,1.0) -- (3.9,1.0) node[right] {$P_3$};
    \pic[rotate=-30] at (3.3,0.2) {pin}; \pic[rotate=20] at (1.3,-0.9) {roller};
    \node[below] at (1.6,-1.6) {\itshape real $(N, M, T)$};
  \end{scope}
  \begin{scope}[xshift=7cm]
    \draw[thick] plot[smooth cycle,tension=0.9] coordinates {(0,0) (1.5,-0.9) (3.2,-0.3) (3.4,1.2) (2.0,2.0) (0.4,1.4)};
    \draw[load] (-0.4,2.2) -- (0.8,1.0) node[pos=0,above] {$\delta P$};
    \pic[rotate=-30] at (3.3,1.2) {pin}; \pic at (1.5,-0.9) {roller};
    \node[below] at (1.6,-1.6) {\itshape virtual $(n, m, t)$};
  \end{scope}
\end{tikzpicture}

\end{fig}
\noindent$***$ Support conditions can be different!

\section{Nonlinear system}
How can we use the unit load method to deal with nonlinear systems? Assume that
\begin{equation*}
  U^* = \int\frac{aM^b}{EI}\,dx
\end{equation*}
where $a$ and $b$ are constants. Then, $\dU^* = \int\dfrac{abM^{b-1}}{EI}\delta M\,dx$. Therefore the unit-load method takes the following form:
\begin{equation}
  \Delta_i = \int\frac{abM^{b-1}}{EI}m\,dx.
\end{equation}

\begin{example}[Application of P.V.F.]
\leavevmode
\begin{fig}
  % Real system (cantilever with tip load) and virtual system 1 (simply supported, load at C)
\begin{tikzpicture}[scale=0.9,font=\small]
  \begin{scope}
    \pic at (0,0) {wallL}; \draw[member] (0,0) -- (6,0);
    \node[above right] at (0,0.05) {$A$}; \node[above] at (3,0.05) {$C$}; \node[above left] at (6,0.05) {$B$};
    \draw[load] (6,0.1) -- (6,1.3) node[above] {$P$};
    \draw[dim] (0,-0.6) -- (3,-0.6) node[midway,fill=white,inner sep=1pt] {$\frac L2$}; \draw[dim] (3,-0.6) -- (6,-0.6) node[midway,fill=white,inner sep=1pt] {$\frac L2$};
    \node[below] at (3,-1.0) {\itshape real system};
  \end{scope}
  \begin{scope}[xshift=8.5cm]
    \draw[member] (0,0) -- (6,0); \pic at (0,0) {roller}; \pic at (6,0) {pin}; \node[node dot] at (0,0) {}; \node[node dot] at (6,0) {};
    \node[above] at (0.3,0.05) {$A$}; \node[above right] at (3,0.05) {$C$}; \node[above] at (5.7,0.05) {$B$};
    \draw[load] (3,0.1) -- (3,1.3) node[above] {$\delta P$};
    \draw[dim] (0,-0.8) -- (3,-0.8) node[midway,fill=white,inner sep=1pt] {$\frac L2$}; \draw[dim] (3,-0.8) -- (6,-0.8) node[midway,fill=white,inner sep=1pt] {$\frac L2$};
    \node[below] at (3,-1.2) {\itshape virtual system 1};
  \end{scope}
\end{tikzpicture}

\end{fig}
For real system: $M(x) = P(L-x)$

\emph{Case 1: virtual system 1}

\emph{The boundary conditions for the real system and the virtual system are totally different, but the P.V.F.\ still works as we shall show below.}

The virtual internal moment distribution is
\begin{align*}
  m(x) &= -\frac{\delta P}{2}x,\quad 0 < x < \frac L2\\
       &= -\frac{\delta P}{2}(L-x),\quad \frac L2 < x < L
\end{align*}
The complementary virtual work and the complementary virtual strain energy are given respectively by
\begin{gather*}
  \dW^* = -\frac{\delta P}{2}\times0 + \delta P\times v_c - \frac{\delta P}{2}\times v_B\\
  \begin{aligned}
    \dU^* &= \int_0^L\frac{M(x)m(x)}{EI}\,dx\\
          &= \frac{1}{EI}\int_0^{L/2}\Bigl[P(L-x)\Bigl(-\frac{\delta P}{2}x\Bigr)\Bigr]dx + \frac{1}{EI}\int_{L/2}^LP(L-x)\Bigl(-\frac{\delta P}{2}\Bigr)(L-x)\,dx\\
          &= -\frac{\delta P}{16}\frac{PL^3}{EI}
  \end{aligned}
\end{gather*}
Then, by the P.V.F.\ $\dW^* = \dU^*$, the following ``compatibility equation'' can be obtained.
\begin{equation*}
  v_C - \frac{v_B}{2} = -\frac{1}{16}\frac{PL^3}{EI}
\end{equation*}
Although this relation may not be useful, it is valid since the exact values are $v_B = \dfrac{PL^3}{3EI}$ and $v_C = \dfrac{5PL^3}{48EI}$.

\emph{Case 2: virtual system 2}
\begin{fig}
  % Virtual system 2: fixed at B, load at A
\begin{tikzpicture}[scale=0.9,font=\small]
  \draw[member] (0,0) -- (6,0); \pic at (6,0) {wallR};
  \node[above right] at (0,0.05) {$A$}; \node[above] at (3,0.05) {$C$}; \node[above left] at (6,0.05) {$B$};
  \draw[load] (0,0.1) -- (0,1.3) node[above] {$\delta P$};
  \draw[dim] (0,-0.6) -- (3,-0.6) node[midway,fill=white,inner sep=1pt] {$\frac L2$}; \draw[dim] (3,-0.6) -- (6,-0.6) node[midway,fill=white,inner sep=1pt] {$\frac L2$};
  \node[below] at (3,-1.0) {\itshape virtual system 2};
\end{tikzpicture}

\end{fig}
\begin{gather*}
  m(x) = \delta P\cdot x\\
  \dW^* = \delta P\cdot L\times\theta_B - \delta P\times v_B\\
  \dU^* = \int_0^L\frac{M(x)m(x)}{EI}\,dx = \frac{\delta P}{6}\frac{PL^3}{EI}
\end{gather*}
Then, by P.V.F.\ $\dW^* = \dU^*$, the following compatibility equation can be obtained.
\begin{equation*}
  -v_B + \theta_BL = \frac16\frac{PL^3}{EI}
\end{equation*}
which is true (but may not be so much useful) since $v_B = \dfrac{PL^3}{3EI}$ and $\theta_B = \dfrac{PL^2}{2EI}$.

\emph{Case 3: virtual system 3: solve the vertical displacement of $B$}
\begin{fig}
  % Virtual system 3: fixed at A, load at B
\begin{tikzpicture}[scale=0.9,font=\small]
  \pic at (0,0) {wallL}; \draw[member] (0,0) -- (6,0);
  \node[above right] at (0,0.05) {$A$}; \node[above] at (3,0.05) {$C$}; \node[above left] at (6,0.05) {$B$};
  \draw[load] (6,0.1) -- (6,1.3) node[above] {$\delta P$};
  \draw[dim] (0,-0.6) -- (3,-0.6) node[midway,fill=white,inner sep=1pt] {$\frac L2$}; \draw[dim] (3,-0.6) -- (6,-0.6) node[midway,fill=white,inner sep=1pt] {$\frac L2$};
  \node[below] at (3,-1.0) {\itshape virtual system 3};
\end{tikzpicture}

\end{fig}
\begin{gather*}
  m(x) = \delta P(L-x)\\
  \dW^* = \delta P\times v_B\\
  \dU^* = \int_0^L\frac{M(x)m(x)}{EI}\,dx = \frac{\delta P}{3}\frac{PL^3}{EI}
\end{gather*}
Then, by P.V.F.\ $\dW^* = \dU^*$, the vertical displacement of $B$ can be obtained as
\begin{equation*}
  v_B = \frac13\frac{PL^3}{EI}
\end{equation*}

\emph{Case 4: virtual system 4: solve the vertical displacement of $C$}
\begin{fig}
  % Virtual system 4: fixed at A, load at C
\begin{tikzpicture}[scale=0.9,font=\small]
  \pic at (0,0) {wallL}; \draw[member] (0,0) -- (6,0);
  \node[above right] at (0,0.05) {$A$}; \node[above right] at (3,0.05) {$C$}; \node[above left] at (6,0.05) {$B$};
  \draw[load] (3,0.1) -- (3,1.3) node[above] {$\delta P$};
  \draw[dim] (0,-0.6) -- (3,-0.6) node[midway,fill=white,inner sep=1pt] {$\frac L2$}; \draw[dim] (3,-0.6) -- (6,-0.6) node[midway,fill=white,inner sep=1pt] {$\frac L2$};
  \node[below] at (3,-1.0) {\itshape virtual system 4};
\end{tikzpicture}

\end{fig}
The virtual force system shall be chosen as that shown.

Note that the following formula is very useful when carrying out the integration of $\int_0^LM(x)m(x)\,dx$. Some other specific cases can be seen in the last page of this chapter (Table~\ref{tab:integral}), which can be considered special cases of this formula.
\begin{fig}
  % Simpson-type formula for the integral of M(x) m(x)
\begin{tikzpicture}[scale=0.9,font=\small]
  \begin{scope}
    \draw (0,0) -- (5,0);
    \draw[thick] (0,0) -- (0,0.9) .. controls (1.5,1.6) and (3.5,1.9) .. (5,2.0) -- (5,0);
    \draw (2.5,0) -- (2.5,1.63);
    \node[right] at (0.05,0.6) {$a_1$}; \node[right] at (2.55,1.2) {$a_2$}; \node[right] at (5.05,1.6) {$a_3$};
    \node[left] at (-0.1,0.5) {$M(x)$};
    \node[below] at (0,-0.05) {0}; \node[below] at (2.5,-0.05) {$\frac L2$}; \node[below] at (5,-0.05) {$L$};
  \end{scope}
  \begin{scope}[xshift=7.5cm]
    \draw (0,0) -- (5,0);
    \draw[thick] (0,0) -- (0,0.5) -- (5,2.2) -- (5,0);
    \draw (2.5,0) -- (2.5,1.35);
    \node[right] at (0.05,0.3) {$b_1$}; \node[right] at (2.55,1.0) {$b_2$}; \node[right] at (5.05,1.8) {$b_3$};
    \node[left] at (-0.1,0.5) {$m(x)$};
    \node[below] at (0,-0.05) {0}; \node[below] at (2.5,-0.05) {$\frac L2$}; \node[below] at (5,-0.05) {$L$};
  \end{scope}
\end{tikzpicture}

\end{fig}
\begin{equation}
  \int_0^LM(x)m(x)\,dx = \frac L6\bigl(a_1b_1 + 4a_2b_2 + a_3b_3\bigr)
\end{equation}
\end{example}

\begin{example}[Application of the unit load method]
\leavevmode

\begin{minipage}{0.55\linewidth}\centering% Propped cantilever with an overhang: 30 kN at the tip
\begin{tikzpicture}[scale=0.9,font=\small]
  \pic at (0,0) {wallL}; \draw[member] (0,0) -- (9,0); \pic at (6,0) {roller};
  \node[above right] at (0,0.05) {$A$}; \node[above] at (6,0.05) {$C$}; \node[above left] at (9,0.05) {$B$};
  \draw[load] (9,0.1) -- (9,1.4) node[above] {30 kN};
  \draw[dim] (0,-1.0) -- (6,-1.0) node[midway,fill=white,inner sep=1pt] {6 m}; \draw[dim] (6,-1.0) -- (9,-1.0) node[midway,fill=white,inner sep=1pt] {3 m};
  \foreach \x in {0,6,9}{\draw (\x,-0.8) -- (\x,-1.2);}
\end{tikzpicture}
\end{minipage}\hfill
\begin{minipage}{0.42\linewidth}
  $EI =$ constant
  \begin{align*}
    M(x) &= -45 + ax,\quad 0\leq x\leq 6\\
         &= c - bx,\quad 6\leq x\leq 9
  \end{align*}
\end{minipage}
\begin{enumerate}[label=(\alph*)]
  \item Determine the values of the coefficients $a$, $b$ and $c$.
  \item Determine $\theta_B$.
\end{enumerate}
\begin{solution}
\begin{enumerate}[label=(\alph*)]
  \item The value of $b = 30$ is obvious since the slope in the moment diagram is equal to the shear force. Knowing that $b = 30$ and $M(9) = 0$, we obtain $c = 270$. Finally, the moment at $x = 6$ shall be continuous, i.e.
    \begin{equation*}
      M(6) = c - bx = 270 - 30\times6 = 90 = -45 + 6a
    \end{equation*}
    Therefore, $a = 22.5$.
  \item Since the virtual system can be of any kind of boundary conditions, the simplest one (a cantilever) is
    \begin{fig}
      % Virtual system: cantilever with a unit moment at B
\begin{tikzpicture}[scale=0.9,font=\small]
  \pic at (0,0) {wallL}; \draw[member] (0,0) -- (9,0);
  \node[above right] at (0,0.05) {$A$}; \node[above] at (6,0.05) {$C$}; \node[above left] at (9,0.05) {$B$};
  \draw[->,thick] (9.5,0.5) arc (60:-120:0.5); \node[right] at (9.7,0) {$m = 1$};
  \draw[dim] (0,-1.0) -- (6,-1.0) node[midway,fill=white,inner sep=1pt] {6 m}; \draw[dim] (6,-1.0) -- (9,-1.0) node[midway,fill=white,inner sep=1pt] {3 m};
  \foreach \x in {0,6,9}{\draw (\x,-0.8) -- (\x,-1.2);}
\end{tikzpicture}

    \end{fig}
    The virtual moment distribution is $m(x) = 1$. Then, $\dW^* = \dU^*$ becomes
    \begin{equation*}
      1\times\theta_B = \frac{1}{EI}\int_0^6(-45 + 22.5x)\,dx + \frac{1}{EI}\int_6^9(270 - 30x)\,dx = \frac{270}{EI}
    \end{equation*}
\end{enumerate}
\end{solution}
\end{example}

\begin{example}[Application of the unit load method]
\leavevmode
\begin{fig}
  % Propped cantilever with a central load: real system (redundant R_B) and virtual system
\begin{tikzpicture}[scale=0.9,font=\small]
  \begin{scope}
    \pic at (0,0) {wallL}; \draw[member] (0,0) -- (6,0); \pic at (6,0) {roller}; \node[node dot] at (6,0) {};
    \node[above right] at (0,0.05) {$A$}; \node[above left] at (3,0.05) {$C$}; \node[above left] at (6,0.05) {$B$};
    \draw[load] (3,1.3) -- (3,0.1) node[pos=0,above] {$P$};
    \draw[load] (6,-1.6) -- (6,-0.7) node[pos=0,below] {$R_B$};
    \draw[dim] (0,-1.1) -- (3,-1.1) node[midway,fill=white,inner sep=1pt] {$\frac L2$}; \draw[dim] (3,-1.1) -- (6,-1.1) node[midway,fill=white,inner sep=1pt] {$\frac L2$};
    \node[below] at (2.5,-1.6) {\itshape real system};
  \end{scope}
  \begin{scope}[xshift=8.5cm]
    \pic at (0,0) {wallL}; \draw[member] (0,0) -- (6,0);
    \node[above right] at (0,0.05) {$A$}; \node[above] at (3,0.05) {$C$}; \node[above left] at (6,0.05) {$B$};
    \draw[load] (6,0.1) -- (6,1.3) node[above] {1};
    \draw[dim] (0,-1.1) -- (3,-1.1) node[midway,fill=white,inner sep=1pt] {$\frac L2$}; \draw[dim] (3,-1.1) -- (6,-1.1) node[midway,fill=white,inner sep=1pt] {$\frac L2$};
    \node[below] at (3,-1.6) {\itshape virtual system};
  \end{scope}
\end{tikzpicture}

\end{fig}
Take the reaction at the roller as the redundant force. Then the moment distribution for the real system is
\begin{align*}
  M(x) &= R_B(L-x),\quad \frac L2 < x < L\\
       &= R_B(L-x) - P\Bigl(\frac L2 - x\Bigr),\quad 0 < x < \frac L2
\end{align*}
From the virtual system, we have $m(x) = 1\times(L-x)$. According to the unit load method
\begin{equation*}
  v_B = \int_0^L\frac{Mm}{EI}\,dx = 0 \;\Rightarrow\; 0 = \frac{5PL^3}{48EI} - \frac{R_BL^3}{3EI} \;\Rightarrow\; R_B = \frac{5}{16}P
\end{equation*}
\begin{fig}
  % Reactions of the propped cantilever
\begin{tikzpicture}[scale=0.9,font=\small]
  \pic at (0,0) {wallL}; \draw[member] (0,0) -- (6,0); \pic at (6,0) {pin}; \node[node dot] at (6,0) {};
  \draw[load] (3,1.3) -- (3,0.1) node[pos=0,above] {$P$};
  \draw[load] (0,-1.6) -- (0,-0.5) node[pos=0,below] {$\frac{11}{16}P$};
  \draw[load] (6,-1.6) -- (6,-0.7) node[pos=0,below] {$\frac{5}{16}P$};
  \draw[->,thick] (-0.9,0.6) arc (110:-40:0.5); \node[left] at (-0.9,0.5) {$\frac{PL}{8}\times\frac32$};
\end{tikzpicture}

\end{fig}
\end{example}

\begin{example}[Application of Engesser's 1st theorem]
\leavevmode
\begin{fig}
  % Truss for Engesser's first theorem and the nonlinear member response
\begin{tikzpicture}[scale=0.9,font=\small]
  \begin{scope}
    \coordinate (a) at (0,0); \coordinate (c) at (3.2,0); \coordinate (b) at (1.6,-2.77); \coordinate (d) at (4.8,-2.77);
    \draw[member] (a) -- (c) -- (b) -- (a); \draw[member] (c) -- (d);
    \pic at (b) {pin}; \pic at (d) {pin};
    \node[above] at (a) {$a$}; \node[above] at (c) {$c$}; \node[below left] at (b) {$b$}; \node[below right] at (d) {$d$};
    \draw[load] (a) -- ++(-1.4,0) node[left] {$P$};
    \node[above] at (1.6,0.05) {$(P)$}; \node[left] at (0.6,-1.4) {$(0)$}; \node[right] at (2.5,-1.4) {$(-P)$}; \node[right] at (4.1,-1.4) {$(P)$};
    \draw[->] ($(a)+(1.0,0)$) arc (0:-60:1.0); \node[right] at ($(a)+(0.9,-0.6)$) {$60^\circ$};
  \end{scope}
  \begin{scope}[xshift=8cm,yshift=-2.8cm]
    \draw[->] (0,0) -- (5.2,0) node[right] {$\delta$}; \draw[->] (0,0) -- (0,4.2) node[above] {$f$};
    \draw[thick,domain=0:4.6,samples=40] plot (\x,{2.3*sqrt(\x/1.0)*0.85});
    \pgfmathsetmacro{\F}{2.3*sqrt(1.4)*0.85}
    \fill[pattern={Lines[angle=45,distance=3pt]}] (0,0) -- plot[domain=0:1.4,samples=30] (\x,{2.3*sqrt(\x)*0.85}) -- (0,\F) -- cycle;
    \draw[dashed] (0,\F) -- (1.4,\F); \node[left] at (0,\F) {$F$};
    \draw[dashed] (0,{\F*0.62}) -- (0.55,{\F*0.62}); \draw[dashed] (0,{\F*0.62+0.25}) -- (0.62,{\F*0.62+0.25});
    \draw[dim] (-0.25,{\F*0.62}) -- (-0.25,{\F*0.62+0.25}); \node[left] at (-0.3,{\F*0.62+0.12}) {$df$};
    \draw[dim] (0,-0.3) -- (0.55,-0.3); \node[below] at (0.3,-0.35) {$\delta$};
    \node[right] at (1.8,1.6) {$\delta = K_1f + K_2f^2$};
  \end{scope}
\end{tikzpicture}

\end{fig}
The complementary strain energy for a member:
\begin{equation*}
  U^*(F) = \int_0^F\delta(f)\,df = \int_0^F(K_1f + K_2f^2)\,df = \tfrac12K_1F^2 + \tfrac13K_2F^3\quad\text{(the shaded area)}
\end{equation*}
The structure is statically determinate and the member forces are already shown in the left figure. The complementary strain energy of the system is
\begin{equation*}
  U^* = \sum U^*_j = \tfrac12K_1\bigl(P^2 + (-P)^2 + P^2\bigr) + \tfrac13K_2\bigl(P^3 + (-P)^3 + P^3\bigr) = \tfrac32K_1P^2 + \tfrac13K_2P^3
\end{equation*}
By Engesser's 1st theorem, the horizontal displacement ($\leftarrow$) at the loaded point is
\begin{equation*}
  \Delta_a = \pd{U^*}{P} = 3K_1P + K_2P^2\quad(\leftarrow)
\end{equation*}
\end{example}

\begin{example}[3D structure]
\leavevmode
\begin{fig}
  % 3D L-shaped structure: cantilever AB with a perpendicular arm BC loaded by P at C
\begin{tikzpicture}[scale=1.0,font=\small,x={(1cm,0cm)},y={(0cm,1cm)},z={(-0.55cm,-0.85cm)}]
  \begin{scope}
    \def\h{0.5} \def\w{0.5}
    % member AB along +x from A(0,0,0) to B(3,0,0)
    \draw (0,0,0) -- (3,0,0) -- (3,\h,0) -- (0,\h,0) -- cycle;
    \draw (3,0,0) -- (3,0,\w) -- (3,\h,\w) -- (3,\h,0); \draw (0,\h,0) -- (0,\h,\w) -- (3,\h,\w);
    % member BC along +z from B to C(3,0,4.5)
    \draw (3,0,\w) -- (3,0,4.5) -- (3,\h,4.5) -- (3,\h,\w); \draw (3,\h,4.5) -- (2.5,\h,4.5) -- (2.5,0,4.5) -- (3,0,4.5);
    \draw (2.5,\h,4.5) -- (2.5,\h,\w);
    % wall at A
    \fill[hatch] (-0.25,-0.3,0) -- (0,-0.3,0) -- (0,\h+0.3,0) -- (-0.25,\h+0.3,0) -- cycle;
    \node[above left] at (0,\h,0) {$A$}; \node[above] at (3,\h,0) {$B$}; \node[below] at (2.75,0,4.6) {$C$};
    \draw[->] (2.2,\h+0.15,0) -- (1.0,\h+0.15,0) node[above] {$x'$};
    \draw[->] (2.75,\h+0.1,3.0) -- (2.75,\h+0.1,1.8) node[above right] {$x'$};
    \draw[load] (2.75,\h+1.5,4.2) -- (2.75,\h+0.15,4.2) node[pos=0,left] {$P$};
    \draw[dim] (0,\h+0.9,0) -- (3,\h+0.9,0) node[midway,above] {$L$};
    \draw[dim] (3.6,0,\w) -- (3.6,0,4.5) node[midway,right] {$2L$};
  \end{scope}
  % section
  \begin{scope}[shift={(7.5,0.5,0)}]
    \draw (-0.5,-0.5,0) rectangle (0.5,0.5,0); \draw (-0.35,-0.35,0) rectangle (0.35,0.35,0);
    \draw[->,thick] (0,0,0) -- (0,2.0,0) node[above] {$I_2\,(z')$};
    \draw[->,thick] (0,0,0) -- (2.0,0,0) node[right] {$I_1\,(y')$};
    \draw[->,thick] (0,0,0) -- (0,0,2.2) node[below] {$J\,(x')$};
    \node[right] at (0.4,-1.0,0) {section};
  \end{scope}
\end{tikzpicture}

\end{fig}
Determine the vertical displacement at node $C$.
\begin{solution}
\leavevmode
\begin{fig}
  % Free bodies of members AB and BC
\begin{tikzpicture}[scale=0.9,font=\small]
  \begin{scope}
    \node[left] at (-0.8,0) {Member $AB$:};
    \pic at (0,0) {wallL}; \draw (0,-0.25) rectangle (5,0.25);
    \node[above right] at (0,0.25) {$A$}; \node[above left] at (5,0.25) {$B$};
    \draw[->] (4.3,0) -- (2.8,0) node[below] {$x'$};
    \draw[load] (5,1.4) -- (5,0.3) node[pos=0,above] {$P$};
    \draw[->,thick,double] (5.2,0) -- (6.4,0) node[right] {$2PL$};
    \draw[->] (0,-0.3) -- (0,-1.5) node[below] {$z'$};
  \end{scope}
  \begin{scope}[yshift=-4.8cm]
    \node[left] at (-0.8,0) {Member $BC$:};
    \draw (0,-0.25) rectangle (7,0.25);
    \node[above right] at (0,0.25) {$C$}; \node[above left] at (7,0.25) {$B$};
    \draw[->] (0.2,0) -- (1.7,0) node[below] {$x'$};
    \draw[load] (0,1.4) -- (0,0.3) node[pos=0,above] {$P$};
    \draw[->] (0,-0.3) -- (0,-1.5) node[below] {$z'$};
  \end{scope}
\end{tikzpicture}

\end{fig}
Member $AB$: $M(x') = Px'$ (w.r.t.\ the $y'$ axis), $T(x') = 2PL$.\\
Member $BC$: $M(x') = Px'$ (w.r.t.\ the $y'$ axis), $T(x') = 0$.

According to the unit load method
\begin{equation*}
  1\cdot v_C = \int_0^L\frac{(Px)x}{EI_1}\,dx + \int_0^L\frac{(2PL)2L}{GJ}\,dx + \int_0^{2L}\frac{(Px)x}{EI_1}\,dx = 3\frac{PL^3}{EI_1} + 4\frac{PL^3}{GJ}
\end{equation*}
\end{solution}
\end{example}

\begin{example}[Deflection of truss due to external loads]
Calculate the horizontal displacement of node $C$, $D_1$, and the \emph{relative displacement} of nodes $B$ and $E$, $D_2$. Note that the cross section for members 6 and 7 $= 1.25A$ and that for other members is $A$. (Ghali and Neville 1997, \emph{Structural Analysis})
\begin{fig}
  % Truss example (Ghali & Neville): geometry, real forces, and virtual forces
\begin{tikzpicture}[scale=0.62,font=\scriptsize]
  \newcommand{\trussgeo}{\coordinate (A) at (0,0); \coordinate (B) at (4,0); \coordinate (C) at (8,0); \coordinate (D) at (4,3); \coordinate (E) at (0,3);
    \draw[thinmember] (A) -- (C); \draw[thinmember] (E) -- (D) -- (C); \draw[thinmember] (A) -- (E); \draw[thinmember] (A) -- (D); \draw[thinmember] (B) -- (D);
    \foreach \p in {A,B,C,D,E}{\fill (\p) circle (0.06);}}
  % (a)
  \begin{scope}
    \trussgeo \pic at (A) {pin}; \pic at (C) {roller};
    \node[left] at (A) {$A$}; \node[below] at (B) {$B$}; \node[below right] at (C) {$C$}; \node[above right] at (D) {$D$}; \node[above left] at (E) {$E$};
    \newcommand{\mem}[2]{\node[circle,draw,fill=white,inner sep=0.8pt,font=\tiny] at #1 {#2};}
    \mem{(2,0)}{1} \mem{(6,0)}{2} \mem{(0,1.5)}{3} \mem{(4,1.5)}{4} \mem{(2,3)}{5} \mem{(6,1.5)}{6} \mem{(1.5,1.4)}{7}
    \draw[load] (-1.3,3) -- (-0.1,3) node[pos=0,above] {$P$}; \draw[load] (4,4.3) -- (4,3.1) node[pos=0,right] {$P$};
    \draw[load] (8.1,0) -- (9.1,0) node[right] {$D_1$}; \node[below] at (8.6,0) {1};
    \draw[load] (E) -- ++(-40:0.9) node[below right] {$\frac{D_2}{2}$}; \draw[load] (B) -- ++(140:0.9) node[above left] {$\frac{D_2}{2}$};
    \draw[dim] (0,-1.0) -- (4,-1.0) node[midway,fill=white,inner sep=1pt] {$l$}; \draw[dim] (4,-1.0) -- (8,-1.0) node[midway,fill=white,inner sep=1pt] {$l$};
    \node[below] at (4,-1.5) {(a) 2D truss};
  \end{scope}
  % (b)
  \begin{scope}[xshift=12cm]
    \trussgeo \pic at (A) {pin}; \pic at (C) {roller};
    \draw[load] (-1.3,3) -- (-0.1,3) node[pos=0,above] {$P$}; \draw[load] (4,4.3) -- (4,3.1) node[pos=0,right] {$P$};
    \draw[load] (-1.3,0) -- (-0.1,0) node[pos=0,above] {$P$};
    \node[above] at (2,3) {$-P$}; \node[left] at (0,1.5) {$0$}; \node[right] at (4,1.5) {$0$}; \node[below] at (2,0) {$1.167P$}; \node[below] at (6,0) {$1.167P$};
    \node[above,sloped] at (2,1.5) {$-0.208P$}; \node[above,sloped] at (6,1.5) {$-1.458P$};
    \draw[load] (0,-1.4) -- (0,-0.6) node[pos=0,below] {$0.125P$}; \draw[load] (8,-1.4) -- (8,-0.7) node[pos=0,below] {$0.875P$};
    \draw[dim] (-1.6,0) -- (-1.6,3) node[midway,left] {$0.75l$};
    \node[below] at (4,-2.0) {(b) (real) internal forces $N_i$};
  \end{scope}
  % (c)
  \begin{scope}[yshift=-7.5cm]
    \trussgeo \pic at (A) {pin}; \pic at (C) {roller};
    \draw[load] (8.1,0) -- (9.1,0) node[above] {1};
    \node[above] at (2,3) {$0$}; \node[left] at (0,1.5) {$0$}; \node[right] at (4,1.5) {$0$}; \node[below] at (2,0) {$-1$}; \node[below] at (6,0) {$-1$};
    \node[above,sloped] at (2,1.5) {$0$}; \node[above,sloped] at (6,1.5) {$0$};
    \node[below] at (4,-1.5) {(c) virtual forces $n_i$ (due to $\delta P_{D1} = 1$)};
  \end{scope}
  % (d)
  \begin{scope}[xshift=12cm,yshift=-7.5cm]
    \trussgeo \pic at (A) {pin}; \pic at (C) {roller};
    \draw[load] (E) -- ++(-36.87:1.0) node[below right] {1}; \draw[load] (B) -- ++(143.13:1.0) node[above left] {1};
    \node[above] at (2,3) {$-0.8$}; \node[left] at (0,1.5) {$-0.6$}; \node[right] at (4,1.5) {$-0.6$}; \node[below] at (2,0) {$-0.8$}; \node[below] at (6,0) {$0$};
    \node[above,sloped] at (2.9,2.2) {$1$}; \node[above,sloped] at (6,1.5) {$0$};
    \node[below] at (4,-1.5) {(d) virtual forces $n_i$ (due to $\delta P_{D2} = 1$)};
  \end{scope}
\end{tikzpicture}

\end{fig}
\begin{gather*}
  \dU^* = \sum_{i=1}^{elem}n_i\delta_i\quad\text{where } \delta_i = \frac{N_i}{(EA/L)_i} = \frac{N_iL_i}{E_iA_i}\ \text{(real axial deformation)}\\
  \dW^* = \delta P_i\Delta_i\ \text{(no summation)}\\
  \text{PVF: } \dW^* = \dU^*
\end{gather*}
Calculation of $D_1$:
\begin{align*}
  (1)D_1 &= (-1)\frac{1.167P(L)}{EA} + (-1)\frac{1.167P(L)}{EA} + (0)(0) + (0)(0) + (0)\frac{-P(L)}{EA}\\
         &\quad + (0)\frac{-1.458P(1.25L)}{E(1.25A)} + (0)\frac{-0.208P(1.25L)}{E(1.25A)} = -2.334\frac{PL}{EA}
  \;\rightarrow\; D_1 = -2.334\frac{PL}{EA}
\end{align*}
Calculation of $D_2$:
\begin{align*}
  (1)D_2 &= (-0.8)\frac{1.167P(L)}{EA} + (0)\frac{1.167P(L)}{EA} + (-0.6)(0) + (-0.6)(0) + (-0.8)\frac{-P(L)}{EA}\\
         &\quad + (0)\frac{-1.458P(1.25L)}{E(1.25A)} + (1)\frac{-0.208P(1.25L)}{E(1.25A)} = -0.342\frac{PL}{EA}
  \;\rightarrow\; D_2 = -0.342\frac{PL}{EA}
\end{align*}
\end{example}

\begin{example}[Deflection of truss due to temperature/fabrication errors]
Consider the above truss but \emph{without applied loads}. Assume member 2 has a temperature rise of 30 degrees and members 5 and 6 have a temperature rise of 20 degrees. Also, assume the coefficient of thermal expansion $\alpha = 0.5\times10^{-6}/{}^\circ C$. Calculate the displacements $D_1$ and $D_2$.
\begin{gather*}
  \text{PVF: } \dW^* = \dU^*\\
  \dW^* = \delta P_i\Delta_i\quad\text{and}\quad \dU^* = \sum_{i=1}^{elem}n_i\delta_i
\end{gather*}
where $\delta = \Delta L = (\alpha\Delta T)L$ for temperature change and $\delta = \Delta L =$ fabrication error. Therefore, in this example
\begin{align*}
  \delta_2 &= \alpha\Delta T_2L_2 = (0.5\times10^{-5})(30)(L) = 15\times10^{-5}L\\
  \delta_5 &= \alpha\Delta T_5L_5 = (0.5\times10^{-5})(20)(L) = 10\times10^{-5}L\\
  \delta_6 &= \alpha\Delta T_6L_6 = (0.5\times10^{-5})(20)(1.25L) = 12.5\times10^{-5}L
\end{align*}
Calculation of $D_1$:
\begin{gather*}
  (1)(D_1) = n_2\delta_2 + n_5\delta_5 + n_6\delta_6 = (-1)(15\times10^{-5}L) + (0)(10\times10^{-5}L) + (0)(12.5\times10^{-5}L)\\
  \rightarrow\; D_1 = -15\times10^{-5}L
\end{gather*}
Calculation of $D_2$:
\begin{gather*}
  (1)(D_2) = n_2\delta_2 + n_5\delta_5 + n_6\delta_6 = (0)(15\times10^{-5}L) + (-0.8)(10\times10^{-5}L) + (0)(12.5\times10^{-5}L)\\
  \rightarrow\; D_2 = -8\times10^{-5}L
\end{gather*}
(The virtual forces $n_i$ due to $\delta P_{D1} = 1$ and $\delta P_{D2} = 1$ are those shown in (c) and (d) of the previous figure.)
\end{example}

\begin{example}[Deflection of frame due to support settlement]
\begin{equation*}
  \dW^* = (\delta P)(\Delta) + \sum_{Support}(\delta R_x\Delta_x + \delta R_y\Delta_y + \delta M\theta)
\end{equation*}
Consider the following frame with $E = 2\times10^8$~kN/m$^2$ and $I = 6\times10^{-4}$~m$^4$. Assume node $E$ has a horizontal settlement of 0.02~m to the right and node $A$ has a downward settlement of 0.03~m. Calculate the horizontal displacement of node $B$.
\begin{fig}
  % Portal frame with support settlements and a horizontal load at B
\begin{tikzpicture}[scale=0.7,font=\small]
  \draw[member] (0,0) -- (0,5) -- (10,5) -- (10,0);
  \pic at (0,0) {pin}; \pic at (10,0) {pin}; \node[node dot] at (0,0) {}; \node[node dot] at (10,0) {}; \node[node dot] at (5,5) {};
  \node[left] at (-0.3,0.1) {$A$}; \node[above left] at (0,5.05) {$B$}; \node[above] at (5,5.1) {$C$}; \node[above right] at (10,5.05) {$D$}; \node[right] at (10.3,0.1) {$E$};
  \draw[load] (-2.5,5) -- (-0.15,5) node[pos=0,left] {100 kN};
  \node[below] at (2.5,4.9) {5 m}; \node[below] at (7.5,4.9) {5 m}; \node[left] at (-0.1,2.5) {5 m};
\end{tikzpicture}

\end{fig}
\begin{fig}
  % Real moments M (100 kN) and virtual moments m (unit load) for the portal frame
\begin{tikzpicture}[scale=0.48,font=\scriptsize]
  \newcommand{\framemom}[5]{% #1 load label, #2 corner moment, #3 mid moment label, #4 reaction H, #5 reaction V
    \draw[thinmember] (0,0) -- (0,5) -- (10,5) -- (10,0);
    % column AB: triangular from 0 at A to #2 at B (drawn to the left)
    \draw (0,0) -- (-1.5,5) -- (0,5);
    \draw (0,5) -- (0,6.5) -- (10,3.5) -- (10,5);
    \draw (10,0) -- (11.5,5) -- (10,5);
    \node[above] at (-0.2,6.5) {#2}; \node[left] at (-1.5,4.7) {#2}; \node[right] at (10.4,3.5) {#2}; \node[right] at (11.5,4.7) {#2};
    \draw[load] (-3.0,5) -- (-1.6,5) node[pos=0,left] {#1};
    \draw[load] (-0.3,-0.9) -- (-1.5,-0.9) node[pos=0.5,below] {#4}; \draw[load] (0,-1.3) -- (0,-2.3) node[below] {#4};
    \draw[load] (11.5,-0.9) -- (10.3,-0.9) node[pos=0.5,below] {#4}; \draw[load] (10,-2.3) -- (10,-1.3) node[pos=0,below] {#4};
    \node at (5,2.5) {\large #3};}
  \begin{scope} \framemom{100}{250}{$M$}{50}{50} \node[below] at (5,-3.2) {(a) real moments}; \end{scope}
  \begin{scope}[xshift=17cm] \framemom{1}{2.5}{$m$}{0.5}{0.5} \node[below] at (5,-3.2) {(b) virtual moments due to $\delta P = 1$}; \end{scope}
\end{tikzpicture}

\end{fig}
\begin{gather*}
  \dW^* = 1\cdot u_B + (-0.5)(-0.03) + (-0.5)(0.02) = u_B + 0.005\\
  \dU^* = \int\frac{Mm\,dx}{EI} = \frac{1041.67}{EI} + \frac{2083.33}{EI} + \frac{1041.67}{EI} = \frac{4166.7}{EI} = \frac{4166.7}{12\times10^4} = 0.0347\\
  \text{(can be calculated using Table~\ref{tab:integral})}\\
  \dW^* = \dU^* \;\Rightarrow\; u_B = 0.0297\text{ m}
\end{gather*}
\end{example}

\clearpage
\begin{table}[H]
  \centering
  \caption{Values of the integral $\int M_uM\,dl$ (Ghali and Neville 1997, \emph{Structural Analysis}).}
  \label{tab:integral}
  % Values of the integral \int M_u M dl (after Ghali & Neville 1997)
% Small shape macros (M diagrams over length l)
\newcommand{\shpRect}[1]{\begin{tikzpicture}[scale=0.55,baseline=(current bounding box.center),font=\tiny]
  \draw[thick] (0,0) rectangle (2.4,0.7); \node[right] at (2.4,0.35) {$#1$}; \draw[dim] (0,-0.35) -- (2.4,-0.35) node[midway,fill=white,inner sep=0.5pt] {$l$};\end{tikzpicture}}
\newcommand{\shpTriR}[1]{\begin{tikzpicture}[scale=0.55,baseline=(current bounding box.center),font=\tiny]
  \draw[thick] (0,0) -- (2.4,0) -- (2.4,0.8) -- cycle; \node[right] at (2.4,0.4) {$#1$}; \draw[dim] (0,-0.35) -- (2.4,-0.35) node[midway,fill=white,inner sep=0.5pt] {$l$};\end{tikzpicture}}
\newcommand{\shpTriL}[1]{\begin{tikzpicture}[scale=0.55,baseline=(current bounding box.center),font=\tiny]
  \draw[thick] (0,0) -- (2.4,0) -- (0,0.8) -- cycle; \node[left] at (0,0.4) {$#1$}; \draw[dim] (0,-0.35) -- (2.4,-0.35) node[midway,fill=white,inner sep=0.5pt] {$l$};\end{tikzpicture}}
\newcommand{\shpTrap}[2]{\begin{tikzpicture}[scale=0.55,baseline=(current bounding box.center),font=\tiny]
  \draw[thick] (0,0) -- (2.4,0) -- (2.4,1.0) -- (0,0.45) -- cycle; \node[left] at (0,0.25) {$#1$}; \node[right] at (2.4,0.6) {$#2$}; \draw[dim] (0,-0.35) -- (2.4,-0.35) node[midway,fill=white,inner sep=0.5pt] {$l$};\end{tikzpicture}}
\newcommand{\shpPeak}[1]{\begin{tikzpicture}[scale=0.55,baseline=(current bounding box.center),font=\tiny]
  \draw[thick] (0,0) -- (2.4,0) -- (0.9,0.9) -- cycle; \draw (0.9,0) -- (0.9,0.9); \node[right] at (0.9,0.45) {$#1$};
  \draw[dim] (0,1.2) -- (0.9,1.2) node[midway,fill=white,inner sep=0.5pt] {$\alpha l$}; \draw[dim] (0.9,1.2) -- (2.4,1.2) node[midway,fill=white,inner sep=0.5pt] {$\beta l$};
  \draw[dim] (0,-0.35) -- (2.4,-0.35) node[midway,fill=white,inner sep=0.5pt] {$l$};\end{tikzpicture}}
\newcommand{\shpPara}[1]{\begin{tikzpicture}[scale=0.55,baseline=(current bounding box.center),font=\tiny]
  \draw[thick] (0,0) -- (2.4,0); \draw[thick,domain=0:2.4,samples=30] plot (\x,{0.9*(1-((\x-1.2)/1.2)^2)}); \draw (1.2,0) -- (1.2,0.9); \node[right] at (1.2,0.45) {$#1$}; \node at (2.6,0.9) {$*$};
  \draw[dim] (0,-0.35) -- (2.4,-0.35) node[midway,fill=white,inner sep=0.5pt] {$l$};\end{tikzpicture}}
\newcommand{\shpParaA}[1]{\begin{tikzpicture}[scale=0.55,baseline=(current bounding box.center),font=\tiny]
  \draw[thick] (0,0) -- (2.4,0) -- (2.4,0.9); \draw[thick,domain=0:2.4,samples=30] plot (\x,{0.9*(\x/2.4)^2}); \draw[thin] (-0.6,0) -- (0.8,0); \node[right] at (2.4,0.45) {$#1$}; \node at (2.6,1.0) {$*$};
  \node[above,font=\tiny] at (0.1,0.05) {tangent}; \draw[dim] (0,-0.35) -- (2.4,-0.35) node[midway,fill=white,inner sep=0.5pt] {$l$};\end{tikzpicture}}
\newcommand{\shpParaB}[1]{\begin{tikzpicture}[scale=0.55,baseline=(current bounding box.center),font=\tiny]
  \draw[thick] (0,0) -- (2.4,0) -- (2.4,0.9); \draw[thick,domain=0:2.4,samples=30] plot (\x,{0.9*(1-(1-\x/2.4)^2)}); \draw[thin] (1.4,0.9) -- (3.0,0.9); \node[right] at (2.4,0.45) {$#1$}; \node at (2.6,1.15) {$*$};
  \node[above,font=\tiny] at (0.5,0.55) {tangent}; \draw[dim] (0,-0.35) -- (2.4,-0.35) node[midway,fill=white,inner sep=0.5pt] {$l$};\end{tikzpicture}}
{\footnotesize
\renewcommand{\arraystretch}{1.6}
\begin{tabular}{|c|c|c|c|c|c|}\hline
  \diagbox[width=2.6cm]{$M$}{$M_u$} & \shpRect{b} & \shpTriR{b} & \shpTriL{b} & \shpTrap{b_1}{b_2} & \shpPeak{b}\\\hline
  \shpRect{a} & $abl$ & $\frac12abl$ & $\frac12abl$ & $\frac{al}{2}(b_1+b_2)$ & $\frac12abl$\\\hline
  \shpTrap{a_1}{a_2} & $\frac{bl}{2}(a_1+a_2)$ & $\frac{bl}{6}(a_1+2a_2)$ & $\frac{bl}{6}(2a_1+a_2)$ & \makecell{$\frac l6(2a_1b_1+a_1b_2+$\\$a_2b_1+2a_2b_2)$} & \makecell{$\frac{bl}{6}[(1+\beta)a_1$\\$+(1+\alpha)a_2]$}\\\hline
  \shpPeak{a} & $\frac12abl$ & $\frac{abl}{6}(1+\alpha)$ & $\frac{abl}{6}(1+\beta)$ & \makecell{$\frac{al}{6}[(1+\beta)b_1+$\\$(1+\alpha)b_2]$} & $\frac13abl$\\\hline
  \shpPara{a} & $\frac23abl$ & $\frac13abl$ & $\frac13abl$ & $\frac13al(b_1+b_2)$ & $\frac{abl}{3}(1+\alpha\beta)$\\\hline
  \shpParaA{a} & $\frac13abl$ & $\frac14abl$ & $\frac1{12}abl$ & $\frac{al}{12}(b_1+3b_2)$ & $\frac{abl}{12}(1+\alpha+\alpha^2)$\\\hline
  \shpParaB{a} & $\frac23abl$ & $\frac5{12}abl$ & $\frac14abl$ & $\frac{al}{12}(3b_1+5b_2)$ & $\frac{abl}{12}(5-\beta-\beta^2)$\\\hline
\end{tabular}\\[2pt]
{\scriptsize $*$ Second-degree parabola.}}

\end{table}

\chapter{Flexibility Method}
\label{ch:flexmethod}

\section{Complementary strain energy}
\begin{equation}
  U^* = \sum_{i=1}^m\tfrac12\F_i^\tr\mat{r}_i = \sum_{i=1}^m\tfrac12\F_i^\tr\mat{d}_i\F_i = \tfrac12\F'^\tr\mat{d}'\F'
\end{equation}
where $\F_i$ is the independent force vector (we use the notation $\F_f$ in Ch.~\ref{ch:flex}) for the $i$-th member, $\mat{r}_i$ is the associated deformation vector, and $\mat{d}_i$ is the flexibility matrix. The definitions of $\F'$ and $\mat{d}'$ are given as follows:
\begin{equation}
  \F' = \Bmat{\F_1\\ \F_2\\ \vdots\\ \F_m},\quad\text{and}\quad \mat{d}' = \bmat{\mat{d}_1 & & & \\ & \mat{d}_2 & & \\ & & \ddots & \\ & & & \mat{d}_m}
\end{equation}
Note that $\F_i$ may include the reaction(s) of a support; in this case, $\mat{d}_i = \mat{0}$ if the support is rigid.

In the flexibility method, certain redundant forces, with the number being equal to redundancy, need to be chosen from $\F'$ first. Let these redundant forces be denoted by $\F_x$. Group $\F'$ into $\F_x$ and the other forces to yield
\begin{equation}
  \F = \Bmat{\F_0\\ \F_x}
\end{equation}
where $\F_0 = \F'$ excludes $\F_x$.

The complementary strain energy, written in terms of $\F$, takes the form
\begin{equation}
  U^* = \tfrac12\F^\tr\mat{d}\F
\end{equation}
The flexibility matrix $\mat{d}$ corresponding to $\F$ can also be obtained by rearranging $\mat{d}'$. Let $\mat{d}$ be expressed as
\begin{equation}
  \mat{d} = \bmat{\mat{d}_{00} & \mat{d}_{0x}\\ \mat{d}_{x0} & \mat{d}_{xx}}
\end{equation}
Note that:
\begin{enumerate}
  \item $\mat{d}$ is symmetric.
  \item $\mat{d}_{0x} = \mat{d}_{x0} = \mat{0}$ for trusses.
  \item In general, $\mat{d}_{0x} = \mat{d}_{x0}\neq\mat{0}$ for beams.
\end{enumerate}

\section{Equilibrium equations}
On the basis of the equilibrium equations for the primary structure (a statically determinate structure), we can write $\F_0$ in terms of $\Pvec_0$ and $\F_x$ as follows.
\begin{equation}
  \F_0 = \mat{C}_1\Pvec_0 + \mat{C}_2\F_x
\end{equation}
where $\Pvec_0 =$ external load vector, and $\mat{C}_1$ and $\mat{C}_2$ are related to the geometry of structures. The $i$-th column of $\mat{C}_1$ represents the forces $\F_0$ due to $(\Pvec_0)_i = 1$ and the $i$-th column of $\mat{C}_2$ represents the forces $\F_0$ due to $(\F_x)_i = 1$.

With the above equilibrium equations, we can write $\F$ in the following form.
\begin{equation}
  \F = \Bmat{\F_0\\ \F_x} = \bmat{\mat{C}_1 & \mat{C}_2\\ \mat{0} & \mat{I}}\Bmat{\Pvec_0\\ \F_x} = \mat{C}\Pvec
\end{equation}
where
\begin{equation}
  \mat{C} = \bmat{\mat{C}_1 & \mat{C}_2\\ \mat{0} & \mat{I}}\quad\text{and}\quad \Pvec = \Bmat{\Pvec_0\\ \F_x}
\end{equation}

\section{Compatibility equations}
Substituting the above expression for $\F$ into $U^*$ results in
\begin{equation}
  U^* = \tfrac12\F^\tr\mat{d}\F = \tfrac12(\mat{C}\Pvec)^\tr\mat{d}(\mat{C}\Pvec) = \tfrac12\Pvec^\tr\mat{D}\Pvec
\end{equation}
where the structural flexibility matrix $\mat{D}$ is defined as
\begin{equation}
  \mat{D} = \mat{C}^\tr\mat{d}\mat{C} = \bmat{\mat{C}_1^\tr & \mat{0}\\ \mat{C}_2^\tr & \mat{I}}\bmat{\mat{d}_{00} & \mat{d}_{0x}\\ \mat{d}_{x0} & \mat{d}_{xx}}\bmat{\mat{C}_1 & \mat{C}_2\\ \mat{0} & \mat{I}} = \bmat{\mat{D}_{pp} & \mat{D}_{px}\\ \mat{D}_{xp} & \mat{D}_{xx}}
\end{equation}
where
\begin{align}
  \mat{D}_{pp} &= \mat{C}_1^\tr\mat{d}_{00}\mat{C}_1\\
  \mat{D}_{px} &= \mat{C}_1^\tr(\mat{d}_{00}\mat{C}_2 + \mat{d}_{0x}) = \mat{D}_{xp}^\tr\\
  \mat{D}_{xx} &= \mat{C}_2^\tr(\mat{d}_{00}\mat{C}_2 + \mat{d}_{0x}) + \mat{d}_{x0}\mat{C}_2 + \mat{d}_{xx}
\end{align}
Note that $\mat{D}$ is symmetric.

According to Castigliano's second theorem, we obtain the following compatibility equations:
\begin{equation}
  \pd{U^*}{P_{0i}} = \Delta_{0i},\quad i = 1, n
\end{equation}
where $n =$ dimension of $\Pvec_0$.

On the other hand we obtain, according to the theorem of least work, the following compatibility equations:
\begin{equation}
  \pd{U^*}{F_{xi}} = 0,\quad i = 1, r
\end{equation}
where $r =$ the number of redundant forces.

The above two sets of equations can be expressed, in vector form, as
\begin{equation}
  \pd{U^*}{\Pvec} = \Bmat{\bs{\Delta}_0\\ \mat{0}}
\end{equation}
Knowing that $U^* = \frac12\Pvec^\tr\mat{D}\Pvec$, therefore $\partial U^*/\partial\Pvec = \mat{D}\Pvec$, and the above equations become
\begin{equation}
  \mat{D}\Pvec = \Bmat{\bs{\Delta}_0\\ \mat{0}},\quad\text{i.e.}\quad
  \bmat{\mat{D}_{pp} & \mat{D}_{px}\\ \mat{D}_{xp} & \mat{D}_{xx}}\Bmat{\Pvec_0\\ \F_x} = \Bmat{\bs{\Delta}_0\\ \mat{0}}
\end{equation}
Solving for $\F_x$ from the lower part of equations results in
\begin{equation}
  \F_x = -\mat{D}_{xx}^{-1}\mat{D}_{xp}\Pvec_0
\end{equation}
Substituting the above expression for $\F_x$ into the upper part of equations yields
\begin{equation}
  \bs{\Delta}_0 = (\mat{D}_{pp} - \mat{D}_{px}\mat{D}_{xx}^{-1}\mat{D}_{xp})\Pvec_0 = \widehat{\mat{D}}_{pp}\Pvec_0
\end{equation}
where
\begin{equation}
  \widehat{\mat{D}}_{pp} = \mat{D}_{pp} - \mat{D}_{px}\mat{D}_{xx}^{-1}\mat{D}_{xp}
\end{equation}
Inverting the relation, $\bs{\Delta}_0 = \widehat{\mat{D}}_{pp}\Pvec_0$, yields the stiffness equations:
\begin{equation}
  \widehat{\Kmat}_{pp}\bs{\Delta}_0 = \Pvec_0
\end{equation}
where $\widehat{\Kmat}_{pp} = \widehat{\mat{D}}_{pp}^{-1}$. Note that $\widehat{\Kmat}_{pp}$ is the condensed stiffness matrix corresponding to $\Pvec_0$ if the stiffness method is employed.

\begin{example}[Truss]
Fig.~(a) shows a truss structure, and Fig.~(b) shows the primary structure, where $F_{bd}$ is the redundant force. Assuming all bars have the same $EA$, determine $F_{bd}$ and $u_c$.
\begin{fig}
  % Truss (a) and primary structure (b) with redundant F_bd
\begin{tikzpicture}[scale=0.8,font=\small]
  \begin{scope}
    \coordinate (a) at (0,0); \coordinate (b) at (0,4); \coordinate (c) at (4,4); \coordinate (d) at (4,0);
    \draw[member] (a) -- (b) -- (c) -- (d); \draw[member] (a) -- (c); \draw[member] (b) -- (d);
    \pic at (a) {pin}; \pic at (d) {pin}; \foreach \p in {a,b,c,d}{\node[node dot] at (\p) {};}
    \node[left] at (a) {$a$}; \node[above left] at (b) {$b$}; \node[above right] at (c) {$c$}; \node[right] at (d) {$d$};
    \draw[load] (c) -- ++(1.4,0) node[right] {$P$};
    \draw[dim] (-1.0,0) -- (-1.0,4) node[midway,left] {$L$}; \draw[dim] (0,-1.0) -- (4,-1.0) node[midway,fill=white,inner sep=1pt] {$L$};
    \node[below] at (2,-1.5) {(a)};
  \end{scope}
  \begin{scope}[xshift=8cm]
    \coordinate (a) at (0,0); \coordinate (b) at (0,4); \coordinate (c) at (4,4); \coordinate (d) at (4,0);
    \draw[member] (a) -- (b) -- (c) -- (d); \draw[member] (a) -- (c);
    \pic at (a) {pin}; \pic at (d) {pin}; \foreach \p in {a,b,c,d}{\node[node dot] at (\p) {};}
    \node[left] at (a) {$a$}; \node[above left] at (b) {$b$}; \node[above right] at (c) {$c$}; \node[right] at (d) {$d$};
    \draw[load] (c) -- ++(1.4,0) node[right] {$P$};
    \draw[load] (b) -- ++(-45:1.4); \draw[load] (d) -- ++(135:1.4); \node at (1.9,2.1) {$F_{bd}$};
    \draw[dim] (-1.0,0) -- (-1.0,4) node[midway,left] {$L$}; \draw[dim] (0,-1.0) -- (4,-1.0) node[midway,fill=white,inner sep=1pt] {$L$};
    \node[below] at (2,-1.5) {(b) Primary structure};
  \end{scope}
\end{tikzpicture}

\end{fig}
First, find the member forces for the primary structure in terms of $P$ and $F_{bd}$.
\begin{equation*}
  \F_0 = \Bmat{F_{ab}\\ F_{bc}\\ F_{cd}\\ F_{ac}} = \Bmat{0\\ 0\\ -1\\ \sqrt2}P + \Bmat{-1\\ -1\\ -1\\ \sqrt2}\frac{F_{bd}}{\sqrt2} = \mat{C}_1\Pvec_0 + \mat{C}_2\F_x,
\end{equation*}
where $\Pvec_0 = P$ and $\F_x = F_{bd}$.

Therefore,
\begin{equation*}
  \mat{C} = \bmat{\mat{C}_1 & \mat{C}_2\\ \mat{0} & \mat{I}} =
  \begin{bmatrix}
    0 & -\dfrac{1}{\sqrt2}\\[6pt] 0 & -\dfrac{1}{\sqrt2}\\[6pt] -1 & -\dfrac{1}{\sqrt2}\\[6pt] \sqrt2 & 1\\[4pt] 0 & 1
  \end{bmatrix}
\end{equation*}
The compatibility equations are
\begin{equation*}
  \mat{D}\Pvec = \Bmat{\bs{\Delta}_0\\ 0} \Rightarrow \mat{D}\Bmat{P\\ F_{bd}} = \Bmat{u_c\\ 0}
\end{equation*}
where $\mat{D} = \mat{C}^\tr\mat{d}\mat{C} = \dfrac{L}{EA}\begin{bmatrix}1+2\sqrt2 & 2+\dfrac{1}{\sqrt2}\\[6pt] 2+\dfrac{1}{\sqrt2} & \dfrac32+2\sqrt2\end{bmatrix}$ since
\begin{equation*}
  \mat{d} = \begin{bmatrix}
    \dfrac{L}{EA} & & & & \\[4pt] & \dfrac{L}{EA} & & & \\[4pt] & & \dfrac{L}{EA} & & \\[4pt] & & & \dfrac{\sqrt2L}{EA} & \\[4pt] & & & & \dfrac{\sqrt2L}{EA}
  \end{bmatrix}
\end{equation*}
Hence, $F_{bd}$ and $u_c$ can be solved from the compatibility equations; the results are:
\begin{equation*}
  \Bmat{F_{bd}\\ u_c} = \Bmat{\dfrac{-(4+\sqrt2)}{3+4\sqrt2}\\[8pt] \dfrac{L}{EA}\dfrac{10+6\sqrt2}{3+4\sqrt2}}P = \Bmat{-0.625P\\[4pt] 2.135\dfrac{PL}{EA}}
\end{equation*}
The axial forces in the other members:
\begin{equation*}
  \Bmat{F_{ab}\\ F_{bc}\\ F_{cd}\\ F_{ac}} = \mat{C}_1\Pvec_0 + \mat{C}_2\F_x = \Bmat{0.442\\ 0.442\\ -0.558\\ 0.789}P
\end{equation*}
Since only a horizontal load is applied at node $C$, the displacement that we can find corresponds to this load. If the displacements of other points or $v_c$ are needed, we shall add artificial loads first and then let them equal zero.
\end{example}

\begin{example}[Beam]
The beam is subjected to a load $P$ at the central point $b$; determine the vertical displacement of node $b$.
\begin{fig}
  % Propped cantilever with central load P and the element end moments
\begin{tikzpicture}[scale=0.7,font=\small]
  \begin{scope}
    \pic at (0,0) {wallL}; \draw[member] (0,0) -- (6,0); \pic at (6,0) {roller}; \node[node dot] at (6,0) {};
    \node[above right] at (0,0.05) {$a$}; \node[above right] at (3,0.05) {$b$}; \node[above left] at (6,0.05) {$C$};
    \draw[load] (3,1.3) -- (3,0.1) node[pos=0,above] {$P$};
    \draw[dim] (0,-1.0) -- (6,-1.0) node[midway,fill=white,inner sep=1pt] {$L$};
  \end{scope}
  \begin{scope}[xshift=8cm]
    \draw[thinmember] (0,0) -- (3.5,0); \draw[thinmember] (6.5,0) -- (10,0);
    \node[above right] at (0.05,0) {$a$}; \node[above left] at (3.5,0) {$b$}; \node[above right] at (6.55,0) {$b$}; \node[above left] at (10,0) {$C$};
    \draw[->,thick] (-0.15,0.45) arc (110:-110:0.45); \node[above] at (-0.2,0.6) {$M_1^{ab}$};
    \draw[->,thick] (3.65,0.45) arc (70:290:0.45); \node[above] at (3.7,0.6) {$M_2^{ab}$};
    \draw[->,thick] (6.35,0.45) arc (110:-110:0.45); \node[above] at (6.4,0.6) {$M_1^{bc}$};
    \draw[->,thick] (10.15,0.45) arc (70:290:0.45); \node[above] at (10.2,0.6) {$M_2^{bc}$};
  \end{scope}
\end{tikzpicture}

\end{fig}
\emph{Case 1: Take the reaction at the roller as the redundant force.}
\begin{fig}
  % Case 1: reaction at the roller as the redundant force
\begin{tikzpicture}[scale=0.7,font=\small]
  \begin{scope}
    \pic at (0,0) {wallL}; \draw[member] (0,0) -- (6,0);
    \node[above right] at (0,0.05) {$a$}; \node[above right] at (3,0.05) {$b$}; \node[above left] at (6,0.05) {$C$};
    \draw[load] (3,1.3) -- (3,0.1) node[pos=0,above] {$P$};
    \draw[load] (6,-1.2) -- (6,-0.1) node[pos=0,below] {$R$};
  \end{scope}
  \begin{scope}[xshift=8cm]
    \draw[thinmember] (0,0) -- (3.5,0); \draw[thinmember] (6.5,0) -- (10,0);
    \node[above right] at (0.05,0) {$a$}; \node[above left] at (3.5,0) {$b$}; \node[above right] at (6.55,0) {$b$}; \node[above left] at (10,0) {$C$};
    \draw[->,thick] (-0.15,0.45) arc (110:-110:0.45); \node[above] at (-0.2,0.6) {$\frac{PL}{2}-RL$};
    \draw[->,thick] (3.65,0.45) arc (70:290:0.45); \node[above] at (3.7,0.6) {$\frac{RL}{2}$};
    \draw[->,thick] (6.35,0.45) arc (110:-110:0.45); \node[above] at (6.4,0.6) {$-\frac{RL}{2}$};
    \draw[->,thick] (10.15,0.45) arc (70:290:0.45); \node[above] at (10.2,0.6) {$0$};
  \end{scope}
\end{tikzpicture}

\end{fig}
First let $\F = \Bmat{\F_0\\ \F_x} = \Bmat{M_1^{ab}\\ M_2^{ab}\\ M_1^{bc}\\ M_2^{bc}\\ R}$. Then analyze this beam and express the element end moments in terms of $P$ and $R$; namely
\begin{equation*}
  \Bmat{M_1^{ab}\\ M_2^{ab}\\ M_1^{bc}\\ M_2^{bc}} = \Bmat{L/2\\ 0\\ 0\\ 0}P + \Bmat{-L\\ L/2\\ -L/2\\ 0}R = \mat{C}_1\Pvec_0 + \mat{C}_2\F_x,\quad\text{where } \mat{C}_1 = \Bmat{L/2\\ 0\\ 0\\ 0}\text{ and } \mat{C}_2 = \Bmat{-L\\ L/2\\ -L/2\\ 0}.
\end{equation*}
Note that $\Pvec_0 = P$ and $\F_x = R$.

We can obtain the structural flexibility matrix as
\begin{equation*}
  \mat{D} = \mat{C}^\tr\mat{d}\mat{C} = \frac{L^3}{EI}\bmat{1/24 & -5/48\\ -5/48 & 1/3}
\end{equation*}
where
\begin{equation*}
  \mat{d} = \bmat{\mat{d}_{ab} & & \\ & \mat{d}_{bc} & \\ & & 0} = \frac{L'}{6EI}\bmat{2 & -1 & & & \\ -1 & 2 & & & \\ & & 2 & -1 & \\ & & -1 & 2 & \\ & & & & 0}\quad\text{with } L' = \frac L2
\end{equation*}
and
\begin{equation*}
  \mat{C}_{5\times2} = \bmat{\mat{C}_1 & \mat{C}_2\\ 0 & 1}
\end{equation*}
Note that the flexibility coefficient for the roller ($d_{55}$) is 0 since it is rigid. Finally, the system compatibility equations can be written as
\begin{equation*}
  \mat{D}\Bmat{P\\ R} = \Bmat{v_b\\ 0},
\end{equation*}
from which we obtain
\begin{equation*}
  \Bmat{R\\ v_b} = \Bmat{\dfrac{5}{16}P\\[6pt] 9.11\times10^{-3}\dfrac{PL^3}{EI}}
\end{equation*}

\emph{Case 2: Take the moment at the fixed end as the redundant force.}
\begin{fig}
  % Case 2: moment at the fixed end as the redundant force
\begin{tikzpicture}[scale=0.7,font=\small]
  \begin{scope}
    \draw[member] (0,0) -- (6,0); \pic at (0,0) {pin}; \pic at (6,0) {roller}; \node[node dot] at (0,0) {}; \node[node dot] at (6,0) {};
    \node[above right] at (0.1,0.05) {$a$}; \node[above right] at (3,0.05) {$b$}; \node[above left] at (5.9,0.05) {$C$};
    \draw[load] (3,1.3) -- (3,0.1) node[pos=0,above] {$P$};
    \draw[->,thick] (-0.5,0.7) arc (100:260:0.5); \node[above] at (-0.6,0.75) {$R$};
    \draw[dim] (0,-1.0) -- (6,-1.0) node[midway,fill=white,inner sep=1pt] {$L$};
  \end{scope}
  \begin{scope}[xshift=8cm]
    \draw[thinmember] (0,0) -- (3.5,0); \draw[thinmember] (6.5,0) -- (10,0);
    \node[above right] at (0.05,0) {$a$}; \node[above left] at (3.5,0) {$b$}; \node[above right] at (6.55,0) {$b$}; \node[above left] at (10,0) {$C$};
    \draw[->,thick] (-0.15,0.45) arc (110:-110:0.45); \node[above] at (-0.2,0.6) {$R$};
    \draw[->,thick] (3.65,0.45) arc (70:290:0.45); \node[above] at (3.5,0.6) {$\frac{PL}{4}-\frac R2$};
    \draw[->,thick] (6.35,0.45) arc (110:-110:0.45); \node[above] at (6.8,0.6) {$-\frac{PL}{4}+\frac R2$};
    \draw[->,thick] (10.15,0.45) arc (70:290:0.45); \node[above] at (10.2,0.6) {$0$};
  \end{scope}
\end{tikzpicture}

\end{fig}
Following the presented procedure, we will have
\begin{equation*}
  \mat{C}_1 = \Bmat{0\\ \dfrac L4\\[4pt] -\dfrac L4\\[4pt] 0},\quad \mat{C}_2 = \Bmat{1\\ -\dfrac12\\[4pt] \dfrac12\\[4pt] 0}\quad\text{and}\quad \mat{D} = \frac{L^3}{EI}\begin{bmatrix}\dfrac{1}{48} & -\dfrac{1}{16L}\\[8pt] -\dfrac{1}{16L} & \dfrac{1}{3L^2}\end{bmatrix}.
\end{equation*}
Finally, we have
\begin{equation*}
  \Bmat{R\\ v_b} = \Bmat{\dfrac{3}{16}PL\\[6pt] 9.11\times10^{-3}\dfrac{PL^3}{EI}}
\end{equation*}
The above two cases use end moments as the DOFs. When the moment and force of the right end are chosen to be the DOFs, we can obtain the same solutions. For example, when the fixed end moment is chosen as the redundant force (case 2), we have
\begin{equation*}
  \mat{d} = \bmat{\mat{d}_{ab} & & \\ & \mat{d}_{bc} & \\ & & 0},\qquad \mat{d}_{ab} = \mat{d}_{bc} = \frac{L'}{6EI}\bmat{2L'^2 & 3L'\\ 3L' & 6},\ \text{where } L' = \frac L2.
\end{equation*}
\begin{equation*}
  \mat{C}_1 = \Bmat{-\dfrac12\\[4pt] \dfrac L4\\[4pt] \dfrac12\\[4pt] 0},\quad \mat{C}_2 = \Bmat{-\dfrac1L\\[4pt] -\dfrac12\\[4pt] -\dfrac1L\\[4pt] 0},\quad\text{and}\quad \mat{D} = \frac{L^3}{EI}\begin{bmatrix}\dfrac{1}{48} & -\dfrac{1}{16L}\\[8pt] -\dfrac{1}{16L} & \dfrac{1}{3L^2}\end{bmatrix}
\end{equation*}
Note that the same $\mat{D}$ is obtained.
\end{example}

\begin{example}[Beam with an internal spring]
The structure is similar to the previous example except that an internal rotational spring is added at the loaded point.
\begin{fig}
  % Propped cantilever with an internal rotational spring at the loaded point
\begin{tikzpicture}[scale=0.8,font=\small]
  \pic at (0,0) {wallL}; \draw[member] (0,0) -- (6,0); \pic at (6,0) {roller}; \node[node dot] at (6,0) {};
  \draw[thick] (2,0) circle (0.18); \draw[thick] (2,0) circle (0.09);
  \node[above right] at (0,0.05) {$a$}; \node[above right] at (2.1,0.1) {$b$}; \node[above left] at (6,0.05) {$C$};
  \draw[load] (2,0.2) -- (2,1.3) node[above] {$P$};
  \draw[dim] (0,-1.0) -- (2,-1.0) node[midway,fill=white,inner sep=1pt] {$L$}; \draw[dim] (2,-1.0) -- (6,-1.0) node[midway,fill=white,inner sep=1pt] {$2L$};
\end{tikzpicture}

\end{fig}
Take the reaction at the roller as the redundant force.
\begin{fig}
  % Primary structure (redundant R at the roller) and element end moments incl. the spring moment
\begin{tikzpicture}[scale=0.8,font=\small]
  \begin{scope}
    \pic at (0,0) {wallL}; \draw[member] (0,0) -- (6,0);
    \draw[thick] (2,0) circle (0.18); \draw[thick] (2,0) circle (0.09);
    \node[above right] at (0,0.05) {$a$}; \node[above right] at (2.1,0.1) {$b$}; \node[above left] at (6,0.05) {$C$};
    \draw[load] (2,0.2) -- (2,1.3) node[above] {$P$};
    \draw[load] (6,-1.2) -- (6,-0.1) node[pos=0,left] {$R$};
  \end{scope}
  \begin{scope}[xshift=7.5cm]
    \draw[thinmember] (0,0) -- (2.2,0); \draw[thick] (3.1,0) circle (0.16); \draw[thick] (3.1,0) circle (0.08); \draw[thinmember] (4.0,0) -- (8.0,0);
    \node[above right] at (0.05,0) {$a$}; \node[above left] at (2.2,0) {$b$}; \node[above right] at (4.05,0) {$b$}; \node[above left] at (8.0,0) {$C$};
    \draw[->,thick] (-0.15,0.45) arc (110:-110:0.45); \node[above] at (-0.2,0.6) {$M_1^{ab}$};
    \draw[->,thick] (2.35,0.45) arc (70:290:0.45); \node[above] at (2.4,0.6) {$M_2^{ab}$};
    \draw[->,thick] (3.6,0.45) arc (70:290:0.45); \node[above] at (3.4,0.6) {$M_s$};
    \draw[->,thick] (3.85,0.45) arc (110:-110:0.45); \node[above] at (4.4,0.6) {$M_1^{bc}$};
    \draw[->,thick] (8.15,0.45) arc (70:290:0.45); \node[above] at (8.2,0.6) {$M_2^{bc}$};
  \end{scope}
\end{tikzpicture}

\end{fig}
\begin{equation*}
  \F = \Bmat{\F_0\\ \F_x} = \Bmat{M_1^{ab}\\ M_2^{ab}\\ M_1^{bc}\\ M_2^{bc}\\ M_s\\ R},\quad\text{and}\quad
  \mat{d} = \bmat{\mat{d}_{ab} & & & \\ & \mat{d}_{bc} & & \\ & & \dfrac{1}{k_\theta} & \\ & & & 0}
\end{equation*}
Using the same procedure, we will obtain $v_b = \dfrac19\dfrac{PL^3}{EI}\dfrac{11\alpha+12}{9\alpha+4}$. Try it.
\end{example}

\begin{example}[2D frame]
\leavevmode

\begin{minipage}{0.6\linewidth}
Even for 2D frames, the same procedure can be used. The DOFs for one member are shown in the figure, and the flexibility matrix is
\begin{equation*}
  \mat{d} = \begin{bmatrix}\dfrac{L}{EA} & & \\[6pt] & \dfrac{L}{3EI} & -\dfrac{L}{6EI}\\[8pt] & -\dfrac{L}{6EI} & \dfrac{L}{3EI}\end{bmatrix}
\end{equation*}
\end{minipage}\hfill
\begin{minipage}{0.35\linewidth}\centering% 2D frame: the three member DOFs (axial + two end rotations) of the left column
\begin{tikzpicture}[scale=0.8,font=\small]
  \draw[member] (0,0) -- (0,4) -- (4,4) -- (4,0);
  \pic at (0,0) {pin}; \pic at (4,0) {fixed}; \node[node dot] at (0,0) {};
  \draw[load] (0,4.1) -- (0,5.2);
  \draw[->,thick] (-0.5,4.4) arc (150:-30:0.5);
  \draw[->,thick] (0.5,0.4) arc (30:210:0.5);
\end{tikzpicture}
\end{minipage}

This problem was solved by the stiffness method in Sec.~\ref{sec:axial-portal}.
\begin{fig}
  % Portal frame solved by the flexibility method: redundants at d
\begin{tikzpicture}[scale=0.8,font=\small]
  \draw[member] (0,0) -- (0,4) -- (4,4) -- (4,0);
  \pic at (0,0) {fixed};
  \node[left] at (0,0.3) {$a$}; \node[above right] at (0,4) {$b$}; \node[above left] at (4,4) {$c$}; \node[above left] at (4,0.3) {$d$};
  \draw[load] (-1.4,4) -- (-0.1,4) node[pos=0,below] {$P$};
  \draw[load] (4.1,0) -- (5.4,0) node[right] {$F_{x1}$};
  \draw[load] (4,-1.3) -- (4,-0.1) node[pos=0,right] {$F_{x2}$};
  \draw[->,thick] (4.6,-0.4) arc (-40:200:0.6); \node[left] at (3.4,-0.4) {$F_{x3}$};
\end{tikzpicture}

\end{fig}
\begin{equation*}
  \mat{C} = \bigl[(\mat{C}_1)_{9\times1}, (\mat{C}_2)_{9\times3}\bigr] =
  \begin{bmatrix}
    0 & 0 & 1 & 0\\ L & 0 & -L & -1\\ 0 & L & L & 1\\ 0 & 1 & 0 & 0\\ 0 & -L & -L & -1\\ 0 & L & 0 & 1\\ 0 & 0 & -1 & 0\\ 0 & -L & 0 & -1\\ 0 & 0 & 0 & 1
  \end{bmatrix}
\end{equation*}
\begin{equation*}
  \F_0 = \Bmat{\F^{ab}\\ \F^{bc}\\ \F^{cd}};\quad\text{member forces of each member: } \F = \Bmat{F\\ M_1\\ M_2}
\end{equation*}
{\small
\begin{equation*}
  \mat{d} = \bmat{\mat{d}_{ab} & \mat{0} & \mat{0}\\ \mat{0} & \mat{d}_{bc} & \mat{0}\\ \mat{0} & \mat{0} & \mat{d}_{cd}} =
  \begin{bmatrix}
    \frac{L}{EA} & 0 & 0 & 0 & 0 & 0 & 0 & 0 & 0\\
    0 & \frac{L}{3EI} & \frac{-L}{6EI} & 0 & 0 & 0 & 0 & 0 & 0\\
    0 & \frac{-L}{6EI} & \frac{L}{3EI} & 0 & 0 & 0 & 0 & 0 & 0\\
    0 & 0 & 0 & \frac{L}{EA} & 0 & 0 & 0 & 0 & 0\\
    0 & 0 & 0 & 0 & \frac{L}{3EI} & \frac{-L}{6EI} & 0 & 0 & 0\\
    0 & 0 & 0 & 0 & \frac{-L}{6EI} & \frac{L}{3EI} & 0 & 0 & 0\\
    0 & 0 & 0 & 0 & 0 & 0 & \frac{L}{EA} & 0 & 0\\
    0 & 0 & 0 & 0 & 0 & 0 & 0 & \frac{L}{3EI} & \frac{-L}{6EI}\\
    0 & 0 & 0 & 0 & 0 & 0 & 0 & \frac{-L}{6EI} & \frac{L}{3EI}
  \end{bmatrix}
\end{equation*}
\begin{equation*}
  \mat{D} = \mat{C}^\tr\mat{d}\mat{C} =
  \begin{bmatrix}
    \frac13\frac{L^3}{EI} & -\frac16\frac{L^3}{EI} & -\frac12\frac{L^3}{EI} & -\frac12\frac{L^2}{EI}\\[4pt]
    -\frac16\frac{L^3}{EI} & \frac53\frac{L^3}{EI}+\frac{L}{EA} & \frac{L^3}{EI} & 2\frac{L^2}{EI}\\[4pt]
    -\frac12\frac{L^3}{EI} & \frac{L^3}{EI} & 2\frac{L}{EA}+\frac43\frac{L^3}{EI} & \frac32\frac{L^2}{EI}\\[4pt]
    -\frac12\frac{L^2}{EI} & 2\frac{L^2}{EI} & \frac32\frac{L^2}{EI} & 3\frac{L}{EI}
  \end{bmatrix};\qquad
  \Pvec = \Bmat{P\\ F_{x1}\\ F_{x2}\\ F_{x3}}
\end{equation*}}
System flexibility equations:
\begin{equation*}
  \mat{D}\Pvec = \Bmat{u_b\\ 0\\ 0\\ 0}\quad\rightarrow\quad\text{the last three equations:}
\end{equation*}
\begin{equation*}
  \begin{bmatrix}
    \frac53\frac{L^3}{EI}+\frac{L}{EA} & \frac{L^3}{EI} & 2\frac{L^2}{EI}\\[4pt]
    \frac{L^3}{EI} & 2\frac{L}{EA}+\frac43\frac{L^3}{EI} & \frac32\frac{L^2}{EI}\\[4pt]
    2\frac{L^2}{EI} & \frac32\frac{L^2}{EI} & 3\frac{L}{EI}
  \end{bmatrix}\Bmat{F_{x1}\\ F_{x2}\\ F_{x3}} = \Bmat{\frac16\frac{L^3}{EI}\\[4pt] \frac12\frac{L^3}{EI}\\[4pt] \frac12\frac{L^2}{EI}}P
\end{equation*}
\begin{equation*}
  \rightarrow\quad \Bmat{F_{x1}\\ F_{x2}\\ F_{x3}} = \Bmat{-\frac12L^2\frac{A}{L^2A+3I}\\[6pt] 3L^2\frac{A}{24I+7L^2A}\\[6pt] L\frac{11L^2AI+2L^4A^2+12I^2}{(24I+7L^2A)(L^2A+3I)}}P
\end{equation*}
$\rightarrow$ Substituting them into the first equation obtains:
\begin{equation*}
  u_b = \frac{1}{12}L^3\frac{84L^2AI+216I^2+5L^4A^2}{EI(45L^2AI+72I^2+7L^4A^2)}
\end{equation*}
The above solution is the same as that obtained in Sec.~\ref{sec:axial-portal}.

\topic{Neglect of axial deformation:}
\begin{equation*}
  \frac{L}{EA} = 0\ \rightarrow\ \mat{d} =
  \begin{bmatrix}
    \frac{L}{3EI} & \frac{-L}{6EI} & 0 & 0 & 0 & 0\\
    \frac{-L}{6EI} & \frac{L}{3EI} & 0 & 0 & 0 & 0\\
    0 & 0 & \frac{L}{3EI} & \frac{-L}{6EI} & 0 & 0\\
    0 & 0 & \frac{-L}{6EI} & \frac{L}{3EI} & 0 & 0\\
    0 & 0 & 0 & 0 & \frac{L}{3EI} & \frac{-L}{6EI}\\
    0 & 0 & 0 & 0 & \frac{-L}{6EI} & \frac{L}{3EI}
  \end{bmatrix};
\end{equation*}
\begin{equation*}
  \mat{C} = \bigl[(\mat{C}_1)_{6\times1}, (\mat{C}_2)_{6\times3}\bigr] =
  \begin{bmatrix}
    L & 0 & -L & -1\\ 0 & L & L & 1\\ 0 & -L & -L & -1\\ 0 & L & 0 & 1\\ 0 & -L & 0 & -1\\ 0 & 0 & 0 & 1
  \end{bmatrix}
\end{equation*}
\begin{equation*}
  \mat{D} = \begin{bmatrix}
    \frac13\frac{L^3}{EI} & -\frac16\frac{L^3}{EI} & -\frac12\frac{L^3}{EI} & -\frac12\frac{L^2}{EI}\\[4pt]
    -\frac16\frac{L^3}{EI} & \frac53\frac{L^3}{EI} & \frac{L^3}{EI} & 2\frac{L^2}{EI}\\[4pt]
    -\frac12\frac{L^3}{EI} & \frac{L^3}{EI} & \frac43\frac{L^3}{EI} & \frac32\frac{L^2}{EI}\\[4pt]
    -\frac12\frac{L^2}{EI} & 2\frac{L^2}{EI} & \frac32\frac{L^2}{EI} & 3\frac{L}{EI}
  \end{bmatrix}
\end{equation*}
The last three equations:
\begin{equation*}
  \begin{bmatrix}
    \frac53\frac{L^3}{EI} & \frac{L^3}{EI} & 2\frac{L^2}{EI}\\[4pt]
    \frac{L^3}{EI} & \frac43\frac{L^3}{EI} & \frac32\frac{L^2}{EI}\\[4pt]
    2\frac{L^2}{EI} & \frac32\frac{L^2}{EI} & 3\frac{L}{EI}
  \end{bmatrix}\Bmat{F_{x1}\\ F_{x2}\\ F_{x3}} = \Bmat{\frac16\frac{L^3}{EI}\\[4pt] \frac12\frac{L^3}{EI}\\[4pt] \frac12\frac{L^2}{EI}}P
  \quad\rightarrow\quad \Bmat{F_{x1}\\ F_{x2}\\ F_{x3}} = \Bmat{-\frac12\\[4pt] \frac37\\[4pt] \frac27L}P
\end{equation*}
$\rightarrow$ Substituting them into the first equation obtains: $u_b = \dfrac{5}{84}\dfrac{PL^3}{EI}$.

This value is the same as that given in Sec.~\ref{sec:axial-portal}.
\end{example}


\end{document}
